\documentclass[11pt]{article}

    
\usepackage[breakable]{tcolorbox}
    \usepackage{parskip} % Stop auto-indenting (to mimic markdown behaviour)
    

    % Basic figure setup, for now with no caption control since it's done
    % automatically by Pandoc (which extracts ![](path) syntax from Markdown).
    \usepackage{graphicx}
    % Keep aspect ratio if custom image width or height is specified
    \setkeys{Gin}{keepaspectratio}
    % Maintain compatibility with old templates. Remove in nbconvert 6.0
    \let\Oldincludegraphics\includegraphics
    % Ensure that by default, figures have no caption (until we provide a
    % proper Figure object with a Caption API and a way to capture that
    % in the conversion process - todo).
    \usepackage{caption}
    \DeclareCaptionFormat{nocaption}{}
    \captionsetup{format=nocaption,aboveskip=0pt,belowskip=0pt}

    \usepackage{float}
    \floatplacement{figure}{H} % forces figures to be placed at the correct location
    \usepackage{xcolor} % Allow colors to be defined
    \usepackage{enumerate} % Needed for markdown enumerations to work
    \usepackage{geometry} % Used to adjust the document margins
    \usepackage{amsmath} % Equations
    \usepackage{amssymb} % Equations
    \usepackage{textcomp} % defines textquotesingle
    % Hack from http://tex.stackexchange.com/a/47451/13684:
    \AtBeginDocument{%
        \def\PYZsq{\textquotesingle}% Upright quotes in Pygmentized code
    }
    \usepackage{upquote} % Upright quotes for verbatim code
    \usepackage{eurosym} % defines \euro

    \usepackage{iftex}
    \ifPDFTeX
        \usepackage[T1]{fontenc}
        \IfFileExists{alphabeta.sty}{
              \usepackage{alphabeta}
          }{
              \usepackage[mathletters]{ucs}
              \usepackage[utf8x]{inputenc}
          }
    \else
        \usepackage{fontspec}
        \usepackage{unicode-math}
    \fi

    \usepackage{fancyvrb} % verbatim replacement that allows latex
    \usepackage{grffile} % extends the file name processing of package graphics
                         % to support a larger range
    \makeatletter % fix for old versions of grffile with XeLaTeX
    \@ifpackagelater{grffile}{2019/11/01}
    {
      % Do nothing on new versions
    }
    {
      \def\Gread@@xetex#1{%
        \IfFileExists{"\Gin@base".bb}%
        {\Gread@eps{\Gin@base.bb}}%
        {\Gread@@xetex@aux#1}%
      }
    }
    \makeatother
    \usepackage[Export]{adjustbox} % Used to constrain images to a maximum size
    \adjustboxset{max size={0.9\linewidth}{0.9\paperheight}}

    % The hyperref package gives us a pdf with properly built
    % internal navigation ('pdf bookmarks' for the table of contents,
    % internal cross-reference links, web links for URLs, etc.)
    \usepackage{hyperref}
    % The default LaTeX title has an obnoxious amount of whitespace. By default,
    % titling removes some of it. It also provides customization options.
    \usepackage{titling}
    \usepackage{longtable} % longtable support required by pandoc >1.10
    \usepackage{booktabs}  % table support for pandoc > 1.12.2
    \usepackage{array}     % table support for pandoc >= 2.11.3
    \usepackage{calc}      % table minipage width calculation for pandoc >= 2.11.1
    \usepackage[inline]{enumitem} % IRkernel/repr support (it uses the enumerate* environment)
    \usepackage[normalem]{ulem} % ulem is needed to support strikethroughs (\sout)
                                % normalem makes italics be italics, not underlines
    \usepackage{soul}      % strikethrough (\st) support for pandoc >= 3.0.0
    \usepackage{mathrsfs}
    
\usepackage[australian]{babel}
\usepackage{datetime}
\usepackage{fmtcount}
% Row shading for pandas Styler LaTeX tables (xcolor is loaded by nbconvert).
\usepackage{colortbl}
\usepackage{notebook_pdf/templates/labstyle}

% Australian date format: 26 October 2026
%\renewcommand{\today}{\number\day\space\monthname\space\number\year}
% Example: 22nd October 2026
\renewcommand{\today}{
    \ordinalnum{\day}\space\monthname\space\number\year%
}


    
    % Colors for the hyperref package
    \definecolor{urlcolor}{rgb}{0,.145,.698}
    \definecolor{linkcolor}{rgb}{.71,0.21,0.01}
    \definecolor{citecolor}{rgb}{.12,.54,.11}

    % ANSI colors
    \definecolor{ansi-black}{HTML}{3E424D}
    \definecolor{ansi-black-intense}{HTML}{282C36}
    \definecolor{ansi-red}{HTML}{E75C58}
    \definecolor{ansi-red-intense}{HTML}{B22B31}
    \definecolor{ansi-green}{HTML}{00A250}
    \definecolor{ansi-green-intense}{HTML}{007427}
    \definecolor{ansi-yellow}{HTML}{DDB62B}
    \definecolor{ansi-yellow-intense}{HTML}{B27D12}
    \definecolor{ansi-blue}{HTML}{208FFB}
    \definecolor{ansi-blue-intense}{HTML}{0065CA}
    \definecolor{ansi-magenta}{HTML}{D160C4}
    \definecolor{ansi-magenta-intense}{HTML}{A03196}
    \definecolor{ansi-cyan}{HTML}{60C6C8}
    \definecolor{ansi-cyan-intense}{HTML}{258F8F}
    \definecolor{ansi-white}{HTML}{C5C1B4}
    \definecolor{ansi-white-intense}{HTML}{A1A6B2}
    \definecolor{ansi-default-inverse-fg}{HTML}{FFFFFF}
    \definecolor{ansi-default-inverse-bg}{HTML}{000000}

    % common color for the border for error outputs.
    \definecolor{outerrorbackground}{HTML}{FFDFDF}

    % commands and environments needed by pandoc snippets
    % extracted from the output of `pandoc -s`
    \providecommand{\tightlist}{%
      \setlength{\itemsep}{0pt}\setlength{\parskip}{0pt}}
    \DefineVerbatimEnvironment{Highlighting}{Verbatim}{commandchars=\\\{\}}
    % Add ',fontsize=\small' for more characters per line
    \newenvironment{Shaded}{}{}
    \newcommand{\KeywordTok}[1]{\textcolor[rgb]{0.00,0.44,0.13}{\textbf{{#1}}}}
    \newcommand{\DataTypeTok}[1]{\textcolor[rgb]{0.56,0.13,0.00}{{#1}}}
    \newcommand{\DecValTok}[1]{\textcolor[rgb]{0.25,0.63,0.44}{{#1}}}
    \newcommand{\BaseNTok}[1]{\textcolor[rgb]{0.25,0.63,0.44}{{#1}}}
    \newcommand{\FloatTok}[1]{\textcolor[rgb]{0.25,0.63,0.44}{{#1}}}
    \newcommand{\CharTok}[1]{\textcolor[rgb]{0.25,0.44,0.63}{{#1}}}
    \newcommand{\StringTok}[1]{\textcolor[rgb]{0.25,0.44,0.63}{{#1}}}
    \newcommand{\CommentTok}[1]{\textcolor[rgb]{0.38,0.63,0.69}{\textit{{#1}}}}
    \newcommand{\OtherTok}[1]{\textcolor[rgb]{0.00,0.44,0.13}{{#1}}}
    \newcommand{\AlertTok}[1]{\textcolor[rgb]{1.00,0.00,0.00}{\textbf{{#1}}}}
    \newcommand{\FunctionTok}[1]{\textcolor[rgb]{0.02,0.16,0.49}{{#1}}}
    \newcommand{\RegionMarkerTok}[1]{{#1}}
    \newcommand{\ErrorTok}[1]{\textcolor[rgb]{1.00,0.00,0.00}{\textbf{{#1}}}}
    \newcommand{\NormalTok}[1]{{#1}}

    % Additional commands for more recent versions of Pandoc
    \newcommand{\ConstantTok}[1]{\textcolor[rgb]{0.53,0.00,0.00}{{#1}}}
    \newcommand{\SpecialCharTok}[1]{\textcolor[rgb]{0.25,0.44,0.63}{{#1}}}
    \newcommand{\VerbatimStringTok}[1]{\textcolor[rgb]{0.25,0.44,0.63}{{#1}}}
    \newcommand{\SpecialStringTok}[1]{\textcolor[rgb]{0.73,0.40,0.53}{{#1}}}
    \newcommand{\ImportTok}[1]{{#1}}
    \newcommand{\DocumentationTok}[1]{\textcolor[rgb]{0.73,0.13,0.13}{\textit{{#1}}}}
    \newcommand{\AnnotationTok}[1]{\textcolor[rgb]{0.38,0.63,0.69}{\textbf{\textit{{#1}}}}}
    \newcommand{\CommentVarTok}[1]{\textcolor[rgb]{0.38,0.63,0.69}{\textbf{\textit{{#1}}}}}
    \newcommand{\VariableTok}[1]{\textcolor[rgb]{0.10,0.09,0.49}{{#1}}}
    \newcommand{\ControlFlowTok}[1]{\textcolor[rgb]{0.00,0.44,0.13}{\textbf{{#1}}}}
    \newcommand{\OperatorTok}[1]{\textcolor[rgb]{0.40,0.40,0.40}{{#1}}}
    \newcommand{\BuiltInTok}[1]{{#1}}
    \newcommand{\ExtensionTok}[1]{{#1}}
    \newcommand{\PreprocessorTok}[1]{\textcolor[rgb]{0.74,0.48,0.00}{{#1}}}
    \newcommand{\AttributeTok}[1]{\textcolor[rgb]{0.49,0.56,0.16}{{#1}}}
    \newcommand{\InformationTok}[1]{\textcolor[rgb]{0.38,0.63,0.69}{\textbf{\textit{{#1}}}}}
    \newcommand{\WarningTok}[1]{\textcolor[rgb]{0.38,0.63,0.69}{\textbf{\textit{{#1}}}}}
    \makeatletter
    \newsavebox\pandoc@box
    \newcommand*\pandocbounded[1]{%
      \sbox\pandoc@box{#1}%
      % scaling factors for width and height
      \Gscale@div\@tempa\textheight{\dimexpr\ht\pandoc@box+\dp\pandoc@box\relax}%
      \Gscale@div\@tempb\linewidth{\wd\pandoc@box}%
      % select the smaller of both
      \ifdim\@tempb\p@<\@tempa\p@
        \let\@tempa\@tempb
      \fi
      % scaling accordingly (\@tempa < 1)
      \ifdim\@tempa\p@<\p@
        \scalebox{\@tempa}{\usebox\pandoc@box}%
      % scaling not needed, use as it is
      \else
        \usebox{\pandoc@box}%
      \fi
    }
    \makeatother

    % Define a nice break command that doesn't care if a line doesn't already
    % exist.
    \def\br{\hspace*{\fill} \\* }
    % Math Jax compatibility definitions
    \def\gt{>}
    \def\lt{<}
    \let\Oldtex\TeX
    \let\Oldlatex\LaTeX
    \renewcommand{\TeX}{\textrm{\Oldtex}}
    \renewcommand{\LaTeX}{\textrm{\Oldlatex}}
    % Document parameters
    % Document title
    

\title{Week 4 Lesson 5 — Variance, Covariance, Eigenvalues and Eigenvectors}

    
    
    
    
    
    \author{Álvaro Morán Gago}
    
    
    
% Pygments definitions
\makeatletter
\def\PY@reset{\let\PY@it=\relax \let\PY@bf=\relax%
    \let\PY@ul=\relax \let\PY@tc=\relax%
    \let\PY@bc=\relax \let\PY@ff=\relax}
\def\PY@tok#1{\csname PY@tok@#1\endcsname}
\def\PY@toks#1+{\ifx\relax#1\empty\else%
    \PY@tok{#1}\expandafter\PY@toks\fi}
\def\PY@do#1{\PY@bc{\PY@tc{\PY@ul{%
    \PY@it{\PY@bf{\PY@ff{#1}}}}}}}
\def\PY#1#2{\PY@reset\PY@toks#1+\relax+\PY@do{#2}}

\@namedef{PY@tok@w}{\def\PY@tc##1{\textcolor[rgb]{0.73,0.73,0.73}{##1}}}
\@namedef{PY@tok@c}{\let\PY@it=\textit\def\PY@tc##1{\textcolor[rgb]{0.24,0.48,0.48}{##1}}}
\@namedef{PY@tok@cp}{\def\PY@tc##1{\textcolor[rgb]{0.61,0.40,0.00}{##1}}}
\@namedef{PY@tok@k}{\let\PY@bf=\textbf\def\PY@tc##1{\textcolor[rgb]{0.00,0.50,0.00}{##1}}}
\@namedef{PY@tok@kp}{\def\PY@tc##1{\textcolor[rgb]{0.00,0.50,0.00}{##1}}}
\@namedef{PY@tok@kt}{\def\PY@tc##1{\textcolor[rgb]{0.69,0.00,0.25}{##1}}}
\@namedef{PY@tok@o}{\def\PY@tc##1{\textcolor[rgb]{0.40,0.40,0.40}{##1}}}
\@namedef{PY@tok@ow}{\let\PY@bf=\textbf\def\PY@tc##1{\textcolor[rgb]{0.67,0.13,1.00}{##1}}}
\@namedef{PY@tok@nb}{\def\PY@tc##1{\textcolor[rgb]{0.00,0.50,0.00}{##1}}}
\@namedef{PY@tok@nf}{\def\PY@tc##1{\textcolor[rgb]{0.00,0.00,1.00}{##1}}}
\@namedef{PY@tok@nc}{\let\PY@bf=\textbf\def\PY@tc##1{\textcolor[rgb]{0.00,0.00,1.00}{##1}}}
\@namedef{PY@tok@nn}{\let\PY@bf=\textbf\def\PY@tc##1{\textcolor[rgb]{0.00,0.00,1.00}{##1}}}
\@namedef{PY@tok@ne}{\let\PY@bf=\textbf\def\PY@tc##1{\textcolor[rgb]{0.80,0.25,0.22}{##1}}}
\@namedef{PY@tok@nv}{\def\PY@tc##1{\textcolor[rgb]{0.10,0.09,0.49}{##1}}}
\@namedef{PY@tok@no}{\def\PY@tc##1{\textcolor[rgb]{0.53,0.00,0.00}{##1}}}
\@namedef{PY@tok@nl}{\def\PY@tc##1{\textcolor[rgb]{0.46,0.46,0.00}{##1}}}
\@namedef{PY@tok@ni}{\let\PY@bf=\textbf\def\PY@tc##1{\textcolor[rgb]{0.44,0.44,0.44}{##1}}}
\@namedef{PY@tok@na}{\def\PY@tc##1{\textcolor[rgb]{0.41,0.47,0.13}{##1}}}
\@namedef{PY@tok@nt}{\let\PY@bf=\textbf\def\PY@tc##1{\textcolor[rgb]{0.00,0.50,0.00}{##1}}}
\@namedef{PY@tok@nd}{\def\PY@tc##1{\textcolor[rgb]{0.67,0.13,1.00}{##1}}}
\@namedef{PY@tok@s}{\def\PY@tc##1{\textcolor[rgb]{0.73,0.13,0.13}{##1}}}
\@namedef{PY@tok@sd}{\let\PY@it=\textit\def\PY@tc##1{\textcolor[rgb]{0.73,0.13,0.13}{##1}}}
\@namedef{PY@tok@si}{\let\PY@bf=\textbf\def\PY@tc##1{\textcolor[rgb]{0.64,0.35,0.47}{##1}}}
\@namedef{PY@tok@se}{\let\PY@bf=\textbf\def\PY@tc##1{\textcolor[rgb]{0.67,0.36,0.12}{##1}}}
\@namedef{PY@tok@sr}{\def\PY@tc##1{\textcolor[rgb]{0.64,0.35,0.47}{##1}}}
\@namedef{PY@tok@ss}{\def\PY@tc##1{\textcolor[rgb]{0.10,0.09,0.49}{##1}}}
\@namedef{PY@tok@sx}{\def\PY@tc##1{\textcolor[rgb]{0.00,0.50,0.00}{##1}}}
\@namedef{PY@tok@m}{\def\PY@tc##1{\textcolor[rgb]{0.40,0.40,0.40}{##1}}}
\@namedef{PY@tok@gh}{\let\PY@bf=\textbf\def\PY@tc##1{\textcolor[rgb]{0.00,0.00,0.50}{##1}}}
\@namedef{PY@tok@gu}{\let\PY@bf=\textbf\def\PY@tc##1{\textcolor[rgb]{0.50,0.00,0.50}{##1}}}
\@namedef{PY@tok@gd}{\def\PY@tc##1{\textcolor[rgb]{0.63,0.00,0.00}{##1}}}
\@namedef{PY@tok@gi}{\def\PY@tc##1{\textcolor[rgb]{0.00,0.52,0.00}{##1}}}
\@namedef{PY@tok@gr}{\def\PY@tc##1{\textcolor[rgb]{0.89,0.00,0.00}{##1}}}
\@namedef{PY@tok@ge}{\let\PY@it=\textit}
\@namedef{PY@tok@gs}{\let\PY@bf=\textbf}
\@namedef{PY@tok@ges}{\let\PY@bf=\textbf\let\PY@it=\textit}
\@namedef{PY@tok@gp}{\let\PY@bf=\textbf\def\PY@tc##1{\textcolor[rgb]{0.00,0.00,0.50}{##1}}}
\@namedef{PY@tok@go}{\def\PY@tc##1{\textcolor[rgb]{0.44,0.44,0.44}{##1}}}
\@namedef{PY@tok@gt}{\def\PY@tc##1{\textcolor[rgb]{0.00,0.27,0.87}{##1}}}
\@namedef{PY@tok@err}{\def\PY@bc##1{{\setlength{\fboxsep}{\string -\fboxrule}\fcolorbox[rgb]{1.00,0.00,0.00}{1,1,1}{\strut ##1}}}}
\@namedef{PY@tok@kc}{\let\PY@bf=\textbf\def\PY@tc##1{\textcolor[rgb]{0.00,0.50,0.00}{##1}}}
\@namedef{PY@tok@kd}{\let\PY@bf=\textbf\def\PY@tc##1{\textcolor[rgb]{0.00,0.50,0.00}{##1}}}
\@namedef{PY@tok@kn}{\let\PY@bf=\textbf\def\PY@tc##1{\textcolor[rgb]{0.00,0.50,0.00}{##1}}}
\@namedef{PY@tok@kr}{\let\PY@bf=\textbf\def\PY@tc##1{\textcolor[rgb]{0.00,0.50,0.00}{##1}}}
\@namedef{PY@tok@bp}{\def\PY@tc##1{\textcolor[rgb]{0.00,0.50,0.00}{##1}}}
\@namedef{PY@tok@fm}{\def\PY@tc##1{\textcolor[rgb]{0.00,0.00,1.00}{##1}}}
\@namedef{PY@tok@vc}{\def\PY@tc##1{\textcolor[rgb]{0.10,0.09,0.49}{##1}}}
\@namedef{PY@tok@vg}{\def\PY@tc##1{\textcolor[rgb]{0.10,0.09,0.49}{##1}}}
\@namedef{PY@tok@vi}{\def\PY@tc##1{\textcolor[rgb]{0.10,0.09,0.49}{##1}}}
\@namedef{PY@tok@vm}{\def\PY@tc##1{\textcolor[rgb]{0.10,0.09,0.49}{##1}}}
\@namedef{PY@tok@sa}{\def\PY@tc##1{\textcolor[rgb]{0.73,0.13,0.13}{##1}}}
\@namedef{PY@tok@sb}{\def\PY@tc##1{\textcolor[rgb]{0.73,0.13,0.13}{##1}}}
\@namedef{PY@tok@sc}{\def\PY@tc##1{\textcolor[rgb]{0.73,0.13,0.13}{##1}}}
\@namedef{PY@tok@dl}{\def\PY@tc##1{\textcolor[rgb]{0.73,0.13,0.13}{##1}}}
\@namedef{PY@tok@s2}{\def\PY@tc##1{\textcolor[rgb]{0.73,0.13,0.13}{##1}}}
\@namedef{PY@tok@sh}{\def\PY@tc##1{\textcolor[rgb]{0.73,0.13,0.13}{##1}}}
\@namedef{PY@tok@s1}{\def\PY@tc##1{\textcolor[rgb]{0.73,0.13,0.13}{##1}}}
\@namedef{PY@tok@mb}{\def\PY@tc##1{\textcolor[rgb]{0.40,0.40,0.40}{##1}}}
\@namedef{PY@tok@mf}{\def\PY@tc##1{\textcolor[rgb]{0.40,0.40,0.40}{##1}}}
\@namedef{PY@tok@mh}{\def\PY@tc##1{\textcolor[rgb]{0.40,0.40,0.40}{##1}}}
\@namedef{PY@tok@mi}{\def\PY@tc##1{\textcolor[rgb]{0.40,0.40,0.40}{##1}}}
\@namedef{PY@tok@il}{\def\PY@tc##1{\textcolor[rgb]{0.40,0.40,0.40}{##1}}}
\@namedef{PY@tok@mo}{\def\PY@tc##1{\textcolor[rgb]{0.40,0.40,0.40}{##1}}}
\@namedef{PY@tok@ch}{\let\PY@it=\textit\def\PY@tc##1{\textcolor[rgb]{0.24,0.48,0.48}{##1}}}
\@namedef{PY@tok@cm}{\let\PY@it=\textit\def\PY@tc##1{\textcolor[rgb]{0.24,0.48,0.48}{##1}}}
\@namedef{PY@tok@cpf}{\let\PY@it=\textit\def\PY@tc##1{\textcolor[rgb]{0.24,0.48,0.48}{##1}}}
\@namedef{PY@tok@c1}{\let\PY@it=\textit\def\PY@tc##1{\textcolor[rgb]{0.24,0.48,0.48}{##1}}}
\@namedef{PY@tok@cs}{\let\PY@it=\textit\def\PY@tc##1{\textcolor[rgb]{0.24,0.48,0.48}{##1}}}

\def\PYZbs{\char`\\}
\def\PYZus{\char`\_}
\def\PYZob{\char`\{}
\def\PYZcb{\char`\}}
\def\PYZca{\char`\^}
\def\PYZam{\char`\&}
\def\PYZlt{\char`\<}
\def\PYZgt{\char`\>}
\def\PYZsh{\char`\#}
\def\PYZpc{\char`\%}
\def\PYZdl{\char`\$}
\def\PYZhy{\char`\-}
\def\PYZsq{\char`\'}
\def\PYZdq{\char`\"}
\def\PYZti{\char`\~}
% for compatibility with earlier versions
\def\PYZat{@}
\def\PYZlb{[}
\def\PYZrb{]}
\makeatother


    % For linebreaks inside Verbatim environment from package fancyvrb.
    \makeatletter
        \newbox\Wrappedcontinuationbox
        \newbox\Wrappedvisiblespacebox
        \newcommand*\Wrappedvisiblespace {\textcolor{red}{\textvisiblespace}}
        \newcommand*\Wrappedcontinuationsymbol {\textcolor{red}{\llap{\tiny$\m@th\hookrightarrow$}}}
        \newcommand*\Wrappedcontinuationindent {3ex }
        \newcommand*\Wrappedafterbreak {\kern\Wrappedcontinuationindent\copy\Wrappedcontinuationbox}
        % Take advantage of the already applied Pygments mark-up to insert
        % potential linebreaks for TeX processing.
        %        {, <, #, %, $, ' and ": go to next line.
        %        _, }, ^, &, >, - and ~: stay at end of broken line.
        % Use of \textquotesingle for straight quote.
        \newcommand*\Wrappedbreaksatspecials {%
            \def\PYGZus{\discretionary{\char`\_}{\Wrappedafterbreak}{\char`\_}}%
            \def\PYGZob{\discretionary{}{\Wrappedafterbreak\char`\{}{\char`\{}}%
            \def\PYGZcb{\discretionary{\char`\}}{\Wrappedafterbreak}{\char`\}}}%
            \def\PYGZca{\discretionary{\char`\^}{\Wrappedafterbreak}{\char`\^}}%
            \def\PYGZam{\discretionary{\char`\&}{\Wrappedafterbreak}{\char`\&}}%
            \def\PYGZlt{\discretionary{}{\Wrappedafterbreak\char`\<}{\char`\<}}%
            \def\PYGZgt{\discretionary{\char`\>}{\Wrappedafterbreak}{\char`\>}}%
            \def\PYGZsh{\discretionary{}{\Wrappedafterbreak\char`\#}{\char`\#}}%
            \def\PYGZpc{\discretionary{}{\Wrappedafterbreak\char`\%}{\char`\%}}%
            \def\PYGZdl{\discretionary{}{\Wrappedafterbreak\char`\$}{\char`\$}}%
            \def\PYGZhy{\discretionary{\char`\-}{\Wrappedafterbreak}{\char`\-}}%
            \def\PYGZsq{\discretionary{}{\Wrappedafterbreak\textquotesingle}{\textquotesingle}}%
            \def\PYGZdq{\discretionary{}{\Wrappedafterbreak\char`\"}{\char`\"}}%
            \def\PYGZti{\discretionary{\char`\~}{\Wrappedafterbreak}{\char`\~}}%
        }
        % Some characters . , ; ? ! / are not pygmentized.
        % This macro makes them "active" and they will insert potential linebreaks
        \newcommand*\Wrappedbreaksatpunct {%
            \lccode`\~`\.\lowercase{\def~}{\discretionary{\hbox{\char`\.}}{\Wrappedafterbreak}{\hbox{\char`\.}}}%
            \lccode`\~`\,\lowercase{\def~}{\discretionary{\hbox{\char`\,}}{\Wrappedafterbreak}{\hbox{\char`\,}}}%
            \lccode`\~`\;\lowercase{\def~}{\discretionary{\hbox{\char`\;}}{\Wrappedafterbreak}{\hbox{\char`\;}}}%
            \lccode`\~`\:\lowercase{\def~}{\discretionary{\hbox{\char`\:}}{\Wrappedafterbreak}{\hbox{\char`\:}}}%
            \lccode`\~`\?\lowercase{\def~}{\discretionary{\hbox{\char`\?}}{\Wrappedafterbreak}{\hbox{\char`\?}}}%
            \lccode`\~`\!\lowercase{\def~}{\discretionary{\hbox{\char`\!}}{\Wrappedafterbreak}{\hbox{\char`\!}}}%
            \lccode`\~`\/\lowercase{\def~}{\discretionary{\hbox{\char`\/}}{\Wrappedafterbreak}{\hbox{\char`\/}}}%
            \catcode`\.\active
            \catcode`\,\active
            \catcode`\;\active
            \catcode`\:\active
            \catcode`\?\active
            \catcode`\!\active
            \catcode`\/\active
            \lccode`\~`\~
        }
    \makeatother

    \let\OriginalVerbatim=\Verbatim
    \makeatletter
    \renewcommand{\Verbatim}[1][1]{%
        %\parskip\z@skip
        \sbox\Wrappedcontinuationbox {\Wrappedcontinuationsymbol}%
        \sbox\Wrappedvisiblespacebox {\FV@SetupFont\Wrappedvisiblespace}%
        \def\FancyVerbFormatLine ##1{\hsize\linewidth
            \vtop{\raggedright\hyphenpenalty\z@\exhyphenpenalty\z@
                \doublehyphendemerits\z@\finalhyphendemerits\z@
                \strut ##1\strut}%
        }%
        % If the linebreak is at a space, the latter will be displayed as visible
        % space at end of first line, and a continuation symbol starts next line.
        % Stretch/shrink are however usually zero for typewriter font.
        \def\FV@Space {%
            \nobreak\hskip\z@ plus\fontdimen3\font minus\fontdimen4\font
            \discretionary{\copy\Wrappedvisiblespacebox}{\Wrappedafterbreak}
            {\kern\fontdimen2\font}%
        }%

        % Allow breaks at special characters using \PYG... macros.
        \Wrappedbreaksatspecials
        % Breaks at punctuation characters . , ; ? ! and / need catcode=\active
        \OriginalVerbatim[#1,codes*=\Wrappedbreaksatpunct]%
    }
    \makeatother

    % Exact colors from NB
    \definecolor{incolor}{HTML}{303F9F}
    \definecolor{outcolor}{HTML}{D84315}
    \definecolor{cellborder}{HTML}{CFCFCF}
    \definecolor{cellbackground}{HTML}{F7F7F7}

    % prompt
    \makeatletter
    \newcommand{\boxspacing}{\kern\kvtcb@left@rule\kern\kvtcb@boxsep}
    \makeatother
    \newcommand{\prompt}[4]{
        {\ttfamily\llap{{\color{#2}[#3]:\hspace{3pt}#4}}\vspace{-\baselineskip}}
    }
    

    
    % Prevent overflowing lines due to hard-to-break entities
    \sloppy
    % Setup hyperref package
    \hypersetup{
      breaklinks=true,  % so long urls are correctly broken across lines
      colorlinks=true,
      urlcolor=urlcolor,
      linkcolor=linkcolor,
      citecolor=citecolor,
      }
    % Slightly bigger margins than the latex defaults
    
    \geometry{verbose,tmargin=1in,bmargin=1in,lmargin=1in,rmargin=1in}
    
    

\begin{document}
    
    \maketitle
    
    

    

    \section{Variance, Covariance, Eigenvalues and
Eigenvectors}\label{variance-covariance-eigenvalues-and-eigenvectors}

\subsubsection{From statistics to Principal Component Analysis
(PCA)}\label{from-statistics-to-principal-component-analysis-pca}

This notebook develops the ideas step by step using a small dataset and
\textbf{NumPy}, \textbf{pandas} and \textbf{Matplotlib}.

\textbf{Learning objectives}

\begin{enumerate}
\def\labelenumi{\arabic{enumi}.}
\tightlist
\item
  Calculate and interpret sample variance and standard deviation.
\item
  Calculate covariance and correlation, and read a covariance matrix.
\item
  Understand an eigenvector as a direction preserved by a matrix
  transformation, and an eigenvalue as its scaling factor.
\item
  Compute eigenvalues and eigenvectors of a covariance matrix.
\item
  Understand how PCA identifies directions of maximum variance.
\end{enumerate}

Run the cells from top to bottom. All datasets are generated locally, so
no downloads are required.


    \begin{tcolorbox}[breakable, size=fbox, boxrule=1pt, pad at break*=1mm,colback=cellbackground, colframe=cellborder]
\prompt{In}{incolor}{1}{\boxspacing}
\begin{Verbatim}[commandchars=\\\{\}]
\PY{k+kn}{import}\PY{+w}{ }\PY{n+nn}{numpy}\PY{+w}{ }\PY{k}{as}\PY{+w}{ }\PY{n+nn}{np}
\PY{k+kn}{import}\PY{+w}{ }\PY{n+nn}{pandas}\PY{+w}{ }\PY{k}{as}\PY{+w}{ }\PY{n+nn}{pd}
\PY{k+kn}{import}\PY{+w}{ }\PY{n+nn}{matplotlib}\PY{n+nn}{.}\PY{n+nn}{pyplot}\PY{+w}{ }\PY{k}{as}\PY{+w}{ }\PY{n+nn}{plt}
\PY{k+kn}{from}\PY{+w}{ }\PY{n+nn}{matplotlib}\PY{n+nn}{.}\PY{n+nn}{lines}\PY{+w}{ }\PY{k+kn}{import} \PY{n}{Line2D}
\PY{k+kn}{from}\PY{+w}{ }\PY{n+nn}{IPython}\PY{n+nn}{.}\PY{n+nn}{display}\PY{+w}{ }\PY{k+kn}{import} \PY{n}{display}\PY{p}{,} \PY{n}{Math}\PY{p}{,} \PY{n}{Markdown}
\PY{k+kn}{from}\PY{+w}{ }\PY{n+nn}{sympy}\PY{+w}{ }\PY{k+kn}{import} \PY{n}{Matrix}\PY{p}{,} \PY{n}{latex}

\PY{n}{np}\PY{o}{.}\PY{n}{set\PYZus{}printoptions}\PY{p}{(}\PY{n}{precision}\PY{o}{=}\PY{l+m+mi}{4}\PY{p}{,} \PY{n}{suppress}\PY{o}{=}\PY{k+kc}{True}\PY{p}{)}
\PY{n+nb}{print}\PY{p}{(}\PY{l+s+s2}{\PYZdq{}}\PY{l+s+s2}{NumPy version:}\PY{l+s+s2}{\PYZdq{}}\PY{p}{,} \PY{n}{np}\PY{o}{.}\PY{n}{\PYZus{}\PYZus{}version\PYZus{}\PYZus{}}\PY{p}{)}
\end{Verbatim}
\end{tcolorbox}

    \begin{Verbatim}[commandchars=\\\{\}]
NumPy version: 1.26.4
    \end{Verbatim}


    \subsection{1. Variance: how much does one variable spread
out?}\label{variance-how-much-does-one-variable-spread-out}

Given observations \(x_1,\ldots,x_n\) with mean \(\bar{x}\), the
\textbf{sample variance} is

\[s_X^2 = \frac{1}{n-1}\sum_{i=1}^n(x_i-\bar{x})^2.\]

The sample standard deviation is \(s_X=\sqrt{s_X^2}\), returning the
spread to the original units.

For a \textbf{population} (rather than a sample), use \(N\) instead of
\(n-1\). NumPy uses the population denominator by default, so specify
\texttt{ddof=1} for sample variance.


    \begin{tcolorbox}[breakable, size=fbox, boxrule=1pt, pad at break*=1mm,colback=cellbackground, colframe=cellborder]
\prompt{In}{incolor}{2}{\boxspacing}
\begin{Verbatim}[commandchars=\\\{\}]
\PY{n}{scores} \PY{o}{=} \PY{n}{np}\PY{o}{.}\PY{n}{array}\PY{p}{(}\PY{p}{[}\PY{l+m+mi}{60}\PY{p}{,} \PY{l+m+mi}{70}\PY{p}{,} \PY{l+m+mi}{80}\PY{p}{,} \PY{l+m+mi}{90}\PY{p}{,} \PY{l+m+mi}{100}\PY{p}{]}\PY{p}{,} \PY{n}{dtype}\PY{o}{=}\PY{n+nb}{float}\PY{p}{)}

\PY{n}{n} \PY{o}{=} \PY{n+nb}{len}\PY{p}{(}\PY{n}{scores}\PY{p}{)}
\PY{n}{mean} \PY{o}{=} \PY{n}{scores}\PY{o}{.}\PY{n}{mean}\PY{p}{(}\PY{p}{)}
\PY{n}{deviations} \PY{o}{=} \PY{n}{scores} \PY{o}{\PYZhy{}} \PY{n}{mean}
\PY{n}{squared\PYZus{}deviations} \PY{o}{=} \PY{n}{deviations}\PY{o}{*}\PY{o}{*}\PY{l+m+mi}{2}
\PY{n}{manual\PYZus{}variance} \PY{o}{=} \PY{n}{squared\PYZus{}deviations}\PY{o}{.}\PY{n}{sum}\PY{p}{(}\PY{p}{)} \PY{o}{/} \PY{p}{(}\PY{n}{n} \PY{o}{\PYZhy{}} \PY{l+m+mi}{1}\PY{p}{)}
\PY{n}{sample\PYZus{}std} \PY{o}{=} \PY{n}{np}\PY{o}{.}\PY{n}{std}\PY{p}{(}\PY{n}{scores}\PY{p}{,} \PY{n}{ddof}\PY{o}{=}\PY{l+m+mi}{1}\PY{p}{)}

\PY{n}{summary} \PY{o}{=} \PY{n}{pd}\PY{o}{.}\PY{n}{DataFrame}\PY{p}{(}\PY{p}{\PYZob{}}
    \PY{l+s+s2}{\PYZdq{}}\PY{l+s+s2}{Score}\PY{l+s+s2}{\PYZdq{}}\PY{p}{:} \PY{n}{scores}\PY{p}{,}
    \PY{l+s+s2}{\PYZdq{}}\PY{l+s+s2}{Deviation from mean}\PY{l+s+s2}{\PYZdq{}}\PY{p}{:} \PY{n}{deviations}\PY{p}{,}
    \PY{l+s+s2}{\PYZdq{}}\PY{l+s+s2}{Squared deviation}\PY{l+s+s2}{\PYZdq{}}\PY{p}{:} \PY{n}{squared\PYZus{}deviations}
\PY{p}{\PYZcb{}}\PY{p}{)}
\PY{n}{display}\PY{p}{(}\PY{n}{summary}\PY{p}{)}


\PY{l+s+sd}{\PYZsq{}\PYZsq{}\PYZsq{}print(f\PYZdq{}Mean = \PYZob{}mean:.1f\PYZcb{}\PYZdq{})}
\PY{l+s+sd}{print(f\PYZdq{}Sample variance = \PYZob{}manual\PYZus{}variance:.1f\PYZcb{}\PYZdq{})}
\PY{l+s+sd}{print(f\PYZdq{}NumPy sample variance = \PYZob{}np.var(scores, ddof=1):.1f\PYZcb{}\PYZdq{})}
\PY{l+s+sd}{print(f\PYZdq{}Sample standard deviation = \PYZob{}np.std(scores, ddof=1):.3f\PYZcb{}\PYZdq{})\PYZsq{}\PYZsq{}\PYZsq{}}

\PY{c+c1}{\PYZsh{} Format the observations for LaTeX}
\PY{n}{values} \PY{o}{=} \PY{l+s+s2}{\PYZdq{}}\PY{l+s+s2}{+}\PY{l+s+s2}{\PYZdq{}}\PY{o}{.}\PY{n}{join}\PY{p}{(}\PY{l+s+sa}{f}\PY{l+s+s2}{\PYZdq{}}\PY{l+s+si}{\PYZob{}}\PY{n}{x}\PY{l+s+si}{:}\PY{l+s+s2}{g}\PY{l+s+si}{\PYZcb{}}\PY{l+s+s2}{\PYZdq{}} \PY{k}{for} \PY{n}{x} \PY{o+ow}{in} \PY{n}{scores}\PY{p}{)}
\PY{n}{deviations} \PY{o}{=} \PY{l+s+s2}{\PYZdq{}}\PY{l+s+s2}{+}\PY{l+s+s2}{\PYZdq{}}\PY{o}{.}\PY{n}{join}\PY{p}{(}
    \PY{l+s+sa}{f}\PY{l+s+s2}{\PYZdq{}}\PY{l+s+s2}{(}\PY{l+s+si}{\PYZob{}}\PY{n}{x}\PY{l+s+si}{:}\PY{l+s+s2}{g}\PY{l+s+si}{\PYZcb{}}\PY{l+s+s2}{\PYZhy{}}\PY{l+s+si}{\PYZob{}}\PY{n}{mean}\PY{l+s+si}{:}\PY{l+s+s2}{g}\PY{l+s+si}{\PYZcb{}}\PY{l+s+s2}{)\PYZca{}2}\PY{l+s+s2}{\PYZdq{}} \PY{k}{for} \PY{n}{x} \PY{o+ow}{in} \PY{n}{scores}
\PY{p}{)}
\PY{n}{squared\PYZus{}deviations} \PY{o}{=} \PY{l+s+s2}{\PYZdq{}}\PY{l+s+s2}{+}\PY{l+s+s2}{\PYZdq{}}\PY{o}{.}\PY{n}{join}\PY{p}{(}
    \PY{l+s+sa}{f}\PY{l+s+s2}{\PYZdq{}}\PY{l+s+si}{\PYZob{}}\PY{p}{(}\PY{n}{x}\PY{+w}{ }\PY{o}{\PYZhy{}}\PY{+w}{ }\PY{n}{mean}\PY{p}{)}\PY{o}{*}\PY{o}{*}\PY{l+m+mi}{2}\PY{l+s+si}{:}\PY{l+s+s2}{g}\PY{l+s+si}{\PYZcb{}}\PY{l+s+s2}{\PYZdq{}} \PY{k}{for} \PY{n}{x} \PY{o+ow}{in} \PY{n}{scores}
\PY{p}{)}

\PY{c+c1}{\PYZsh{} 1. Mean}
\PY{n}{display}\PY{p}{(}\PY{n}{Math}\PY{p}{(}
    \PY{l+s+sa}{rf}\PY{l+s+s2}{\PYZdq{}}\PY{l+s+s2}{\PYZbs{}}\PY{l+s+s2}{bar}\PY{l+s+se}{\PYZob{}\PYZob{}}\PY{l+s+s2}{x}\PY{l+s+se}{\PYZcb{}\PYZcb{}}\PY{l+s+s2}{ = }\PY{l+s+s2}{\PYZbs{}}\PY{l+s+s2}{frac}\PY{l+s+se}{\PYZob{}\PYZob{}}\PY{l+s+si}{\PYZob{}}\PY{n}{values}\PY{l+s+si}{\PYZcb{}}\PY{l+s+se}{\PYZcb{}\PYZcb{}}\PY{l+s+se}{\PYZob{}\PYZob{}}\PY{l+s+si}{\PYZob{}}\PY{n}{n}\PY{l+s+si}{\PYZcb{}}\PY{l+s+se}{\PYZcb{}\PYZcb{}}\PY{l+s+s2}{ = }\PY{l+s+si}{\PYZob{}}\PY{n}{mean}\PY{l+s+si}{:}\PY{l+s+s2}{.1f}\PY{l+s+si}{\PYZcb{}}\PY{l+s+s2}{\PYZdq{}}
\PY{p}{)}\PY{p}{)}

\PY{c+c1}{\PYZsh{} 2. Sample variance}
\PY{n}{display}\PY{p}{(}\PY{n}{Math}\PY{p}{(}
    \PY{l+s+sa}{rf}\PY{l+s+s2}{\PYZdq{}}\PY{l+s+s2}{s\PYZca{}2 = }\PY{l+s+s2}{\PYZbs{}}\PY{l+s+s2}{frac}\PY{l+s+se}{\PYZob{}\PYZob{}}\PY{l+s+si}{\PYZob{}}\PY{n}{deviations}\PY{l+s+si}{\PYZcb{}}\PY{l+s+se}{\PYZcb{}\PYZcb{}}\PY{l+s+se}{\PYZob{}\PYZob{}}\PY{l+s+si}{\PYZob{}}\PY{n}{n}\PY{l+s+si}{\PYZcb{}}\PY{l+s+s2}{\PYZhy{}1}\PY{l+s+se}{\PYZcb{}\PYZcb{}}\PY{l+s+s2}{\PYZdq{}}
    \PY{l+s+sa}{rf}\PY{l+s+s2}{\PYZdq{}}\PY{l+s+s2}{ = }\PY{l+s+s2}{\PYZbs{}}\PY{l+s+s2}{frac}\PY{l+s+se}{\PYZob{}\PYZob{}}\PY{l+s+si}{\PYZob{}}\PY{n}{squared\PYZus{}deviations}\PY{l+s+si}{\PYZcb{}}\PY{l+s+se}{\PYZcb{}\PYZcb{}}\PY{l+s+se}{\PYZob{}\PYZob{}}\PY{l+s+si}{\PYZob{}}\PY{n}{n}\PY{o}{\PYZhy{}}\PY{l+m+mi}{1}\PY{l+s+si}{\PYZcb{}}\PY{l+s+se}{\PYZcb{}\PYZcb{}}\PY{l+s+s2}{\PYZdq{}}
    \PY{l+s+sa}{rf}\PY{l+s+s2}{\PYZdq{}}\PY{l+s+s2}{ = }\PY{l+s+si}{\PYZob{}}\PY{n}{manual\PYZus{}variance}\PY{l+s+si}{:}\PY{l+s+s2}{.1f}\PY{l+s+si}{\PYZcb{}}\PY{l+s+s2}{\PYZdq{}}
\PY{p}{)}\PY{p}{)}

\PY{c+c1}{\PYZsh{} 3. NumPy sample variance}
\PY{n}{display}\PY{p}{(}\PY{n}{Math}\PY{p}{(}
    \PY{l+s+sa}{rf}\PY{l+s+s2}{\PYZdq{}}\PY{l+s+s2}{\PYZbs{}}\PY{l+s+s2}{operatorname}\PY{l+s+se}{\PYZob{}\PYZob{}}\PY{l+s+s2}{Var}\PY{l+s+se}{\PYZcb{}\PYZcb{}}\PY{l+s+s2}{\PYZus{}}\PY{l+s+se}{\PYZob{}\PYZob{}}\PY{l+s+s2}{\PYZbs{}}\PY{l+s+s2}{mathrm}\PY{l+s+se}{\PYZob{}\PYZob{}}\PY{l+s+s2}{NumPy}\PY{l+s+se}{\PYZcb{}\PYZcb{}}\PY{l+s+se}{\PYZcb{}\PYZcb{}}\PY{l+s+s2}{(X)}\PY{l+s+s2}{\PYZdq{}}
    \PY{l+s+sa}{rf}\PY{l+s+s2}{\PYZdq{}}\PY{l+s+s2}{ = }\PY{l+s+si}{\PYZob{}}\PY{n}{np}\PY{o}{.}\PY{n}{var}\PY{p}{(}\PY{n}{scores}\PY{p}{,}\PY{+w}{ }\PY{n}{ddof}\PY{o}{=}\PY{l+m+mi}{1}\PY{p}{)}\PY{l+s+si}{:}\PY{l+s+s2}{.1f}\PY{l+s+si}{\PYZcb{}}\PY{l+s+s2}{\PYZdq{}}
\PY{p}{)}\PY{p}{)}

\PY{c+c1}{\PYZsh{} 4. Sample standard deviation}
\PY{n}{display}\PY{p}{(}\PY{n}{Math}\PY{p}{(}
    \PY{l+s+sa}{rf}\PY{l+s+s2}{\PYZdq{}}\PY{l+s+s2}{s = }\PY{l+s+s2}{\PYZbs{}}\PY{l+s+s2}{sqrt}\PY{l+s+se}{\PYZob{}\PYZob{}}\PY{l+s+s2}{s\PYZca{}2}\PY{l+s+se}{\PYZcb{}\PYZcb{}}\PY{l+s+s2}{\PYZdq{}}
    \PY{l+s+sa}{rf}\PY{l+s+s2}{\PYZdq{}}\PY{l+s+s2}{ = }\PY{l+s+s2}{\PYZbs{}}\PY{l+s+s2}{sqrt}\PY{l+s+se}{\PYZob{}\PYZob{}}\PY{l+s+si}{\PYZob{}}\PY{n}{manual\PYZus{}variance}\PY{l+s+si}{:}\PY{l+s+s2}{.1f}\PY{l+s+se}{\PYZcb{}\PYZcb{}}\PY{l+s+si}{\PYZcb{}}\PY{l+s+s2}{\PYZdq{}}
    \PY{l+s+sa}{rf}\PY{l+s+s2}{\PYZdq{}}\PY{l+s+s2}{ = }\PY{l+s+si}{\PYZob{}}\PY{n}{sample\PYZus{}std}\PY{l+s+si}{:}\PY{l+s+s2}{.3f}\PY{l+s+si}{\PYZcb{}}\PY{l+s+s2}{\PYZdq{}}
\PY{p}{)}\PY{p}{)}
\end{Verbatim}
\end{tcolorbox}

    
    \begin{Verbatim}[commandchars=\\\{\}]
   Score  Deviation from mean  Squared deviation
0   60.0                -20.0              400.0
1   70.0                -10.0              100.0
2   80.0                  0.0                0.0
3   90.0                 10.0              100.0
4  100.0                 20.0              400.0
    \end{Verbatim}

    
    $\displaystyle \bar{x} = \frac{60+70+80+90+100}{5} = 80.0$

    
    $\displaystyle s^2 = \frac{(60-80)^2+(70-80)^2+(80-80)^2+(90-80)^2+(100-80)^2}{5-1} = \frac{400+100+0+100+400}{4} = 250.0$

    
    $\displaystyle \operatorname{Var}_{\mathrm{NumPy}}(X) = 250.0$

    
    $\displaystyle s = \sqrt{s^2} = \sqrt{250.0} = 15.811$

    

    \subsection{2. Sample Standard
Deviation}\label{sample-standard-deviation}

The \textbf{sample standard deviation} measures how spread out the
observations in a sample are around their mean.

It is the \textbf{square root of the sample variance}, expressed in the
same units as the original data.

\subsubsection{2.1. Mathematical formula}\label{mathematical-formula}

\[
\boxed{s = \sqrt{\frac{\sum_{i=1}^{n}(x_i-\bar{x})^2}{n-1}}}
\]

Where:

\begin{itemize}
\tightlist
\item
  \(s\): Sample standard deviation.
\item
  \(x_i\): Individual observation.
\item
  \(\bar{x}\): Sample mean.
\item
  \(n\): Number of observations.
\item
  \(n-1\): Degrees of freedom (Bessel's correction).
\end{itemize}

\subsubsection{2.2. Worked example}\label{worked-example}

Consider five exam scores:

\[
X = [60,70,80,90,100]
\]

\textbf{Step 1: Calculate the sample mean}

\[
\bar{x}=\frac{60+70+80+90+100}{5}=80
\]

\textbf{Step 2: Calculate the squared deviations}

{\def\LTcaptype{none} % do not increment counter
\begin{longtable}[]{@{}
  >{\raggedright\arraybackslash}p{(\linewidth - 4\tabcolsep) * \real{0.3333}}
  >{\raggedright\arraybackslash}p{(\linewidth - 4\tabcolsep) * \real{0.3333}}
  >{\raggedright\arraybackslash}p{(\linewidth - 4\tabcolsep) * \real{0.3333}}@{}}
\toprule\noalign{}
\begin{minipage}[b]{\linewidth}\raggedright
Score
\end{minipage} & \begin{minipage}[b]{\linewidth}\raggedright
Deviation \((x_i-\bar{x})\)
\end{minipage} & \begin{minipage}[b]{\linewidth}\raggedright
Squared deviation \((x_i-\bar{x})^2\)
\end{minipage} \\
\midrule\noalign{}
\endhead
\bottomrule\noalign{}
\endlastfoot
60 & -20 & 400 \\
70 & -10 & 100 \\
80 & 0 & 0 \\
90 & 10 & 100 \\
100 & 20 & 400 \\
\textbf{Total} & \textbf{0} & \textbf{1000} \\
\end{longtable}
}

\textbf{Step 3: Calculate the sample variance}

\[
s^2=\frac{1000}{5-1}=\boxed{250}
\]

\textbf{Step 4: Calculate the sample standard deviation}

\[
s=\sqrt{250}=\boxed{15.81}
\]

The sample standard deviation is approximately \textbf{15.81 points}.

\subsubsection{2.3. Sample vs Population Standard
Deviation}\label{sample-vs-population-standard-deviation}

\textbf{Sample standard deviation:}

\[
s=\sqrt{\frac{\sum_{i=1}^{n}(x_i-\bar{x})^2}{n-1}}
\]

\textbf{Population standard deviation:}

\[
\sigma=\sqrt{\frac{\sum_{i=1}^{N}(x_i-\mu)^2}{N}}
\]

{\def\LTcaptype{none} % do not increment counter
\begin{longtable}[]{@{}
  >{\raggedright\arraybackslash}p{(\linewidth - 4\tabcolsep) * \real{0.3333}}
  >{\raggedright\arraybackslash}p{(\linewidth - 4\tabcolsep) * \real{0.3333}}
  >{\raggedright\arraybackslash}p{(\linewidth - 4\tabcolsep) * \real{0.3333}}@{}}
\toprule\noalign{}
\begin{minipage}[b]{\linewidth}\raggedright
Property
\end{minipage} & \begin{minipage}[b]{\linewidth}\raggedright
Sample
\end{minipage} & \begin{minipage}[b]{\linewidth}\raggedright
Population
\end{minipage} \\
\midrule\noalign{}
\endhead
\bottomrule\noalign{}
\endlastfoot
Symbol & \(s\) & \(\sigma\) \\
Denominator & \(n-1\) & \(N\) \\
Purpose & Estimate population variability & Describe the entire
population \\
Example result & 15.81 & 14.14 \\
\end{longtable}
}

\subsubsection{\texorpdfstring{2.4. Why divide by
\(n-1\)?}{2.4. Why divide by n-1?}}\label{why-divide-by-n-1}

When calculating the sample mean from the same observations, the
deviations must sum to zero:

\[
\sum_{i=1}^{n}(x_i-\bar{x})=0
\]

Therefore, only \(n-1\) deviations can vary independently.

Dividing by \(n-1\) instead of \(n\) is known as \textbf{Bessel's
correction}.

It makes sample variance an unbiased estimator of population variance
under independent, identically distributed sampling with finite
variance.

However, the sample standard deviation itself is not generally an
unbiased estimator of the population standard deviation.

\subsubsection{2.5. Python implementation}\label{python-implementation}

\begin{Shaded}
\begin{Highlighting}[]
\ImportTok{import}\NormalTok{ numpy }\ImportTok{as}\NormalTok{ np}

\NormalTok{scores }\OperatorTok{=}\NormalTok{ np.array([}\DecValTok{60}\NormalTok{, }\DecValTok{70}\NormalTok{, }\DecValTok{80}\NormalTok{, }\DecValTok{90}\NormalTok{, }\DecValTok{100}\NormalTok{])}

\CommentTok{\# Sample variance}
\NormalTok{sample\_variance }\OperatorTok{=}\NormalTok{ np.var(scores, ddof}\OperatorTok{=}\DecValTok{1}\NormalTok{)}

\CommentTok{\# Sample standard deviation}
\NormalTok{sample\_std }\OperatorTok{=}\NormalTok{ np.std(scores, ddof}\OperatorTok{=}\DecValTok{1}\NormalTok{)}

\CommentTok{\# Population standard deviation}
\NormalTok{population\_std }\OperatorTok{=}\NormalTok{ np.std(scores, ddof}\OperatorTok{=}\DecValTok{0}\NormalTok{)}

\BuiltInTok{print}\NormalTok{(}\SpecialStringTok{f"Sample variance: }\SpecialCharTok{\{}\NormalTok{sample\_variance}\SpecialCharTok{:.2f\}}\SpecialStringTok{"}\NormalTok{)}
\BuiltInTok{print}\NormalTok{(}\SpecialStringTok{f"Sample standard deviation: }\SpecialCharTok{\{}\NormalTok{sample\_std}\SpecialCharTok{:.2f\}}\SpecialStringTok{"}\NormalTok{)}
\BuiltInTok{print}\NormalTok{(}\SpecialStringTok{f"Population standard deviation: }\SpecialCharTok{\{}\NormalTok{population\_std}\SpecialCharTok{:.2f\}}\SpecialStringTok{"}\NormalTok{)}
\end{Highlighting}
\end{Shaded}

Expected output:

\begin{Shaded}
\begin{Highlighting}[]
\NormalTok{Sample variance: 250.00}
\NormalTok{Sample standard deviation: 15.81}
\NormalTok{Population standard deviation: 14.14}
\end{Highlighting}
\end{Shaded}

\textbf{Note:} In NumPy, \texttt{ddof=1} applies Bessel's correction by
dividing by \(n-1\). The default \texttt{ddof=0} divides by \(n\).

\subsubsection{2.6. Variance vs Standard
Deviation}\label{variance-vs-standard-deviation}

{\def\LTcaptype{none} % do not increment counter
\begin{longtable}[]{@{}lll@{}}
\toprule\noalign{}
Property & Variance & Standard deviation \\
\midrule\noalign{}
\endhead
\bottomrule\noalign{}
\endlastfoot
Measures & Spread around the mean & Spread around the mean \\
Formula & \(s^2\) & \(\sqrt{s^2}\) \\
Units & Squared original units & Original units \\
Example & 250 points² & 15.81 points \\
\end{longtable}
}

\textbf{Key takeaway:}

\[
\boxed{\text{Standard Deviation}=\sqrt{\text{Variance}}}
\]

Standard deviation is generally easier to interpret because it uses the
same units as the original observations.


    \begin{tcolorbox}[breakable, size=fbox, boxrule=1pt, pad at break*=1mm,colback=cellbackground, colframe=cellborder]
\prompt{In}{incolor}{3}{\boxspacing}
\begin{Verbatim}[commandchars=\\\{\}]
\PY{n}{fig}\PY{p}{,} \PY{n}{ax} \PY{o}{=} \PY{n}{plt}\PY{o}{.}\PY{n}{subplots}\PY{p}{(}\PY{n}{figsize}\PY{o}{=}\PY{p}{(}\PY{l+m+mi}{8}\PY{p}{,} \PY{l+m+mf}{3.3}\PY{p}{)}\PY{p}{)}
\PY{n}{ax}\PY{o}{.}\PY{n}{scatter}\PY{p}{(}\PY{n}{scores}\PY{p}{,} \PY{n}{np}\PY{o}{.}\PY{n}{zeros\PYZus{}like}\PY{p}{(}\PY{n}{scores}\PY{p}{)}\PY{p}{,} \PY{n}{s}\PY{o}{=}\PY{l+m+mi}{65}\PY{p}{,} \PY{n}{zorder}\PY{o}{=}\PY{l+m+mi}{3}\PY{p}{,} \PY{n}{label}\PY{o}{=}\PY{l+s+s2}{\PYZdq{}}\PY{l+s+s2}{Observations}\PY{l+s+s2}{\PYZdq{}}\PY{p}{)}
\PY{n}{ax}\PY{o}{.}\PY{n}{axvline}\PY{p}{(}\PY{n}{mean}\PY{p}{,} \PY{n}{linestyle}\PY{o}{=}\PY{l+s+s2}{\PYZdq{}}\PY{l+s+s2}{\PYZhy{}\PYZhy{}}\PY{l+s+s2}{\PYZdq{}}\PY{p}{,} \PY{n}{label}\PY{o}{=}\PY{l+s+sa}{f}\PY{l+s+s2}{\PYZdq{}}\PY{l+s+s2}{Mean = }\PY{l+s+si}{\PYZob{}}\PY{n}{mean}\PY{l+s+si}{:}\PY{l+s+s2}{.0f}\PY{l+s+si}{\PYZcb{}}\PY{l+s+s2}{\PYZdq{}}\PY{p}{)}
\PY{k}{for} \PY{n}{x} \PY{o+ow}{in} \PY{n}{scores}\PY{p}{:}
    \PY{n}{ax}\PY{o}{.}\PY{n}{plot}\PY{p}{(}\PY{p}{[}\PY{n}{mean}\PY{p}{,} \PY{n}{x}\PY{p}{]}\PY{p}{,} \PY{p}{[}\PY{l+m+mi}{0}\PY{p}{,} \PY{l+m+mi}{0}\PY{p}{]}\PY{p}{,} \PY{n}{alpha}\PY{o}{=}\PY{l+m+mf}{0.35}\PY{p}{,} \PY{n}{linewidth}\PY{o}{=}\PY{l+m+mi}{3}\PY{p}{)}
\PY{n}{ax}\PY{o}{.}\PY{n}{set}\PY{p}{(}\PY{n}{xlabel}\PY{o}{=}\PY{l+s+s2}{\PYZdq{}}\PY{l+s+s2}{Exam score}\PY{l+s+s2}{\PYZdq{}}\PY{p}{,} \PY{n}{yticks}\PY{o}{=}\PY{p}{[}\PY{p}{]}\PY{p}{,} \PY{n}{title}\PY{o}{=}\PY{l+s+s2}{\PYZdq{}}\PY{l+s+s2}{Variance measures spread around the mean}\PY{l+s+s2}{\PYZdq{}}\PY{p}{)}
\PY{n}{ax}\PY{o}{.}\PY{n}{legend}\PY{p}{(}\PY{p}{)}
\PY{n}{plt}\PY{o}{.}\PY{n}{show}\PY{p}{(}\PY{p}{)}
\end{Verbatim}
\end{tcolorbox}

    \begin{center}
    \adjustimage{max size={0.9\linewidth}{0.9\paperheight}}{week4_pdf_build_pca_files/week4_pdf_build_pca_5_0.png}
    \end{center}
    { \hspace*{\fill} \\}
    

    \subsection{3. Covariance: how do two variables move
together?}\label{covariance-how-do-two-variables-move-together}

The \textbf{sample covariance} of variables \(X\) and \(Y\) is

\[\operatorname{Cov}(X,Y)=\frac{1}{n-1}\sum_{i=1}^n(x_i-\bar{x})(y_i-\bar{y}).\]

\begin{itemize}
\tightlist
\item
  \textbf{Positive:} as one variable increases, the other tends to
  increase.
\item
  \textbf{Negative:} as one increases, the other tends to decrease.
\item
  \textbf{Near zero:} little \emph{linear} co-movement; nonlinear
  relationships can still exist.
\end{itemize}

Notice: \(\operatorname{Var}(X)=\operatorname{Cov}(X,X)\).


    \begin{tcolorbox}[breakable, size=fbox, boxrule=1pt, pad at break*=1mm,colback=cellbackground, colframe=cellborder]
\prompt{In}{incolor}{4}{\boxspacing}
\begin{Verbatim}[commandchars=\\\{\}]
\PY{n}{hours} \PY{o}{=} \PY{n}{np}\PY{o}{.}\PY{n}{array}\PY{p}{(}\PY{p}{[}\PY{l+m+mi}{1}\PY{p}{,} \PY{l+m+mi}{2}\PY{p}{,} \PY{l+m+mi}{3}\PY{p}{,} \PY{l+m+mi}{4}\PY{p}{,} \PY{l+m+mi}{5}\PY{p}{]}\PY{p}{,} \PY{n}{dtype}\PY{o}{=}\PY{n+nb}{float}\PY{p}{)}
\PY{n}{exam\PYZus{}scores} \PY{o}{=} \PY{n}{np}\PY{o}{.}\PY{n}{array}\PY{p}{(}\PY{p}{[}\PY{l+m+mi}{50}\PY{p}{,} \PY{l+m+mi}{60}\PY{p}{,} \PY{l+m+mi}{70}\PY{p}{,} \PY{l+m+mi}{80}\PY{p}{,} \PY{l+m+mi}{90}\PY{p}{]}\PY{p}{,} \PY{n}{dtype}\PY{o}{=}\PY{n+nb}{float}\PY{p}{)}

\PY{n}{df} \PY{o}{=} \PY{n}{pd}\PY{o}{.}\PY{n}{DataFrame}\PY{p}{(}\PY{p}{\PYZob{}}
    \PY{l+s+s2}{\PYZdq{}}\PY{l+s+s2}{Study Hours}\PY{l+s+s2}{\PYZdq{}}\PY{p}{:} \PY{n}{hours}\PY{p}{,}
    \PY{l+s+s2}{\PYZdq{}}\PY{l+s+s2}{Exam Scores}\PY{l+s+s2}{\PYZdq{}}\PY{p}{:} \PY{n}{exam\PYZus{}scores}
\PY{p}{\PYZcb{}}\PY{p}{)}

\PY{n}{display}\PY{p}{(}\PY{n}{df}\PY{p}{)}

\PY{n}{df}\PY{o}{.}\PY{n}{describe}\PY{p}{(}\PY{p}{)}
\end{Verbatim}
\end{tcolorbox}

    
    \begin{Verbatim}[commandchars=\\\{\}]
   Study Hours  Exam Scores
0          1.0         50.0
1          2.0         60.0
2          3.0         70.0
3          4.0         80.0
4          5.0         90.0
    \end{Verbatim}

    
            \begin{tcolorbox}[breakable, size=fbox, boxrule=.5pt, pad at break*=1mm, opacityfill=0]
\prompt{Out}{outcolor}{4}{\boxspacing}
\begin{Verbatim}[commandchars=\\\{\}]
       Study Hours  Exam Scores
count     5.000000     5.000000
mean      3.000000    70.000000
std       1.581139    15.811388
min       1.000000    50.000000
25\%       2.000000    60.000000
50\%       3.000000    70.000000
75\%       4.000000    80.000000
max       5.000000    90.000000
\end{Verbatim}
\end{tcolorbox}
        
    \begin{tcolorbox}[breakable, size=fbox, boxrule=1pt, pad at break*=1mm,colback=cellbackground, colframe=cellborder]
\prompt{In}{incolor}{5}{\boxspacing}
\begin{Verbatim}[commandchars=\\\{\}]
\PY{n}{products} \PY{o}{=} \PY{p}{(}\PY{n}{hours} \PY{o}{\PYZhy{}} \PY{n}{hours}\PY{o}{.}\PY{n}{mean}\PY{p}{(}\PY{p}{)}\PY{p}{)} \PY{o}{*} \PY{p}{(}\PY{n}{exam\PYZus{}scores} \PY{o}{\PYZhy{}} \PY{n}{exam\PYZus{}scores}\PY{o}{.}\PY{n}{mean}\PY{p}{(}\PY{p}{)}\PY{p}{)}
\PY{n}{manual\PYZus{}covariance} \PY{o}{=} \PY{n}{products}\PY{o}{.}\PY{n}{sum}\PY{p}{(}\PY{p}{)} \PY{o}{/} \PY{p}{(}\PY{n+nb}{len}\PY{p}{(}\PY{n}{hours}\PY{p}{)}\PY{o}{\PYZhy{}}\PY{l+m+mi}{1}\PY{p}{)}
\PY{n}{calculation} \PY{o}{=} \PY{n}{pd}\PY{o}{.}\PY{n}{DataFrame}\PY{p}{(}\PY{p}{\PYZob{}}
    \PY{l+s+s2}{\PYZdq{}}\PY{l+s+s2}{Hours}\PY{l+s+s2}{\PYZdq{}}\PY{p}{:} \PY{n}{hours}\PY{p}{,}
    \PY{l+s+s2}{\PYZdq{}}\PY{l+s+s2}{Score}\PY{l+s+s2}{\PYZdq{}}\PY{p}{:} \PY{n}{exam\PYZus{}scores}\PY{p}{,}
    \PY{l+s+s2}{\PYZdq{}}\PY{l+s+s2}{Hours deviation}\PY{l+s+s2}{\PYZdq{}}\PY{p}{:} \PY{n}{hours}\PY{o}{\PYZhy{}}\PY{n}{hours}\PY{o}{.}\PY{n}{mean}\PY{p}{(}\PY{p}{)}\PY{p}{,}
    \PY{l+s+s2}{\PYZdq{}}\PY{l+s+s2}{Score deviation}\PY{l+s+s2}{\PYZdq{}}\PY{p}{:} \PY{n}{exam\PYZus{}scores}\PY{o}{\PYZhy{}}\PY{n}{exam\PYZus{}scores}\PY{o}{.}\PY{n}{mean}\PY{p}{(}\PY{p}{)}\PY{p}{,}
    \PY{l+s+s2}{\PYZdq{}}\PY{l+s+s2}{Product}\PY{l+s+s2}{\PYZdq{}}\PY{p}{:} \PY{n}{products}
\PY{p}{\PYZcb{}}\PY{p}{)}
\PY{n}{display}\PY{p}{(}\PY{n}{calculation}\PY{p}{)}

\PY{l+s+sd}{\PYZsq{}\PYZsq{}\PYZsq{}print(\PYZdq{}Sample variance of hours:\PYZdq{}, np.var(hours, ddof=1))}
\PY{l+s+sd}{print(\PYZdq{}Sample variance of scores:\PYZdq{}, np.var(exam\PYZus{}scores, ddof=1))}
\PY{l+s+sd}{print(\PYZdq{}Sample covariance:\PYZdq{}, manual\PYZus{}covariance)\PYZsq{}\PYZsq{}\PYZsq{}}




\PY{c+c1}{\PYZsh{} Helper function to format values for LaTeX}
\PY{k}{def}\PY{+w}{ }\PY{n+nf}{latex\PYZus{}sum}\PY{p}{(}\PY{n}{values}\PY{p}{)}\PY{p}{:}
    \PY{k}{return} \PY{l+s+s2}{\PYZdq{}}\PY{l+s+s2}{ + }\PY{l+s+s2}{\PYZdq{}}\PY{o}{.}\PY{n}{join}\PY{p}{(}\PY{l+s+sa}{f}\PY{l+s+s2}{\PYZdq{}}\PY{l+s+si}{\PYZob{}}\PY{n}{v}\PY{l+s+si}{:}\PY{l+s+s2}{g}\PY{l+s+si}{\PYZcb{}}\PY{l+s+s2}{\PYZdq{}} \PY{k}{for} \PY{n}{v} \PY{o+ow}{in} \PY{n}{values}\PY{p}{)}

\PY{n}{n} \PY{o}{=} \PY{n+nb}{len}\PY{p}{(}\PY{n}{hours}\PY{p}{)}

\PY{n}{hours\PYZus{}mean} \PY{o}{=} \PY{n}{hours}\PY{o}{.}\PY{n}{mean}\PY{p}{(}\PY{p}{)}
\PY{n}{scores\PYZus{}mean} \PY{o}{=} \PY{n}{exam\PYZus{}scores}\PY{o}{.}\PY{n}{mean}\PY{p}{(}\PY{p}{)}

\PY{n}{hours\PYZus{}deviations} \PY{o}{=} \PY{n}{hours} \PY{o}{\PYZhy{}} \PY{n}{hours\PYZus{}mean}
\PY{n}{scores\PYZus{}deviations} \PY{o}{=} \PY{n}{exam\PYZus{}scores} \PY{o}{\PYZhy{}} \PY{n}{scores\PYZus{}mean}

\PY{n}{hours\PYZus{}squared} \PY{o}{=} \PY{n}{hours\PYZus{}deviations}\PY{o}{*}\PY{o}{*}\PY{l+m+mi}{2}
\PY{n}{scores\PYZus{}squared} \PY{o}{=} \PY{n}{scores\PYZus{}deviations}\PY{o}{*}\PY{o}{*}\PY{l+m+mi}{2}

\PY{n}{n} \PY{o}{=} \PY{n+nb}{len}\PY{p}{(}\PY{n}{hours}\PY{p}{)}

\PY{c+c1}{\PYZsh{} 1. Mean of study hours}
\PY{n}{hours\PYZus{}mean} \PY{o}{=} \PY{n}{hours}\PY{o}{.}\PY{n}{mean}\PY{p}{(}\PY{p}{)}
\PY{n}{hours\PYZus{}values} \PY{o}{=} \PY{l+s+s2}{\PYZdq{}}\PY{l+s+s2}{ + }\PY{l+s+s2}{\PYZdq{}}\PY{o}{.}\PY{n}{join}\PY{p}{(}\PY{l+s+sa}{f}\PY{l+s+s2}{\PYZdq{}}\PY{l+s+si}{\PYZob{}}\PY{n}{x}\PY{l+s+si}{:}\PY{l+s+s2}{g}\PY{l+s+si}{\PYZcb{}}\PY{l+s+s2}{\PYZdq{}} \PY{k}{for} \PY{n}{x} \PY{o+ow}{in} \PY{n}{hours}\PY{p}{)}

\PY{n}{display}\PY{p}{(}\PY{n}{Math}\PY{p}{(}
    \PY{l+s+sa}{rf}\PY{l+s+s2}{\PYZdq{}}\PY{l+s+s2}{\PYZbs{}}\PY{l+s+s2}{bar}\PY{l+s+se}{\PYZob{}\PYZob{}}\PY{l+s+s2}{x}\PY{l+s+se}{\PYZcb{}\PYZcb{}}\PY{l+s+s2}{\PYZus{}}\PY{l+s+se}{\PYZob{}\PYZob{}}\PY{l+s+s2}{\PYZbs{}}\PY{l+s+s2}{mathrm}\PY{l+s+se}{\PYZob{}\PYZob{}}\PY{l+s+s2}{hours}\PY{l+s+se}{\PYZcb{}\PYZcb{}}\PY{l+s+se}{\PYZcb{}\PYZcb{}}\PY{l+s+s2}{\PYZdq{}}
    \PY{l+s+sa}{rf}\PY{l+s+s2}{\PYZdq{}}\PY{l+s+s2}{ = }\PY{l+s+s2}{\PYZbs{}}\PY{l+s+s2}{frac}\PY{l+s+se}{\PYZob{}\PYZob{}}\PY{l+s+si}{\PYZob{}}\PY{n}{hours\PYZus{}values}\PY{l+s+si}{\PYZcb{}}\PY{l+s+se}{\PYZcb{}\PYZcb{}}\PY{l+s+se}{\PYZob{}\PYZob{}}\PY{l+s+si}{\PYZob{}}\PY{n}{n}\PY{l+s+si}{\PYZcb{}}\PY{l+s+se}{\PYZcb{}\PYZcb{}}\PY{l+s+s2}{\PYZdq{}}
    \PY{l+s+sa}{rf}\PY{l+s+s2}{\PYZdq{}}\PY{l+s+s2}{ = }\PY{l+s+si}{\PYZob{}}\PY{n}{hours\PYZus{}mean}\PY{l+s+si}{:}\PY{l+s+s2}{.1f}\PY{l+s+si}{\PYZcb{}}\PY{l+s+s2}{\PYZdq{}}
\PY{p}{)}\PY{p}{)}

\PY{c+c1}{\PYZsh{} 2. Mean of exam scores}
\PY{n}{scores\PYZus{}mean} \PY{o}{=} \PY{n}{exam\PYZus{}scores}\PY{o}{.}\PY{n}{mean}\PY{p}{(}\PY{p}{)}
\PY{n}{scores\PYZus{}values} \PY{o}{=} \PY{l+s+s2}{\PYZdq{}}\PY{l+s+s2}{ + }\PY{l+s+s2}{\PYZdq{}}\PY{o}{.}\PY{n}{join}\PY{p}{(}\PY{l+s+sa}{f}\PY{l+s+s2}{\PYZdq{}}\PY{l+s+si}{\PYZob{}}\PY{n}{x}\PY{l+s+si}{:}\PY{l+s+s2}{g}\PY{l+s+si}{\PYZcb{}}\PY{l+s+s2}{\PYZdq{}} \PY{k}{for} \PY{n}{x} \PY{o+ow}{in} \PY{n}{exam\PYZus{}scores}\PY{p}{)}

\PY{n}{display}\PY{p}{(}\PY{n}{Math}\PY{p}{(}
    \PY{l+s+sa}{rf}\PY{l+s+s2}{\PYZdq{}}\PY{l+s+s2}{\PYZbs{}}\PY{l+s+s2}{bar}\PY{l+s+se}{\PYZob{}\PYZob{}}\PY{l+s+s2}{x}\PY{l+s+se}{\PYZcb{}\PYZcb{}}\PY{l+s+s2}{\PYZus{}}\PY{l+s+se}{\PYZob{}\PYZob{}}\PY{l+s+s2}{\PYZbs{}}\PY{l+s+s2}{mathrm}\PY{l+s+se}{\PYZob{}\PYZob{}}\PY{l+s+s2}{scores}\PY{l+s+se}{\PYZcb{}\PYZcb{}}\PY{l+s+se}{\PYZcb{}\PYZcb{}}\PY{l+s+s2}{\PYZdq{}}
    \PY{l+s+sa}{rf}\PY{l+s+s2}{\PYZdq{}}\PY{l+s+s2}{ = }\PY{l+s+s2}{\PYZbs{}}\PY{l+s+s2}{frac}\PY{l+s+se}{\PYZob{}\PYZob{}}\PY{l+s+si}{\PYZob{}}\PY{n}{scores\PYZus{}values}\PY{l+s+si}{\PYZcb{}}\PY{l+s+se}{\PYZcb{}\PYZcb{}}\PY{l+s+se}{\PYZob{}\PYZob{}}\PY{l+s+si}{\PYZob{}}\PY{n}{n}\PY{l+s+si}{\PYZcb{}}\PY{l+s+se}{\PYZcb{}\PYZcb{}}\PY{l+s+s2}{\PYZdq{}}
    \PY{l+s+sa}{rf}\PY{l+s+s2}{\PYZdq{}}\PY{l+s+s2}{ = }\PY{l+s+si}{\PYZob{}}\PY{n}{scores\PYZus{}mean}\PY{l+s+si}{:}\PY{l+s+s2}{.1f}\PY{l+s+si}{\PYZcb{}}\PY{l+s+s2}{\PYZdq{}}
\PY{p}{)}\PY{p}{)}

\PY{c+c1}{\PYZsh{} Helper function to show deviations before squaring}
\PY{k}{def}\PY{+w}{ }\PY{n+nf}{latex\PYZus{}squared\PYZus{}deviations}\PY{p}{(}\PY{n}{data}\PY{p}{,} \PY{n}{mean}\PY{p}{)}\PY{p}{:}
    \PY{k}{return} \PY{l+s+s2}{\PYZdq{}}\PY{l+s+s2}{ + }\PY{l+s+s2}{\PYZdq{}}\PY{o}{.}\PY{n}{join}\PY{p}{(}
        \PY{l+s+sa}{rf}\PY{l+s+s2}{\PYZdq{}}\PY{l+s+s2}{(}\PY{l+s+si}{\PYZob{}}\PY{n}{x}\PY{l+s+si}{:}\PY{l+s+s2}{g}\PY{l+s+si}{\PYZcb{}}\PY{l+s+s2}{ \PYZhy{} }\PY{l+s+si}{\PYZob{}}\PY{n}{mean}\PY{l+s+si}{:}\PY{l+s+s2}{g}\PY{l+s+si}{\PYZcb{}}\PY{l+s+s2}{)\PYZca{}2}\PY{l+s+s2}{\PYZdq{}}
        \PY{k}{for} \PY{n}{x} \PY{o+ow}{in} \PY{n}{data}
    \PY{p}{)}

\PY{c+c1}{\PYZsh{} 3. Sample variance of hours}
\PY{n}{display}\PY{p}{(}\PY{n}{Math}\PY{p}{(}
    \PY{l+s+sa}{rf}\PY{l+s+s2}{\PYZdq{}}\PY{l+s+s2}{\PYZbs{}}\PY{l+s+s2}{begin}\PY{l+s+se}{\PYZob{}\PYZob{}}\PY{l+s+s2}{aligned}\PY{l+s+se}{\PYZcb{}\PYZcb{}}\PY{l+s+s2}{\PYZdq{}}
    \PY{l+s+sa}{rf}\PY{l+s+s2}{\PYZdq{}}\PY{l+s+s2}{s\PYZus{}}\PY{l+s+se}{\PYZob{}\PYZob{}}\PY{l+s+s2}{\PYZbs{}}\PY{l+s+s2}{mathrm}\PY{l+s+se}{\PYZob{}\PYZob{}}\PY{l+s+s2}{hours}\PY{l+s+se}{\PYZcb{}\PYZcb{}}\PY{l+s+se}{\PYZcb{}\PYZcb{}}\PY{l+s+s2}{\PYZca{}2}\PY{l+s+s2}{\PYZdq{}}
    \PY{l+s+sa}{rf}\PY{l+s+s2}{\PYZdq{}}\PY{l+s+s2}{ \PYZam{}= }\PY{l+s+s2}{\PYZbs{}}\PY{l+s+s2}{frac}\PY{l+s+se}{\PYZob{}\PYZob{}}\PY{l+s+si}{\PYZob{}}\PY{n}{latex\PYZus{}squared\PYZus{}deviations}\PY{p}{(}\PY{n}{hours}\PY{p}{,}\PY{+w}{ }\PY{n}{hours\PYZus{}mean}\PY{p}{)}\PY{l+s+si}{\PYZcb{}}\PY{l+s+se}{\PYZcb{}\PYZcb{}}\PY{l+s+se}{\PYZob{}\PYZob{}}\PY{l+s+si}{\PYZob{}}\PY{n}{n}\PY{l+s+si}{\PYZcb{}}\PY{l+s+s2}{\PYZhy{}1}\PY{l+s+se}{\PYZcb{}\PYZcb{}}\PY{l+s+s2}{ }\PY{l+s+se}{\PYZbs{}\PYZbs{}}\PY{l+s+s2}{[6pt]}\PY{l+s+s2}{\PYZdq{}}
    \PY{l+s+sa}{rf}\PY{l+s+s2}{\PYZdq{}}\PY{l+s+s2}{ \PYZam{}= }\PY{l+s+s2}{\PYZbs{}}\PY{l+s+s2}{frac}\PY{l+s+se}{\PYZob{}\PYZob{}}\PY{l+s+si}{\PYZob{}}\PY{n}{latex\PYZus{}sum}\PY{p}{(}\PY{n}{hours\PYZus{}squared}\PY{p}{)}\PY{l+s+si}{\PYZcb{}}\PY{l+s+se}{\PYZcb{}\PYZcb{}}\PY{l+s+se}{\PYZob{}\PYZob{}}\PY{l+s+si}{\PYZob{}}\PY{n}{n}\PY{o}{\PYZhy{}}\PY{l+m+mi}{1}\PY{l+s+si}{\PYZcb{}}\PY{l+s+se}{\PYZcb{}\PYZcb{}}\PY{l+s+s2}{\PYZdq{}}
    \PY{l+s+sa}{rf}\PY{l+s+s2}{\PYZdq{}}\PY{l+s+s2}{ = }\PY{l+s+s2}{\PYZbs{}}\PY{l+s+s2}{frac}\PY{l+s+se}{\PYZob{}\PYZob{}}\PY{l+s+si}{\PYZob{}}\PY{n}{hours\PYZus{}squared}\PY{o}{.}\PY{n}{sum}\PY{p}{(}\PY{p}{)}\PY{l+s+si}{:}\PY{l+s+s2}{g}\PY{l+s+se}{\PYZcb{}\PYZcb{}}\PY{l+s+si}{\PYZcb{}}\PY{l+s+se}{\PYZob{}\PYZob{}}\PY{l+s+si}{\PYZob{}}\PY{n}{n}\PY{o}{\PYZhy{}}\PY{l+m+mi}{1}\PY{l+s+si}{\PYZcb{}}\PY{l+s+se}{\PYZcb{}\PYZcb{}}\PY{l+s+s2}{\PYZdq{}}
    \PY{l+s+sa}{rf}\PY{l+s+s2}{\PYZdq{}}\PY{l+s+s2}{ = }\PY{l+s+s2}{\PYZbs{}}\PY{l+s+s2}{boxed}\PY{l+s+se}{\PYZob{}\PYZob{}}\PY{l+s+si}{\PYZob{}}\PY{n}{np}\PY{o}{.}\PY{n}{var}\PY{p}{(}\PY{n}{hours}\PY{p}{,}\PY{+w}{ }\PY{n}{ddof}\PY{o}{=}\PY{l+m+mi}{1}\PY{p}{)}\PY{l+s+si}{:}\PY{l+s+s2}{g}\PY{l+s+se}{\PYZcb{}\PYZcb{}}\PY{l+s+si}{\PYZcb{}}\PY{l+s+s2}{\PYZdq{}}
    \PY{l+s+sa}{rf}\PY{l+s+s2}{\PYZdq{}}\PY{l+s+s2}{\PYZbs{}}\PY{l+s+s2}{end}\PY{l+s+se}{\PYZob{}\PYZob{}}\PY{l+s+s2}{aligned}\PY{l+s+se}{\PYZcb{}\PYZcb{}}\PY{l+s+s2}{\PYZdq{}}
\PY{p}{)}\PY{p}{)}

\PY{c+c1}{\PYZsh{} 4. Sample variance of scores}
\PY{n}{display}\PY{p}{(}\PY{n}{Math}\PY{p}{(}
    \PY{l+s+sa}{rf}\PY{l+s+s2}{\PYZdq{}}\PY{l+s+s2}{\PYZbs{}}\PY{l+s+s2}{begin}\PY{l+s+se}{\PYZob{}\PYZob{}}\PY{l+s+s2}{aligned}\PY{l+s+se}{\PYZcb{}\PYZcb{}}\PY{l+s+s2}{\PYZdq{}}
    \PY{l+s+sa}{rf}\PY{l+s+s2}{\PYZdq{}}\PY{l+s+s2}{s\PYZus{}}\PY{l+s+se}{\PYZob{}\PYZob{}}\PY{l+s+s2}{\PYZbs{}}\PY{l+s+s2}{mathrm}\PY{l+s+se}{\PYZob{}\PYZob{}}\PY{l+s+s2}{scores}\PY{l+s+se}{\PYZcb{}\PYZcb{}}\PY{l+s+se}{\PYZcb{}\PYZcb{}}\PY{l+s+s2}{\PYZca{}2}\PY{l+s+s2}{\PYZdq{}}
    \PY{l+s+sa}{rf}\PY{l+s+s2}{\PYZdq{}}\PY{l+s+s2}{ \PYZam{}= }\PY{l+s+s2}{\PYZbs{}}\PY{l+s+s2}{frac}\PY{l+s+se}{\PYZob{}\PYZob{}}\PY{l+s+si}{\PYZob{}}\PY{n}{latex\PYZus{}squared\PYZus{}deviations}\PY{p}{(}\PY{n}{exam\PYZus{}scores}\PY{p}{,}\PY{+w}{ }\PY{n}{scores\PYZus{}mean}\PY{p}{)}\PY{l+s+si}{\PYZcb{}}\PY{l+s+se}{\PYZcb{}\PYZcb{}}\PY{l+s+se}{\PYZob{}\PYZob{}}\PY{l+s+si}{\PYZob{}}\PY{n}{n}\PY{l+s+si}{\PYZcb{}}\PY{l+s+s2}{\PYZhy{}1}\PY{l+s+se}{\PYZcb{}\PYZcb{}}\PY{l+s+s2}{ }\PY{l+s+se}{\PYZbs{}\PYZbs{}}\PY{l+s+s2}{[6pt]}\PY{l+s+s2}{\PYZdq{}}
    \PY{l+s+sa}{rf}\PY{l+s+s2}{\PYZdq{}}\PY{l+s+s2}{ \PYZam{}= }\PY{l+s+s2}{\PYZbs{}}\PY{l+s+s2}{frac}\PY{l+s+se}{\PYZob{}\PYZob{}}\PY{l+s+si}{\PYZob{}}\PY{n}{latex\PYZus{}sum}\PY{p}{(}\PY{n}{scores\PYZus{}squared}\PY{p}{)}\PY{l+s+si}{\PYZcb{}}\PY{l+s+se}{\PYZcb{}\PYZcb{}}\PY{l+s+se}{\PYZob{}\PYZob{}}\PY{l+s+si}{\PYZob{}}\PY{n}{n}\PY{o}{\PYZhy{}}\PY{l+m+mi}{1}\PY{l+s+si}{\PYZcb{}}\PY{l+s+se}{\PYZcb{}\PYZcb{}}\PY{l+s+s2}{\PYZdq{}}
    \PY{l+s+sa}{rf}\PY{l+s+s2}{\PYZdq{}}\PY{l+s+s2}{ = }\PY{l+s+s2}{\PYZbs{}}\PY{l+s+s2}{frac}\PY{l+s+se}{\PYZob{}\PYZob{}}\PY{l+s+si}{\PYZob{}}\PY{n}{scores\PYZus{}squared}\PY{o}{.}\PY{n}{sum}\PY{p}{(}\PY{p}{)}\PY{l+s+si}{:}\PY{l+s+s2}{g}\PY{l+s+se}{\PYZcb{}\PYZcb{}}\PY{l+s+si}{\PYZcb{}}\PY{l+s+se}{\PYZob{}\PYZob{}}\PY{l+s+si}{\PYZob{}}\PY{n}{n}\PY{o}{\PYZhy{}}\PY{l+m+mi}{1}\PY{l+s+si}{\PYZcb{}}\PY{l+s+se}{\PYZcb{}\PYZcb{}}\PY{l+s+s2}{\PYZdq{}}
    \PY{l+s+sa}{rf}\PY{l+s+s2}{\PYZdq{}}\PY{l+s+s2}{ = }\PY{l+s+s2}{\PYZbs{}}\PY{l+s+s2}{boxed}\PY{l+s+se}{\PYZob{}\PYZob{}}\PY{l+s+si}{\PYZob{}}\PY{n}{np}\PY{o}{.}\PY{n}{var}\PY{p}{(}\PY{n}{exam\PYZus{}scores}\PY{p}{,}\PY{+w}{ }\PY{n}{ddof}\PY{o}{=}\PY{l+m+mi}{1}\PY{p}{)}\PY{l+s+si}{:}\PY{l+s+s2}{g}\PY{l+s+se}{\PYZcb{}\PYZcb{}}\PY{l+s+si}{\PYZcb{}}\PY{l+s+s2}{\PYZdq{}}
    \PY{l+s+sa}{rf}\PY{l+s+s2}{\PYZdq{}}\PY{l+s+s2}{\PYZbs{}}\PY{l+s+s2}{end}\PY{l+s+se}{\PYZob{}\PYZob{}}\PY{l+s+s2}{aligned}\PY{l+s+se}{\PYZcb{}\PYZcb{}}\PY{l+s+s2}{\PYZdq{}}
\PY{p}{)}\PY{p}{)}

\PY{c+c1}{\PYZsh{} Sample standard deviations}
\PY{n}{hours\PYZus{}std} \PY{o}{=} \PY{n}{np}\PY{o}{.}\PY{n}{std}\PY{p}{(}\PY{n}{hours}\PY{p}{,} \PY{n}{ddof}\PY{o}{=}\PY{l+m+mi}{1}\PY{p}{)}
\PY{n}{scores\PYZus{}std} \PY{o}{=} \PY{n}{np}\PY{o}{.}\PY{n}{std}\PY{p}{(}\PY{n}{exam\PYZus{}scores}\PY{p}{,} \PY{n}{ddof}\PY{o}{=}\PY{l+m+mi}{1}\PY{p}{)}

\PY{c+c1}{\PYZsh{} 5. Standard deviation of study hours}
\PY{n}{display}\PY{p}{(}\PY{n}{Math}\PY{p}{(}
    \PY{l+s+sa}{rf}\PY{l+s+s2}{\PYZdq{}}\PY{l+s+s2}{\PYZbs{}}\PY{l+s+s2}{begin}\PY{l+s+se}{\PYZob{}\PYZob{}}\PY{l+s+s2}{aligned}\PY{l+s+se}{\PYZcb{}\PYZcb{}}\PY{l+s+s2}{\PYZdq{}}
    \PY{l+s+sa}{rf}\PY{l+s+s2}{\PYZdq{}}\PY{l+s+s2}{s\PYZus{}}\PY{l+s+se}{\PYZob{}\PYZob{}}\PY{l+s+s2}{\PYZbs{}}\PY{l+s+s2}{mathrm}\PY{l+s+se}{\PYZob{}\PYZob{}}\PY{l+s+s2}{hours}\PY{l+s+se}{\PYZcb{}\PYZcb{}}\PY{l+s+se}{\PYZcb{}\PYZcb{}}\PY{l+s+s2}{\PYZdq{}}
    \PY{l+s+sa}{rf}\PY{l+s+s2}{\PYZdq{}}\PY{l+s+s2}{ \PYZam{}= }\PY{l+s+s2}{\PYZbs{}}\PY{l+s+s2}{sqrt}\PY{l+s+se}{\PYZob{}\PYZob{}}\PY{l+s+s2}{s\PYZus{}}\PY{l+s+se}{\PYZob{}\PYZob{}}\PY{l+s+s2}{\PYZbs{}}\PY{l+s+s2}{mathrm}\PY{l+s+se}{\PYZob{}\PYZob{}}\PY{l+s+s2}{hours}\PY{l+s+se}{\PYZcb{}\PYZcb{}}\PY{l+s+se}{\PYZcb{}\PYZcb{}}\PY{l+s+s2}{\PYZca{}2}\PY{l+s+se}{\PYZcb{}\PYZcb{}}\PY{l+s+s2}{\PYZdq{}}
    \PY{l+s+sa}{rf}\PY{l+s+s2}{\PYZdq{}}\PY{l+s+s2}{ = }\PY{l+s+s2}{\PYZbs{}}\PY{l+s+s2}{sqrt}\PY{l+s+se}{\PYZob{}\PYZob{}}\PY{l+s+si}{\PYZob{}}\PY{n}{np}\PY{o}{.}\PY{n}{var}\PY{p}{(}\PY{n}{hours}\PY{p}{,}\PY{+w}{ }\PY{n}{ddof}\PY{o}{=}\PY{l+m+mi}{1}\PY{p}{)}\PY{l+s+si}{:}\PY{l+s+s2}{g}\PY{l+s+se}{\PYZcb{}\PYZcb{}}\PY{l+s+si}{\PYZcb{}}\PY{l+s+s2}{\PYZdq{}}
    \PY{l+s+sa}{rf}\PY{l+s+s2}{\PYZdq{}}\PY{l+s+s2}{ = }\PY{l+s+s2}{\PYZbs{}}\PY{l+s+s2}{boxed}\PY{l+s+se}{\PYZob{}\PYZob{}}\PY{l+s+si}{\PYZob{}}\PY{n}{hours\PYZus{}std}\PY{l+s+si}{:}\PY{l+s+s2}{.3f}\PY{l+s+se}{\PYZcb{}\PYZcb{}}\PY{l+s+si}{\PYZcb{}}\PY{l+s+s2}{\PYZdq{}}
    \PY{l+s+sa}{rf}\PY{l+s+s2}{\PYZdq{}}\PY{l+s+s2}{\PYZbs{}}\PY{l+s+s2}{end}\PY{l+s+se}{\PYZob{}\PYZob{}}\PY{l+s+s2}{aligned}\PY{l+s+se}{\PYZcb{}\PYZcb{}}\PY{l+s+s2}{\PYZdq{}}
\PY{p}{)}\PY{p}{)}

\PY{c+c1}{\PYZsh{} 6. Standard deviation of exam scores}
\PY{n}{display}\PY{p}{(}\PY{n}{Math}\PY{p}{(}
    \PY{l+s+sa}{rf}\PY{l+s+s2}{\PYZdq{}}\PY{l+s+s2}{\PYZbs{}}\PY{l+s+s2}{begin}\PY{l+s+se}{\PYZob{}\PYZob{}}\PY{l+s+s2}{aligned}\PY{l+s+se}{\PYZcb{}\PYZcb{}}\PY{l+s+s2}{\PYZdq{}}
    \PY{l+s+sa}{rf}\PY{l+s+s2}{\PYZdq{}}\PY{l+s+s2}{s\PYZus{}}\PY{l+s+se}{\PYZob{}\PYZob{}}\PY{l+s+s2}{\PYZbs{}}\PY{l+s+s2}{mathrm}\PY{l+s+se}{\PYZob{}\PYZob{}}\PY{l+s+s2}{scores}\PY{l+s+se}{\PYZcb{}\PYZcb{}}\PY{l+s+se}{\PYZcb{}\PYZcb{}}\PY{l+s+s2}{\PYZdq{}}
    \PY{l+s+sa}{rf}\PY{l+s+s2}{\PYZdq{}}\PY{l+s+s2}{ \PYZam{}= }\PY{l+s+s2}{\PYZbs{}}\PY{l+s+s2}{sqrt}\PY{l+s+se}{\PYZob{}\PYZob{}}\PY{l+s+s2}{s\PYZus{}}\PY{l+s+se}{\PYZob{}\PYZob{}}\PY{l+s+s2}{\PYZbs{}}\PY{l+s+s2}{mathrm}\PY{l+s+se}{\PYZob{}\PYZob{}}\PY{l+s+s2}{scores}\PY{l+s+se}{\PYZcb{}\PYZcb{}}\PY{l+s+se}{\PYZcb{}\PYZcb{}}\PY{l+s+s2}{\PYZca{}2}\PY{l+s+se}{\PYZcb{}\PYZcb{}}\PY{l+s+s2}{\PYZdq{}}
    \PY{l+s+sa}{rf}\PY{l+s+s2}{\PYZdq{}}\PY{l+s+s2}{ = }\PY{l+s+s2}{\PYZbs{}}\PY{l+s+s2}{sqrt}\PY{l+s+se}{\PYZob{}\PYZob{}}\PY{l+s+si}{\PYZob{}}\PY{n}{np}\PY{o}{.}\PY{n}{var}\PY{p}{(}\PY{n}{exam\PYZus{}scores}\PY{p}{,}\PY{+w}{ }\PY{n}{ddof}\PY{o}{=}\PY{l+m+mi}{1}\PY{p}{)}\PY{l+s+si}{:}\PY{l+s+s2}{g}\PY{l+s+se}{\PYZcb{}\PYZcb{}}\PY{l+s+si}{\PYZcb{}}\PY{l+s+s2}{\PYZdq{}}
    \PY{l+s+sa}{rf}\PY{l+s+s2}{\PYZdq{}}\PY{l+s+s2}{ = }\PY{l+s+s2}{\PYZbs{}}\PY{l+s+s2}{boxed}\PY{l+s+se}{\PYZob{}\PYZob{}}\PY{l+s+si}{\PYZob{}}\PY{n}{scores\PYZus{}std}\PY{l+s+si}{:}\PY{l+s+s2}{.3f}\PY{l+s+se}{\PYZcb{}\PYZcb{}}\PY{l+s+si}{\PYZcb{}}\PY{l+s+s2}{\PYZdq{}}
    \PY{l+s+sa}{rf}\PY{l+s+s2}{\PYZdq{}}\PY{l+s+s2}{\PYZbs{}}\PY{l+s+s2}{end}\PY{l+s+se}{\PYZob{}\PYZob{}}\PY{l+s+s2}{aligned}\PY{l+s+se}{\PYZcb{}\PYZcb{}}\PY{l+s+s2}{\PYZdq{}}
\PY{p}{)}\PY{p}{)}

\PY{c+c1}{\PYZsh{} 7. Sample covariance}
\PY{c+c1}{\PYZsh{} Show each pair of deviations multiplied together}
\PY{n}{covariance\PYZus{}terms} \PY{o}{=} \PY{l+s+s2}{\PYZdq{}}\PY{l+s+s2}{ + }\PY{l+s+s2}{\PYZdq{}}\PY{o}{.}\PY{n}{join}\PY{p}{(}
    \PY{l+s+sa}{rf}\PY{l+s+s2}{\PYZdq{}}\PY{l+s+s2}{(}\PY{l+s+si}{\PYZob{}}\PY{n}{h}\PY{l+s+si}{:}\PY{l+s+s2}{g}\PY{l+s+si}{\PYZcb{}}\PY{l+s+s2}{\PYZhy{}}\PY{l+s+si}{\PYZob{}}\PY{n}{hours\PYZus{}mean}\PY{l+s+si}{:}\PY{l+s+s2}{g}\PY{l+s+si}{\PYZcb{}}\PY{l+s+s2}{)(}\PY{l+s+si}{\PYZob{}}\PY{n}{s}\PY{l+s+si}{:}\PY{l+s+s2}{g}\PY{l+s+si}{\PYZcb{}}\PY{l+s+s2}{\PYZhy{}}\PY{l+s+si}{\PYZob{}}\PY{n}{scores\PYZus{}mean}\PY{l+s+si}{:}\PY{l+s+s2}{g}\PY{l+s+si}{\PYZcb{}}\PY{l+s+s2}{)}\PY{l+s+s2}{\PYZdq{}}
    \PY{k}{for} \PY{n}{h}\PY{p}{,} \PY{n}{s} \PY{o+ow}{in} \PY{n+nb}{zip}\PY{p}{(}\PY{n}{hours}\PY{p}{,} \PY{n}{exam\PYZus{}scores}\PY{p}{)}
\PY{p}{)}

\PY{n}{display}\PY{p}{(}\PY{n}{Math}\PY{p}{(}
    \PY{l+s+sa}{r}\PY{l+s+s2}{\PYZdq{}}\PY{l+s+s2}{\PYZbs{}}\PY{l+s+s2}{operatorname}\PY{l+s+si}{\PYZob{}Cov\PYZcb{}}\PY{l+s+s2}{(X,Y)}\PY{l+s+s2}{\PYZdq{}}
    \PY{l+s+sa}{r}\PY{l+s+s2}{\PYZdq{}}\PY{l+s+s2}{ = }\PY{l+s+s2}{\PYZbs{}}\PY{l+s+s2}{frac}\PY{l+s+s2}{\PYZob{}}\PY{l+s+s2}{\PYZbs{}}\PY{l+s+s2}{sum\PYZus{}}\PY{l+s+s2}{\PYZob{}}\PY{l+s+s2}{i=1\PYZcb{}\PYZca{}}\PY{l+s+si}{\PYZob{}n\PYZcb{}}\PY{l+s+s2}{(x\PYZus{}i\PYZhy{}}\PY{l+s+s2}{\PYZbs{}}\PY{l+s+s2}{bar}\PY{l+s+si}{\PYZob{}x\PYZcb{}}\PY{l+s+s2}{)(y\PYZus{}i\PYZhy{}}\PY{l+s+s2}{\PYZbs{}}\PY{l+s+s2}{bar}\PY{l+s+si}{\PYZob{}y\PYZcb{}}\PY{l+s+s2}{)\PYZcb{}}\PY{l+s+s2}{\PYZob{}}\PY{l+s+s2}{n\PYZhy{}1\PYZcb{}}\PY{l+s+s2}{\PYZdq{}}
\PY{p}{)}\PY{p}{)}

\PY{n}{display}\PY{p}{(}\PY{n}{Math}\PY{p}{(}
    \PY{l+s+sa}{rf}\PY{l+s+s2}{\PYZdq{}}\PY{l+s+s2}{\PYZbs{}}\PY{l+s+s2}{begin}\PY{l+s+se}{\PYZob{}\PYZob{}}\PY{l+s+s2}{aligned}\PY{l+s+se}{\PYZcb{}\PYZcb{}}\PY{l+s+s2}{\PYZdq{}}
    \PY{l+s+sa}{rf}\PY{l+s+s2}{\PYZdq{}}\PY{l+s+s2}{\PYZbs{}}\PY{l+s+s2}{operatorname}\PY{l+s+se}{\PYZob{}\PYZob{}}\PY{l+s+s2}{Cov}\PY{l+s+se}{\PYZcb{}\PYZcb{}}\PY{l+s+s2}{(}\PY{l+s+s2}{\PYZbs{}}\PY{l+s+s2}{mathrm}\PY{l+s+se}{\PYZob{}\PYZob{}}\PY{l+s+s2}{hours}\PY{l+s+se}{\PYZcb{}\PYZcb{}}\PY{l+s+s2}{,}\PY{l+s+s2}{\PYZbs{}}\PY{l+s+s2}{mathrm}\PY{l+s+se}{\PYZob{}\PYZob{}}\PY{l+s+s2}{scores}\PY{l+s+se}{\PYZcb{}\PYZcb{}}\PY{l+s+s2}{)}\PY{l+s+s2}{\PYZdq{}}
    \PY{l+s+sa}{rf}\PY{l+s+s2}{\PYZdq{}}\PY{l+s+s2}{ \PYZam{}= }\PY{l+s+s2}{\PYZbs{}}\PY{l+s+s2}{frac}\PY{l+s+se}{\PYZob{}\PYZob{}}\PY{l+s+si}{\PYZob{}}\PY{n}{covariance\PYZus{}terms}\PY{l+s+si}{\PYZcb{}}\PY{l+s+se}{\PYZcb{}\PYZcb{}}\PY{l+s+se}{\PYZob{}\PYZob{}}\PY{l+s+si}{\PYZob{}}\PY{n}{n}\PY{l+s+si}{\PYZcb{}}\PY{l+s+s2}{\PYZhy{}1}\PY{l+s+se}{\PYZcb{}\PYZcb{}}\PY{l+s+s2}{ }\PY{l+s+se}{\PYZbs{}\PYZbs{}}\PY{l+s+s2}{[6pt]}\PY{l+s+s2}{\PYZdq{}}
    \PY{l+s+sa}{rf}\PY{l+s+s2}{\PYZdq{}}\PY{l+s+s2}{ \PYZam{}= }\PY{l+s+s2}{\PYZbs{}}\PY{l+s+s2}{frac}\PY{l+s+se}{\PYZob{}\PYZob{}}\PY{l+s+si}{\PYZob{}}\PY{n}{latex\PYZus{}sum}\PY{p}{(}\PY{n}{products}\PY{p}{)}\PY{l+s+si}{\PYZcb{}}\PY{l+s+se}{\PYZcb{}\PYZcb{}}\PY{l+s+se}{\PYZob{}\PYZob{}}\PY{l+s+si}{\PYZob{}}\PY{n}{n}\PY{o}{\PYZhy{}}\PY{l+m+mi}{1}\PY{l+s+si}{\PYZcb{}}\PY{l+s+se}{\PYZcb{}\PYZcb{}}\PY{l+s+s2}{\PYZdq{}}
    \PY{l+s+sa}{rf}\PY{l+s+s2}{\PYZdq{}}\PY{l+s+s2}{ = }\PY{l+s+s2}{\PYZbs{}}\PY{l+s+s2}{frac}\PY{l+s+se}{\PYZob{}\PYZob{}}\PY{l+s+si}{\PYZob{}}\PY{n}{products}\PY{o}{.}\PY{n}{sum}\PY{p}{(}\PY{p}{)}\PY{l+s+si}{:}\PY{l+s+s2}{g}\PY{l+s+se}{\PYZcb{}\PYZcb{}}\PY{l+s+si}{\PYZcb{}}\PY{l+s+se}{\PYZob{}\PYZob{}}\PY{l+s+si}{\PYZob{}}\PY{n}{n}\PY{o}{\PYZhy{}}\PY{l+m+mi}{1}\PY{l+s+si}{\PYZcb{}}\PY{l+s+se}{\PYZcb{}\PYZcb{}}\PY{l+s+s2}{\PYZdq{}}
    \PY{l+s+sa}{rf}\PY{l+s+s2}{\PYZdq{}}\PY{l+s+s2}{ = }\PY{l+s+s2}{\PYZbs{}}\PY{l+s+s2}{boxed}\PY{l+s+se}{\PYZob{}\PYZob{}}\PY{l+s+si}{\PYZob{}}\PY{n}{manual\PYZus{}covariance}\PY{l+s+si}{:}\PY{l+s+s2}{g}\PY{l+s+se}{\PYZcb{}\PYZcb{}}\PY{l+s+si}{\PYZcb{}}\PY{l+s+s2}{\PYZdq{}}
    \PY{l+s+sa}{rf}\PY{l+s+s2}{\PYZdq{}}\PY{l+s+s2}{\PYZbs{}}\PY{l+s+s2}{end}\PY{l+s+se}{\PYZob{}\PYZob{}}\PY{l+s+s2}{aligned}\PY{l+s+se}{\PYZcb{}\PYZcb{}}\PY{l+s+s2}{\PYZdq{}}
\PY{p}{)}\PY{p}{)}
\end{Verbatim}
\end{tcolorbox}

    
    \begin{Verbatim}[commandchars=\\\{\}]
   Hours  Score  Hours deviation  Score deviation  Product
0    1.0   50.0             -2.0            -20.0     40.0
1    2.0   60.0             -1.0            -10.0     10.0
2    3.0   70.0              0.0              0.0      0.0
3    4.0   80.0              1.0             10.0     10.0
4    5.0   90.0              2.0             20.0     40.0
    \end{Verbatim}

    
    $\displaystyle \bar{x}_{\mathrm{hours}} = \frac{1 + 2 + 3 + 4 + 5}{5} = 3.0$

    
    $\displaystyle \bar{x}_{\mathrm{scores}} = \frac{50 + 60 + 70 + 80 + 90}{5} = 70.0$

    
    $\displaystyle \begin{aligned}s_{\mathrm{hours}}^2 &= \frac{(1 - 3)^2 + (2 - 3)^2 + (3 - 3)^2 + (4 - 3)^2 + (5 - 3)^2}{5-1} \\[6pt] &= \frac{4 + 1 + 0 + 1 + 4}{4} = \frac{10}{4} = \boxed{2.5}\end{aligned}$

    
    $\displaystyle \begin{aligned}s_{\mathrm{scores}}^2 &= \frac{(50 - 70)^2 + (60 - 70)^2 + (70 - 70)^2 + (80 - 70)^2 + (90 - 70)^2}{5-1} \\[6pt] &= \frac{400 + 100 + 0 + 100 + 400}{4} = \frac{1000}{4} = \boxed{250}\end{aligned}$

    
    $\displaystyle \begin{aligned}s_{\mathrm{hours}} &= \sqrt{s_{\mathrm{hours}}^2} = \sqrt{2.5} = \boxed{1.581}\end{aligned}$

    
    $\displaystyle \begin{aligned}s_{\mathrm{scores}} &= \sqrt{s_{\mathrm{scores}}^2} = \sqrt{250} = \boxed{15.811}\end{aligned}$

    
    $\displaystyle \operatorname{Cov}(X,Y) = \frac{\sum_{i=1}^{n}(x_i-\bar{x})(y_i-\bar{y})}{n-1}$

    
    $\displaystyle \begin{aligned}\operatorname{Cov}(\mathrm{hours},\mathrm{scores}) &= \frac{(1-3)(50-70) + (2-3)(60-70) + (3-3)(70-70) + (4-3)(80-70) + (5-3)(90-70)}{5-1} \\[6pt] &= \frac{40 + 10 + 0 + 10 + 40}{4} = \frac{100}{4} = \boxed{25}\end{aligned}$

    
    \begin{tcolorbox}[breakable, size=fbox, boxrule=1pt, pad at break*=1mm,colback=cellbackground, colframe=cellborder]
\prompt{In}{incolor}{6}{\boxspacing}
\begin{Verbatim}[commandchars=\\\{\}]
\PY{n}{rng} \PY{o}{=} \PY{n}{np}\PY{o}{.}\PY{n}{random}\PY{o}{.}\PY{n}{default\PYZus{}rng}\PY{p}{(}\PY{l+m+mi}{42}\PY{p}{)}
\PY{n}{x} \PY{o}{=} \PY{n}{rng}\PY{o}{.}\PY{n}{normal}\PY{p}{(}\PY{n}{size}\PY{o}{=}\PY{l+m+mi}{100}\PY{p}{)}
\PY{n}{relationships} \PY{o}{=} \PY{p}{\PYZob{}} \PY{c+c1}{\PYZsh{} This contains the \PYZsq{}y\PYZsq{} values for each relationship}
    \PY{l+s+s2}{\PYZdq{}}\PY{l+s+s2}{Positive}\PY{l+s+s2}{\PYZdq{}}\PY{p}{:} \PY{l+m+mf}{1.8}\PY{o}{*}\PY{n}{x} \PY{o}{+} \PY{n}{rng}\PY{o}{.}\PY{n}{normal}\PY{p}{(}\PY{n}{scale}\PY{o}{=}\PY{l+m+mf}{.7}\PY{p}{,} \PY{n}{size}\PY{o}{=}\PY{n+nb}{len}\PY{p}{(}\PY{n}{x}\PY{p}{)}\PY{p}{)}\PY{p}{,}
    \PY{l+s+s2}{\PYZdq{}}\PY{l+s+s2}{Negative}\PY{l+s+s2}{\PYZdq{}}\PY{p}{:} \PY{o}{\PYZhy{}}\PY{l+m+mf}{1.8}\PY{o}{*}\PY{n}{x} \PY{o}{+} \PY{n}{rng}\PY{o}{.}\PY{n}{normal}\PY{p}{(}\PY{n}{scale}\PY{o}{=}\PY{l+m+mf}{.7}\PY{p}{,} \PY{n}{size}\PY{o}{=}\PY{n+nb}{len}\PY{p}{(}\PY{n}{x}\PY{p}{)}\PY{p}{)}\PY{p}{,}
    \PY{l+s+s2}{\PYZdq{}}\PY{l+s+s2}{Near zero}\PY{l+s+s2}{\PYZdq{}}\PY{p}{:} \PY{n}{rng}\PY{o}{.}\PY{n}{normal}\PY{p}{(}\PY{n}{size}\PY{o}{=}\PY{n+nb}{len}\PY{p}{(}\PY{n}{x}\PY{p}{)}\PY{p}{)}
\PY{p}{\PYZcb{}}
\PY{n}{fig}\PY{p}{,} \PY{n}{axes} \PY{o}{=} \PY{n}{plt}\PY{o}{.}\PY{n}{subplots}\PY{p}{(}\PY{l+m+mi}{1}\PY{p}{,} \PY{l+m+mi}{3}\PY{p}{,} \PY{n}{figsize}\PY{o}{=}\PY{p}{(}\PY{l+m+mi}{12}\PY{p}{,} \PY{l+m+mf}{3.5}\PY{p}{)}\PY{p}{,} \PY{n}{sharex}\PY{o}{=}\PY{k+kc}{True}\PY{p}{,} \PY{n}{sharey}\PY{o}{=}\PY{k+kc}{True}\PY{p}{)}
\PY{k}{for} \PY{n}{ax}\PY{p}{,} \PY{p}{(}\PY{n}{name}\PY{p}{,} \PY{n}{y}\PY{p}{)} \PY{o+ow}{in} \PY{n+nb}{zip}\PY{p}{(}\PY{n}{axes}\PY{p}{,} \PY{n}{relationships}\PY{o}{.}\PY{n}{items}\PY{p}{(}\PY{p}{)}\PY{p}{)}\PY{p}{:}
    \PY{n}{ax}\PY{o}{.}\PY{n}{scatter}\PY{p}{(}\PY{n}{x}\PY{p}{,} \PY{n}{y}\PY{p}{,} \PY{n}{alpha}\PY{o}{=}\PY{l+m+mf}{.65}\PY{p}{,} \PY{n}{s}\PY{o}{=}\PY{l+m+mi}{18}\PY{p}{)}
    \PY{n}{covariance} \PY{o}{=} \PY{n}{np}\PY{o}{.}\PY{n}{cov}\PY{p}{(}\PY{n}{x}\PY{p}{,} \PY{n}{y}\PY{p}{)}\PY{p}{[}\PY{l+m+mi}{0}\PY{p}{,} \PY{l+m+mi}{1}\PY{p}{]}
    \PY{n}{ax}\PY{o}{.}\PY{n}{set\PYZus{}title}\PY{p}{(}\PY{l+s+sa}{f}\PY{l+s+s2}{\PYZdq{}}\PY{l+s+si}{\PYZob{}}\PY{n}{name}\PY{l+s+si}{\PYZcb{}}\PY{l+s+se}{\PYZbs{}n}\PY{l+s+s2}{Covariance = }\PY{l+s+si}{\PYZob{}}\PY{n}{covariance}\PY{l+s+si}{:}\PY{l+s+s2}{+.2f}\PY{l+s+si}{\PYZcb{}}\PY{l+s+s2}{\PYZdq{}}\PY{p}{)}
    \PY{n}{ax}\PY{o}{.}\PY{n}{set\PYZus{}xlabel}\PY{p}{(}\PY{l+s+s2}{\PYZdq{}}\PY{l+s+s2}{X}\PY{l+s+s2}{\PYZdq{}}\PY{p}{)}
\PY{n}{axes}\PY{p}{[}\PY{l+m+mi}{0}\PY{p}{]}\PY{o}{.}\PY{n}{set\PYZus{}ylabel}\PY{p}{(}\PY{l+s+s2}{\PYZdq{}}\PY{l+s+s2}{Y}\PY{l+s+s2}{\PYZdq{}}\PY{p}{)}
\PY{n}{fig}\PY{o}{.}\PY{n}{suptitle}\PY{p}{(}\PY{l+s+s2}{\PYZdq{}}\PY{l+s+s2}{Three possible signs of covariance}\PY{l+s+s2}{\PYZdq{}}\PY{p}{)}
\PY{n}{fig}\PY{o}{.}\PY{n}{tight\PYZus{}layout}\PY{p}{(}\PY{p}{)}
\PY{n}{plt}\PY{o}{.}\PY{n}{show}\PY{p}{(}\PY{p}{)}
\end{Verbatim}
\end{tcolorbox}

    \begin{center}
    \adjustimage{max size={0.9\linewidth}{0.9\paperheight}}{week4_pdf_build_pca_files/week4_pdf_build_pca_9_0.png}
    \end{center}
    { \hspace*{\fill} \\}
    

    \subsection{4. Covariance matrix and
correlation}\label{covariance-matrix-and-correlation}

For two variables the sample covariance matrix is

\[
\mathbf{S}=\begin{bmatrix}
\operatorname{Var}(X) & \operatorname{Cov}(X,Y)\\
\operatorname{Cov}(Y,X) & \operatorname{Var}(Y)
\end{bmatrix}.
\]

The diagonal contains \textbf{variances}. Off-diagonal terms contain
\textbf{covariances}. The matrix is symmetric.

Covariance depends on units and scales, so correlation standardises it:

\[r_{XY}=\frac{\operatorname{Cov}(X,Y)}{s_Xs_Y},\qquad -1\le r_{XY}\le1.\]

A correlation of zero does \textbf{not} prove independence.

\subsubsection{Pearson's correlation}\label{pearsons-correlation}

Pearson's correlation coefficient divides the covariance by the product
of the standard deviations of the two variables, not by their variances.

\[
\boxed{r_{XY}=\frac{\operatorname{Cov}(X,Y)}{s_Xs_Y}}
\]

Since standard deviation is the square root of variance:

\[
s_X=\sqrt{\operatorname{Var}(X)}
\]

We can also express correlation as:

\[
\boxed{r_{XY}=\frac{\operatorname{Cov}(X,Y)}{\sqrt{\operatorname{Var}(X)\operatorname{Var}(Y)}}}
\]

These two formulas are mathematically equivalent.

\subsubsection{Why standard deviation instead of
variance?}\label{why-standard-deviation-instead-of-variance}

Because the purpose of correlation is to remove the units and scales of
the original variables.

Let's use your previous example:

{\def\LTcaptype{none} % do not increment counter
\begin{longtable}[]{@{}lrr@{}}
\toprule\noalign{}
Statistic & Study hours & Exam scores \\
\midrule\noalign{}
\endhead
\bottomrule\noalign{}
\endlastfoot
Sample variance & 2.5 & 250 \\
Sample standard deviation & 1.581 & 15.811 \\
Sample covariance (hours, scores) & \textbf{25} & \\
\end{longtable}
}

The covariance has units of hours × score points.

Dividing by the standard deviations (\(s_{x} \times s_{y}\)):

\[
r=\frac{\operatorname{Cov}(hours,scores)}{s_{hours} \times s_{scores}}=\frac{25}{\sqrt{2.5}\sqrt{250}}=\frac{25}{\sqrt{2.5\times 250}}=\frac{25}{25}=\boxed{1}
\]

The units cancel:

\[
\frac{\text{hours}\times\text{points}}{\text{hours}\times\text{points}}=1
\]

The result is dimensionless and always falls between -1 and +1, provided
both variables have non-zero standard deviations.

If you divided by the product of the variances instead
(\(s_{x}^2 \times s_{y}^2\)), you would not obtain a scale-independent
correlation coefficient.


    \begin{tcolorbox}[breakable, size=fbox, boxrule=1pt, pad at break*=1mm,colback=cellbackground, colframe=cellborder]
\prompt{In}{incolor}{7}{\boxspacing}
\begin{Verbatim}[commandchars=\\\{\}]
\PY{n}{data} \PY{o}{=} \PY{n}{pd}\PY{o}{.}\PY{n}{DataFrame}\PY{p}{(}\PY{p}{\PYZob{}}\PY{l+s+s2}{\PYZdq{}}\PY{l+s+s2}{Study hours}\PY{l+s+s2}{\PYZdq{}}\PY{p}{:} \PY{n}{hours}\PY{p}{,} \PY{l+s+s2}{\PYZdq{}}\PY{l+s+s2}{Exam score}\PY{l+s+s2}{\PYZdq{}}\PY{p}{:} \PY{n}{exam\PYZus{}scores}\PY{p}{\PYZcb{}}\PY{p}{)}
\PY{n+nb}{print}\PY{p}{(}\PY{l+s+s2}{\PYZdq{}}\PY{l+s+s2}{Sample covariance matrix:}\PY{l+s+s2}{\PYZdq{}}\PY{p}{)}
\PY{n}{display}\PY{p}{(}\PY{n}{data}\PY{o}{.}\PY{n}{cov}\PY{p}{(}\PY{p}{)}\PY{p}{)}
\PY{n+nb}{print}\PY{p}{(}\PY{l+s+s2}{\PYZdq{}}\PY{l+s+s2}{Correlation matrix:}\PY{l+s+s2}{\PYZdq{}}\PY{p}{)}
\PY{n}{display}\PY{p}{(}\PY{n}{data}\PY{o}{.}\PY{n}{corr}\PY{p}{(}\PY{p}{)}\PY{p}{)}
\PY{n+nb}{print}\PY{p}{(}\PY{l+s+s2}{\PYZdq{}}\PY{l+s+s2}{np.cov agrees with pandas:}\PY{l+s+se}{\PYZbs{}n}\PY{l+s+s2}{\PYZdq{}}\PY{p}{,} \PY{n}{np}\PY{o}{.}\PY{n}{cov}\PY{p}{(}\PY{n}{hours}\PY{p}{,} \PY{n}{exam\PYZus{}scores}\PY{p}{)}\PY{p}{)}
\PY{n+nb}{print}\PY{p}{(}\PY{l+s+s2}{\PYZdq{}}\PY{l+s+s2}{np.corrcoef agrees with pandas:}\PY{l+s+se}{\PYZbs{}n}\PY{l+s+s2}{\PYZdq{}}\PY{p}{,} \PY{n}{np}\PY{o}{.}\PY{n}{corrcoef}\PY{p}{(}\PY{n}{hours}\PY{p}{,} \PY{n}{exam\PYZus{}scores}\PY{p}{)}\PY{p}{)}
\end{Verbatim}
\end{tcolorbox}

    \begin{Verbatim}[commandchars=\\\{\}]
Sample covariance matrix:
    \end{Verbatim}

    
    \begin{Verbatim}[commandchars=\\\{\}]
             Study hours  Exam score
Study hours          2.5        25.0
Exam score          25.0       250.0
    \end{Verbatim}

    
    \begin{Verbatim}[commandchars=\\\{\}]
Correlation matrix:
    \end{Verbatim}

    
    \begin{Verbatim}[commandchars=\\\{\}]
             Study hours  Exam score
Study hours          1.0         1.0
Exam score           1.0         1.0
    \end{Verbatim}

    
    \begin{Verbatim}[commandchars=\\\{\}]
np.cov agrees with pandas:
 [[  2.5  25. ]
 [ 25.  250. ]]
np.corrcoef agrees with pandas:
 [[1. 1.]
 [1. 1.]]
    \end{Verbatim}

    \begin{tcolorbox}[breakable, size=fbox, boxrule=1pt, pad at break*=1mm,colback=cellbackground, colframe=cellborder]
\prompt{In}{incolor}{8}{\boxspacing}
\begin{Verbatim}[commandchars=\\\{\}]
\PY{c+c1}{\PYZsh{} Sample covariance matrix}
\PY{n}{cov\PYZus{}matrix} \PY{o}{=} \PY{n}{np}\PY{o}{.}\PY{n}{cov}\PY{p}{(}\PY{n}{hours}\PY{p}{,} \PY{n}{exam\PYZus{}scores}\PY{p}{)}

\PY{c+c1}{\PYZsh{} Standard deviations from the covariance matrix}
\PY{n}{std} \PY{o}{=} \PY{n}{np}\PY{o}{.}\PY{n}{sqrt}\PY{p}{(}\PY{n}{np}\PY{o}{.}\PY{n}{diag}\PY{p}{(}\PY{n}{cov\PYZus{}matrix}\PY{p}{)}\PY{p}{)}

\PY{c+c1}{\PYZsh{} Convert covariance matrix to correlation matrix}
\PY{n}{corr\PYZus{}from\PYZus{}cov} \PY{o}{=} \PY{n}{cov\PYZus{}matrix} \PY{o}{/} \PY{n}{np}\PY{o}{.}\PY{n}{outer}\PY{p}{(}\PY{n}{std}\PY{p}{,} \PY{n}{std}\PY{p}{)}

\PY{n+nb}{print}\PY{p}{(}\PY{l+s+s2}{\PYZdq{}}\PY{l+s+s2}{Covariance matrix:}\PY{l+s+s2}{\PYZdq{}}\PY{p}{)}
\PY{n+nb}{print}\PY{p}{(}\PY{n}{cov\PYZus{}matrix}\PY{p}{)}

\PY{n+nb}{print}\PY{p}{(}\PY{l+s+s2}{\PYZdq{}}\PY{l+s+se}{\PYZbs{}n}\PY{l+s+s2}{Correlation calculated from covariance:}\PY{l+s+s2}{\PYZdq{}}\PY{p}{)}
\PY{n+nb}{print}\PY{p}{(}\PY{n}{corr\PYZus{}from\PYZus{}cov}\PY{p}{)}

\PY{n+nb}{print}\PY{p}{(}\PY{l+s+s2}{\PYZdq{}}\PY{l+s+se}{\PYZbs{}n}\PY{l+s+s2}{Correlation calculated directly:}\PY{l+s+s2}{\PYZdq{}}\PY{p}{)}
\PY{n+nb}{print}\PY{p}{(}\PY{n}{np}\PY{o}{.}\PY{n}{corrcoef}\PY{p}{(}\PY{n}{hours}\PY{p}{,} \PY{n}{exam\PYZus{}scores}\PY{p}{)}\PY{p}{)}

\PY{c+c1}{\PYZsh{} Verify equivalence}
\PY{n+nb}{print}\PY{p}{(}\PY{l+s+s2}{\PYZdq{}}\PY{l+s+se}{\PYZbs{}n}\PY{l+s+s2}{Both methods agree:}\PY{l+s+s2}{\PYZdq{}}\PY{p}{,}
      \PY{n}{np}\PY{o}{.}\PY{n}{allclose}\PY{p}{(}\PY{n}{corr\PYZus{}from\PYZus{}cov}\PY{p}{,} \PY{n}{np}\PY{o}{.}\PY{n}{corrcoef}\PY{p}{(}\PY{n}{hours}\PY{p}{,} \PY{n}{exam\PYZus{}scores}\PY{p}{)}\PY{p}{)}\PY{p}{)}
\end{Verbatim}
\end{tcolorbox}

    \begin{Verbatim}[commandchars=\\\{\}]
Covariance matrix:
[[  2.5  25. ]
 [ 25.  250. ]]

Correlation calculated from covariance:
[[1. 1.]
 [1. 1.]]

Correlation calculated directly:
[[1. 1.]
 [1. 1.]]

Both methods agree: True
    \end{Verbatim}

    \begin{tcolorbox}[breakable, size=fbox, boxrule=1pt, pad at break*=1mm,colback=cellbackground, colframe=cellborder]
\prompt{In}{incolor}{9}{\boxspacing}
\begin{Verbatim}[commandchars=\\\{\}]
\PY{c+c1}{\PYZsh{} Calculate covariance and correlation matrices}
\PY{n}{cov\PYZus{}matrix} \PY{o}{=} \PY{n}{np}\PY{o}{.}\PY{n}{cov}\PY{p}{(}\PY{n}{hours}\PY{p}{,} \PY{n}{exam\PYZus{}scores}\PY{p}{)}
\PY{n}{corr\PYZus{}matrix} \PY{o}{=} \PY{n}{np}\PY{o}{.}\PY{n}{corrcoef}\PY{p}{(}\PY{n}{hours}\PY{p}{,} \PY{n}{exam\PYZus{}scores}\PY{p}{)}

\PY{c+c1}{\PYZsh{} Calculate standard deviations}
\PY{n}{std} \PY{o}{=} \PY{n}{np}\PY{o}{.}\PY{n}{sqrt}\PY{p}{(}\PY{n}{np}\PY{o}{.}\PY{n}{diag}\PY{p}{(}\PY{n}{cov\PYZus{}matrix}\PY{p}{)}\PY{p}{)}

\PY{c+c1}{\PYZsh{} Convert a NumPy matrix to LaTeX}
\PY{k}{def}\PY{+w}{ }\PY{n+nf}{latex\PYZus{}matrix}\PY{p}{(}\PY{n}{matrix}\PY{p}{,} \PY{n}{precision}\PY{o}{=}\PY{l+m+mi}{3}\PY{p}{)}\PY{p}{:}
    \PY{n}{rows} \PY{o}{=} \PY{p}{[}
        \PY{l+s+s2}{\PYZdq{}}\PY{l+s+s2}{ \PYZam{} }\PY{l+s+s2}{\PYZdq{}}\PY{o}{.}\PY{n}{join}\PY{p}{(}\PY{l+s+sa}{f}\PY{l+s+s2}{\PYZdq{}}\PY{l+s+si}{\PYZob{}}\PY{n}{v}\PY{l+s+si}{:}\PY{l+s+s2}{.}\PY{l+s+si}{\PYZob{}}\PY{n}{precision}\PY{l+s+si}{\PYZcb{}}\PY{l+s+s2}{f}\PY{l+s+si}{\PYZcb{}}\PY{l+s+s2}{\PYZdq{}} \PY{k}{for} \PY{n}{v} \PY{o+ow}{in} \PY{n}{row}\PY{p}{)}
        \PY{k}{for} \PY{n}{row} \PY{o+ow}{in} \PY{n}{matrix}
    \PY{p}{]}
    \PY{k}{return} \PY{l+s+sa}{r}\PY{l+s+s2}{\PYZdq{}}\PY{l+s+s2}{\PYZbs{}}\PY{l+s+s2}{begin}\PY{l+s+si}{\PYZob{}bmatrix\PYZcb{}}\PY{l+s+s2}{\PYZdq{}} \PY{o}{+} \PY{l+s+sa}{r}\PY{l+s+s2}{\PYZdq{}}\PY{l+s+s2}{ }\PY{l+s+se}{\PYZbs{}\PYZbs{}}\PY{l+s+s2}{ }\PY{l+s+s2}{\PYZdq{}}\PY{o}{.}\PY{n}{join}\PY{p}{(}\PY{n}{rows}\PY{p}{)} \PY{o}{+} \PY{l+s+sa}{r}\PY{l+s+s2}{\PYZdq{}}\PY{l+s+s2}{\PYZbs{}}\PY{l+s+s2}{end}\PY{l+s+si}{\PYZob{}bmatrix\PYZcb{}}\PY{l+s+s2}{\PYZdq{}}

\PY{c+c1}{\PYZsh{} Build the intermediate matrix using covariance / standard deviations}
\PY{n}{fractions} \PY{o}{=} \PY{p}{[}
    \PY{p}{[}
        \PY{l+s+sa}{rf}\PY{l+s+s2}{\PYZdq{}}\PY{l+s+s2}{\PYZbs{}}\PY{l+s+s2}{frac}\PY{l+s+se}{\PYZob{}\PYZob{}}\PY{l+s+si}{\PYZob{}}\PY{n}{cov\PYZus{}matrix}\PY{p}{[}\PY{n}{i}\PY{p}{,}\PY{+w}{ }\PY{n}{j}\PY{p}{]}\PY{l+s+si}{:}\PY{l+s+s2}{g}\PY{l+s+se}{\PYZcb{}\PYZcb{}}\PY{l+s+si}{\PYZcb{}}\PY{l+s+s2}{\PYZdq{}}
        \PY{l+s+sa}{rf}\PY{l+s+s2}{\PYZdq{}}\PY{l+s+se}{\PYZob{}\PYZob{}}\PY{l+s+s2}{\PYZbs{}}\PY{l+s+s2}{sqrt}\PY{l+s+se}{\PYZob{}\PYZob{}}\PY{l+s+si}{\PYZob{}}\PY{n}{cov\PYZus{}matrix}\PY{p}{[}\PY{n}{i}\PY{p}{,}\PY{+w}{ }\PY{n}{i}\PY{p}{]}\PY{l+s+si}{:}\PY{l+s+s2}{g}\PY{l+s+se}{\PYZcb{}\PYZcb{}}\PY{l+s+si}{\PYZcb{}}\PY{l+s+s2}{\PYZdq{}}
        \PY{l+s+sa}{rf}\PY{l+s+s2}{\PYZdq{}}\PY{l+s+s2}{\PYZbs{}}\PY{l+s+s2}{sqrt}\PY{l+s+se}{\PYZob{}\PYZob{}}\PY{l+s+si}{\PYZob{}}\PY{n}{cov\PYZus{}matrix}\PY{p}{[}\PY{n}{j}\PY{p}{,}\PY{+w}{ }\PY{n}{j}\PY{p}{]}\PY{l+s+si}{:}\PY{l+s+s2}{g}\PY{l+s+se}{\PYZcb{}\PYZcb{}}\PY{l+s+se}{\PYZcb{}\PYZcb{}}\PY{l+s+si}{\PYZcb{}}\PY{l+s+s2}{\PYZdq{}}
        \PY{k}{for} \PY{n}{j} \PY{o+ow}{in} \PY{n+nb}{range}\PY{p}{(}\PY{l+m+mi}{2}\PY{p}{)}
    \PY{p}{]}
    \PY{k}{for} \PY{n}{i} \PY{o+ow}{in} \PY{n+nb}{range}\PY{p}{(}\PY{l+m+mi}{2}\PY{p}{)}
\PY{p}{]}

\PY{n}{fractions\PYZus{}latex} \PY{o}{=} \PY{p}{(}
    \PY{l+s+sa}{r}\PY{l+s+s2}{\PYZdq{}}\PY{l+s+s2}{\PYZbs{}}\PY{l+s+s2}{begin}\PY{l+s+si}{\PYZob{}bmatrix\PYZcb{}}\PY{l+s+s2}{\PYZdq{}}
    \PY{o}{+} \PY{l+s+sa}{r}\PY{l+s+s2}{\PYZdq{}}\PY{l+s+s2}{ }\PY{l+s+se}{\PYZbs{}\PYZbs{}}\PY{l+s+s2}{ }\PY{l+s+s2}{\PYZdq{}}\PY{o}{.}\PY{n}{join}\PY{p}{(}\PY{l+s+s2}{\PYZdq{}}\PY{l+s+s2}{ \PYZam{} }\PY{l+s+s2}{\PYZdq{}}\PY{o}{.}\PY{n}{join}\PY{p}{(}\PY{n}{row}\PY{p}{)} \PY{k}{for} \PY{n}{row} \PY{o+ow}{in} \PY{n}{fractions}\PY{p}{)}
    \PY{o}{+} \PY{l+s+sa}{r}\PY{l+s+s2}{\PYZdq{}}\PY{l+s+s2}{\PYZbs{}}\PY{l+s+s2}{end}\PY{l+s+si}{\PYZob{}bmatrix\PYZcb{}}\PY{l+s+s2}{\PYZdq{}}
\PY{p}{)}

\PY{c+c1}{\PYZsh{} Display the full derivation}
\PY{n}{display}\PY{p}{(}\PY{n}{Math}\PY{p}{(}
    \PY{l+s+sa}{rf}\PY{l+s+s2}{\PYZdq{}}\PY{l+s+s2}{\PYZbs{}}\PY{l+s+s2}{begin}\PY{l+s+se}{\PYZob{}\PYZob{}}\PY{l+s+s2}{aligned}\PY{l+s+se}{\PYZcb{}\PYZcb{}}\PY{l+s+s2}{\PYZdq{}}
    \PY{l+s+sa}{rf}\PY{l+s+s2}{\PYZdq{}}\PY{l+s+s2}{\PYZbs{}}\PY{l+s+s2}{mathbf}\PY{l+s+se}{\PYZob{}\PYZob{}}\PY{l+s+s2}{R}\PY{l+s+se}{\PYZcb{}\PYZcb{}}\PY{l+s+s2}{\PYZdq{}}
    \PY{l+s+sa}{rf}\PY{l+s+s2}{\PYZdq{}}\PY{l+s+s2}{ \PYZam{}= }\PY{l+s+s2}{\PYZbs{}}\PY{l+s+s2}{begin}\PY{l+s+se}{\PYZob{}\PYZob{}}\PY{l+s+s2}{bmatrix}\PY{l+s+se}{\PYZcb{}\PYZcb{}}\PY{l+s+s2}{\PYZdq{}}
    \PY{l+s+sa}{rf}\PY{l+s+s2}{\PYZdq{}}\PY{l+s+s2}{\PYZbs{}}\PY{l+s+s2}{frac}\PY{l+s+se}{\PYZob{}\PYZob{}}\PY{l+s+s2}{s\PYZus{}X\PYZca{}2}\PY{l+s+se}{\PYZcb{}\PYZcb{}}\PY{l+s+se}{\PYZob{}\PYZob{}}\PY{l+s+s2}{s\PYZus{}Xs\PYZus{}X}\PY{l+s+se}{\PYZcb{}\PYZcb{}}\PY{l+s+s2}{ \PYZam{} }\PY{l+s+s2}{\PYZbs{}}\PY{l+s+s2}{frac}\PY{l+s+se}{\PYZob{}\PYZob{}}\PY{l+s+s2}{s\PYZus{}}\PY{l+s+se}{\PYZob{}\PYZob{}}\PY{l+s+s2}{XY}\PY{l+s+se}{\PYZcb{}\PYZcb{}}\PY{l+s+se}{\PYZcb{}\PYZcb{}}\PY{l+s+se}{\PYZob{}\PYZob{}}\PY{l+s+s2}{s\PYZus{}Xs\PYZus{}Y}\PY{l+s+se}{\PYZcb{}\PYZcb{}}\PY{l+s+s2}{ }\PY{l+s+se}{\PYZbs{}\PYZbs{}}\PY{l+s+s2}{\PYZdq{}}
    \PY{l+s+sa}{rf}\PY{l+s+s2}{\PYZdq{}}\PY{l+s+s2}{\PYZbs{}}\PY{l+s+s2}{frac}\PY{l+s+se}{\PYZob{}\PYZob{}}\PY{l+s+s2}{s\PYZus{}}\PY{l+s+se}{\PYZob{}\PYZob{}}\PY{l+s+s2}{XY}\PY{l+s+se}{\PYZcb{}\PYZcb{}}\PY{l+s+se}{\PYZcb{}\PYZcb{}}\PY{l+s+se}{\PYZob{}\PYZob{}}\PY{l+s+s2}{s\PYZus{}Ys\PYZus{}X}\PY{l+s+se}{\PYZcb{}\PYZcb{}}\PY{l+s+s2}{ \PYZam{} }\PY{l+s+s2}{\PYZbs{}}\PY{l+s+s2}{frac}\PY{l+s+se}{\PYZob{}\PYZob{}}\PY{l+s+s2}{s\PYZus{}Y\PYZca{}2}\PY{l+s+se}{\PYZcb{}\PYZcb{}}\PY{l+s+se}{\PYZob{}\PYZob{}}\PY{l+s+s2}{s\PYZus{}Ys\PYZus{}Y}\PY{l+s+se}{\PYZcb{}\PYZcb{}}\PY{l+s+s2}{\PYZdq{}}
    \PY{l+s+sa}{rf}\PY{l+s+s2}{\PYZdq{}}\PY{l+s+s2}{\PYZbs{}}\PY{l+s+s2}{end}\PY{l+s+se}{\PYZob{}\PYZob{}}\PY{l+s+s2}{bmatrix}\PY{l+s+se}{\PYZcb{}\PYZcb{}}\PY{l+s+s2}{ }\PY{l+s+se}{\PYZbs{}\PYZbs{}}\PY{l+s+s2}{[12pt]}\PY{l+s+s2}{\PYZdq{}}
    \PY{l+s+sa}{rf}\PY{l+s+s2}{\PYZdq{}}\PY{l+s+s2}{ \PYZam{}= }\PY{l+s+si}{\PYZob{}}\PY{n}{fractions\PYZus{}latex}\PY{l+s+si}{\PYZcb{}}\PY{l+s+s2}{ }\PY{l+s+se}{\PYZbs{}\PYZbs{}}\PY{l+s+s2}{[12pt]}\PY{l+s+s2}{\PYZdq{}}
    \PY{l+s+sa}{rf}\PY{l+s+s2}{\PYZdq{}}\PY{l+s+s2}{ \PYZam{}= }\PY{l+s+s2}{\PYZbs{}}\PY{l+s+s2}{boxed}\PY{l+s+se}{\PYZob{}\PYZob{}}\PY{l+s+si}{\PYZob{}}\PY{n}{latex\PYZus{}matrix}\PY{p}{(}\PY{n}{corr\PYZus{}matrix}\PY{p}{)}\PY{l+s+si}{\PYZcb{}}\PY{l+s+se}{\PYZcb{}\PYZcb{}}\PY{l+s+s2}{\PYZdq{}}
    \PY{l+s+sa}{rf}\PY{l+s+s2}{\PYZdq{}}\PY{l+s+s2}{\PYZbs{}}\PY{l+s+s2}{end}\PY{l+s+se}{\PYZob{}\PYZob{}}\PY{l+s+s2}{aligned}\PY{l+s+se}{\PYZcb{}\PYZcb{}}\PY{l+s+s2}{\PYZdq{}}
\PY{p}{)}\PY{p}{)}
\end{Verbatim}
\end{tcolorbox}

    $\displaystyle \begin{aligned}\mathbf{R} &= \begin{bmatrix}\frac{s_X^2}{s_Xs_X} & \frac{s_{XY}}{s_Xs_Y} \\\frac{s_{XY}}{s_Ys_X} & \frac{s_Y^2}{s_Ys_Y}\end{bmatrix} \\[12pt] &= \begin{bmatrix}\frac{2.5}{\sqrt{2.5}\sqrt{2.5}} & \frac{25}{\sqrt{2.5}\sqrt{250}} \\ \frac{25}{\sqrt{250}\sqrt{2.5}} & \frac{250}{\sqrt{250}\sqrt{250}}\end{bmatrix} \\[12pt] &= \boxed{\begin{bmatrix}1.000 & 1.000 \\ 1.000 & 1.000\end{bmatrix}}\end{aligned}$

    
    \begin{tcolorbox}[breakable, size=fbox, boxrule=1pt, pad at break*=1mm,colback=cellbackground, colframe=cellborder]
\prompt{In}{incolor}{10}{\boxspacing}
\begin{Verbatim}[commandchars=\\\{\}]
\PY{c+c1}{\PYZsh{} Convert covariance to correlation manually}
\PY{n}{manual\PYZus{}corr\PYZus{}matrix} \PY{o}{=} \PY{n}{cov\PYZus{}matrix} \PY{o}{/} \PY{n}{np}\PY{o}{.}\PY{n}{outer}\PY{p}{(}\PY{n}{std}\PY{p}{,} \PY{n}{std}\PY{p}{)}

\PY{c+c1}{\PYZsh{} Verify against NumPy}
\PY{k}{assert} \PY{n}{np}\PY{o}{.}\PY{n}{allclose}\PY{p}{(}\PY{n}{manual\PYZus{}corr\PYZus{}matrix}\PY{p}{,} \PY{n}{corr\PYZus{}matrix}\PY{p}{)}
\end{Verbatim}
\end{tcolorbox}


    \subsubsection{Covariance Matrix vs Correlation
Matrix}\label{covariance-matrix-vs-correlation-matrix}

Both the \textbf{covariance matrix} and the \textbf{correlation matrix}
describe relationships between variables. However, they represent these
relationships differently.

\paragraph{1. Covariance Matrix}\label{covariance-matrix}

For two variables \(X\) and \(Y\), the covariance matrix is defined as:

\[
S =
\begin{bmatrix}
\operatorname{Var}(X) & \operatorname{Cov}(X,Y) \\
\operatorname{Cov}(Y,X) & \operatorname{Var}(Y)
\end{bmatrix}
\]

\begin{itemize}
\tightlist
\item
  \textbf{Diagonal:} Variance of each variable.
\item
  \textbf{Off-diagonal:} Covariance between different variables.
\end{itemize}

\textbf{Important property: Covariance is symmetric.}

\[
\boxed{\operatorname{Cov}(X,Y)=\operatorname{Cov}(Y,X)}
\]

This follows from the sample covariance formula:

\[
\operatorname{Cov}(X,Y)
=
\frac{1}{n-1}
\sum_{i=1}^{n}
(x_i-\bar{x})(y_i-\bar{y})
\]

Since multiplication is commutative, exchanging \(X\) and \(Y\) does not
change the result.

Therefore, the covariance matrix can also be expressed as:

\[
S =
\begin{bmatrix}
s_X^2 & s_{XY} \\
s_{XY} & s_Y^2
\end{bmatrix}
\]

And it is always \textbf{symmetric}:

\[
\boxed{S=S^T}
\]

\paragraph{2. Correlation Matrix}\label{correlation-matrix}

The correlation matrix contains the \textbf{Pearson correlation
coefficients} between variables.

For two variables \(X\) and \(Y\):

\[
R =
\begin{bmatrix}
r_{XX} & r_{XY} \\
r_{YX} & r_{YY}
\end{bmatrix}
\]

The Pearson correlation coefficient is defined as:

\[
r_{XY} =
\frac{\operatorname{Cov}(X,Y)}{s_Xs_Y}
\]

Where \(s_X\) and \(s_Y\) are the standard deviations of \(X\) and
\(Y\).

\begin{itemize}
\tightlist
\item
  \textbf{Diagonal:} Always 1 for variables with nonzero variance,
  because each variable is perfectly correlated with itself.
\item
  \textbf{Off-diagonal:} Correlation between different variables,
  ranging from -1 to +1.
\end{itemize}

For the diagonal elements:

\[
\begin{aligned}
r_{XX}
&=\frac{\operatorname{Cov}(X,X)}{s_Xs_X}\\[6pt]
&=\frac{\operatorname{Var}(X)}{s_X^2}\\[6pt]
&=\boxed{1}
\end{aligned}
\]

Similarly:

\[
r_{YY}=1
\]

\textbf{Important property: Correlation is also symmetric.}

\[
\boxed{r_{XY}=r_{YX}}
\]

Therefore, the correlation matrix can be written as:

\[
R =
\begin{bmatrix}
1 & r_{XY} \\
r_{XY} & 1
\end{bmatrix}
\]

And:

\[
\boxed{R=R^T}
\]

\textbf{Important:} If a variable has zero variance, its correlation is
undefined because its standard deviation is zero.

\paragraph{3. Relationship Between the Two
Matrices}\label{relationship-between-the-two-matrices}

The correlation matrix is obtained by standardising the covariance
matrix using the standard deviations of the variables:

\[
\boxed{
r_{XY} =
\frac{\operatorname{Cov}(X,Y)}{s_Xs_Y}
}
\]

In matrix notation:

\[
\boxed{R=D^{-1}SD^{-1}}
\]

Where:

\[
D =
\begin{bmatrix}
s_X & 0 \\
0 & s_Y
\end{bmatrix}
\]

\begin{itemize}
\tightlist
\item
  \(R\) is the correlation matrix.
\item
  \(S\) is the covariance matrix.
\item
  \(D\) is the diagonal matrix containing the standard deviations.
\item
  \(D^{-1}\) is the inverse of \(D\).
\end{itemize}

This relationship requires nonzero standard deviations.

For example, consider the covariance matrix:

\[
S =
\begin{bmatrix}
4 & 2 \\
2 & 9
\end{bmatrix}
\]

The standard deviations are:

\[
s_X=\sqrt{4}=2,
\qquad
s_Y=\sqrt{9}=3
\]

Therefore:

\[
r_{XY}=\frac{2}{2\times3}=\frac{1}{3}
\]

And the correlation matrix becomes:

\[
R =
\begin{bmatrix}
1 & \frac{1}{3} \\
\frac{1}{3} & 1
\end{bmatrix}
\]

\paragraph{4. Key Differences}\label{key-differences}

{\def\LTcaptype{none} % do not increment counter
\begin{longtable}[]{@{}
  >{\raggedright\arraybackslash}p{(\linewidth - 4\tabcolsep) * \real{0.3333}}
  >{\raggedright\arraybackslash}p{(\linewidth - 4\tabcolsep) * \real{0.3333}}
  >{\raggedright\arraybackslash}p{(\linewidth - 4\tabcolsep) * \real{0.3333}}@{}}
\toprule\noalign{}
\begin{minipage}[b]{\linewidth}\raggedright
Property
\end{minipage} & \begin{minipage}[b]{\linewidth}\raggedright
Covariance Matrix
\end{minipage} & \begin{minipage}[b]{\linewidth}\raggedright
Correlation Matrix
\end{minipage} \\
\midrule\noalign{}
\endhead
\bottomrule\noalign{}
\endlastfoot
Diagonal & Variances & Ones (1) \\
Off-diagonal & Covariances & Correlation coefficients \\
Units & Depends on the variables & Dimensionless \\
Range & Unbounded in general & Between -1 and +1 \\
Symmetric & Yes & Yes \\
Used in PCA & PCA on original feature scales & PCA on standardised
features \\
\end{longtable}
}

\textbf{Key takeaway:} The correlation matrix is the covariance matrix
of variables standardised to have zero mean and unit variance.

Both matrices are symmetric, but the correlation matrix removes
differences caused by measurement units and individual feature scales.


    \begin{tcolorbox}[breakable, size=fbox, boxrule=1pt, pad at break*=1mm,colback=cellbackground, colframe=cellborder]
\prompt{In}{incolor}{11}{\boxspacing}
\begin{Verbatim}[commandchars=\\\{\}]
\PY{k+kn}{import}\PY{+w}{ }\PY{n+nn}{numpy}\PY{+w}{ }\PY{k}{as}\PY{+w}{ }\PY{n+nn}{np}

\PY{c+c1}{\PYZsh{} S is the covariance matrix}
\PY{n}{S} \PY{o}{=} \PY{n}{np}\PY{o}{.}\PY{n}{array}\PY{p}{(}\PY{p}{[}\PY{p}{[}\PY{l+m+mi}{4}\PY{p}{,} \PY{l+m+mi}{2}\PY{p}{]}\PY{p}{,} \PY{p}{[}\PY{l+m+mi}{2}\PY{p}{,} \PY{l+m+mi}{9}\PY{p}{]}\PY{p}{]}\PY{p}{)}

\PY{c+c1}{\PYZsh{} Standard deviations from covariance diagonal}
\PY{n}{D} \PY{o}{=} \PY{n}{np}\PY{o}{.}\PY{n}{diag}\PY{p}{(}\PY{n}{np}\PY{o}{.}\PY{n}{sqrt}\PY{p}{(}\PY{n}{np}\PY{o}{.}\PY{n}{diag}\PY{p}{(}\PY{n}{S}\PY{p}{)}\PY{p}{)}\PY{p}{)}

\PY{c+c1}{\PYZsh{} diagonal elements}
\PY{n}{diag} \PY{o}{=} \PY{n}{np}\PY{o}{.}\PY{n}{diag}\PY{p}{(}\PY{n}{S}\PY{p}{)}
\PY{n+nb}{print}\PY{p}{(}\PY{l+s+sa}{f}\PY{l+s+s2}{\PYZdq{}}\PY{l+s+s2}{The diagonal elements of the covariance matrix are }\PY{l+s+si}{\PYZob{}}\PY{n}{diag}\PY{l+s+si}{\PYZcb{}}\PY{l+s+s2}{ which represent the variances.}\PY{l+s+s2}{\PYZdq{}}\PY{p}{)}
\PY{n}{diag\PYZus{}sqr\PYZus{}root} \PY{o}{=} \PY{n}{np}\PY{o}{.}\PY{n}{sqrt}\PY{p}{(}\PY{n}{diag}\PY{p}{)}
\PY{n+nb}{print}\PY{p}{(}\PY{l+s+sa}{f}\PY{l+s+s2}{\PYZdq{}}\PY{l+s+s2}{The diagonal elements of the correlation matrix are }\PY{l+s+si}{\PYZob{}}\PY{n}{diag\PYZus{}sqr\PYZus{}root}\PY{l+s+si}{\PYZcb{}}\PY{l+s+s2}{ which represent the standard deviations.}\PY{l+s+s2}{\PYZdq{}}\PY{p}{)}

\PY{c+c1}{\PYZsh{} Inverse of D}
\PY{n}{D\PYZus{}inv} \PY{o}{=} \PY{n}{np}\PY{o}{.}\PY{n}{linalg}\PY{o}{.}\PY{n}{inv}\PY{p}{(}\PY{n}{D}\PY{p}{)}

\PY{c+c1}{\PYZsh{} Correlation matrix}
\PY{n}{R} \PY{o}{=} \PY{n}{D\PYZus{}inv} \PY{o}{@} \PY{n}{S} \PY{o}{@} \PY{n}{D\PYZus{}inv}

\PY{n+nb}{print}\PY{p}{(}\PY{l+s+sa}{f}\PY{l+s+s2}{\PYZdq{}}\PY{l+s+s2}{The correlation matrix is:}\PY{l+s+se}{\PYZbs{}n}\PY{l+s+si}{\PYZob{}}\PY{n}{R}\PY{l+s+si}{\PYZcb{}}\PY{l+s+s2}{\PYZdq{}}\PY{p}{)}
\end{Verbatim}
\end{tcolorbox}

    \begin{Verbatim}[commandchars=\\\{\}]
The diagonal elements of the covariance matrix are [4 9] which represent the
variances.
The diagonal elements of the correlation matrix are [2. 3.] which represent the
standard deviations.
The correlation matrix is:
[[1.     0.3333]
 [0.3333 1.    ]]
    \end{Verbatim}


    \subsubsection{Why doesn't zero correlation prove
independence?}\label{why-doesnt-zero-correlation-prove-independence}

\textbf{Pearson correlation measures the strength and direction of a
linear relationship between two variables.}

However, \textbf{independence} means that knowing the value of one
variable provides no information about the probability distribution of
the other.

Therefore:

\begin{itemize}
\tightlist
\item
  \textbf{Zero correlation} means there is no linear association between
  the variables.
\item
  \textbf{Independence} means there is no statistical dependence between
  the variables, whether linear or nonlinear.
\end{itemize}

\paragraph{Example: A nonlinear
relationship}\label{example-a-nonlinear-relationship}

Consider two variables:

\[
X=[-2,-1,0,1,2]
\]

\[
Y=X^2=[4,1,0,1,4]
\]

Although \(Y\) is completely determined by \(X\), their Pearson
correlation is zero.

\textbf{Step 1: Calculate the means}

\[
\bar{X}=\frac{-2-1+0+1+2}{5}=0
\]

\[
\bar{Y}=\frac{4+1+0+1+4}{5}=2
\]

\textbf{Step 2: Calculate the sample covariance}

\[
\begin{aligned}
\operatorname{Cov}(X,Y)
&=\frac{\sum_{i=1}^{n}(x_i-\bar{X})(y_i-\bar{Y})}{n-1}\\[6pt]
&=\frac{(-2)(2)+(-1)(-1)+(0)(-2)+(1)(-1)+(2)(2)}{5-1}\\[6pt]
&=\frac{-4+1+0-1+4}{4}\\[6pt]
&=\boxed{0}
\end{aligned}
\]

\textbf{Step 3: Calculate Pearson correlation}

\[
r_{XY}=\frac{\operatorname{Cov}(X,Y)}{s_Xs_Y}
=\frac{0}{s_Xs_Y}
=\boxed{0}
\]

The correlation is zero because the positive and negative contributions
to the covariance cancel each other.

Nevertheless, the variables are \textbf{not independent}, because \(Y\)
is completely determined by \(X\) through the relationship \(Y=X^2\).

\paragraph{Independence vs zero
correlation}\label{independence-vs-zero-correlation}

{\def\LTcaptype{none} % do not increment counter
\begin{longtable}[]{@{}
  >{\raggedright\arraybackslash}p{(\linewidth - 2\tabcolsep) * \real{0.5000}}
  >{\raggedright\arraybackslash}p{(\linewidth - 2\tabcolsep) * \real{0.5000}}@{}}
\toprule\noalign{}
\begin{minipage}[b]{\linewidth}\raggedright
Concept
\end{minipage} & \begin{minipage}[b]{\linewidth}\raggedright
Interpretation
\end{minipage} \\
\midrule\noalign{}
\endhead
\bottomrule\noalign{}
\endlastfoot
Positive correlation & Variables tend to increase together \\
Negative correlation & One variable tends to decrease as the other
increases \\
Zero correlation & No linear association \\
Independence & Knowing one variable provides no information about the
distribution of the other \\
\end{longtable}
}

\textbf{Important mathematical property:}

\[
X\perp Y \implies \operatorname{Cov}(X,Y)=0
\]

However:

\[
\operatorname{Cov}(X,Y)=0 \not\Longrightarrow X\perp Y
\]

The first implication assumes that the covariance exists. Zero
covariance is equivalent to zero Pearson correlation when both standard
deviations are finite and non-zero.

\textbf{Exception:} For variables with a \textbf{joint Gaussian
distribution}, zero correlation does imply independence.

NB: The symbol \(\perp\) means perpendicular in geometry, but in
statistics it is commonly used to represent independence.

\paragraph{Key takeaway}\label{key-takeaway}

\textbf{Zero correlation does not necessarily mean there is no
relationship.} It only means Pearson's correlation coefficient does not
detect a linear relationship.

This is important in machine learning because features with zero or
near-zero linear correlation may still contain useful information for
nonlinear models such as decision trees and random forests.


    \begin{tcolorbox}[breakable, size=fbox, boxrule=1pt, pad at break*=1mm,colback=cellbackground, colframe=cellborder]
\prompt{In}{incolor}{12}{\boxspacing}
\begin{Verbatim}[commandchars=\\\{\}]
\PY{c+c1}{\PYZsh{} Dataset}
\PY{n}{x} \PY{o}{=} \PY{n}{np}\PY{o}{.}\PY{n}{array}\PY{p}{(}\PY{p}{[}\PY{o}{\PYZhy{}}\PY{l+m+mi}{2}\PY{p}{,} \PY{o}{\PYZhy{}}\PY{l+m+mi}{1}\PY{p}{,} \PY{l+m+mi}{0}\PY{p}{,} \PY{l+m+mi}{1}\PY{p}{,} \PY{l+m+mi}{2}\PY{p}{]}\PY{p}{,} \PY{n}{dtype}\PY{o}{=}\PY{n+nb}{float}\PY{p}{)}
\PY{n}{y} \PY{o}{=} \PY{n}{x}\PY{o}{*}\PY{o}{*}\PY{l+m+mi}{2}

\PY{n}{df} \PY{o}{=} \PY{n}{pd}\PY{o}{.}\PY{n}{DataFrame}\PY{p}{(}\PY{p}{\PYZob{}}
    \PY{l+s+s2}{\PYZdq{}}\PY{l+s+s2}{X}\PY{l+s+s2}{\PYZdq{}}\PY{p}{:} \PY{n}{x}\PY{p}{,}
    \PY{l+s+s2}{\PYZdq{}}\PY{l+s+s2}{Y = X²}\PY{l+s+s2}{\PYZdq{}}\PY{p}{:} \PY{n}{y}
\PY{p}{\PYZcb{}}\PY{p}{)}

\PY{n}{display}\PY{p}{(}\PY{n}{df}\PY{p}{)}

\PY{n+nb}{print}\PY{p}{(}\PY{l+s+s2}{\PYZdq{}}\PY{l+s+s2}{Sample covariance:}\PY{l+s+s2}{\PYZdq{}}\PY{p}{,} \PY{n}{np}\PY{o}{.}\PY{n}{cov}\PY{p}{(}\PY{n}{x}\PY{p}{,} \PY{n}{y}\PY{p}{)}\PY{p}{[}\PY{l+m+mi}{0}\PY{p}{,} \PY{l+m+mi}{1}\PY{p}{]}\PY{p}{)}
\PY{n+nb}{print}\PY{p}{(}\PY{l+s+s2}{\PYZdq{}}\PY{l+s+s2}{Pearson correlation:}\PY{l+s+s2}{\PYZdq{}}\PY{p}{,} \PY{n}{np}\PY{o}{.}\PY{n}{corrcoef}\PY{p}{(}\PY{n}{x}\PY{p}{,} \PY{n}{y}\PY{p}{)}\PY{p}{[}\PY{l+m+mi}{0}\PY{p}{,} \PY{l+m+mi}{1}\PY{p}{]}\PY{p}{)}
\end{Verbatim}
\end{tcolorbox}

    
    \begin{Verbatim}[commandchars=\\\{\}]
     X  Y = X²
0 -2.0     4.0
1 -1.0     1.0
2  0.0     0.0
3  1.0     1.0
4  2.0     4.0
    \end{Verbatim}

    
    \begin{Verbatim}[commandchars=\\\{\}]
Sample covariance: 0.0
Pearson correlation: 0.0
    \end{Verbatim}

    \begin{tcolorbox}[breakable, size=fbox, boxrule=1pt, pad at break*=1mm,colback=cellbackground, colframe=cellborder]
\prompt{In}{incolor}{13}{\boxspacing}
\begin{Verbatim}[commandchars=\\\{\}]
\PY{c+c1}{\PYZsh{} Display DataFrame}
\PY{c+c1}{\PYZsh{}display(df.style.format(\PYZdq{}\PYZob{}:.0f\PYZcb{}\PYZdq{}))}

\PY{k}{def}\PY{+w}{ }\PY{n+nf}{display\PYZus{}dataframe}\PY{p}{(}\PY{n}{df}\PY{p}{,} \PY{n}{precision}\PY{o}{=}\PY{l+m+mi}{0}\PY{p}{,} \PY{n}{show\PYZus{}index}\PY{o}{=}\PY{k+kc}{True}\PY{p}{)}\PY{p}{:}
\PY{+w}{    }\PY{l+s+sd}{\PYZdq{}\PYZdq{}\PYZdq{}}
\PY{l+s+sd}{    Display a DataFrame in Jupyter and export it to LaTeX.}

\PY{l+s+sd}{    Jupyter uses the HTML representation.}
\PY{l+s+sd}{    nbconvert uses the LaTeX representation.}

\PY{l+s+sd}{    Parameters}
\PY{l+s+sd}{    \PYZhy{}\PYZhy{}\PYZhy{}\PYZhy{}\PYZhy{}\PYZhy{}\PYZhy{}\PYZhy{}\PYZhy{}\PYZhy{}}
\PY{l+s+sd}{    df : pd.DataFrame}
\PY{l+s+sd}{        DataFrame to display.}
\PY{l+s+sd}{    precision : int}
\PY{l+s+sd}{        Number of decimal places (default: 0).}
\PY{l+s+sd}{    show\PYZus{}index : bool}
\PY{l+s+sd}{        Whether to include the DataFrame index.}
\PY{l+s+sd}{    \PYZdq{}\PYZdq{}\PYZdq{}}

    \PY{n}{number\PYZus{}format} \PY{o}{=} \PY{l+s+sa}{f}\PY{l+s+s2}{\PYZdq{}}\PY{l+s+se}{\PYZob{}\PYZob{}}\PY{l+s+s2}{:.}\PY{l+s+si}{\PYZob{}}\PY{n}{precision}\PY{l+s+si}{\PYZcb{}}\PY{l+s+s2}{f}\PY{l+s+se}{\PYZcb{}\PYZcb{}}\PY{l+s+s2}{\PYZdq{}}

    \PY{c+c1}{\PYZsh{} HTML representation for Jupyter}
    \PY{n}{html\PYZus{}output} \PY{o}{=} \PY{p}{(}
        \PY{n}{df}\PY{o}{.}\PY{n}{style}
        \PY{o}{.}\PY{n}{format}\PY{p}{(}\PY{n}{number\PYZus{}format}\PY{p}{)}
        \PY{o}{.}\PY{n}{hide}\PY{p}{(}\PY{n}{axis}\PY{o}{=}\PY{l+s+s2}{\PYZdq{}}\PY{l+s+s2}{index}\PY{l+s+s2}{\PYZdq{}}\PY{p}{)} \PY{k}{if} \PY{o+ow}{not} \PY{n}{show\PYZus{}index}
        \PY{k}{else} \PY{n}{df}\PY{o}{.}\PY{n}{style}\PY{o}{.}\PY{n}{format}\PY{p}{(}\PY{n}{number\PYZus{}format}\PY{p}{)}
    \PY{p}{)}\PY{o}{.}\PY{n}{to\PYZus{}html}\PY{p}{(}\PY{p}{)}

    \PY{c+c1}{\PYZsh{} LaTeX representation for PDF export}
    \PY{n}{latex\PYZus{}output} \PY{o}{=} \PY{n}{df}\PY{o}{.}\PY{n}{to\PYZus{}latex}\PY{p}{(}
        \PY{n}{index}\PY{o}{=}\PY{n}{show\PYZus{}index}\PY{p}{,}
        \PY{n}{float\PYZus{}format}\PY{o}{=}\PY{l+s+sa}{f}\PY{l+s+s2}{\PYZdq{}}\PY{l+s+s2}{\PYZpc{}.}\PY{l+s+si}{\PYZob{}}\PY{n}{precision}\PY{l+s+si}{\PYZcb{}}\PY{l+s+s2}{f}\PY{l+s+s2}{\PYZdq{}}
    \PY{p}{)}

    \PY{c+c1}{\PYZsh{} One output containing both representations}
    \PY{n}{display}\PY{p}{(}\PY{p}{\PYZob{}}
        \PY{l+s+s2}{\PYZdq{}}\PY{l+s+s2}{text/html}\PY{l+s+s2}{\PYZdq{}}\PY{p}{:} \PY{n}{html\PYZus{}output}\PY{p}{,}
        \PY{l+s+s2}{\PYZdq{}}\PY{l+s+s2}{text/latex}\PY{l+s+s2}{\PYZdq{}}\PY{p}{:} \PY{n}{latex\PYZus{}output}
    \PY{p}{\PYZcb{}}\PY{p}{,} \PY{n}{raw}\PY{o}{=}\PY{k+kc}{True}\PY{p}{)}

\PY{n}{display\PYZus{}dataframe}\PY{p}{(}\PY{n}{df}\PY{p}{)}

\PY{c+c1}{\PYZsh{} Calculate covariance and correlation}
\PY{n}{covariance} \PY{o}{=} \PY{n}{np}\PY{o}{.}\PY{n}{cov}\PY{p}{(}\PY{n}{x}\PY{p}{,} \PY{n}{y}\PY{p}{)}\PY{p}{[}\PY{l+m+mi}{0}\PY{p}{,} \PY{l+m+mi}{1}\PY{p}{]}
\PY{n}{correlation} \PY{o}{=} \PY{n}{np}\PY{o}{.}\PY{n}{corrcoef}\PY{p}{(}\PY{n}{x}\PY{p}{,} \PY{n}{y}\PY{p}{)}\PY{p}{[}\PY{l+m+mi}{0}\PY{p}{,} \PY{l+m+mi}{1}\PY{p}{]}

\PY{c+c1}{\PYZsh{} Display mathematical results}
\PY{n}{display}\PY{p}{(}\PY{n}{Math}\PY{p}{(}
    \PY{l+s+sa}{rf}\PY{l+s+s2}{\PYZdq{}}\PY{l+s+s2}{\PYZbs{}}\PY{l+s+s2}{operatorname}\PY{l+s+se}{\PYZob{}\PYZob{}}\PY{l+s+s2}{Cov}\PY{l+s+se}{\PYZcb{}\PYZcb{}}\PY{l+s+s2}{(X,Y) = }\PY{l+s+s2}{\PYZbs{}}\PY{l+s+s2}{boxed}\PY{l+s+se}{\PYZob{}\PYZob{}}\PY{l+s+si}{\PYZob{}}\PY{n}{covariance}\PY{l+s+si}{:}\PY{l+s+s2}{.3f}\PY{l+s+se}{\PYZcb{}\PYZcb{}}\PY{l+s+si}{\PYZcb{}}\PY{l+s+s2}{\PYZdq{}}
\PY{p}{)}\PY{p}{)}

\PY{n}{display}\PY{p}{(}\PY{n}{Math}\PY{p}{(}
    \PY{l+s+sa}{rf}\PY{l+s+s2}{\PYZdq{}}\PY{l+s+s2}{r\PYZus{}}\PY{l+s+se}{\PYZob{}\PYZob{}}\PY{l+s+s2}{XY}\PY{l+s+se}{\PYZcb{}\PYZcb{}}\PY{l+s+s2}{ = }\PY{l+s+s2}{\PYZbs{}}\PY{l+s+s2}{boxed}\PY{l+s+se}{\PYZob{}\PYZob{}}\PY{l+s+si}{\PYZob{}}\PY{n}{correlation}\PY{l+s+si}{:}\PY{l+s+s2}{.3f}\PY{l+s+se}{\PYZcb{}\PYZcb{}}\PY{l+s+si}{\PYZcb{}}\PY{l+s+s2}{\PYZdq{}}
\PY{p}{)}\PY{p}{)}

\PY{c+c1}{\PYZsh{} Generate a smooth parabola}
\PY{n}{x\PYZus{}curve} \PY{o}{=} \PY{n}{np}\PY{o}{.}\PY{n}{linspace}\PY{p}{(}\PY{n}{x}\PY{o}{.}\PY{n}{min}\PY{p}{(}\PY{p}{)}\PY{p}{,} \PY{n}{x}\PY{o}{.}\PY{n}{max}\PY{p}{(}\PY{p}{)}\PY{p}{,} \PY{l+m+mi}{200}\PY{p}{)}
\PY{n}{y\PYZus{}curve} \PY{o}{=} \PY{n}{x\PYZus{}curve}\PY{o}{*}\PY{o}{*}\PY{l+m+mi}{2}

\PY{c+c1}{\PYZsh{} Plot}
\PY{n}{fig}\PY{p}{,} \PY{n}{ax} \PY{o}{=} \PY{n}{plt}\PY{o}{.}\PY{n}{subplots}\PY{p}{(}\PY{n}{figsize}\PY{o}{=}\PY{p}{(}\PY{l+m+mi}{8}\PY{p}{,} \PY{l+m+mi}{5}\PY{p}{)}\PY{p}{)}

\PY{n}{ax}\PY{o}{.}\PY{n}{scatter}\PY{p}{(}
    \PY{n}{x}\PY{p}{,} \PY{n}{y}\PY{p}{,}
    \PY{n}{s}\PY{o}{=}\PY{l+m+mi}{90}\PY{p}{,}
    \PY{n}{color}\PY{o}{=}\PY{l+s+s2}{\PYZdq{}}\PY{l+s+s2}{teal}\PY{l+s+s2}{\PYZdq{}}\PY{p}{,}
    \PY{n}{edgecolor}\PY{o}{=}\PY{l+s+s2}{\PYZdq{}}\PY{l+s+s2}{black}\PY{l+s+s2}{\PYZdq{}}\PY{p}{,}
    \PY{n}{zorder}\PY{o}{=}\PY{l+m+mi}{3}\PY{p}{,}
    \PY{n}{label}\PY{o}{=}\PY{l+s+s2}{\PYZdq{}}\PY{l+s+s2}{Observations}\PY{l+s+s2}{\PYZdq{}}
\PY{p}{)}

\PY{n}{ax}\PY{o}{.}\PY{n}{plot}\PY{p}{(}
    \PY{n}{x\PYZus{}curve}\PY{p}{,} \PY{n}{y\PYZus{}curve}\PY{p}{,}
    \PY{n}{color}\PY{o}{=}\PY{l+s+s2}{\PYZdq{}}\PY{l+s+s2}{darkorange}\PY{l+s+s2}{\PYZdq{}}\PY{p}{,}
    \PY{n}{linewidth}\PY{o}{=}\PY{l+m+mi}{2}\PY{p}{,}
    \PY{n}{label}\PY{o}{=}\PY{l+s+sa}{r}\PY{l+s+s2}{\PYZdq{}}\PY{l+s+s2}{\PYZdl{}Y = X\PYZca{}2\PYZdl{}}\PY{l+s+s2}{\PYZdq{}}
\PY{p}{)}

\PY{n}{ax}\PY{o}{.}\PY{n}{set}\PY{p}{(}
    \PY{n}{title}\PY{o}{=}\PY{l+s+sa}{f}\PY{l+s+s2}{\PYZdq{}}\PY{l+s+s2}{Nonlinear Dependence Despite Zero Correlation (r = }\PY{l+s+si}{\PYZob{}}\PY{n}{correlation}\PY{l+s+si}{:}\PY{l+s+s2}{.2f}\PY{l+s+si}{\PYZcb{}}\PY{l+s+s2}{)}\PY{l+s+s2}{\PYZdq{}}\PY{p}{,}
    \PY{n}{xlabel}\PY{o}{=}\PY{l+s+s2}{\PYZdq{}}\PY{l+s+s2}{X}\PY{l+s+s2}{\PYZdq{}}\PY{p}{,}
    \PY{n}{ylabel}\PY{o}{=}\PY{l+s+s2}{\PYZdq{}}\PY{l+s+s2}{Y}\PY{l+s+s2}{\PYZdq{}}
\PY{p}{)}

\PY{n}{ax}\PY{o}{.}\PY{n}{set\PYZus{}xticks}\PY{p}{(}\PY{n}{x}\PY{p}{)}
\PY{n}{ax}\PY{o}{.}\PY{n}{grid}\PY{p}{(}\PY{k+kc}{True}\PY{p}{,} \PY{n}{alpha}\PY{o}{=}\PY{l+m+mf}{0.3}\PY{p}{)}
\PY{n}{ax}\PY{o}{.}\PY{n}{legend}\PY{p}{(}\PY{p}{)}

\PY{n}{plt}\PY{o}{.}\PY{n}{show}\PY{p}{(}\PY{p}{)}
\end{Verbatim}
\end{tcolorbox}

    \begin{tabular}{lrr}
\toprule
 & X & Y = X² \\
\midrule
0 & -2 & 4 \\
1 & -1 & 1 \\
2 & 0 & 0 \\
3 & 1 & 1 \\
4 & 2 & 4 \\
\bottomrule
\end{tabular}


    
    $\displaystyle \operatorname{Cov}(X,Y) = \boxed{0.000}$

    
    $\displaystyle r_{XY} = \boxed{0.000}$

    
    \begin{center}
    \adjustimage{max size={0.9\linewidth}{0.9\paperheight}}{week4_pdf_build_pca_files/week4_pdf_build_pca_19_3.png}
    \end{center}
    { \hspace*{\fill} \\}
    

    \subsection{5. Eigenvectors and eigenvalues: directions and
scaling}\label{eigenvectors-and-eigenvalues-directions-and-scaling}

For a square matrix \(\mathbf A\), an \textbf{eigenvector}
\(\mathbf v\ne \mathbf 0\) and \textbf{eigenvalue} \(\lambda\) satisfy

\[\boxed{\mathbf A\mathbf v=\lambda\mathbf v}.\]

In words: transforming \(\mathbf v\) by \(\mathbf A\) changes its
\textbf{length} and possibly its sign, but not its underlying line of
direction.

\begin{itemize}
\tightlist
\item
  \(\lambda>1\): scaling/expansion along that eigenvector.
\item
  \(0<\lambda<1\): shrinking along that eigenvector.
\item
  \(\lambda<0\): the direction is also reversed.
\item
  \(\lambda=0\): the vector is mapped to zero.
\end{itemize}

For example,
\(\mathbf A=\left[\begin{smallmatrix}3&0\\0&1\end{smallmatrix}\right]\)
stretches the horizontal axis by three while leaving the vertical axis
unchanged.


    \begin{tcolorbox}[breakable, size=fbox, boxrule=1pt, pad at break*=1mm,colback=cellbackground, colframe=cellborder]
\prompt{In}{incolor}{14}{\boxspacing}
\begin{Verbatim}[commandchars=\\\{\}]
\PY{n}{A} \PY{o}{=} \PY{n}{np}\PY{o}{.}\PY{n}{array}\PY{p}{(}\PY{p}{[}\PY{p}{[}\PY{l+m+mf}{3.}\PY{p}{,} \PY{l+m+mf}{0.}\PY{p}{]}\PY{p}{,} \PY{p}{[}\PY{l+m+mf}{0.}\PY{p}{,} \PY{l+m+mf}{1.}\PY{p}{]}\PY{p}{]}\PY{p}{)}
\PY{n}{eigenvalues}\PY{p}{,} \PY{n}{eigenvectors} \PY{o}{=} \PY{n}{np}\PY{o}{.}\PY{n}{linalg}\PY{o}{.}\PY{n}{eig}\PY{p}{(}\PY{n}{A}\PY{p}{)}
\PY{n+nb}{print}\PY{p}{(}\PY{l+s+s2}{\PYZdq{}}\PY{l+s+s2}{Matrix A:}\PY{l+s+se}{\PYZbs{}n}\PY{l+s+s2}{\PYZdq{}}\PY{p}{,} \PY{n}{A}\PY{p}{)}
\PY{n+nb}{print}\PY{p}{(}\PY{l+s+s2}{\PYZdq{}}\PY{l+s+s2}{Eigenvalues:}\PY{l+s+s2}{\PYZdq{}}\PY{p}{,} \PY{n}{eigenvalues}\PY{p}{)}
\PY{n+nb}{print}\PY{p}{(}\PY{l+s+s2}{\PYZdq{}}\PY{l+s+s2}{Eigenvectors (COLUMNS):}\PY{l+s+se}{\PYZbs{}n}\PY{l+s+s2}{\PYZdq{}}\PY{p}{,} \PY{n}{eigenvectors}\PY{p}{)}
\PY{k}{for} \PY{n}{i} \PY{o+ow}{in} \PY{n+nb}{range}\PY{p}{(}\PY{l+m+mi}{2}\PY{p}{)}\PY{p}{:}
    \PY{n}{v} \PY{o}{=} \PY{n}{eigenvectors}\PY{p}{[}\PY{p}{:}\PY{p}{,} \PY{n}{i}\PY{p}{]}
    \PY{n}{lam} \PY{o}{=} \PY{n}{eigenvalues}\PY{p}{[}\PY{n}{i}\PY{p}{]}
    \PY{n+nb}{print}\PY{p}{(}\PY{l+s+sa}{f}\PY{l+s+s2}{\PYZdq{}}\PY{l+s+s2}{Check eigenpair }\PY{l+s+si}{\PYZob{}}\PY{n}{i}\PY{o}{+}\PY{l+m+mi}{1}\PY{l+s+si}{\PYZcb{}}\PY{l+s+s2}{: A @ v = }\PY{l+s+si}{\PYZob{}}\PY{n}{A}\PY{+w}{ }\PY{o}{@}\PY{+w}{ }\PY{n}{v}\PY{l+s+si}{\PYZcb{}}\PY{l+s+s2}{, lambda * v = }\PY{l+s+si}{\PYZob{}}\PY{n}{lam}\PY{+w}{ }\PY{o}{*}\PY{+w}{ }\PY{n}{v}\PY{l+s+si}{\PYZcb{}}\PY{l+s+s2}{\PYZdq{}}\PY{p}{)}
    \PY{k}{assert} \PY{n}{np}\PY{o}{.}\PY{n}{allclose}\PY{p}{(}\PY{n}{A} \PY{o}{@} \PY{n}{v}\PY{p}{,} \PY{n}{lam}\PY{o}{*}\PY{n}{v}\PY{p}{)}
\end{Verbatim}
\end{tcolorbox}

    \begin{Verbatim}[commandchars=\\\{\}]
Matrix A:
 [[3. 0.]
 [0. 1.]]
Eigenvalues: [3. 1.]
Eigenvectors (COLUMNS):
 [[1. 0.]
 [0. 1.]]
Check eigenpair 1: A @ v = [3. 0.], lambda * v = [3. 0.]
Check eigenpair 2: A @ v = [0. 1.], lambda * v = [0. 1.]
    \end{Verbatim}


    \subsubsection{5.1. What happens when we
multiply?}\label{what-happens-when-we-multiply}

The fundamental eigenvalue equation is:

\[
Av=\lambda v
\]

For the first eigenvector:

\[
\begin{aligned}
A\mathbf{v}_1
&=\begin{bmatrix}3&0\\0&1\end{bmatrix}
\begin{bmatrix}1\\0\end{bmatrix}
=\begin{bmatrix}3\\0\end{bmatrix}
=3\begin{bmatrix}1\\0\end{bmatrix}
\end{aligned}
\]

For the second:

\[
\begin{aligned}
A\mathbf{v}_2
&=\begin{bmatrix}3&0\\0&1\end{bmatrix}
\begin{bmatrix}0\\1\end{bmatrix}
=\begin{bmatrix}0\\1\end{bmatrix}
=1\begin{bmatrix}0\\1\end{bmatrix}
\end{aligned}
\]

\subsubsection{5.2. Visual interpretation}\label{visual-interpretation}


    \begin{tcolorbox}[breakable, size=fbox, boxrule=1pt, pad at break*=1mm,colback=cellbackground, colframe=cellborder]
\prompt{In}{incolor}{15}{\boxspacing}
\begin{Verbatim}[commandchars=\\\{\}]
\PY{n}{fig}\PY{p}{,} \PY{n}{ax} \PY{o}{=} \PY{n}{plt}\PY{o}{.}\PY{n}{subplots}\PY{p}{(}\PY{n}{figsize}\PY{o}{=}\PY{p}{(}\PY{l+m+mi}{6}\PY{p}{,} \PY{l+m+mi}{5}\PY{p}{)}\PY{p}{)}

\PY{n}{original\PYZus{}vectors} \PY{o}{=} \PY{p}{[}
    \PY{n}{np}\PY{o}{.}\PY{n}{array}\PY{p}{(}\PY{p}{[}\PY{l+m+mf}{1.}\PY{p}{,} \PY{l+m+mf}{0.}\PY{p}{]}\PY{p}{)}\PY{p}{,}   \PY{c+c1}{\PYZsh{} Eigenvector 1}
    \PY{n}{np}\PY{o}{.}\PY{n}{array}\PY{p}{(}\PY{p}{[}\PY{l+m+mf}{0.}\PY{p}{,} \PY{l+m+mf}{1.}\PY{p}{]}\PY{p}{)}\PY{p}{,}   \PY{c+c1}{\PYZsh{} Eigenvector 2}
    \PY{n}{np}\PY{o}{.}\PY{n}{array}\PY{p}{(}\PY{p}{[}\PY{l+m+mf}{.7}\PY{p}{,} \PY{l+m+mf}{.7}\PY{p}{]}\PY{p}{)}    \PY{c+c1}{\PYZsh{} Not an eigenvector}
\PY{p}{]}

\PY{k}{for} \PY{n}{v} \PY{o+ow}{in} \PY{n}{original\PYZus{}vectors}\PY{p}{:}
    \PY{n}{transformed} \PY{o}{=} \PY{n}{A} \PY{o}{@} \PY{n}{v}
    \PY{n}{ax}\PY{o}{.}\PY{n}{quiver}\PY{p}{(}\PY{l+m+mi}{0}\PY{p}{,} \PY{l+m+mi}{0}\PY{p}{,} \PY{o}{*}\PY{n}{v}\PY{p}{,} \PY{n}{angles}\PY{o}{=}\PY{l+s+s1}{\PYZsq{}}\PY{l+s+s1}{xy}\PY{l+s+s1}{\PYZsq{}}\PY{p}{,} \PY{n}{scale\PYZus{}units}\PY{o}{=}\PY{l+s+s1}{\PYZsq{}}\PY{l+s+s1}{xy}\PY{l+s+s1}{\PYZsq{}}\PY{p}{,} \PY{n}{scale}\PY{o}{=}\PY{l+m+mi}{1}\PY{p}{,} \PY{n}{alpha}\PY{o}{=}\PY{l+m+mf}{.65}\PY{p}{)}
    \PY{n}{ax}\PY{o}{.}\PY{n}{quiver}\PY{p}{(}\PY{l+m+mi}{0}\PY{p}{,} \PY{l+m+mi}{0}\PY{p}{,} \PY{o}{*}\PY{n}{transformed}\PY{p}{,} \PY{n}{angles}\PY{o}{=}\PY{l+s+s1}{\PYZsq{}}\PY{l+s+s1}{xy}\PY{l+s+s1}{\PYZsq{}}\PY{p}{,} \PY{n}{scale\PYZus{}units}\PY{o}{=}\PY{l+s+s1}{\PYZsq{}}\PY{l+s+s1}{xy}\PY{l+s+s1}{\PYZsq{}}\PY{p}{,} \PY{n}{scale}\PY{o}{=}\PY{l+m+mi}{1}\PY{p}{,} \PY{n}{alpha}\PY{o}{=}\PY{l+m+mf}{.85}\PY{p}{)}
\PY{n}{ax}\PY{o}{.}\PY{n}{set}\PY{p}{(}\PY{n}{xlim}\PY{o}{=}\PY{p}{(}\PY{o}{\PYZhy{}}\PY{l+m+mf}{.2}\PY{p}{,} \PY{l+m+mf}{3.4}\PY{p}{)}\PY{p}{,} \PY{n}{ylim}\PY{o}{=}\PY{p}{(}\PY{o}{\PYZhy{}}\PY{l+m+mf}{.2}\PY{p}{,} \PY{l+m+mf}{1.5}\PY{p}{)}\PY{p}{,} \PY{n}{xlabel}\PY{o}{=}\PY{l+s+s2}{\PYZdq{}}\PY{l+s+s2}{x}\PY{l+s+s2}{\PYZdq{}}\PY{p}{,} \PY{n}{ylabel}\PY{o}{=}\PY{l+s+s2}{\PYZdq{}}\PY{l+s+s2}{y}\PY{l+s+s2}{\PYZdq{}}\PY{p}{,}
       \PY{n}{title}\PY{o}{=}\PY{l+s+s2}{\PYZdq{}}\PY{l+s+s2}{Original vectors and their transformed versions}\PY{l+s+s2}{\PYZdq{}}\PY{p}{)}
\PY{n}{ax}\PY{o}{.}\PY{n}{set\PYZus{}aspect}\PY{p}{(}\PY{l+s+s1}{\PYZsq{}}\PY{l+s+s1}{equal}\PY{l+s+s1}{\PYZsq{}}\PY{p}{)}
\PY{n}{ax}\PY{o}{.}\PY{n}{grid}\PY{p}{(}\PY{n}{alpha}\PY{o}{=}\PY{l+m+mf}{.2}\PY{p}{)}
\PY{n}{ax}\PY{o}{.}\PY{n}{text}\PY{p}{(}\PY{l+m+mf}{.45}\PY{p}{,} \PY{l+m+mf}{1.35}\PY{p}{,} \PY{l+s+s2}{\PYZdq{}}\PY{l+s+s2}{Each pair: shorter original → transformed}\PY{l+s+s2}{\PYZdq{}}\PY{p}{,} \PY{n}{fontsize}\PY{o}{=}\PY{l+m+mi}{9}\PY{p}{)}
\PY{n}{plt}\PY{o}{.}\PY{n}{show}\PY{p}{(}\PY{p}{)}
\end{Verbatim}
\end{tcolorbox}

    \begin{center}
    \adjustimage{max size={0.9\linewidth}{0.9\paperheight}}{week4_pdf_build_pca_files/week4_pdf_build_pca_23_0.png}
    \end{center}
    { \hspace*{\fill} \\}
    
    \begin{tcolorbox}[breakable, size=fbox, boxrule=1pt, pad at break*=1mm,colback=cellbackground, colframe=cellborder]
\prompt{In}{incolor}{16}{\boxspacing}
\begin{Verbatim}[commandchars=\\\{\}]
\PY{n}{colors} \PY{o}{=} \PY{p}{[}\PY{l+s+s2}{\PYZdq{}}\PY{l+s+s2}{royalblue}\PY{l+s+s2}{\PYZdq{}}\PY{p}{,} \PY{l+s+s2}{\PYZdq{}}\PY{l+s+s2}{seagreen}\PY{l+s+s2}{\PYZdq{}}\PY{p}{,} \PY{l+s+s2}{\PYZdq{}}\PY{l+s+s2}{darkorange}\PY{l+s+s2}{\PYZdq{}}\PY{p}{]}

\PY{n}{labels} \PY{o}{=} \PY{p}{[}
    \PY{l+s+sa}{r}\PY{l+s+s2}{\PYZdq{}}\PY{l+s+s2}{\PYZdl{}v\PYZus{}1\PYZdl{} (\PYZdl{}}\PY{l+s+s2}{\PYZbs{}}\PY{l+s+s2}{lambda\PYZus{}1=3\PYZdl{})}\PY{l+s+s2}{\PYZdq{}}\PY{p}{,}
    \PY{l+s+sa}{r}\PY{l+s+s2}{\PYZdq{}}\PY{l+s+s2}{\PYZdl{}v\PYZus{}2\PYZdl{} (\PYZdl{}}\PY{l+s+s2}{\PYZbs{}}\PY{l+s+s2}{lambda\PYZus{}2=1\PYZdl{})}\PY{l+s+s2}{\PYZdq{}}\PY{p}{,}
    \PY{l+s+sa}{r}\PY{l+s+s2}{\PYZdq{}}\PY{l+s+s2}{\PYZdl{}v\PYZus{}3\PYZdl{} (not an eigenvector)}\PY{l+s+s2}{\PYZdq{}}
\PY{p}{]}

\PY{n}{fig}\PY{p}{,} \PY{n}{ax} \PY{o}{=} \PY{n}{plt}\PY{o}{.}\PY{n}{subplots}\PY{p}{(}\PY{n}{figsize}\PY{o}{=}\PY{p}{(}\PY{l+m+mi}{9}\PY{p}{,} \PY{l+m+mi}{6}\PY{p}{)}\PY{p}{)}

\PY{k}{for} \PY{n}{i}\PY{p}{,} \PY{n}{v} \PY{o+ow}{in} \PY{n+nb}{enumerate}\PY{p}{(}\PY{n}{original\PYZus{}vectors}\PY{p}{)}\PY{p}{:}
    \PY{n}{transformed} \PY{o}{=} \PY{n}{A} \PY{o}{@} \PY{n}{v}
    \PY{n}{color} \PY{o}{=} \PY{n}{colors}\PY{p}{[}\PY{n}{i}\PY{p}{]}

    \PY{c+c1}{\PYZsh{} Original vector: dashed}
    \PY{n}{ax}\PY{o}{.}\PY{n}{quiver}\PY{p}{(}
        \PY{l+m+mi}{0}\PY{p}{,} \PY{l+m+mi}{0}\PY{p}{,} \PY{o}{*}\PY{n}{v}\PY{p}{,}
        \PY{n}{angles}\PY{o}{=}\PY{l+s+s2}{\PYZdq{}}\PY{l+s+s2}{xy}\PY{l+s+s2}{\PYZdq{}}\PY{p}{,} \PY{n}{scale\PYZus{}units}\PY{o}{=}\PY{l+s+s2}{\PYZdq{}}\PY{l+s+s2}{xy}\PY{l+s+s2}{\PYZdq{}}\PY{p}{,} \PY{n}{scale}\PY{o}{=}\PY{l+m+mi}{1}\PY{p}{,}
        \PY{n}{color}\PY{o}{=}\PY{n}{color}\PY{p}{,} \PY{n}{alpha}\PY{o}{=}\PY{l+m+mf}{0.45}\PY{p}{,}
        \PY{n}{linestyle}\PY{o}{=}\PY{l+s+s2}{\PYZdq{}}\PY{l+s+s2}{\PYZhy{}\PYZhy{}}\PY{l+s+s2}{\PYZdq{}}\PY{p}{,} \PY{n}{width}\PY{o}{=}\PY{l+m+mf}{0.006}
    \PY{p}{)}

    \PY{c+c1}{\PYZsh{} Transformed vector: solid}
    \PY{n}{ax}\PY{o}{.}\PY{n}{quiver}\PY{p}{(}
        \PY{l+m+mi}{0}\PY{p}{,} \PY{l+m+mi}{0}\PY{p}{,} \PY{o}{*}\PY{n}{transformed}\PY{p}{,}
        \PY{n}{angles}\PY{o}{=}\PY{l+s+s2}{\PYZdq{}}\PY{l+s+s2}{xy}\PY{l+s+s2}{\PYZdq{}}\PY{p}{,} \PY{n}{scale\PYZus{}units}\PY{o}{=}\PY{l+s+s2}{\PYZdq{}}\PY{l+s+s2}{xy}\PY{l+s+s2}{\PYZdq{}}\PY{p}{,} \PY{n}{scale}\PY{o}{=}\PY{l+m+mi}{1}\PY{p}{,}
        \PY{n}{color}\PY{o}{=}\PY{n}{color}\PY{p}{,} \PY{n}{alpha}\PY{o}{=}\PY{l+m+mf}{0.95}\PY{p}{,} \PY{n}{width}\PY{o}{=}\PY{l+m+mf}{0.008}
    \PY{p}{)}

    \PY{c+c1}{\PYZsh{} Original label}
    \PY{n}{ax}\PY{o}{.}\PY{n}{annotate}\PY{p}{(}
        \PY{l+s+sa}{f}\PY{l+s+s2}{\PYZdq{}}\PY{l+s+s2}{Original }\PY{l+s+si}{\PYZob{}}\PY{n}{labels}\PY{p}{[}\PY{n}{i}\PY{p}{]}\PY{l+s+si}{\PYZcb{}}\PY{l+s+s2}{\PYZdq{}}\PY{p}{,}
        \PY{n}{xy}\PY{o}{=}\PY{n}{v}\PY{p}{,}
        \PY{n}{xytext}\PY{o}{=}\PY{p}{(}\PY{l+m+mi}{5}\PY{p}{,} \PY{l+m+mi}{10} \PY{o}{+} \PY{n}{i} \PY{o}{*} \PY{l+m+mi}{5}\PY{p}{)}\PY{p}{,}
        \PY{n}{textcoords}\PY{o}{=}\PY{l+s+s2}{\PYZdq{}}\PY{l+s+s2}{offset points}\PY{l+s+s2}{\PYZdq{}}\PY{p}{,}
        \PY{n}{color}\PY{o}{=}\PY{n}{color}\PY{p}{,}
        \PY{n}{fontsize}\PY{o}{=}\PY{l+m+mi}{10}
    \PY{p}{)}

    \PY{c+c1}{\PYZsh{} Transformed label}
    \PY{n}{ax}\PY{o}{.}\PY{n}{annotate}\PY{p}{(}
        \PY{l+s+sa}{rf}\PY{l+s+s2}{\PYZdq{}}\PY{l+s+s2}{\PYZdl{}Av\PYZus{}}\PY{l+s+si}{\PYZob{}}\PY{n}{i}\PY{o}{+}\PY{l+m+mi}{1}\PY{l+s+si}{\PYZcb{}}\PY{l+s+s2}{\PYZdl{}}\PY{l+s+s2}{\PYZdq{}}\PY{p}{,}
        \PY{n}{xy}\PY{o}{=}\PY{n}{transformed}\PY{p}{,}
        \PY{n}{xytext}\PY{o}{=}\PY{p}{(}\PY{l+m+mi}{7}\PY{p}{,} \PY{o}{\PYZhy{}}\PY{l+m+mi}{18} \PY{o}{\PYZhy{}} \PY{n}{i} \PY{o}{*} \PY{l+m+mi}{3}\PY{p}{)}\PY{p}{,}
        \PY{n}{textcoords}\PY{o}{=}\PY{l+s+s2}{\PYZdq{}}\PY{l+s+s2}{offset points}\PY{l+s+s2}{\PYZdq{}}\PY{p}{,}
        \PY{n}{color}\PY{o}{=}\PY{n}{color}\PY{p}{,}
        \PY{n}{fontsize}\PY{o}{=}\PY{l+m+mi}{11}\PY{p}{,}
        \PY{n}{fontweight}\PY{o}{=}\PY{l+s+s2}{\PYZdq{}}\PY{l+s+s2}{bold}\PY{l+s+s2}{\PYZdq{}}
    \PY{p}{)}

\PY{c+c1}{\PYZsh{} Custom legend}
\PY{n}{legend\PYZus{}elements} \PY{o}{=} \PY{p}{[}
    \PY{n}{Line2D}\PY{p}{(}\PY{p}{[}\PY{l+m+mi}{0}\PY{p}{]}\PY{p}{,} \PY{p}{[}\PY{l+m+mi}{0}\PY{p}{]}\PY{p}{,} \PY{n}{color}\PY{o}{=}\PY{l+s+s2}{\PYZdq{}}\PY{l+s+s2}{royalblue}\PY{l+s+s2}{\PYZdq{}}\PY{p}{,} \PY{n}{lw}\PY{o}{=}\PY{l+m+mi}{2}\PY{p}{,}
           \PY{n}{label}\PY{o}{=}\PY{l+s+sa}{r}\PY{l+s+s2}{\PYZdq{}}\PY{l+s+s2}{\PYZdl{}v\PYZus{}1\PYZdl{}: Eigenvector (\PYZdl{}}\PY{l+s+s2}{\PYZbs{}}\PY{l+s+s2}{lambda\PYZus{}1=3\PYZdl{})}\PY{l+s+s2}{\PYZdq{}}\PY{p}{)}\PY{p}{,}
    \PY{n}{Line2D}\PY{p}{(}\PY{p}{[}\PY{l+m+mi}{0}\PY{p}{]}\PY{p}{,} \PY{p}{[}\PY{l+m+mi}{0}\PY{p}{]}\PY{p}{,} \PY{n}{color}\PY{o}{=}\PY{l+s+s2}{\PYZdq{}}\PY{l+s+s2}{seagreen}\PY{l+s+s2}{\PYZdq{}}\PY{p}{,} \PY{n}{lw}\PY{o}{=}\PY{l+m+mi}{2}\PY{p}{,}
           \PY{n}{label}\PY{o}{=}\PY{l+s+sa}{r}\PY{l+s+s2}{\PYZdq{}}\PY{l+s+s2}{\PYZdl{}v\PYZus{}2\PYZdl{}: Eigenvector (\PYZdl{}}\PY{l+s+s2}{\PYZbs{}}\PY{l+s+s2}{lambda\PYZus{}2=1\PYZdl{})}\PY{l+s+s2}{\PYZdq{}}\PY{p}{)}\PY{p}{,}
    \PY{n}{Line2D}\PY{p}{(}\PY{p}{[}\PY{l+m+mi}{0}\PY{p}{]}\PY{p}{,} \PY{p}{[}\PY{l+m+mi}{0}\PY{p}{]}\PY{p}{,} \PY{n}{color}\PY{o}{=}\PY{l+s+s2}{\PYZdq{}}\PY{l+s+s2}{darkorange}\PY{l+s+s2}{\PYZdq{}}\PY{p}{,} \PY{n}{lw}\PY{o}{=}\PY{l+m+mi}{2}\PY{p}{,}
           \PY{n}{label}\PY{o}{=}\PY{l+s+sa}{r}\PY{l+s+s2}{\PYZdq{}}\PY{l+s+s2}{\PYZdl{}v\PYZus{}3\PYZdl{}: Not an eigenvector}\PY{l+s+s2}{\PYZdq{}}\PY{p}{)}\PY{p}{,}
    \PY{n}{Line2D}\PY{p}{(}\PY{p}{[}\PY{l+m+mi}{0}\PY{p}{]}\PY{p}{,} \PY{p}{[}\PY{l+m+mi}{0}\PY{p}{]}\PY{p}{,} \PY{n}{color}\PY{o}{=}\PY{l+s+s2}{\PYZdq{}}\PY{l+s+s2}{grey}\PY{l+s+s2}{\PYZdq{}}\PY{p}{,} \PY{n}{lw}\PY{o}{=}\PY{l+m+mi}{2}\PY{p}{,} \PY{n}{linestyle}\PY{o}{=}\PY{l+s+s2}{\PYZdq{}}\PY{l+s+s2}{\PYZhy{}\PYZhy{}}\PY{l+s+s2}{\PYZdq{}}\PY{p}{,}
           \PY{n}{label}\PY{o}{=}\PY{l+s+s2}{\PYZdq{}}\PY{l+s+s2}{Original}\PY{l+s+s2}{\PYZdq{}}\PY{p}{)}\PY{p}{,}
    \PY{n}{Line2D}\PY{p}{(}\PY{p}{[}\PY{l+m+mi}{0}\PY{p}{]}\PY{p}{,} \PY{p}{[}\PY{l+m+mi}{0}\PY{p}{]}\PY{p}{,} \PY{n}{color}\PY{o}{=}\PY{l+s+s2}{\PYZdq{}}\PY{l+s+s2}{grey}\PY{l+s+s2}{\PYZdq{}}\PY{p}{,} \PY{n}{lw}\PY{o}{=}\PY{l+m+mi}{2}\PY{p}{,}
           \PY{n}{label}\PY{o}{=}\PY{l+s+s2}{\PYZdq{}}\PY{l+s+s2}{Transformed}\PY{l+s+s2}{\PYZdq{}}\PY{p}{)}\PY{p}{,}
\PY{p}{]}

\PY{n}{ax}\PY{o}{.}\PY{n}{legend}\PY{p}{(}\PY{n}{handles}\PY{o}{=}\PY{n}{legend\PYZus{}elements}\PY{p}{,} \PY{n}{loc}\PY{o}{=}\PY{l+s+s2}{\PYZdq{}}\PY{l+s+s2}{upper right}\PY{l+s+s2}{\PYZdq{}}\PY{p}{)}

\PY{n}{ax}\PY{o}{.}\PY{n}{set}\PY{p}{(}
    \PY{n}{xlim}\PY{o}{=}\PY{p}{(}\PY{o}{\PYZhy{}}\PY{l+m+mf}{0.2}\PY{p}{,} \PY{l+m+mf}{3.4}\PY{p}{)}\PY{p}{,}
    \PY{n}{ylim}\PY{o}{=}\PY{p}{(}\PY{o}{\PYZhy{}}\PY{l+m+mf}{0.2}\PY{p}{,} \PY{l+m+mf}{1.5}\PY{p}{)}\PY{p}{,}
    \PY{n}{xlabel}\PY{o}{=}\PY{l+s+s2}{\PYZdq{}}\PY{l+s+s2}{X}\PY{l+s+s2}{\PYZdq{}}\PY{p}{,}
    \PY{n}{ylabel}\PY{o}{=}\PY{l+s+s2}{\PYZdq{}}\PY{l+s+s2}{Y}\PY{l+s+s2}{\PYZdq{}}\PY{p}{,}
    \PY{n}{title}\PY{o}{=}\PY{l+s+s2}{\PYZdq{}}\PY{l+s+s2}{Eigenvectors: Original vs Transformed}\PY{l+s+s2}{\PYZdq{}}
\PY{p}{)}

\PY{n}{ax}\PY{o}{.}\PY{n}{set\PYZus{}aspect}\PY{p}{(}\PY{l+s+s2}{\PYZdq{}}\PY{l+s+s2}{equal}\PY{l+s+s2}{\PYZdq{}}\PY{p}{)}
\PY{n}{ax}\PY{o}{.}\PY{n}{grid}\PY{p}{(}\PY{n}{alpha}\PY{o}{=}\PY{l+m+mf}{0.2}\PY{p}{)}

\PY{n}{plt}\PY{o}{.}\PY{n}{tight\PYZus{}layout}\PY{p}{(}\PY{p}{)}
\PY{n}{plt}\PY{o}{.}\PY{n}{show}\PY{p}{(}\PY{p}{)}
\end{Verbatim}
\end{tcolorbox}

    \begin{center}
    \adjustimage{max size={0.9\linewidth}{0.9\paperheight}}{week4_pdf_build_pca_files/week4_pdf_build_pca_24_0.png}
    \end{center}
    { \hspace*{\fill} \\}
    

    \subsubsection{What the plot
demonstrates}\label{what-the-plot-demonstrates}

The matrix

\[
A=\begin{bmatrix}3&0\\0&1\end{bmatrix}
\]

scales the X-axis by a factor of 3 while leaving the Y-axis unchanged.

{\def\LTcaptype{none} % do not increment counter
\begin{longtable}[]{@{}
  >{\raggedright\arraybackslash}p{(\linewidth - 6\tabcolsep) * \real{0.2500}}
  >{\raggedright\arraybackslash}p{(\linewidth - 6\tabcolsep) * \real{0.2500}}
  >{\raggedright\arraybackslash}p{(\linewidth - 6\tabcolsep) * \real{0.2500}}
  >{\raggedright\arraybackslash}p{(\linewidth - 6\tabcolsep) * \real{0.2500}}@{}}
\toprule\noalign{}
\begin{minipage}[b]{\linewidth}\raggedright
Vector
\end{minipage} & \begin{minipage}[b]{\linewidth}\raggedright
Original
\end{minipage} & \begin{minipage}[b]{\linewidth}\raggedright
Transformed
\end{minipage} & \begin{minipage}[b]{\linewidth}\raggedright
Interpretation
\end{minipage} \\
\midrule\noalign{}
\endhead
\bottomrule\noalign{}
\endlastfoot
\(v_1\) & \((1,0)\) & \((3,0)\) & Same direction, stretched by
\(\lambda_1=3\) \\
\(v_2\) & \((0,1)\) & \((0,1)\) & Unchanged, \(\lambda_2=1\) \\
\(v_3\) & \((0.7,0.7)\) & \((2.1,0.7)\) & Changes direction, so it is
not an eigenvector \\
\end{longtable}
}

\textbf{Key observations:}

\begin{itemize}
\tightlist
\item
  \textbf{Eigenvector \(v_1\):} Its direction remains unchanged, but its
  magnitude is multiplied by 3.
\item
  \textbf{Eigenvector \(v_2\):} Both its direction and magnitude remain
  unchanged because its eigenvalue is 1.
\item
  \textbf{Vector \(v_3\):} Its direction changes because the X component
  is multiplied by 3 while the Y component remains unchanged. Therefore,
  it is not an eigenvector.
\item
  \textbf{Overlapping vectors:} \(v_2\) and \(Av_2\) overlap perfectly
  because \(\lambda_2=1\).
\end{itemize}

\textbf{Conclusion:}

An eigenvector is a non-zero vector whose direction remains on the same
line when transformed by a matrix. Its corresponding eigenvalue
determines the scaling factor:

\[
\boxed{Av=\lambda v}
\]

\begin{itemize}
\tightlist
\item
  \(|\lambda|>1\): The vector is stretched.
\item
  \(0<|\lambda|<1\): The vector is compressed.
\item
  \(\lambda=1\): The vector remains unchanged.
\item
  \(\lambda=-1\): The vector reverses direction without changing
  magnitude.
\item
  \(\lambda=0\): The vector collapses to the zero vector.
\end{itemize}

\textbf{Connection to PCA:} Eigenvectors of the covariance matrix
identify the principal directions of variation in the data, while their
corresponding eigenvalues quantify the variance along those directions.



    \subsubsection{Why are eigenvalues and eigenvectors useful in
PCA?}\label{why-are-eigenvalues-and-eigenvectors-useful-in-pca}

For \textbf{symmetric covariance matrices}, NumPy's
\href{https://numpy.org/doc/stable/reference/generated/numpy.linalg.eigh.html}{\texttt{np.linalg.eigh}}
is preferable to
\href{https://numpy.org/doc/stable/reference/generated/numpy.linalg.eig.html}{\texttt{eig}}:
the eigenvalues are real and the eigenvectors can be chosen orthonormal.

\begin{itemize}
\tightlist
\item
  A covariance matrix's eigenvectors give orthogonal directions in
  feature space.
\item
  Its eigenvalues give the variance along each corresponding unit
  eigenvector.
\item
  The eigenvector corresponding to the largest eigenvalue identifies the
  \textbf{principal direction of greatest variance (PC1).}
\end{itemize}

Eigenvector signs are arbitrary: \(\mathbf v\) and \(-\mathbf v\)
represent the same axis.



    \subsection{6. From covariance to PCA}\label{section-6}

We'll generate \textbf{200 observations with two correlated features}
(the same dataset size you mentioned for clustering). The points form an
elongated cloud. PCA finds the cloud's principal axes.

\textbf{Important:} This example works in the original feature scales.
When features have very different units, standardising before PCA is
often appropriate.


    \begin{tcolorbox}[breakable, size=fbox, boxrule=1pt, pad at break*=1mm,colback=cellbackground, colframe=cellborder]
\prompt{In}{incolor}{17}{\boxspacing}
\begin{Verbatim}[commandchars=\\\{\}]
\PY{n}{rng} \PY{o}{=} \PY{n}{np}\PY{o}{.}\PY{n}{random}\PY{o}{.}\PY{n}{default\PYZus{}rng}\PY{p}{(}\PY{l+m+mi}{7}\PY{p}{)}
\PY{n}{X} \PY{o}{=} \PY{n}{rng}\PY{o}{.}\PY{n}{multivariate\PYZus{}normal}\PY{p}{(}
    \PY{n}{mean}\PY{o}{=}\PY{p}{[}\PY{l+m+mf}{2.0}\PY{p}{,} \PY{l+m+mf}{3.0}\PY{p}{]}\PY{p}{,}
    \PY{n}{cov}\PY{o}{=}\PY{p}{[}\PY{p}{[}\PY{l+m+mf}{3.0}\PY{p}{,} \PY{l+m+mf}{2.1}\PY{p}{]}\PY{p}{,} \PY{p}{[}\PY{l+m+mf}{2.1}\PY{p}{,} \PY{l+m+mf}{2.0}\PY{p}{]}\PY{p}{]}\PY{p}{,}
    \PY{n}{size}\PY{o}{=}\PY{l+m+mi}{200}
\PY{p}{)}
\PY{n}{X\PYZus{}centered} \PY{o}{=} \PY{n}{X} \PY{o}{\PYZhy{}} \PY{n}{X}\PY{o}{.}\PY{n}{mean}\PY{p}{(}\PY{n}{axis}\PY{o}{=}\PY{l+m+mi}{0}\PY{p}{)}
\PY{n}{S} \PY{o}{=} \PY{n}{np}\PY{o}{.}\PY{n}{cov}\PY{p}{(}\PY{n}{X\PYZus{}centered}\PY{p}{,} \PY{n}{rowvar}\PY{o}{=}\PY{k+kc}{False}\PY{p}{,} \PY{n}{ddof}\PY{o}{=}\PY{l+m+mi}{1}\PY{p}{)}

\PY{c+c1}{\PYZsh{} eigh returns eigenvalues in ascending order}
\PY{n}{X\PYZus{}eigen\PYZus{}values}\PY{p}{,} \PY{n}{X\PYZus{}eigen\PYZus{}vectors} \PY{o}{=} \PY{n}{np}\PY{o}{.}\PY{n}{linalg}\PY{o}{.}\PY{n}{eigh}\PY{p}{(}\PY{n}{S}\PY{p}{)}
\PY{n}{order} \PY{o}{=} \PY{n}{np}\PY{o}{.}\PY{n}{argsort}\PY{p}{(}\PY{n}{X\PYZus{}eigen\PYZus{}values}\PY{p}{)}\PY{p}{[}\PY{p}{:}\PY{p}{:}\PY{o}{\PYZhy{}}\PY{l+m+mi}{1}\PY{p}{]}
\PY{n}{X\PYZus{}eigen\PYZus{}values} \PY{o}{=} \PY{n}{X\PYZus{}eigen\PYZus{}values}\PY{p}{[}\PY{n}{order}\PY{p}{]}
\PY{n}{X\PYZus{}eigen\PYZus{}vectors} \PY{o}{=} \PY{n}{X\PYZus{}eigen\PYZus{}vectors}\PY{p}{[}\PY{p}{:}\PY{p}{,} \PY{n}{order}\PY{p}{]}
\PY{n}{explained\PYZus{}variance\PYZus{}ratio} \PY{o}{=} \PY{n}{X\PYZus{}eigen\PYZus{}values} \PY{o}{/} \PY{n}{X\PYZus{}eigen\PYZus{}values}\PY{o}{.}\PY{n}{sum}\PY{p}{(}\PY{p}{)}

\PY{n+nb}{print}\PY{p}{(}\PY{l+s+s2}{\PYZdq{}}\PY{l+s+s2}{Means:}\PY{l+s+s2}{\PYZdq{}}\PY{p}{,} \PY{n}{X}\PY{o}{.}\PY{n}{mean}\PY{p}{(}\PY{n}{axis}\PY{o}{=}\PY{l+m+mi}{0}\PY{p}{)}\PY{p}{)} \PY{c+c1}{\PYZsh{} axis=0 (Column\PYZhy{}wise). Collapses the rows vertically, calculating the mean of each column.}
\PY{n+nb}{print}\PY{p}{(}\PY{l+s+s2}{\PYZdq{}}\PY{l+s+s2}{Covariance matrix:}\PY{l+s+se}{\PYZbs{}n}\PY{l+s+s2}{\PYZdq{}}\PY{p}{,} \PY{n}{S}\PY{p}{)}
\PY{n+nb}{print}\PY{p}{(}\PY{l+s+s2}{\PYZdq{}}\PY{l+s+s2}{Eigenvalues (descending):}\PY{l+s+s2}{\PYZdq{}}\PY{p}{,} \PY{n}{X\PYZus{}eigen\PYZus{}values}\PY{p}{)}
\PY{n+nb}{print}\PY{p}{(}\PY{l+s+s2}{\PYZdq{}}\PY{l+s+s2}{Eigenvectors (columns):}\PY{l+s+se}{\PYZbs{}n}\PY{l+s+s2}{\PYZdq{}}\PY{p}{,} \PY{n}{X\PYZus{}eigen\PYZus{}vectors}\PY{p}{)}
\PY{n+nb}{print}\PY{p}{(}\PY{l+s+s2}{\PYZdq{}}\PY{l+s+s2}{Explained variance ratios:}\PY{l+s+s2}{\PYZdq{}}\PY{p}{,} \PY{n}{explained\PYZus{}variance\PYZus{}ratio}\PY{p}{)}
\PY{n+nb}{print}\PY{p}{(}\PY{l+s+s2}{\PYZdq{}}\PY{l+s+s2}{Total variance = trace(S) = sum(eigenvalues):}\PY{l+s+s2}{\PYZdq{}}\PY{p}{,} \PY{n}{np}\PY{o}{.}\PY{n}{trace}\PY{p}{(}\PY{n}{S}\PY{p}{)}\PY{p}{,} \PY{n}{X\PYZus{}eigen\PYZus{}values}\PY{o}{.}\PY{n}{sum}\PY{p}{(}\PY{p}{)}\PY{p}{)}
\PY{k}{assert} \PY{n}{np}\PY{o}{.}\PY{n}{allclose}\PY{p}{(}\PY{n}{S} \PY{o}{@} \PY{n}{X\PYZus{}eigen\PYZus{}vectors}\PY{p}{,} \PY{n}{X\PYZus{}eigen\PYZus{}vectors} \PY{o}{@} \PY{n}{np}\PY{o}{.}\PY{n}{diag}\PY{p}{(}\PY{n}{X\PYZus{}eigen\PYZus{}values}\PY{p}{)}\PY{p}{)}
\end{Verbatim}
\end{tcolorbox}

    \begin{Verbatim}[commandchars=\\\{\}]
Means: [2.2273 3.1063]
Covariance matrix:
 [[2.6729 1.7538]
 [1.7538 1.5808]]
Eigenvalues (descending): [3.9636 0.29  ]
Eigenvectors (columns):
 [[-0.8054  0.5928]
 [-0.5928 -0.8054]]
Explained variance ratios: [0.9318 0.0682]
Total variance = trace(S) = sum(eigenvalues): 4.253641263852073
4.253641263852072
    \end{Verbatim}


    \subsubsection{Explained Variance Ratio
(PCA)}\label{explained-variance-ratio-pca}

The \textbf{explained variance ratio} measures the proportion of the
total variance in a dataset captured by each Principal Component (PC).

In Principal Component Analysis (PCA):

\begin{itemize}
\tightlist
\item
  \textbf{Eigenvectors} identify the directions of maximum variance
  (principal components).
\item
  \textbf{Eigenvalues} quantify the variance along each principal
  component.
\item
  \textbf{Explained variance ratios} express each eigenvalue as a
  proportion of the total variance.
\end{itemize}

\paragraph{Mathematical formula}\label{mathematical-formula}

\[
\boxed{\text{Explained Variance Ratio}_i = \frac{\lambda_i}{\sum_{j=1}^{n}\lambda_j}}
\]

Where:

\begin{itemize}
\tightlist
\item
  \(\lambda_i\): Eigenvalue of principal component \(i\).
\item
  \(\sum_{j=1}^{n}\lambda_j\): Sum of all eigenvalues (total variance).
\item
  \(n\): Number of principal components.
\end{itemize}

\paragraph{Example}\label{example}

Suppose the eigenvalues of a covariance matrix are:

\[
\lambda_1=4.5,\qquad \lambda_2=0.5
\]

The total variance is:

\[
\lambda_1+\lambda_2=4.5+0.5=5
\]

\textbf{First Principal Component (PC1):}

\[
\text{EVR}_1=\frac{4.5}{5}=\boxed{90\%}
\]

\textbf{Second Principal Component (PC2):}

\[
\text{EVR}_2=\frac{0.5}{5}=\boxed{10\%}
\]

{\def\LTcaptype{none} % do not increment counter
\begin{longtable}[]{@{}lrr@{}}
\toprule\noalign{}
Principal Component & Eigenvalue & Explained Variance Ratio \\
\midrule\noalign{}
\endhead
\bottomrule\noalign{}
\endlastfoot
PC1 & 4.5 & 90\% \\
PC2 & 0.5 & 10\% \\
\textbf{Total} & \textbf{5.0} & \textbf{100\%} \\
\end{longtable}
}

\paragraph{Cumulative Explained
Variance}\label{cumulative-explained-variance}

The \textbf{cumulative explained variance ratio} is the sum of the
explained variance ratios of the first \(k\) principal components:

\[
\boxed{\text{Cumulative EVR}_k=\sum_{i=1}^{k}\text{EVR}_i}
\]

It helps determine how many principal components are needed to retain a
desired proportion of the dataset's variance.

For example, if PC1 explains 90\% of the variance, reducing a
two-dimensional dataset to one principal component retains 90\% of its
total variance.

\paragraph{Key Takeaways}\label{key-takeaways}

\begin{itemize}
\tightlist
\item
  The explained variance ratios sum to \textbf{100\%} when all principal
  components are included and total variance is non-zero.
\item
  A \textbf{larger eigenvalue} corresponds to a principal component
  explaining more variance.
\item
  PCA orders principal components from the largest to the smallest
  eigenvalue.
\item
  Components explaining very little variance can potentially be
  discarded to reduce dimensionality.
\item
  Retaining 95\% of the variance is a common heuristic, but the
  appropriate threshold depends on the application.
\end{itemize}

\textbf{Important:} Explained variance measures retained variability,
not necessarily retained predictive information. A low-variance
component can still contain useful information for a machine learning
model.


    \begin{tcolorbox}[breakable, size=fbox, boxrule=1pt, pad at break*=1mm,colback=cellbackground, colframe=cellborder]
\prompt{In}{incolor}{18}{\boxspacing}
\begin{Verbatim}[commandchars=\\\{\}]
\PY{n}{fig}\PY{p}{,} \PY{n}{axes} \PY{o}{=} \PY{n}{plt}\PY{o}{.}\PY{n}{subplots}\PY{p}{(}\PY{l+m+mi}{1}\PY{p}{,} \PY{l+m+mi}{2}\PY{p}{,} \PY{n}{figsize}\PY{o}{=}\PY{p}{(}\PY{l+m+mi}{12}\PY{p}{,} \PY{l+m+mi}{5}\PY{p}{)}\PY{p}{,} \PY{n}{sharex}\PY{o}{=}\PY{k+kc}{True}\PY{p}{,} \PY{n}{sharey}\PY{o}{=}\PY{k+kc}{True}\PY{p}{)}

\PY{c+c1}{\PYZsh{} Original data}
\PY{n}{axes}\PY{p}{[}\PY{l+m+mi}{0}\PY{p}{]}\PY{o}{.}\PY{n}{scatter}\PY{p}{(}\PY{n}{X}\PY{p}{[}\PY{p}{:}\PY{p}{,} \PY{l+m+mi}{0}\PY{p}{]}\PY{p}{,} \PY{n}{X}\PY{p}{[}\PY{p}{:}\PY{p}{,} \PY{l+m+mi}{1}\PY{p}{]}\PY{p}{,} \PY{n}{alpha}\PY{o}{=}\PY{l+m+mf}{0.65}\PY{p}{,} \PY{n}{color}\PY{o}{=}\PY{l+s+s2}{\PYZdq{}}\PY{l+s+s2}{royalblue}\PY{l+s+s2}{\PYZdq{}}\PY{p}{)}
\PY{n}{axes}\PY{p}{[}\PY{l+m+mi}{0}\PY{p}{]}\PY{o}{.}\PY{n}{scatter}\PY{p}{(}
    \PY{o}{*}\PY{n}{X}\PY{o}{.}\PY{n}{mean}\PY{p}{(}\PY{n}{axis}\PY{o}{=}\PY{l+m+mi}{0}\PY{p}{)}\PY{p}{,} \PY{c+c1}{\PYZsh{} The * unpacks the tuple returned by X.mean(axis=0)}
    \PY{n}{color}\PY{o}{=}\PY{l+s+s2}{\PYZdq{}}\PY{l+s+s2}{red}\PY{l+s+s2}{\PYZdq{}}\PY{p}{,} \PY{n}{marker}\PY{o}{=}\PY{l+s+s2}{\PYZdq{}}\PY{l+s+s2}{x}\PY{l+s+s2}{\PYZdq{}}\PY{p}{,} \PY{n}{s}\PY{o}{=}\PY{l+m+mi}{120}\PY{p}{,} \PY{n}{linewidths}\PY{o}{=}\PY{l+m+mi}{2}\PY{p}{,}
    \PY{n}{label}\PY{o}{=}\PY{l+s+s2}{\PYZdq{}}\PY{l+s+s2}{Sample mean}\PY{l+s+s2}{\PYZdq{}}
\PY{p}{)}
\PY{n}{axes}\PY{p}{[}\PY{l+m+mi}{0}\PY{p}{]}\PY{o}{.}\PY{n}{set\PYZus{}title}\PY{p}{(}\PY{l+s+s2}{\PYZdq{}}\PY{l+s+s2}{Original dataset (X)}\PY{l+s+s2}{\PYZdq{}}\PY{p}{)}
\PY{n}{axes}\PY{p}{[}\PY{l+m+mi}{0}\PY{p}{]}\PY{o}{.}\PY{n}{legend}\PY{p}{(}\PY{p}{)}

\PY{c+c1}{\PYZsh{} Centred data}
\PY{n}{axes}\PY{p}{[}\PY{l+m+mi}{1}\PY{p}{]}\PY{o}{.}\PY{n}{scatter}\PY{p}{(}
    \PY{n}{X\PYZus{}centered}\PY{p}{[}\PY{p}{:}\PY{p}{,} \PY{l+m+mi}{0}\PY{p}{]}\PY{p}{,} \PY{n}{X\PYZus{}centered}\PY{p}{[}\PY{p}{:}\PY{p}{,} \PY{l+m+mi}{1}\PY{p}{]}\PY{p}{,}
    \PY{n}{alpha}\PY{o}{=}\PY{l+m+mf}{0.65}\PY{p}{,} \PY{n}{color}\PY{o}{=}\PY{l+s+s2}{\PYZdq{}}\PY{l+s+s2}{seagreen}\PY{l+s+s2}{\PYZdq{}}
\PY{p}{)}
\PY{n}{axes}\PY{p}{[}\PY{l+m+mi}{1}\PY{p}{]}\PY{o}{.}\PY{n}{scatter}\PY{p}{(}
    \PY{l+m+mi}{0}\PY{p}{,} \PY{l+m+mi}{0}\PY{p}{,}
    \PY{n}{color}\PY{o}{=}\PY{l+s+s2}{\PYZdq{}}\PY{l+s+s2}{red}\PY{l+s+s2}{\PYZdq{}}\PY{p}{,} \PY{n}{marker}\PY{o}{=}\PY{l+s+s2}{\PYZdq{}}\PY{l+s+s2}{x}\PY{l+s+s2}{\PYZdq{}}\PY{p}{,} \PY{n}{s}\PY{o}{=}\PY{l+m+mi}{120}\PY{p}{,} \PY{n}{linewidths}\PY{o}{=}\PY{l+m+mi}{2}\PY{p}{,}
    \PY{n}{label}\PY{o}{=}\PY{l+s+s2}{\PYZdq{}}\PY{l+s+s2}{Centred mean}\PY{l+s+s2}{\PYZdq{}}
\PY{p}{)}
\PY{n}{axes}\PY{p}{[}\PY{l+m+mi}{1}\PY{p}{]}\PY{o}{.}\PY{n}{set\PYZus{}title}\PY{p}{(}\PY{l+s+s2}{\PYZdq{}}\PY{l+s+s2}{Centred dataset (X\PYZus{}centered)}\PY{l+s+s2}{\PYZdq{}}\PY{p}{)}
\PY{n}{axes}\PY{p}{[}\PY{l+m+mi}{1}\PY{p}{]}\PY{o}{.}\PY{n}{legend}\PY{p}{(}\PY{p}{)}

\PY{k}{for} \PY{n}{ax} \PY{o+ow}{in} \PY{n}{axes}\PY{p}{:}
    \PY{n}{ax}\PY{o}{.}\PY{n}{axhline}\PY{p}{(}\PY{l+m+mi}{0}\PY{p}{,} \PY{n}{color}\PY{o}{=}\PY{l+s+s2}{\PYZdq{}}\PY{l+s+s2}{grey}\PY{l+s+s2}{\PYZdq{}}\PY{p}{,} \PY{n}{linewidth}\PY{o}{=}\PY{l+m+mf}{0.8}\PY{p}{)}
    \PY{n}{ax}\PY{o}{.}\PY{n}{axvline}\PY{p}{(}\PY{l+m+mi}{0}\PY{p}{,} \PY{n}{color}\PY{o}{=}\PY{l+s+s2}{\PYZdq{}}\PY{l+s+s2}{grey}\PY{l+s+s2}{\PYZdq{}}\PY{p}{,} \PY{n}{linewidth}\PY{o}{=}\PY{l+m+mf}{0.8}\PY{p}{)}
    \PY{n}{ax}\PY{o}{.}\PY{n}{set\PYZus{}xlabel}\PY{p}{(}\PY{l+s+s2}{\PYZdq{}}\PY{l+s+s2}{Feature 1}\PY{l+s+s2}{\PYZdq{}}\PY{p}{)}
    \PY{n}{ax}\PY{o}{.}\PY{n}{set\PYZus{}ylabel}\PY{p}{(}\PY{l+s+s2}{\PYZdq{}}\PY{l+s+s2}{Feature 2}\PY{l+s+s2}{\PYZdq{}}\PY{p}{)}
    \PY{n}{ax}\PY{o}{.}\PY{n}{grid}\PY{p}{(}\PY{n}{alpha}\PY{o}{=}\PY{l+m+mf}{0.2}\PY{p}{)}
    \PY{n}{ax}\PY{o}{.}\PY{n}{set\PYZus{}aspect}\PY{p}{(}\PY{l+s+s2}{\PYZdq{}}\PY{l+s+s2}{equal}\PY{l+s+s2}{\PYZdq{}}\PY{p}{)}

\PY{n}{plt}\PY{o}{.}\PY{n}{tight\PYZus{}layout}\PY{p}{(}\PY{p}{)}
\PY{n}{plt}\PY{o}{.}\PY{n}{show}\PY{p}{(}\PY{p}{)}
\end{Verbatim}
\end{tcolorbox}

    \begin{center}
    \adjustimage{max size={0.9\linewidth}{0.9\paperheight}}{week4_pdf_build_pca_files/week4_pdf_build_pca_30_0.png}
    \end{center}
    { \hspace*{\fill} \\}
    
    \begin{tcolorbox}[breakable, size=fbox, boxrule=1pt, pad at break*=1mm,colback=cellbackground, colframe=cellborder]
\prompt{In}{incolor}{19}{\boxspacing}
\begin{Verbatim}[commandchars=\\\{\}]
\PY{k}{def}\PY{+w}{ }\PY{n+nf}{plot\PYZus{}pca\PYZus{}components}\PY{p}{(}\PY{n}{data}\PY{p}{,} \PY{n}{eigen\PYZus{}vectors}\PY{p}{,} \PY{n}{eigen\PYZus{}values}\PY{p}{,} \PY{n}{mode}\PY{o}{=}\PY{l+s+s2}{\PYZdq{}}\PY{l+s+s2}{axis}\PY{l+s+s2}{\PYZdq{}}\PY{p}{,} \PY{n}{components}\PY{o}{=}\PY{k+kc}{None}\PY{p}{,} \PY{n}{figsize}\PY{o}{=}\PY{p}{(}\PY{l+m+mi}{7}\PY{p}{,} \PY{l+m+mi}{6}\PY{p}{)}\PY{p}{,} \PY{n}{debug\PYZus{}mode}\PY{o}{=}\PY{k+kc}{False}\PY{p}{)}\PY{p}{:}
    \PY{n}{fig}\PY{p}{,} \PY{n}{ax} \PY{o}{=} \PY{n}{plt}\PY{o}{.}\PY{n}{subplots}\PY{p}{(}\PY{n}{figsize}\PY{o}{=}\PY{n}{figsize}\PY{p}{)}

    \PY{c+c1}{\PYZsh{} Plot data points}
    \PY{n}{ax}\PY{o}{.}\PY{n}{scatter}\PY{p}{(}\PY{n}{data}\PY{p}{[}\PY{p}{:}\PY{p}{,} \PY{l+m+mi}{0}\PY{p}{]}\PY{p}{,} \PY{n}{data}\PY{p}{[}\PY{p}{:}\PY{p}{,} \PY{l+m+mi}{1}\PY{p}{]}\PY{p}{,} \PY{n}{alpha}\PY{o}{=}\PY{l+m+mf}{0.40}\PY{p}{,} \PY{n}{s}\PY{o}{=}\PY{l+m+mi}{19}\PY{p}{)}

    \PY{c+c1}{\PYZsh{} Calculate the mean and explained variance ratio}
    \PY{n}{centre} \PY{o}{=} \PY{n}{data}\PY{o}{.}\PY{n}{mean}\PY{p}{(}\PY{n}{axis}\PY{o}{=}\PY{l+m+mi}{0}\PY{p}{)}
    \PY{n}{explained\PYZus{}variance\PYZus{}ratio} \PY{o}{=} \PY{n}{eigen\PYZus{}values} \PY{o}{/} \PY{n}{eigen\PYZus{}values}\PY{o}{.}\PY{n}{sum}\PY{p}{(}\PY{p}{)}

    \PY{k}{if} \PY{n}{components} \PY{o+ow}{is} \PY{k+kc}{None}\PY{p}{:}
        \PY{n}{components} \PY{o}{=} \PY{n+nb}{range}\PY{p}{(}\PY{l+m+mi}{2}\PY{p}{)}

    \PY{c+c1}{\PYZsh{} Plot principal components}
    \PY{k}{for} \PY{n}{i} \PY{o+ow}{in} \PY{n}{components}\PY{p}{:}
        \PY{n}{direction} \PY{o}{=} \PY{n}{eigen\PYZus{}vectors}\PY{p}{[}\PY{p}{:}\PY{p}{,} \PY{n}{i}\PY{p}{]} \PY{o}{*} \PY{n}{np}\PY{o}{.}\PY{n}{sqrt}\PY{p}{(}\PY{n}{eigen\PYZus{}values}\PY{p}{[}\PY{n}{i}\PY{p}{]}\PY{p}{)} \PY{o}{*} \PY{l+m+mi}{2}

        \PY{k}{if} \PY{n}{mode} \PY{o}{==} \PY{l+s+s2}{\PYZdq{}}\PY{l+s+s2}{axis}\PY{l+s+s2}{\PYZdq{}}\PY{p}{:}
            \PY{n}{start} \PY{o}{=} \PY{n}{centre} \PY{o}{\PYZhy{}} \PY{n}{direction}
            \PY{n}{end} \PY{o}{=} \PY{n}{centre} \PY{o}{+} \PY{n}{direction}
            \PY{n}{arrow\PYZus{}style} \PY{o}{=} \PY{l+s+s2}{\PYZdq{}}\PY{l+s+s2}{\PYZlt{}\PYZhy{}\PYZgt{}}\PY{l+s+s2}{\PYZdq{}}

        \PY{k}{elif} \PY{n}{mode} \PY{o}{==} \PY{l+s+s2}{\PYZdq{}}\PY{l+s+s2}{vector}\PY{l+s+s2}{\PYZdq{}}\PY{p}{:}
            \PY{n}{start} \PY{o}{=} \PY{n}{centre}
            \PY{n}{end} \PY{o}{=} \PY{n}{centre} \PY{o}{+} \PY{n}{direction}
            \PY{n}{arrow\PYZus{}style} \PY{o}{=} \PY{l+s+s2}{\PYZdq{}}\PY{l+s+s2}{\PYZhy{}\PYZgt{}}\PY{l+s+s2}{\PYZdq{}}

        \PY{k}{else}\PY{p}{:}
            \PY{k}{raise} \PY{n+ne}{ValueError}\PY{p}{(}\PY{l+s+s2}{\PYZdq{}}\PY{l+s+s2}{mode must be }\PY{l+s+s2}{\PYZsq{}}\PY{l+s+s2}{axis}\PY{l+s+s2}{\PYZsq{}}\PY{l+s+s2}{ or }\PY{l+s+s2}{\PYZsq{}}\PY{l+s+s2}{vector}\PY{l+s+s2}{\PYZsq{}}\PY{l+s+s2}{\PYZdq{}}\PY{p}{)}

        \PY{k}{if} \PY{n}{debug\PYZus{}mode}\PY{p}{:}
            \PY{n+nb}{print}\PY{p}{(}\PY{l+s+sa}{f}\PY{l+s+s2}{\PYZdq{}}\PY{l+s+s2}{Component }\PY{l+s+si}{\PYZob{}}\PY{n}{i}\PY{o}{+}\PY{l+m+mi}{1}\PY{l+s+si}{\PYZcb{}}\PY{l+s+s2}{: Start = }\PY{l+s+si}{\PYZob{}}\PY{n}{start}\PY{l+s+si}{\PYZcb{}}\PY{l+s+s2}{, End = }\PY{l+s+si}{\PYZob{}}\PY{n}{end}\PY{l+s+si}{\PYZcb{}}\PY{l+s+s2}{, Arrow Style = }\PY{l+s+si}{\PYZob{}}\PY{n}{arrow\PYZus{}style}\PY{l+s+si}{\PYZcb{}}\PY{l+s+s2}{\PYZdq{}}\PY{p}{)}

        \PY{n}{ax}\PY{o}{.}\PY{n}{annotate}\PY{p}{(}
            \PY{l+s+s2}{\PYZdq{}}\PY{l+s+s2}{\PYZdq{}}\PY{p}{,}
            \PY{n}{xy}\PY{o}{=}\PY{n}{end}\PY{p}{,}
            \PY{n}{xytext}\PY{o}{=}\PY{n}{start}\PY{p}{,}
            \PY{n}{arrowprops}\PY{o}{=}\PY{n+nb}{dict}\PY{p}{(}\PY{n}{arrowstyle}\PY{o}{=}\PY{n}{arrow\PYZus{}style}\PY{p}{,} \PY{n}{lw}\PY{o}{=}\PY{l+m+mi}{3}\PY{p}{,} \PY{n}{zorder}\PY{o}{=}\PY{l+m+mi}{2}\PY{p}{)}
        \PY{p}{)}

        \PY{n}{ax}\PY{o}{.}\PY{n}{text}\PY{p}{(}
            \PY{o}{*}\PY{p}{(}\PY{n}{end} \PY{o}{+} \PY{n}{direction} \PY{o}{*} \PY{l+m+mf}{0.08}\PY{p}{)}\PY{p}{,}
            \PY{l+s+sa}{f}\PY{l+s+s2}{\PYZdq{}}\PY{l+s+s2}{PC}\PY{l+s+si}{\PYZob{}}\PY{n}{i}\PY{o}{+}\PY{l+m+mi}{1}\PY{l+s+si}{\PYZcb{}}\PY{l+s+s2}{ (}\PY{l+s+si}{\PYZob{}}\PY{n}{explained\PYZus{}variance\PYZus{}ratio}\PY{p}{[}\PY{n}{i}\PY{p}{]}\PY{l+s+si}{:}\PY{l+s+s2}{.1\PYZpc{}}\PY{l+s+si}{\PYZcb{}}\PY{l+s+s2}{)}\PY{l+s+s2}{\PYZdq{}}
        \PY{p}{)}

    \PY{c+c1}{\PYZsh{} Plot mean}
    \PY{c+c1}{\PYZsh{}ax.scatter(*centre, marker=\PYZdq{}x\PYZdq{}, s=95, label=\PYZdq{}Mean\PYZdq{}) \PYZsh{} With the original code it is not shown properly.}
    \PY{n}{ax}\PY{o}{.}\PY{n}{scatter}\PY{p}{(}
        \PY{o}{*}\PY{n}{centre}\PY{p}{,}
        \PY{n}{marker}\PY{o}{=}\PY{l+s+s2}{\PYZdq{}}\PY{l+s+s2}{X}\PY{l+s+s2}{\PYZdq{}}\PY{p}{,}
        \PY{n}{s}\PY{o}{=}\PY{l+m+mi}{180}\PY{p}{,}
        \PY{n}{color}\PY{o}{=}\PY{l+s+s2}{\PYZdq{}}\PY{l+s+s2}{red}\PY{l+s+s2}{\PYZdq{}}\PY{p}{,}
        \PY{n}{edgecolors}\PY{o}{=}\PY{l+s+s2}{\PYZdq{}}\PY{l+s+s2}{white}\PY{l+s+s2}{\PYZdq{}}\PY{p}{,}
        \PY{n}{linewidths}\PY{o}{=}\PY{l+m+mi}{1}\PY{p}{,}
        \PY{n}{zorder}\PY{o}{=}\PY{l+m+mi}{10}\PY{p}{,}
        \PY{n}{label}\PY{o}{=}\PY{l+s+s2}{\PYZdq{}}\PY{l+s+s2}{Mean}\PY{l+s+s2}{\PYZdq{}}
    \PY{p}{)}

    \PY{k}{if} \PY{n}{debug\PYZus{}mode}\PY{p}{:}
        \PY{n+nb}{print}\PY{p}{(}\PY{l+s+sa}{f}\PY{l+s+s2}{\PYZdq{}}\PY{l+s+s2}{Mean: }\PY{l+s+si}{\PYZob{}}\PY{n}{centre}\PY{l+s+si}{\PYZcb{}}\PY{l+s+s2}{\PYZdq{}}\PY{p}{)}

    \PY{n}{fig}\PY{o}{.}\PY{n}{suptitle}\PY{p}{(}
        \PY{l+s+sa}{f}\PY{l+s+s2}{\PYZdq{}}\PY{l+s+s2}{PCA: }\PY{l+s+si}{\PYZob{}}\PY{n}{mode}\PY{o}{.}\PY{n}{capitalize}\PY{p}{(}\PY{p}{)}\PY{l+s+si}{\PYZcb{}}\PY{l+s+s2}{ Representation}\PY{l+s+s2}{\PYZdq{}}\PY{p}{,}
        \PY{n}{fontsize}\PY{o}{=}\PY{l+m+mi}{14}\PY{p}{,}
        \PY{n}{fontweight}\PY{o}{=}\PY{l+s+s2}{\PYZdq{}}\PY{l+s+s2}{bold}\PY{l+s+s2}{\PYZdq{}}\PY{p}{,}
        \PY{n}{y}\PY{o}{=}\PY{l+m+mf}{0.98}
    \PY{p}{)}

    \PY{n}{ax}\PY{o}{.}\PY{n}{set\PYZus{}title}\PY{p}{(}
        \PY{l+s+s2}{\PYZdq{}}\PY{l+s+s2}{Eigenvectors define directions; eigenvalues quantify variance}\PY{l+s+s2}{\PYZdq{}}\PY{p}{,}
        \PY{n}{fontsize}\PY{o}{=}\PY{l+m+mi}{10}\PY{p}{,}
        \PY{n}{color}\PY{o}{=}\PY{l+s+s2}{\PYZdq{}}\PY{l+s+s2}{grey}\PY{l+s+s2}{\PYZdq{}}\PY{p}{,}
        \PY{n}{pad}\PY{o}{=}\PY{l+m+mi}{8}
    \PY{p}{)}


    \PY{c+c1}{\PYZsh{} Default layout}
    \PY{n}{fig}\PY{o}{.}\PY{n}{subplots\PYZus{}adjust}\PY{p}{(}\PY{n}{top}\PY{o}{=}\PY{l+m+mf}{0.95}\PY{p}{)}

    \PY{c+c1}{\PYZsh{} Preserve equal scaling}
    \PY{n}{ax}\PY{o}{.}\PY{n}{set\PYZus{}aspect}\PY{p}{(}\PY{l+s+s2}{\PYZdq{}}\PY{l+s+s2}{equal}\PY{l+s+s2}{\PYZdq{}}\PY{p}{,} \PY{n}{adjustable}\PY{o}{=}\PY{l+s+s2}{\PYZdq{}}\PY{l+s+s2}{box}\PY{l+s+s2}{\PYZdq{}}\PY{p}{)}

    \PY{c+c1}{\PYZsh{} Check the actual rendered title positions}
    \PY{n}{fig}\PY{o}{.}\PY{n}{canvas}\PY{o}{.}\PY{n}{draw}\PY{p}{(}\PY{p}{)}
    \PY{n}{renderer} \PY{o}{=} \PY{n}{fig}\PY{o}{.}\PY{n}{canvas}\PY{o}{.}\PY{n}{get\PYZus{}renderer}\PY{p}{(}\PY{p}{)}

    \PY{n}{main\PYZus{}title} \PY{o}{=} \PY{n}{fig}\PY{o}{.}\PY{n}{\PYZus{}suptitle}\PY{o}{.}\PY{n}{get\PYZus{}window\PYZus{}extent}\PY{p}{(}\PY{n}{renderer}\PY{p}{)}
    \PY{n}{subtitle} \PY{o}{=} \PY{n}{ax}\PY{o}{.}\PY{n}{title}\PY{o}{.}\PY{n}{get\PYZus{}window\PYZus{}extent}\PY{p}{(}\PY{n}{renderer}\PY{p}{)}

    \PY{n}{gap} \PY{o}{=} \PY{n}{main\PYZus{}title}\PY{o}{.}\PY{n}{y0} \PY{o}{\PYZhy{}} \PY{n}{subtitle}\PY{o}{.}\PY{n}{y1}

    \PY{c+c1}{\PYZsh{} Adjust only when the titles are too close}
    \PY{k}{if} \PY{n}{gap} \PY{o}{\PYZlt{}} \PY{l+m+mi}{5}\PY{p}{:}
        \PY{n}{fig}\PY{o}{.}\PY{n}{subplots\PYZus{}adjust}\PY{p}{(}\PY{n}{top}\PY{o}{=}\PY{l+m+mf}{0.90}\PY{p}{)}
        \PY{n}{fig}\PY{o}{.}\PY{n}{canvas}\PY{o}{.}\PY{n}{draw}\PY{p}{(}\PY{p}{)}

    \PY{n}{ax}\PY{o}{.}\PY{n}{grid}\PY{p}{(}\PY{n}{alpha}\PY{o}{=}\PY{l+m+mf}{0.2}\PY{p}{)}
    \PY{n}{ax}\PY{o}{.}\PY{n}{legend}\PY{p}{(}\PY{p}{)}

    \PY{k}{if} \PY{n}{debug\PYZus{}mode}\PY{p}{:}
        \PY{n}{renderer} \PY{o}{=} \PY{n}{fig}\PY{o}{.}\PY{n}{canvas}\PY{o}{.}\PY{n}{get\PYZus{}renderer}\PY{p}{(}\PY{p}{)}
        \PY{n}{main\PYZus{}title} \PY{o}{=} \PY{n}{fig}\PY{o}{.}\PY{n}{\PYZus{}suptitle}\PY{o}{.}\PY{n}{get\PYZus{}window\PYZus{}extent}\PY{p}{(}\PY{n}{renderer}\PY{p}{)}
        \PY{n}{subtitle} \PY{o}{=} \PY{n}{ax}\PY{o}{.}\PY{n}{title}\PY{o}{.}\PY{n}{get\PYZus{}window\PYZus{}extent}\PY{p}{(}\PY{n}{renderer}\PY{p}{)}

        \PY{n}{final\PYZus{}gap} \PY{o}{=} \PY{n}{main\PYZus{}title}\PY{o}{.}\PY{n}{y0} \PY{o}{\PYZhy{}} \PY{n}{subtitle}\PY{o}{.}\PY{n}{y1}
        \PY{n+nb}{print}\PY{p}{(}\PY{l+s+sa}{f}\PY{l+s+s2}{\PYZdq{}}\PY{l+s+s2}{Final title gap: }\PY{l+s+si}{\PYZob{}}\PY{n}{final\PYZus{}gap}\PY{l+s+si}{:}\PY{l+s+s2}{.2f}\PY{l+s+si}{\PYZcb{}}\PY{l+s+s2}{ px}\PY{l+s+s2}{\PYZdq{}}\PY{p}{)}

    \PY{n}{plt}\PY{o}{.}\PY{n}{show}\PY{p}{(}\PY{p}{)}

    \PY{k}{return} \PY{n}{centre}

\PY{n}{centre} \PY{o}{=} \PY{n}{plot\PYZus{}pca\PYZus{}components}\PY{p}{(}\PY{n}{X}\PY{p}{,} \PY{n}{X\PYZus{}eigen\PYZus{}vectors}\PY{p}{,} \PY{n}{X\PYZus{}eigen\PYZus{}values}\PY{p}{,} \PY{n}{mode}\PY{o}{=}\PY{l+s+s2}{\PYZdq{}}\PY{l+s+s2}{axis}\PY{l+s+s2}{\PYZdq{}}\PY{p}{,} \PY{n}{debug\PYZus{}mode}\PY{o}{=}\PY{k+kc}{False}\PY{p}{)}
\end{Verbatim}
\end{tcolorbox}

    \begin{center}
    \adjustimage{max size={0.9\linewidth}{0.9\paperheight}}{week4_pdf_build_pca_files/week4_pdf_build_pca_31_0.png}
    \end{center}
    { \hspace*{\fill} \\}
    

    \subsection{7. Project the data onto the first principal
component}\label{project-the-data-onto-the-first-principal-component}

Let \(\mathbf V\) be the matrix containing the \textbf{eigenvectors of
the covariance matrix as columns}, ordered by their corresponding
eigenvalues in descending order.

The PCA coordinates are calculated as:

\[
\mathbf Z=(\mathbf X-\bar{\mathbf X})\mathbf V
\]

To reduce the data from two dimensions to one, retain only the
\textbf{first column of \(\mathbf V\)}, corresponding to the largest
eigenvalue.

This preserves the greatest possible sample variance among all
one-dimensional \textbf{linear projections} of the centred data.


    \begin{tcolorbox}[breakable, size=fbox, boxrule=1pt, pad at break*=1mm,colback=cellbackground, colframe=cellborder]
\prompt{In}{incolor}{20}{\boxspacing}
\begin{Verbatim}[commandchars=\\\{\}]
\PY{n}{Z} \PY{o}{=} \PY{n}{X\PYZus{}centered} \PY{o}{@} \PY{n}{X\PYZus{}eigen\PYZus{}vectors}
\PY{n}{pc1\PYZus{}scores} \PY{o}{=} \PY{n}{Z}\PY{p}{[}\PY{p}{:}\PY{p}{,} \PY{l+m+mi}{0}\PY{p}{]} \PY{c+c1}{\PYZsh{} All rows, first column of Z}
\PY{n}{X\PYZus{}reconstructed\PYZus{}1d} \PY{o}{=} \PY{n}{np}\PY{o}{.}\PY{n}{outer}\PY{p}{(}\PY{n}{pc1\PYZus{}scores}\PY{p}{,} \PY{n}{X\PYZus{}eigen\PYZus{}vectors}\PY{p}{[}\PY{p}{:}\PY{p}{,}\PY{l+m+mi}{0}\PY{p}{]}\PY{p}{)} \PY{o}{+} \PY{n}{centre}

\PY{n+nb}{print}\PY{p}{(}\PY{l+s+s2}{\PYZdq{}}\PY{l+s+s2}{Original X shape:}\PY{l+s+s2}{\PYZdq{}}\PY{p}{,} \PY{n}{X}\PY{o}{.}\PY{n}{shape}\PY{p}{)}
\PY{n+nb}{print}\PY{p}{(}\PY{l+s+s2}{\PYZdq{}}\PY{l+s+s2}{Full PCA Z shape:}\PY{l+s+s2}{\PYZdq{}}\PY{p}{,} \PY{n}{Z}\PY{o}{.}\PY{n}{shape}\PY{p}{)}
\PY{n+nb}{print}\PY{p}{(}\PY{l+s+s2}{\PYZdq{}}\PY{l+s+s2}{One\PYZhy{}dimensional PCA shape:}\PY{l+s+s2}{\PYZdq{}}\PY{p}{,} \PY{n}{pc1\PYZus{}scores}\PY{o}{.}\PY{n}{shape}\PY{p}{)}
\PY{n+nb}{print}\PY{p}{(}\PY{l+s+s2}{\PYZdq{}}\PY{l+s+s2}{PC1 sample variance:}\PY{l+s+s2}{\PYZdq{}}\PY{p}{,} \PY{n}{np}\PY{o}{.}\PY{n}{var}\PY{p}{(}\PY{n}{pc1\PYZus{}scores}\PY{p}{,} \PY{n}{ddof}\PY{o}{=}\PY{l+m+mi}{1}\PY{p}{)}\PY{p}{)}
\PY{n+nb}{print}\PY{p}{(}\PY{l+s+s2}{\PYZdq{}}\PY{l+s+s2}{PC1 eigenvalue:}\PY{l+s+s2}{\PYZdq{}}\PY{p}{,} \PY{n}{X\PYZus{}eigen\PYZus{}values}\PY{p}{[}\PY{l+m+mi}{0}\PY{p}{]}\PY{p}{)}
\PY{n+nb}{print}\PY{p}{(}\PY{l+s+s2}{\PYZdq{}}\PY{l+s+s2}{1D explained variance:}\PY{l+s+s2}{\PYZdq{}}\PY{p}{,} \PY{l+s+sa}{f}\PY{l+s+s2}{\PYZdq{}}\PY{l+s+si}{\PYZob{}}\PY{n}{explained\PYZus{}variance\PYZus{}ratio}\PY{p}{[}\PY{l+m+mi}{0}\PY{p}{]}\PY{l+s+si}{:}\PY{l+s+s2}{.1\PYZpc{}}\PY{l+s+si}{\PYZcb{}}\PY{l+s+s2}{\PYZdq{}}\PY{p}{)}
\PY{k}{assert} \PY{n}{np}\PY{o}{.}\PY{n}{isclose}\PY{p}{(}\PY{n}{np}\PY{o}{.}\PY{n}{var}\PY{p}{(}\PY{n}{pc1\PYZus{}scores}\PY{p}{,} \PY{n}{ddof}\PY{o}{=}\PY{l+m+mi}{1}\PY{p}{)}\PY{p}{,} \PY{n}{X\PYZus{}eigen\PYZus{}values}\PY{p}{[}\PY{l+m+mi}{0}\PY{p}{]}\PY{p}{)}

\PY{n}{fig}\PY{p}{,} \PY{n}{ax} \PY{o}{=} \PY{n}{plt}\PY{o}{.}\PY{n}{subplots}\PY{p}{(}\PY{n}{figsize}\PY{o}{=}\PY{p}{(}\PY{l+m+mi}{7}\PY{p}{,} \PY{l+m+mi}{5}\PY{p}{)}\PY{p}{)}
\PY{n}{ax}\PY{o}{.}\PY{n}{scatter}\PY{p}{(}\PY{n}{X}\PY{p}{[}\PY{p}{:}\PY{p}{,}\PY{l+m+mi}{0}\PY{p}{]}\PY{p}{,} \PY{n}{X}\PY{p}{[}\PY{p}{:}\PY{p}{,}\PY{l+m+mi}{1}\PY{p}{]}\PY{p}{,} \PY{n}{alpha}\PY{o}{=}\PY{l+m+mf}{.25}\PY{p}{,} \PY{n}{label}\PY{o}{=}\PY{l+s+s2}{\PYZdq{}}\PY{l+s+s2}{Original data}\PY{l+s+s2}{\PYZdq{}}\PY{p}{)}
\PY{n}{ax}\PY{o}{.}\PY{n}{scatter}\PY{p}{(}\PY{n}{X\PYZus{}reconstructed\PYZus{}1d}\PY{p}{[}\PY{p}{:}\PY{p}{,}\PY{l+m+mi}{0}\PY{p}{]}\PY{p}{,} \PY{n}{X\PYZus{}reconstructed\PYZus{}1d}\PY{p}{[}\PY{p}{:}\PY{p}{,}\PY{l+m+mi}{1}\PY{p}{]}\PY{p}{,} \PY{n}{s}\PY{o}{=}\PY{l+m+mi}{12}\PY{p}{,}
           \PY{n}{alpha}\PY{o}{=}\PY{l+m+mf}{.8}\PY{p}{,} \PY{n}{label}\PY{o}{=}\PY{l+s+s2}{\PYZdq{}}\PY{l+s+s2}{Reconstruction from PC1}\PY{l+s+s2}{\PYZdq{}}\PY{p}{)}
\PY{n}{ax}\PY{o}{.}\PY{n}{set}\PY{p}{(}\PY{n}{xlabel}\PY{o}{=}\PY{l+s+s2}{\PYZdq{}}\PY{l+s+s2}{Feature 1}\PY{l+s+s2}{\PYZdq{}}\PY{p}{,} \PY{n}{ylabel}\PY{o}{=}\PY{l+s+s2}{\PYZdq{}}\PY{l+s+s2}{Feature 2}\PY{l+s+s2}{\PYZdq{}}\PY{p}{,} \PY{n}{title}\PY{o}{=}\PY{l+s+s2}{\PYZdq{}}\PY{l+s+s2}{PCA dimensionality reduction: 2D → 1D}\PY{l+s+s2}{\PYZdq{}}\PY{p}{)}
\PY{n}{ax}\PY{o}{.}\PY{n}{set\PYZus{}aspect}\PY{p}{(}\PY{l+s+s1}{\PYZsq{}}\PY{l+s+s1}{equal}\PY{l+s+s1}{\PYZsq{}}\PY{p}{,} \PY{n}{adjustable}\PY{o}{=}\PY{l+s+s1}{\PYZsq{}}\PY{l+s+s1}{box}\PY{l+s+s1}{\PYZsq{}}\PY{p}{)}
\PY{n}{ax}\PY{o}{.}\PY{n}{grid}\PY{p}{(}\PY{n}{alpha}\PY{o}{=}\PY{l+m+mf}{.2}\PY{p}{)}
\PY{n}{ax}\PY{o}{.}\PY{n}{legend}\PY{p}{(}\PY{p}{)}
\PY{n}{plt}\PY{o}{.}\PY{n}{show}\PY{p}{(}\PY{p}{)}
\end{Verbatim}
\end{tcolorbox}

    \begin{Verbatim}[commandchars=\\\{\}]
Original X shape: (200, 2)
Full PCA Z shape: (200, 2)
One-dimensional PCA shape: (200,)
PC1 sample variance: 3.963641677159418
PC1 eigenvalue: 3.9636416771594183
1D explained variance: 93.2\%
    \end{Verbatim}

    \begin{center}
    \adjustimage{max size={0.9\linewidth}{0.9\paperheight}}{week4_pdf_build_pca_files/week4_pdf_build_pca_33_1.png}
    \end{center}
    { \hspace*{\fill} \\}
    

    \subsubsection{PCA Projection, Dimensionality Reduction and
Reconstruction}\label{pca-projection-dimensionality-reduction-and-reconstruction}

The following three lines perform \textbf{PCA projection, dimensionality
reduction and reconstruction}:

\begin{Shaded}
\begin{Highlighting}[]
\NormalTok{Z }\OperatorTok{=}\NormalTok{ X\_centered }\OperatorTok{@}\NormalTok{ X\_eigen\_vectors}
\NormalTok{pc1\_scores }\OperatorTok{=}\NormalTok{ Z[:, }\DecValTok{0}\NormalTok{]}
\NormalTok{X\_reconstructed\_1d }\OperatorTok{=}\NormalTok{ np.outer(pc1\_scores, X\_eigen\_vectors[:, }\DecValTok{0}\NormalTok{]) }\OperatorTok{+}\NormalTok{ centre}
\end{Highlighting}
\end{Shaded}

The process consists of three stages:

\begin{enumerate}
\def\labelenumi{\arabic{enumi}.}
\tightlist
\item
  \textbf{Projection:} Transform the original features into principal
  component coordinates.
\item
  \textbf{Dimensionality reduction:} Retain only the first principal
  component (PC1).
\item
  \textbf{Reconstruction:} Approximate the original dataset using only
  the information retained by PC1.
\end{enumerate}

The key distinction is that \textbf{projection changes the coordinate
system, dimensionality reduction discards components, and reconstruction
returns the reduced data to the original feature space}.

\paragraph{1. Project the centred data onto the principal
components}\label{project-the-centred-data-onto-the-principal-components}

\begin{Shaded}
\begin{Highlighting}[]
\NormalTok{Z }\OperatorTok{=}\NormalTok{ X\_centered }\OperatorTok{@}\NormalTok{ X\_eigen\_vectors}
\end{Highlighting}
\end{Shaded}

Mathematically:

\[
\boxed{\mathbf Z=\mathbf X_c\mathbf V}
\]

Where:

\begin{itemize}
\tightlist
\item
  \(\mathbf X_c\) is the \textbf{centred dataset}, obtained by
  subtracting the mean of each feature.
\item
  \(\mathbf V\) is the matrix containing the \textbf{eigenvectors of the
  covariance matrix as columns}, ordered by decreasing eigenvalue.
\item
  \(\mathbf Z\) contains the new PCA coordinates, also called
  \textbf{principal component scores}.
\end{itemize}

The centred dataset is calculated as:

\[
\mathbf X_c=\mathbf X-\bar{\mathbf X}
\]

Here, \(\bar{\mathbf X}\) represents the vector of feature means,
subtracted from every observation.

\subparagraph{Numerical example}\label{numerical-example}

Consider one centred observation:

\[
\mathbf x_c=
\begin{bmatrix}
2 & 1
\end{bmatrix}
\]

And the following eigenvector matrix:

\[
\mathbf V=
\begin{bmatrix}
0.8 & -0.6\\
0.6 & 0.8
\end{bmatrix}
\]

Each column represents a principal component direction:

\[
\mathbf v_1=
\begin{bmatrix}
0.8\\
0.6
\end{bmatrix}
\qquad
\mathbf v_2=
\begin{bmatrix}
-0.6\\
0.8
\end{bmatrix}
\]

The projection is:

\[
\begin{aligned}
\mathbf z
&=\mathbf x_c\mathbf V\\[6pt]
&=
\begin{bmatrix}
2 & 1
\end{bmatrix}
\begin{bmatrix}
0.8 & -0.6\\
0.6 & 0.8
\end{bmatrix}\\[6pt]
&=
\begin{bmatrix}
2(0.8)+1(0.6) &
2(-0.6)+1(0.8)
\end{bmatrix}\\[6pt]
&=
\boxed{
\begin{bmatrix}
2.2 & -0.4
\end{bmatrix}
}
\end{aligned}
\]

The observation originally had centred coordinates \((2,1)\). In the PCA
coordinate system, its coordinates are \((2.2,-0.4)\).

\textbf{Important:} At this stage, no dimensionality reduction has
occurred. The observation still has two coordinates, but they now
describe its position along PC1 and PC2 rather than the original feature
axes.

\paragraph{2. Keep only the first principal
component}\label{keep-only-the-first-principal-component}

\begin{Shaded}
\begin{Highlighting}[]
\NormalTok{pc1\_scores }\OperatorTok{=}\NormalTok{ Z[:, }\DecValTok{0}\NormalTok{]}
\end{Highlighting}
\end{Shaded}

In NumPy:

\begin{itemize}
\tightlist
\item
  \texttt{:} selects \textbf{all rows}, representing all observations.
\item
  \texttt{0} selects the \textbf{first column}, representing PC1.
\end{itemize}

Suppose the projected dataset is:

\[
\mathbf Z=
\begin{bmatrix}
2.2 & -0.4\\
1.5 & 0.2\\
-0.8 & 0.1
\end{bmatrix}
\]

After selecting the first column:

\[
\mathbf z_1=
\begin{bmatrix}
2.2\\
1.5\\
-0.8
\end{bmatrix}
\]

The dataset has now been reduced from \textbf{two features to one
principal component}.

Each observation is represented by a single number describing its
position along the PC1 axis.

For the dataset analysed earlier, PC1 explains approximately
\textbf{93.18\% of the total variance}.

Therefore, retaining PC1 preserves most of the original variability,
while discarding PC2 removes approximately 6.82\%.

\paragraph{3. Reconstruct the dataset using only
PC1}\label{reconstruct-the-dataset-using-only-pc1}

\begin{Shaded}
\begin{Highlighting}[]
\NormalTok{X\_reconstructed\_1d }\OperatorTok{=}\NormalTok{ np.outer(pc1\_scores, X\_eigen\_vectors[:, }\DecValTok{0}\NormalTok{]) }\OperatorTok{+}\NormalTok{ centre}
\end{Highlighting}
\end{Shaded}

This reconstructs an approximation of the original observations using
only their PC1 scores.

The mathematical expression is:

\[
\boxed{
\hat{\mathbf X}
=
\mathbf z_1\mathbf v_1^T+\boldsymbol\mu
}
\]

Where:

\begin{itemize}
\tightlist
\item
  \(\hat{\mathbf X}\) is the reconstructed dataset.
\item
  \(\mathbf z_1\) contains the PC1 scores for every observation.
\item
  \(\mathbf v_1^T\) is the transpose of the first eigenvector.
\item
  \(\boldsymbol\mu\) is the vector of original feature means, added to
  every reconstructed observation.
\end{itemize}

\subparagraph{Step 3.1: Extract the first
eigenvector}\label{step-3.1-extract-the-first-eigenvector}

\begin{Shaded}
\begin{Highlighting}[]
\NormalTok{X\_eigen\_vectors[:, }\DecValTok{0}\NormalTok{]}
\end{Highlighting}
\end{Shaded}

This extracts the first column of the eigenvector matrix:

\[
\mathbf v_1=
\begin{bmatrix}
0.8\\
0.6
\end{bmatrix}
\]

Its transpose is:

\[
\mathbf v_1^T=
\begin{bmatrix}
0.8 & 0.6
\end{bmatrix}
\]

\subparagraph{Step 3.2: Calculate the outer
product}\label{step-3.2-calculate-the-outer-product}

\begin{Shaded}
\begin{Highlighting}[]
\NormalTok{np.outer(pc1\_scores, X\_eigen\_vectors[:, }\DecValTok{0}\NormalTok{])}
\end{Highlighting}
\end{Shaded}

The outer product multiplies every PC1 score by the coordinates of the
first eigenvector.

Using our example:

\[
\mathbf z_1=
\begin{bmatrix}
2.2\\
1.5\\
-0.8
\end{bmatrix}
\]

We calculate:

\[
\begin{aligned}
\mathbf z_1\mathbf v_1^T
&=
\begin{bmatrix}
2.2\\
1.5\\
-0.8
\end{bmatrix}
\begin{bmatrix}
0.8 & 0.6
\end{bmatrix}\\[6pt]
&=
\begin{bmatrix}
2.2(0.8) & 2.2(0.6)\\
1.5(0.8) & 1.5(0.6)\\
-0.8(0.8) & -0.8(0.6)
\end{bmatrix}\\[6pt]
&=
\boxed{
\begin{bmatrix}
1.76 & 1.32\\
1.20 & 0.90\\
-0.64 & -0.48
\end{bmatrix}
}
\end{aligned}
\]

We have transformed one score per observation back into two feature
coordinates.

However, these coordinates are still \textbf{centred around zero}.

\subparagraph{Step 3.3: Add the original
mean}\label{step-3.3-add-the-original-mean}

\begin{Shaded}
\begin{Highlighting}[]
\OperatorTok{+}\NormalTok{ centre}
\end{Highlighting}
\end{Shaded}

Remember that PCA was performed on centred data:

\begin{Shaded}
\begin{Highlighting}[]
\NormalTok{centre }\OperatorTok{=}\NormalTok{ X.mean(axis}\OperatorTok{=}\DecValTok{0}\NormalTok{)}
\NormalTok{X\_centered }\OperatorTok{=}\NormalTok{ X }\OperatorTok{{-}}\NormalTok{ centre}
\end{Highlighting}
\end{Shaded}

To reconstruct observations in the original coordinate system, we must
add the feature means back.

For example, if:

\[
\boldsymbol\mu=
\begin{bmatrix}
3 & 4
\end{bmatrix}
\]

Then:

\[
\begin{aligned}
\hat{\mathbf X}
&=
\begin{bmatrix}
1.76 & 1.32\\
1.20 & 0.90\\
-0.64 & -0.48
\end{bmatrix}
+
\begin{bmatrix}
3 & 4
\end{bmatrix}\\[6pt]
&=
\boxed{
\begin{bmatrix}
4.76 & 5.32\\
4.20 & 4.90\\
2.36 & 3.52
\end{bmatrix}
}
\end{aligned}
\]

NumPy uses
\href{https://numpy.org/doc/2.5/user/basics.broadcasting.html}{\textbf{broadcasting}}
to add the mean vector to every row.

The reconstructed dataset again has two features, but the observations
now lie entirely along the PC1 axis.

\paragraph{4. Visualising projection and
reconstruction}\label{visualising-projection-and-reconstruction}

\begin{figure}
\centering
\pandocbounded{\includesvg[keepaspectratio]{images/week_4/projection_and_reconstruction.svg}}
\caption{PCA projection and reconstruction using PC1}
\end{figure}

The diagram illustrates:

\begin{itemize}
\tightlist
\item
  \textbf{Orange points:} Original observations in two-dimensional
  space.
\item
  \textbf{Blue points:} Reconstructed observations using only PC1.
\item
  \textbf{Blue line:} First principal component axis.
\item
  \textbf{Dashed lines:} The differences between original and
  reconstructed observations.
\item
  \textbf{Mean:} The centre through which the PC1 axis passes.
\end{itemize}

Each original observation is projected onto the PC1 axis.

Because PC2 has been discarded, the reconstructed points cannot
reproduce deviations along the second principal component.

The differences between original and reconstructed observations
represent the \textbf{reconstruction error}.

\paragraph{5. Dimensionality at each
stage}\label{dimensionality-at-each-stage}

Suppose the original dataset contains 100 observations and two features.

{\def\LTcaptype{none} % do not increment counter
\begin{longtable}[]{@{}lll@{}}
\toprule\noalign{}
Python expression & Operation & Shape \\
\midrule\noalign{}
\endhead
\bottomrule\noalign{}
\endlastfoot
\texttt{X\_centered} & Original centred dataset & \texttt{(100,\ 2)} \\
\texttt{X\_eigen\_vectors} & Eigenvector matrix & \texttt{(2,\ 2)} \\
\texttt{Z} & PCA coordinates & \texttt{(100,\ 2)} \\
\texttt{pc1\_scores} & Only the PC1 coordinates & \texttt{(100,)} \\
\texttt{X\_reconstructed\_1d} & Reconstructed 2D dataset &
\texttt{(100,\ 2)} \\
\end{longtable}
}

The dimensionality changes as follows:

\[
\underbrace{2D}_{\text{Original}}
\;\longrightarrow\;
\underbrace{2D}_{\text{PCA projection}}
\;\longrightarrow\;
\underbrace{1D}_{\text{PC1 only}}
\;\longrightarrow\;
\underbrace{2D}_{\text{Approximate reconstruction}}
\]

Although the final reconstructed dataset contains two features, it still
has only \textbf{one independent dimension of variation}, because all
observations lie along the PC1 axis.

\paragraph{6. Alternative reconstruction using NumPy
broadcasting}\label{alternative-reconstruction-using-numpy-broadcasting}

Instead of:

\begin{Shaded}
\begin{Highlighting}[]
\NormalTok{X\_reconstructed\_1d }\OperatorTok{=}\NormalTok{ np.outer(}
\NormalTok{    pc1\_scores,}
\NormalTok{    X\_eigen\_vectors[:, }\DecValTok{0}\NormalTok{]}
\NormalTok{) }\OperatorTok{+}\NormalTok{ centre}
\end{Highlighting}
\end{Shaded}

We can use NumPy broadcasting:

\begin{Shaded}
\begin{Highlighting}[]
\NormalTok{X\_reconstructed\_1d }\OperatorTok{=}\NormalTok{ (}
\NormalTok{    pc1\_scores[:, np.newaxis] }\OperatorTok{*}\NormalTok{ X\_eigen\_vectors[:, }\DecValTok{0}\NormalTok{]}
    \OperatorTok{+}\NormalTok{ centre}
\NormalTok{)}
\end{Highlighting}
\end{Shaded}

Here, \texttt{np.newaxis} converts the PC1 scores from shape
\texttt{(n,)} into \texttt{(n,\ 1)}, allowing each score to multiply the
corresponding eigenvector coordinates.

Both implementations produce the same result.

\paragraph{7. Key takeaways}\label{key-takeaways}

\begin{itemize}
\tightlist
\item
  \textbf{Projection:} Converts original feature coordinates into
  principal component coordinates.
\item
  \textbf{Dimensionality reduction:} Discards less important principal
  components, retaining only those explaining the greatest variance.
\item
  \textbf{Reconstruction:} Uses the retained component scores and
  eigenvectors to approximate the original dataset.
\item
  \textbf{Mean restoration:} Adds the original feature means back after
  reconstruction.
\item
  \textbf{Information loss:} Discarded components cannot be recovered
  from the retained scores alone.
\end{itemize}

\textbf{Conclusion:} PCA allows us to represent a dataset using fewer
coordinates while preserving as much variance as possible.
Reconstructing the data from only PC1 produces the closest
one-dimensional linear approximation to the centred dataset in the
least-squares sense.


    \begin{tcolorbox}[breakable, size=fbox, boxrule=1pt, pad at break*=1mm,colback=cellbackground, colframe=cellborder]
\prompt{In}{incolor}{21}{\boxspacing}
\begin{Verbatim}[commandchars=\\\{\}]
\PY{n+nb}{print}\PY{p}{(}\PY{l+s+s2}{\PYZdq{}}\PY{l+s+s2}{Example for the first observation:}\PY{l+s+s2}{\PYZdq{}}\PY{p}{)}
\PY{n}{first\PYZus{}row\PYZus{}of\PYZus{}x\PYZus{}centered} \PY{o}{=} \PY{n}{X\PYZus{}centered}\PY{p}{[}\PY{l+m+mi}{0}\PY{p}{]}

\PY{n+nb}{print}\PY{p}{(}\PY{l+s+sa}{f}\PY{l+s+s2}{\PYZdq{}}\PY{l+s+s2}{First row of X centered: }\PY{l+s+si}{\PYZob{}}\PY{n}{first\PYZus{}row\PYZus{}of\PYZus{}x\PYZus{}centered}\PY{l+s+si}{\PYZcb{}}\PY{l+s+s2}{\PYZdq{}}\PY{p}{)}
\PY{n+nb}{print}\PY{p}{(}\PY{l+s+sa}{f}\PY{l+s+s2}{\PYZdq{}}\PY{l+s+se}{\PYZbs{}n}\PY{l+s+s2}{Eigen vector matrix:}\PY{l+s+se}{\PYZbs{}n}\PY{l+s+si}{\PYZob{}}\PY{n}{X\PYZus{}eigen\PYZus{}vectors}\PY{l+s+si}{\PYZcb{}}\PY{l+s+s2}{\PYZdq{}}\PY{p}{)}

\PY{n}{manual\PYZus{}1} \PY{o}{=} \PY{n}{first\PYZus{}row\PYZus{}of\PYZus{}x\PYZus{}centered}\PY{p}{[}\PY{l+m+mi}{0}\PY{p}{]}\PY{o}{*}\PY{n}{X\PYZus{}eigen\PYZus{}vectors}\PY{p}{[}\PY{l+m+mi}{0}\PY{p}{,}\PY{l+m+mi}{0}\PY{p}{]} \PY{o}{+} \PY{n}{first\PYZus{}row\PYZus{}of\PYZus{}x\PYZus{}centered}\PY{p}{[}\PY{l+m+mi}{1}\PY{p}{]}\PY{o}{*}\PY{n}{X\PYZus{}eigen\PYZus{}vectors}\PY{p}{[}\PY{l+m+mi}{1}\PY{p}{,}\PY{l+m+mi}{0}\PY{p}{]} \PY{c+c1}{\PYZsh{} Using the first column of X\PYZus{}eigen\PYZus{}vectors}
\PY{n}{manual\PYZus{}2} \PY{o}{=} \PY{n}{first\PYZus{}row\PYZus{}of\PYZus{}x\PYZus{}centered}\PY{p}{[}\PY{l+m+mi}{0}\PY{p}{]}\PY{o}{*}\PY{n}{X\PYZus{}eigen\PYZus{}vectors}\PY{p}{[}\PY{l+m+mi}{0}\PY{p}{,}\PY{l+m+mi}{1}\PY{p}{]} \PY{o}{+} \PY{n}{first\PYZus{}row\PYZus{}of\PYZus{}x\PYZus{}centered}\PY{p}{[}\PY{l+m+mi}{1}\PY{p}{]}\PY{o}{*}\PY{n}{X\PYZus{}eigen\PYZus{}vectors}\PY{p}{[}\PY{l+m+mi}{1}\PY{p}{,}\PY{l+m+mi}{1}\PY{p}{]} \PY{c+c1}{\PYZsh{} Using the second column of X\PYZus{}eigen\PYZus{}vectors}
\PY{n+nb}{print}\PY{p}{(}\PY{l+s+sa}{f}\PY{l+s+s2}{\PYZdq{}}\PY{l+s+se}{\PYZbs{}n}\PY{l+s+s2}{Manual values = [}\PY{l+s+si}{\PYZob{}}\PY{n}{manual\PYZus{}1}\PY{l+s+si}{:}\PY{l+s+s2}{.4f}\PY{l+s+si}{\PYZcb{}}\PY{l+s+s2}{, }\PY{l+s+si}{\PYZob{}}\PY{n}{manual\PYZus{}2}\PY{l+s+si}{:}\PY{l+s+s2}{.4f}\PY{l+s+si}{\PYZcb{}}\PY{l+s+s2}{]}\PY{l+s+s2}{\PYZdq{}}\PY{p}{)}
\PY{n}{multiply\PYZus{}first\PYZus{}and\PYZus{}eigen} \PY{o}{=} \PY{n}{first\PYZus{}row\PYZus{}of\PYZus{}x\PYZus{}centered} \PY{o}{@} \PY{n}{X\PYZus{}eigen\PYZus{}vectors}

\PY{n+nb}{print}\PY{p}{(}\PY{l+s+sa}{f}\PY{l+s+s2}{\PYZdq{}}\PY{l+s+se}{\PYZbs{}n}\PY{l+s+s2}{First row of X multiplied by eigen vectors:}\PY{l+s+se}{\PYZbs{}t}\PY{l+s+si}{\PYZob{}}\PY{n}{multiply\PYZus{}first\PYZus{}and\PYZus{}eigen}\PY{l+s+si}{\PYZcb{}}\PY{l+s+s2}{\PYZdq{}}\PY{p}{)}

\PY{n+nb}{print}\PY{p}{(}\PY{l+s+sa}{f}\PY{l+s+s2}{\PYZdq{}}\PY{l+s+s2}{Principal component scores:}\PY{l+s+se}{\PYZbs{}t}\PY{l+s+se}{\PYZbs{}t}\PY{l+s+se}{\PYZbs{}t}\PY{l+s+se}{\PYZbs{}t}\PY{l+s+se}{\PYZbs{}t}\PY{l+s+si}{\PYZob{}}\PY{n}{Z}\PY{p}{[}\PY{l+m+mi}{0}\PY{p}{]}\PY{l+s+si}{\PYZcb{}}\PY{l+s+s2}{\PYZdq{}}\PY{p}{)}

\PY{n+nb}{print}\PY{p}{(}\PY{l+s+s2}{\PYZdq{}}\PY{l+s+se}{\PYZbs{}n}\PY{l+s+s2}{Difference:}\PY{l+s+se}{\PYZbs{}t}\PY{l+s+se}{\PYZbs{}t}\PY{l+s+s2}{\PYZdq{}}\PY{p}{,} \PY{n}{multiply\PYZus{}first\PYZus{}and\PYZus{}eigen} \PY{o}{\PYZhy{}} \PY{n}{Z}\PY{p}{[}\PY{l+m+mi}{0}\PY{p}{]}\PY{p}{)}
\PY{n+nb}{print}\PY{p}{(}\PY{l+s+s2}{\PYZdq{}}\PY{l+s+s2}{Match:}\PY{l+s+se}{\PYZbs{}t}\PY{l+s+se}{\PYZbs{}t}\PY{l+s+se}{\PYZbs{}t}\PY{l+s+s2}{\PYZdq{}}\PY{p}{,} \PY{n}{np}\PY{o}{.}\PY{n}{allclose}\PY{p}{(}\PY{n}{multiply\PYZus{}first\PYZus{}and\PYZus{}eigen}\PY{p}{,} \PY{n}{Z}\PY{p}{[}\PY{l+m+mi}{0}\PY{p}{]}\PY{p}{)}\PY{p}{)}
\end{Verbatim}
\end{tcolorbox}

    \begin{Verbatim}[commandchars=\\\{\}]
Example for the first observation:
First row of X centered: [-0.3375  0.029 ]

Eigen vector matrix:
[[-0.8054  0.5928]
 [-0.5928 -0.8054]]

Manual values = [0.2547, -0.2234]

First row of X multiplied by eigen vectors:     [ 0.2547 -0.2234]
Principal component scores:                                     [ 0.2547
-0.2234]

Difference:              [0. 0.]
Match:                   True
    \end{Verbatim}


    \subsection{8. Check PCA against scikit-learn
(optional)}\label{check-pca-against-scikit-learn-optional}

The preceding calculations require only NumPy. If \texttt{scikit-learn}
is installed, this cell verifies its \texttt{PCA} implementation.
Eigenvector signs can differ, so compare their absolute values.


    \begin{tcolorbox}[breakable, size=fbox, boxrule=1pt, pad at break*=1mm,colback=cellbackground, colframe=cellborder]
\prompt{In}{incolor}{22}{\boxspacing}
\begin{Verbatim}[commandchars=\\\{\}]
\PY{k}{try}\PY{p}{:}
    \PY{k+kn}{from}\PY{+w}{ }\PY{n+nn}{sklearn}\PY{n+nn}{.}\PY{n+nn}{decomposition}\PY{+w}{ }\PY{k+kn}{import} \PY{n}{PCA}
    \PY{n}{pca} \PY{o}{=} \PY{n}{PCA}\PY{p}{(}\PY{n}{n\PYZus{}components}\PY{o}{=}\PY{l+m+mi}{2}\PY{p}{)}
    \PY{n}{pca}\PY{o}{.}\PY{n}{fit}\PY{p}{(}\PY{n}{X}\PY{p}{)}
    \PY{n+nb}{print}\PY{p}{(}\PY{l+s+s2}{\PYZdq{}}\PY{l+s+s2}{scikit\PYZhy{}learn eigenvalues:}\PY{l+s+s2}{\PYZdq{}}\PY{p}{,} \PY{n}{pca}\PY{o}{.}\PY{n}{explained\PYZus{}variance\PYZus{}}\PY{p}{)}
    \PY{n+nb}{print}\PY{p}{(}\PY{l+s+s2}{\PYZdq{}}\PY{l+s+s2}{NumPy eigenvalues:       }\PY{l+s+s2}{\PYZdq{}}\PY{p}{,} \PY{n}{X\PYZus{}eigen\PYZus{}values}\PY{p}{)}
    \PY{n+nb}{print}\PY{p}{(}\PY{l+s+s2}{\PYZdq{}}\PY{l+s+s2}{scikit\PYZhy{}learn explained variance ratios:}\PY{l+s+s2}{\PYZdq{}}\PY{p}{,} \PY{n}{pca}\PY{o}{.}\PY{n}{explained\PYZus{}variance\PYZus{}ratio\PYZus{}}\PY{p}{)}
    \PY{k}{assert} \PY{n}{np}\PY{o}{.}\PY{n}{allclose}\PY{p}{(}\PY{n}{pca}\PY{o}{.}\PY{n}{explained\PYZus{}variance\PYZus{}}\PY{p}{,} \PY{n}{X\PYZus{}eigen\PYZus{}values}\PY{p}{)}
\PY{k}{except} \PY{n+ne}{ImportError}\PY{p}{:}
    \PY{n+nb}{print}\PY{p}{(}\PY{l+s+s2}{\PYZdq{}}\PY{l+s+s2}{scikit\PYZhy{}learn is not installed; skip this optional verification.}\PY{l+s+s2}{\PYZdq{}}\PY{p}{)}
\end{Verbatim}
\end{tcolorbox}

    \begin{Verbatim}[commandchars=\\\{\}]
scikit-learn eigenvalues: [3.9636 0.29  ]
NumPy eigenvalues:        [3.9636 0.29  ]
scikit-learn explained variance ratios: [0.9318 0.0682]
    \end{Verbatim}


    \subsection{9. Quick practice exercises}\label{quick-practice-exercises}

\begin{enumerate}
\def\labelenumi{\arabic{enumi}.}
\tightlist
\item
  For \(X=[2,4,6]\), calculate the \textbf{sample variance} manually,
  then check with NumPy.
\item
  If \(Y=[6,4,2]\), what sign will \(\operatorname{Cov}(X,Y)\) have?
  Why?
\item
  For
  \(\mathbf A=\left[\begin{smallmatrix}4&0\\0&2\end{smallmatrix}\right]\),
  identify both eigenvalues and their corresponding eigenvectors.
\item
  In \hyperref[section-6]{Section 6}, what fraction of the variance does
  PC1 explain? What happens if you reverse the signs of its eigenvector?
\item
  Change the data covariance in \hyperref[section-6]{Section 6} to
  \texttt{{[}{[}2,\ 0{]},\ {[}0,\ 2{]}{]}}. Why might the principal
  directions now be unstable across samples?
\end{enumerate}

\textbf{Bonus:} Try standardising the columns of \texttt{X} and
comparing its covariance matrix, eigenvectors and explained variance
ratios with the unstandardised version.


    \begin{tcolorbox}[breakable, size=fbox, boxrule=1pt, pad at break*=1mm,colback=cellbackground, colframe=cellborder]
\prompt{In}{incolor}{23}{\boxspacing}
\begin{Verbatim}[commandchars=\\\{\}]
\PY{c+c1}{\PYZsh{} Exercise workspace}
\PY{c+c1}{\PYZsh{} Replace these with your own calculations.}
\PY{n}{practice\PYZus{}x} \PY{o}{=} \PY{n}{np}\PY{o}{.}\PY{n}{array}\PY{p}{(}\PY{p}{[}\PY{l+m+mf}{2.}\PY{p}{,} \PY{l+m+mf}{4.}\PY{p}{,} \PY{l+m+mf}{6.}\PY{p}{]}\PY{p}{)}
\PY{n}{practice\PYZus{}y} \PY{o}{=} \PY{n}{np}\PY{o}{.}\PY{n}{array}\PY{p}{(}\PY{p}{[}\PY{l+m+mf}{6.}\PY{p}{,} \PY{l+m+mf}{4.}\PY{p}{,} \PY{l+m+mf}{2.}\PY{p}{]}\PY{p}{)}
\PY{n+nb}{print}\PY{p}{(}\PY{l+s+s2}{\PYZdq{}}\PY{l+s+s2}{Sample variance of X:}\PY{l+s+s2}{\PYZdq{}}\PY{p}{,} \PY{n}{np}\PY{o}{.}\PY{n}{var}\PY{p}{(}\PY{n}{practice\PYZus{}x}\PY{p}{,} \PY{n}{ddof}\PY{o}{=}\PY{l+m+mi}{1}\PY{p}{)}\PY{p}{)}
\PY{n+nb}{print}\PY{p}{(}\PY{l+s+s2}{\PYZdq{}}\PY{l+s+s2}{Sample covariance of X and Y:}\PY{l+s+s2}{\PYZdq{}}\PY{p}{,} \PY{n}{np}\PY{o}{.}\PY{n}{cov}\PY{p}{(}\PY{n}{practice\PYZus{}x}\PY{p}{,} \PY{n}{practice\PYZus{}y}\PY{p}{)}\PY{p}{[}\PY{l+m+mi}{0}\PY{p}{,}\PY{l+m+mi}{1}\PY{p}{]}\PY{p}{)}
\end{Verbatim}
\end{tcolorbox}

    \begin{Verbatim}[commandchars=\\\{\}]
Sample variance of X: 4.0
Sample covariance of X and Y: -4.0
    \end{Verbatim}


    \subsubsection{Exercices completed}\label{exercices-completed}


    \begin{tcolorbox}[breakable, size=fbox, boxrule=1pt, pad at break*=1mm,colback=cellbackground, colframe=cellborder]
\prompt{In}{incolor}{24}{\boxspacing}
\begin{Verbatim}[commandchars=\\\{\}]
\PY{n}{display}\PY{p}{(}\PY{n}{Markdown}\PY{p}{(}\PY{l+s+s2}{\PYZdq{}}\PY{l+s+s2}{\PYZsh{}\PYZsh{}\PYZsh{} Exercice 1}\PY{l+s+s2}{\PYZdq{}}\PY{p}{)}\PY{p}{)}

\PY{n}{n\PYZus{}practice\PYZus{}x} \PY{o}{=} \PY{n+nb}{len}\PY{p}{(}\PY{n}{practice\PYZus{}x}\PY{p}{)}
\PY{n}{mean\PYZus{}practice\PYZus{}x} \PY{o}{=} \PY{n}{practice\PYZus{}x}\PY{o}{.}\PY{n}{mean}\PY{p}{(}\PY{p}{)}
\PY{n}{deviations\PYZus{}practice\PYZus{}x} \PY{o}{=} \PY{n}{practice\PYZus{}x} \PY{o}{\PYZhy{}} \PY{n}{mean\PYZus{}practice\PYZus{}x}
\PY{n}{squared\PYZus{}deviations\PYZus{}practice\PYZus{}x} \PY{o}{=} \PY{n}{deviations\PYZus{}practice\PYZus{}x}\PY{o}{*}\PY{o}{*}\PY{l+m+mi}{2}
\PY{n}{manual\PYZus{}variance\PYZus{}practice\PYZus{}x} \PY{o}{=} \PY{n}{squared\PYZus{}deviations\PYZus{}practice\PYZus{}x}\PY{o}{.}\PY{n}{sum}\PY{p}{(}\PY{p}{)} \PY{o}{/} \PY{p}{(}\PY{n}{n\PYZus{}practice\PYZus{}x} \PY{o}{\PYZhy{}} \PY{l+m+mi}{1}\PY{p}{)}
\PY{n}{sample\PYZus{}std\PYZus{}practice\PYZus{}x} \PY{o}{=} \PY{n}{np}\PY{o}{.}\PY{n}{std}\PY{p}{(}\PY{n}{practice\PYZus{}x}\PY{p}{,} \PY{n}{ddof}\PY{o}{=}\PY{l+m+mi}{1}\PY{p}{)}

\PY{n}{summary} \PY{o}{=} \PY{n}{pd}\PY{o}{.}\PY{n}{DataFrame}\PY{p}{(}\PY{p}{\PYZob{}}
    \PY{l+s+s2}{\PYZdq{}}\PY{l+s+s2}{X value}\PY{l+s+s2}{\PYZdq{}}\PY{p}{:} \PY{n}{practice\PYZus{}x}\PY{p}{,}
    \PY{l+s+s2}{\PYZdq{}}\PY{l+s+s2}{Deviation from mean}\PY{l+s+s2}{\PYZdq{}}\PY{p}{:} \PY{n}{deviations\PYZus{}practice\PYZus{}x}\PY{p}{,}
    \PY{l+s+s2}{\PYZdq{}}\PY{l+s+s2}{Squared deviation}\PY{l+s+s2}{\PYZdq{}}\PY{p}{:} \PY{n}{squared\PYZus{}deviations\PYZus{}practice\PYZus{}x}
\PY{p}{\PYZcb{}}\PY{p}{)}

\PY{c+c1}{\PYZsh{}display(summary.style.format(\PYZdq{}\PYZob{}:.0f\PYZcb{}\PYZdq{}))}
\PY{n}{display\PYZus{}dataframe}\PY{p}{(}\PY{n}{summary}\PY{p}{)}

\PY{c+c1}{\PYZsh{} Format the observations for LaTeX}
\PY{n}{values} \PY{o}{=} \PY{l+s+s2}{\PYZdq{}}\PY{l+s+s2}{+}\PY{l+s+s2}{\PYZdq{}}\PY{o}{.}\PY{n}{join}\PY{p}{(}\PY{l+s+sa}{f}\PY{l+s+s2}{\PYZdq{}}\PY{l+s+si}{\PYZob{}}\PY{n}{x}\PY{l+s+si}{:}\PY{l+s+s2}{g}\PY{l+s+si}{\PYZcb{}}\PY{l+s+s2}{\PYZdq{}} \PY{k}{for} \PY{n}{x} \PY{o+ow}{in} \PY{n}{practice\PYZus{}x}\PY{p}{)}
\PY{n}{deviations\PYZus{}practice\PYZus{}x} \PY{o}{=} \PY{l+s+s2}{\PYZdq{}}\PY{l+s+s2}{+}\PY{l+s+s2}{\PYZdq{}}\PY{o}{.}\PY{n}{join}\PY{p}{(}
    \PY{l+s+sa}{f}\PY{l+s+s2}{\PYZdq{}}\PY{l+s+s2}{(}\PY{l+s+si}{\PYZob{}}\PY{n}{x}\PY{l+s+si}{:}\PY{l+s+s2}{g}\PY{l+s+si}{\PYZcb{}}\PY{l+s+s2}{\PYZhy{}}\PY{l+s+si}{\PYZob{}}\PY{n}{mean\PYZus{}practice\PYZus{}x}\PY{l+s+si}{:}\PY{l+s+s2}{g}\PY{l+s+si}{\PYZcb{}}\PY{l+s+s2}{)\PYZca{}2}\PY{l+s+s2}{\PYZdq{}} \PY{k}{for} \PY{n}{x} \PY{o+ow}{in} \PY{n}{practice\PYZus{}x}
\PY{p}{)}
\PY{n}{squared\PYZus{}deviations\PYZus{}practice\PYZus{}x} \PY{o}{=} \PY{l+s+s2}{\PYZdq{}}\PY{l+s+s2}{+}\PY{l+s+s2}{\PYZdq{}}\PY{o}{.}\PY{n}{join}\PY{p}{(}
    \PY{l+s+sa}{f}\PY{l+s+s2}{\PYZdq{}}\PY{l+s+si}{\PYZob{}}\PY{p}{(}\PY{n}{x}\PY{+w}{ }\PY{o}{\PYZhy{}}\PY{+w}{ }\PY{n}{mean\PYZus{}practice\PYZus{}x}\PY{p}{)}\PY{o}{*}\PY{o}{*}\PY{l+m+mi}{2}\PY{l+s+si}{:}\PY{l+s+s2}{g}\PY{l+s+si}{\PYZcb{}}\PY{l+s+s2}{\PYZdq{}} \PY{k}{for} \PY{n}{x} \PY{o+ow}{in} \PY{n}{practice\PYZus{}x}
\PY{p}{)}

\PY{c+c1}{\PYZsh{} 1. Mean}
\PY{n}{display}\PY{p}{(}\PY{n}{Math}\PY{p}{(}
    \PY{l+s+sa}{rf}\PY{l+s+s2}{\PYZdq{}}\PY{l+s+s2}{\PYZbs{}}\PY{l+s+s2}{bar}\PY{l+s+se}{\PYZob{}\PYZob{}}\PY{l+s+s2}{x}\PY{l+s+se}{\PYZcb{}\PYZcb{}}\PY{l+s+s2}{ = }\PY{l+s+s2}{\PYZbs{}}\PY{l+s+s2}{frac}\PY{l+s+se}{\PYZob{}\PYZob{}}\PY{l+s+si}{\PYZob{}}\PY{n}{values}\PY{l+s+si}{\PYZcb{}}\PY{l+s+se}{\PYZcb{}\PYZcb{}}\PY{l+s+se}{\PYZob{}\PYZob{}}\PY{l+s+si}{\PYZob{}}\PY{n}{n\PYZus{}practice\PYZus{}x}\PY{l+s+si}{\PYZcb{}}\PY{l+s+se}{\PYZcb{}\PYZcb{}}\PY{l+s+s2}{ = }\PY{l+s+si}{\PYZob{}}\PY{n}{mean\PYZus{}practice\PYZus{}x}\PY{l+s+si}{:}\PY{l+s+s2}{.1f}\PY{l+s+si}{\PYZcb{}}\PY{l+s+s2}{\PYZdq{}}
\PY{p}{)}\PY{p}{)}

\PY{c+c1}{\PYZsh{} 2. Sample variance}
\PY{n}{display}\PY{p}{(}\PY{n}{Math}\PY{p}{(}
    \PY{l+s+sa}{rf}\PY{l+s+s2}{\PYZdq{}}\PY{l+s+s2}{\PYZbs{}}\PY{l+s+s2}{boxed}\PY{l+s+se}{\PYZob{}\PYZob{}}\PY{l+s+s2}{s\PYZca{}2 = }\PY{l+s+s2}{\PYZbs{}}\PY{l+s+s2}{frac}\PY{l+s+se}{\PYZob{}\PYZob{}}\PY{l+s+si}{\PYZob{}}\PY{n}{deviations\PYZus{}practice\PYZus{}x}\PY{l+s+si}{\PYZcb{}}\PY{l+s+se}{\PYZcb{}\PYZcb{}}\PY{l+s+se}{\PYZob{}\PYZob{}}\PY{l+s+si}{\PYZob{}}\PY{n}{n\PYZus{}practice\PYZus{}x}\PY{l+s+si}{\PYZcb{}}\PY{l+s+s2}{\PYZhy{}1}\PY{l+s+se}{\PYZcb{}\PYZcb{}}\PY{l+s+s2}{\PYZdq{}}
    \PY{l+s+sa}{rf}\PY{l+s+s2}{\PYZdq{}}\PY{l+s+s2}{ = }\PY{l+s+s2}{\PYZbs{}}\PY{l+s+s2}{frac}\PY{l+s+se}{\PYZob{}\PYZob{}}\PY{l+s+si}{\PYZob{}}\PY{n}{squared\PYZus{}deviations\PYZus{}practice\PYZus{}x}\PY{l+s+si}{\PYZcb{}}\PY{l+s+se}{\PYZcb{}\PYZcb{}}\PY{l+s+se}{\PYZob{}\PYZob{}}\PY{l+s+si}{\PYZob{}}\PY{n}{n\PYZus{}practice\PYZus{}x}\PY{o}{\PYZhy{}}\PY{l+m+mi}{1}\PY{l+s+si}{\PYZcb{}}\PY{l+s+se}{\PYZcb{}\PYZcb{}}\PY{l+s+s2}{\PYZdq{}}
    \PY{l+s+sa}{rf}\PY{l+s+s2}{\PYZdq{}}\PY{l+s+s2}{ = }\PY{l+s+si}{\PYZob{}}\PY{n}{manual\PYZus{}variance\PYZus{}practice\PYZus{}x}\PY{l+s+si}{:}\PY{l+s+s2}{.1f}\PY{l+s+se}{\PYZcb{}\PYZcb{}}\PY{l+s+si}{\PYZcb{}}\PY{l+s+s2}{\PYZdq{}}
\PY{p}{)}\PY{p}{)}

\PY{c+c1}{\PYZsh{} 3. NumPy sample variance}
\PY{n}{display}\PY{p}{(}\PY{n}{Math}\PY{p}{(}
    \PY{l+s+sa}{rf}\PY{l+s+s2}{\PYZdq{}}\PY{l+s+s2}{\PYZbs{}}\PY{l+s+s2}{operatorname}\PY{l+s+se}{\PYZob{}\PYZob{}}\PY{l+s+s2}{Var}\PY{l+s+se}{\PYZcb{}\PYZcb{}}\PY{l+s+s2}{\PYZus{}}\PY{l+s+se}{\PYZob{}\PYZob{}}\PY{l+s+s2}{\PYZbs{}}\PY{l+s+s2}{mathrm}\PY{l+s+se}{\PYZob{}\PYZob{}}\PY{l+s+s2}{NumPy}\PY{l+s+se}{\PYZcb{}\PYZcb{}}\PY{l+s+se}{\PYZcb{}\PYZcb{}}\PY{l+s+s2}{(X)}\PY{l+s+s2}{\PYZdq{}}
    \PY{l+s+sa}{rf}\PY{l+s+s2}{\PYZdq{}}\PY{l+s+s2}{ = }\PY{l+s+si}{\PYZob{}}\PY{n}{np}\PY{o}{.}\PY{n}{var}\PY{p}{(}\PY{n}{practice\PYZus{}x}\PY{p}{,}\PY{+w}{ }\PY{n}{ddof}\PY{o}{=}\PY{l+m+mi}{1}\PY{p}{)}\PY{l+s+si}{:}\PY{l+s+s2}{.1f}\PY{l+s+si}{\PYZcb{}}\PY{l+s+s2}{\PYZdq{}}
\PY{p}{)}\PY{p}{)}

\PY{c+c1}{\PYZsh{} 4. Sample standard deviation}
\PY{n}{display}\PY{p}{(}\PY{n}{Math}\PY{p}{(}
    \PY{l+s+sa}{rf}\PY{l+s+s2}{\PYZdq{}}\PY{l+s+s2}{s = }\PY{l+s+s2}{\PYZbs{}}\PY{l+s+s2}{sqrt}\PY{l+s+se}{\PYZob{}\PYZob{}}\PY{l+s+s2}{s\PYZca{}2}\PY{l+s+se}{\PYZcb{}\PYZcb{}}\PY{l+s+s2}{\PYZdq{}}
    \PY{l+s+sa}{rf}\PY{l+s+s2}{\PYZdq{}}\PY{l+s+s2}{ = }\PY{l+s+s2}{\PYZbs{}}\PY{l+s+s2}{sqrt}\PY{l+s+se}{\PYZob{}\PYZob{}}\PY{l+s+si}{\PYZob{}}\PY{n}{manual\PYZus{}variance\PYZus{}practice\PYZus{}x}\PY{l+s+si}{:}\PY{l+s+s2}{.1f}\PY{l+s+se}{\PYZcb{}\PYZcb{}}\PY{l+s+si}{\PYZcb{}}\PY{l+s+s2}{\PYZdq{}}
    \PY{l+s+sa}{rf}\PY{l+s+s2}{\PYZdq{}}\PY{l+s+s2}{ = }\PY{l+s+si}{\PYZob{}}\PY{n}{sample\PYZus{}std\PYZus{}practice\PYZus{}x}\PY{l+s+si}{:}\PY{l+s+s2}{.3f}\PY{l+s+si}{\PYZcb{}}\PY{l+s+s2}{\PYZdq{}}
\PY{p}{)}\PY{p}{)}
\end{Verbatim}
\end{tcolorbox}

    \subsubsection{Exercice 1}\label{exercice-1}

    
    \begin{tabular}{lrrr}
\toprule
 & X value & Deviation from mean & Squared deviation \\
\midrule
0 & 2 & -2 & 4 \\
1 & 4 & 0 & 0 \\
2 & 6 & 2 & 4 \\
\bottomrule
\end{tabular}


    
    $\displaystyle \bar{x} = \frac{2+4+6}{3} = 4.0$

    
    $\displaystyle \boxed{s^2 = \frac{(2-4)^2+(4-4)^2+(6-4)^2}{3-1} = \frac{4+0+4}{2} = 4.0}$

    
    $\displaystyle \operatorname{Var}_{\mathrm{NumPy}}(X) = 4.0$

    
    $\displaystyle s = \sqrt{s^2} = \sqrt{4.0} = 2.000$

    
    \begin{tcolorbox}[breakable, size=fbox, boxrule=1pt, pad at break*=1mm,colback=cellbackground, colframe=cellborder]
\prompt{In}{incolor}{25}{\boxspacing}
\begin{Verbatim}[commandchars=\\\{\}]
\PY{n}{display}\PY{p}{(}\PY{n}{Markdown}\PY{p}{(}\PY{l+s+s2}{\PYZdq{}}\PY{l+s+s2}{\PYZsh{}\PYZsh{}\PYZsh{} Exercice 2}\PY{l+s+s2}{\PYZdq{}}\PY{p}{)}\PY{p}{)}

\PY{n}{n\PYZus{}practice\PYZus{}y} \PY{o}{=} \PY{n+nb}{len}\PY{p}{(}\PY{n}{practice\PYZus{}y}\PY{p}{)}
\PY{n}{mean\PYZus{}practice\PYZus{}y} \PY{o}{=} \PY{n}{practice\PYZus{}y}\PY{o}{.}\PY{n}{mean}\PY{p}{(}\PY{p}{)}
\PY{n}{deviations\PYZus{}practice\PYZus{}y} \PY{o}{=} \PY{n}{practice\PYZus{}y} \PY{o}{\PYZhy{}} \PY{n}{mean\PYZus{}practice\PYZus{}y}
\PY{n}{squared\PYZus{}deviations\PYZus{}practice\PYZus{}y} \PY{o}{=} \PY{n}{deviations\PYZus{}practice\PYZus{}y}\PY{o}{*}\PY{o}{*}\PY{l+m+mi}{2}
\PY{n}{manual\PYZus{}variance\PYZus{}practice\PYZus{}y} \PY{o}{=} \PY{n}{squared\PYZus{}deviations\PYZus{}practice\PYZus{}y}\PY{o}{.}\PY{n}{sum}\PY{p}{(}\PY{p}{)} \PY{o}{/} \PY{p}{(}\PY{n}{n\PYZus{}practice\PYZus{}y} \PY{o}{\PYZhy{}} \PY{l+m+mi}{1}\PY{p}{)}
\PY{n}{sample\PYZus{}std\PYZus{}practice\PYZus{}y} \PY{o}{=} \PY{n}{np}\PY{o}{.}\PY{n}{std}\PY{p}{(}\PY{n}{practice\PYZus{}y}\PY{p}{,} \PY{n}{ddof}\PY{o}{=}\PY{l+m+mi}{1}\PY{p}{)}

\PY{n}{summary} \PY{o}{=} \PY{n}{pd}\PY{o}{.}\PY{n}{DataFrame}\PY{p}{(}\PY{p}{\PYZob{}}
    \PY{l+s+s2}{\PYZdq{}}\PY{l+s+s2}{X value}\PY{l+s+s2}{\PYZdq{}}\PY{p}{:} \PY{n}{practice\PYZus{}x}\PY{p}{,}
    \PY{l+s+s2}{\PYZdq{}}\PY{l+s+s2}{Deviation from mean X }\PY{l+s+s2}{\PYZdq{}}\PY{p}{:} \PY{p}{(}\PY{n}{practice\PYZus{}x} \PY{o}{\PYZhy{}} \PY{n}{mean\PYZus{}practice\PYZus{}x}\PY{p}{)}\PY{p}{,}
    \PY{l+s+s2}{\PYZdq{}}\PY{l+s+s2}{Squared deviation X}\PY{l+s+s2}{\PYZdq{}}\PY{p}{:} \PY{p}{(}\PY{n}{practice\PYZus{}x} \PY{o}{\PYZhy{}} \PY{n}{mean\PYZus{}practice\PYZus{}x}\PY{p}{)}\PY{o}{*}\PY{o}{*}\PY{l+m+mi}{2}\PY{p}{,}
    \PY{l+s+s2}{\PYZdq{}}\PY{l+s+s2}{Y value}\PY{l+s+s2}{\PYZdq{}}\PY{p}{:} \PY{n}{practice\PYZus{}y}\PY{p}{,}
    \PY{l+s+s2}{\PYZdq{}}\PY{l+s+s2}{Deviation from mean Y}\PY{l+s+s2}{\PYZdq{}}\PY{p}{:} \PY{n}{deviations\PYZus{}practice\PYZus{}y}\PY{p}{,}
    \PY{l+s+s2}{\PYZdq{}}\PY{l+s+s2}{Squared deviation Y}\PY{l+s+s2}{\PYZdq{}}\PY{p}{:} \PY{n}{squared\PYZus{}deviations\PYZus{}practice\PYZus{}y}
\PY{p}{\PYZcb{}}\PY{p}{)}

\PY{c+c1}{\PYZsh{}display(summary.style.format(\PYZdq{}\PYZob{}:.0f\PYZcb{}\PYZdq{}))}
\PY{n}{display\PYZus{}dataframe}\PY{p}{(}\PY{n}{summary}\PY{p}{)}

\PY{c+c1}{\PYZsh{} Format the observations for LaTeX}
\PY{n}{values} \PY{o}{=} \PY{l+s+s2}{\PYZdq{}}\PY{l+s+s2}{+}\PY{l+s+s2}{\PYZdq{}}\PY{o}{.}\PY{n}{join}\PY{p}{(}\PY{l+s+sa}{f}\PY{l+s+s2}{\PYZdq{}}\PY{l+s+si}{\PYZob{}}\PY{n}{x}\PY{l+s+si}{:}\PY{l+s+s2}{g}\PY{l+s+si}{\PYZcb{}}\PY{l+s+s2}{\PYZdq{}} \PY{k}{for} \PY{n}{x} \PY{o+ow}{in} \PY{n}{practice\PYZus{}y}\PY{p}{)}
\PY{n}{deviations\PYZus{}practice\PYZus{}y} \PY{o}{=} \PY{l+s+s2}{\PYZdq{}}\PY{l+s+s2}{+}\PY{l+s+s2}{\PYZdq{}}\PY{o}{.}\PY{n}{join}\PY{p}{(}
    \PY{l+s+sa}{f}\PY{l+s+s2}{\PYZdq{}}\PY{l+s+s2}{(}\PY{l+s+si}{\PYZob{}}\PY{n}{x}\PY{l+s+si}{:}\PY{l+s+s2}{g}\PY{l+s+si}{\PYZcb{}}\PY{l+s+s2}{\PYZhy{}}\PY{l+s+si}{\PYZob{}}\PY{n}{mean\PYZus{}practice\PYZus{}y}\PY{l+s+si}{:}\PY{l+s+s2}{g}\PY{l+s+si}{\PYZcb{}}\PY{l+s+s2}{)\PYZca{}2}\PY{l+s+s2}{\PYZdq{}} \PY{k}{for} \PY{n}{x} \PY{o+ow}{in} \PY{n}{practice\PYZus{}y}
\PY{p}{)}
\PY{n}{squared\PYZus{}deviations\PYZus{}practice\PYZus{}y} \PY{o}{=} \PY{l+s+s2}{\PYZdq{}}\PY{l+s+s2}{+}\PY{l+s+s2}{\PYZdq{}}\PY{o}{.}\PY{n}{join}\PY{p}{(}
    \PY{l+s+sa}{f}\PY{l+s+s2}{\PYZdq{}}\PY{l+s+si}{\PYZob{}}\PY{p}{(}\PY{n}{x}\PY{+w}{ }\PY{o}{\PYZhy{}}\PY{+w}{ }\PY{n}{mean\PYZus{}practice\PYZus{}y}\PY{p}{)}\PY{o}{*}\PY{o}{*}\PY{l+m+mi}{2}\PY{l+s+si}{:}\PY{l+s+s2}{g}\PY{l+s+si}{\PYZcb{}}\PY{l+s+s2}{\PYZdq{}} \PY{k}{for} \PY{n}{x} \PY{o+ow}{in} \PY{n}{practice\PYZus{}y}
\PY{p}{)}

\PY{c+c1}{\PYZsh{} 1. Mean}
\PY{n}{display}\PY{p}{(}\PY{n}{Math}\PY{p}{(}
    \PY{l+s+sa}{rf}\PY{l+s+s2}{\PYZdq{}}\PY{l+s+s2}{\PYZbs{}}\PY{l+s+s2}{bar}\PY{l+s+se}{\PYZob{}\PYZob{}}\PY{l+s+s2}{y}\PY{l+s+se}{\PYZcb{}\PYZcb{}}\PY{l+s+s2}{ = }\PY{l+s+s2}{\PYZbs{}}\PY{l+s+s2}{frac}\PY{l+s+se}{\PYZob{}\PYZob{}}\PY{l+s+si}{\PYZob{}}\PY{n}{values}\PY{l+s+si}{\PYZcb{}}\PY{l+s+se}{\PYZcb{}\PYZcb{}}\PY{l+s+se}{\PYZob{}\PYZob{}}\PY{l+s+si}{\PYZob{}}\PY{n}{n\PYZus{}practice\PYZus{}y}\PY{l+s+si}{\PYZcb{}}\PY{l+s+se}{\PYZcb{}\PYZcb{}}\PY{l+s+s2}{ = }\PY{l+s+si}{\PYZob{}}\PY{n}{mean\PYZus{}practice\PYZus{}y}\PY{l+s+si}{:}\PY{l+s+s2}{.1f}\PY{l+s+si}{\PYZcb{}}\PY{l+s+s2}{\PYZdq{}}
\PY{p}{)}\PY{p}{)}

\PY{c+c1}{\PYZsh{} 2. Sample variance}
\PY{n}{display}\PY{p}{(}\PY{n}{Math}\PY{p}{(}
    \PY{l+s+sa}{rf}\PY{l+s+s2}{\PYZdq{}}\PY{l+s+s2}{s\PYZca{}2 = }\PY{l+s+s2}{\PYZbs{}}\PY{l+s+s2}{frac}\PY{l+s+se}{\PYZob{}\PYZob{}}\PY{l+s+si}{\PYZob{}}\PY{n}{deviations\PYZus{}practice\PYZus{}y}\PY{l+s+si}{\PYZcb{}}\PY{l+s+se}{\PYZcb{}\PYZcb{}}\PY{l+s+se}{\PYZob{}\PYZob{}}\PY{l+s+si}{\PYZob{}}\PY{n}{n\PYZus{}practice\PYZus{}y}\PY{l+s+si}{\PYZcb{}}\PY{l+s+s2}{\PYZhy{}1}\PY{l+s+se}{\PYZcb{}\PYZcb{}}\PY{l+s+s2}{\PYZdq{}}
    \PY{l+s+sa}{rf}\PY{l+s+s2}{\PYZdq{}}\PY{l+s+s2}{ = }\PY{l+s+s2}{\PYZbs{}}\PY{l+s+s2}{frac}\PY{l+s+se}{\PYZob{}\PYZob{}}\PY{l+s+si}{\PYZob{}}\PY{n}{squared\PYZus{}deviations\PYZus{}practice\PYZus{}y}\PY{l+s+si}{\PYZcb{}}\PY{l+s+se}{\PYZcb{}\PYZcb{}}\PY{l+s+se}{\PYZob{}\PYZob{}}\PY{l+s+si}{\PYZob{}}\PY{n}{n\PYZus{}practice\PYZus{}y}\PY{o}{\PYZhy{}}\PY{l+m+mi}{1}\PY{l+s+si}{\PYZcb{}}\PY{l+s+se}{\PYZcb{}\PYZcb{}}\PY{l+s+s2}{\PYZdq{}}
    \PY{l+s+sa}{rf}\PY{l+s+s2}{\PYZdq{}}\PY{l+s+s2}{ = }\PY{l+s+si}{\PYZob{}}\PY{n}{manual\PYZus{}variance\PYZus{}practice\PYZus{}y}\PY{l+s+si}{:}\PY{l+s+s2}{.1f}\PY{l+s+si}{\PYZcb{}}\PY{l+s+s2}{\PYZdq{}}
\PY{p}{)}\PY{p}{)}

\PY{c+c1}{\PYZsh{} 3. NumPy sample variance}
\PY{n}{display}\PY{p}{(}\PY{n}{Math}\PY{p}{(}
    \PY{l+s+sa}{rf}\PY{l+s+s2}{\PYZdq{}}\PY{l+s+s2}{\PYZbs{}}\PY{l+s+s2}{operatorname}\PY{l+s+se}{\PYZob{}\PYZob{}}\PY{l+s+s2}{Var}\PY{l+s+se}{\PYZcb{}\PYZcb{}}\PY{l+s+s2}{\PYZus{}}\PY{l+s+se}{\PYZob{}\PYZob{}}\PY{l+s+s2}{\PYZbs{}}\PY{l+s+s2}{mathrm}\PY{l+s+se}{\PYZob{}\PYZob{}}\PY{l+s+s2}{NumPy}\PY{l+s+se}{\PYZcb{}\PYZcb{}}\PY{l+s+se}{\PYZcb{}\PYZcb{}}\PY{l+s+s2}{(X)}\PY{l+s+s2}{\PYZdq{}}
    \PY{l+s+sa}{rf}\PY{l+s+s2}{\PYZdq{}}\PY{l+s+s2}{ = }\PY{l+s+si}{\PYZob{}}\PY{n}{np}\PY{o}{.}\PY{n}{var}\PY{p}{(}\PY{n}{practice\PYZus{}y}\PY{p}{,}\PY{+w}{ }\PY{n}{ddof}\PY{o}{=}\PY{l+m+mi}{1}\PY{p}{)}\PY{l+s+si}{:}\PY{l+s+s2}{.1f}\PY{l+s+si}{\PYZcb{}}\PY{l+s+s2}{\PYZdq{}}
\PY{p}{)}\PY{p}{)}

\PY{c+c1}{\PYZsh{} 4. Sample standard deviation}
\PY{n}{display}\PY{p}{(}\PY{n}{Math}\PY{p}{(}
    \PY{l+s+sa}{rf}\PY{l+s+s2}{\PYZdq{}}\PY{l+s+s2}{s = }\PY{l+s+s2}{\PYZbs{}}\PY{l+s+s2}{sqrt}\PY{l+s+se}{\PYZob{}\PYZob{}}\PY{l+s+s2}{s\PYZca{}2}\PY{l+s+se}{\PYZcb{}\PYZcb{}}\PY{l+s+s2}{\PYZdq{}}
    \PY{l+s+sa}{rf}\PY{l+s+s2}{\PYZdq{}}\PY{l+s+s2}{ = }\PY{l+s+s2}{\PYZbs{}}\PY{l+s+s2}{sqrt}\PY{l+s+se}{\PYZob{}\PYZob{}}\PY{l+s+si}{\PYZob{}}\PY{n}{manual\PYZus{}variance\PYZus{}practice\PYZus{}y}\PY{l+s+si}{:}\PY{l+s+s2}{.1f}\PY{l+s+se}{\PYZcb{}\PYZcb{}}\PY{l+s+si}{\PYZcb{}}\PY{l+s+s2}{\PYZdq{}}
    \PY{l+s+sa}{rf}\PY{l+s+s2}{\PYZdq{}}\PY{l+s+s2}{ = }\PY{l+s+si}{\PYZob{}}\PY{n}{sample\PYZus{}std\PYZus{}practice\PYZus{}y}\PY{l+s+si}{:}\PY{l+s+s2}{.3f}\PY{l+s+si}{\PYZcb{}}\PY{l+s+s2}{\PYZdq{}}
\PY{p}{)}\PY{p}{)}

\PY{c+c1}{\PYZsh{} 5. Covariance formula}
\PY{n}{display}\PY{p}{(}\PY{n}{Math}\PY{p}{(}
    \PY{l+s+sa}{rf}\PY{l+s+s2}{\PYZdq{}}\PY{l+s+s2}{\PYZbs{}}\PY{l+s+s2}{operatorname}\PY{l+s+se}{\PYZob{}\PYZob{}}\PY{l+s+s2}{Cov}\PY{l+s+se}{\PYZcb{}\PYZcb{}}\PY{l+s+s2}{(X,Y)=}\PY{l+s+s2}{\PYZbs{}}\PY{l+s+s2}{frac}\PY{l+s+se}{\PYZob{}\PYZob{}}\PY{l+s+s2}{1}\PY{l+s+se}{\PYZcb{}\PYZcb{}}\PY{l+s+se}{\PYZob{}\PYZob{}}\PY{l+s+s2}{n\PYZhy{}1}\PY{l+s+se}{\PYZcb{}\PYZcb{}}\PY{l+s+s2}{\PYZbs{}}\PY{l+s+s2}{sum\PYZus{}}\PY{l+s+se}{\PYZob{}\PYZob{}}\PY{l+s+s2}{i=1}\PY{l+s+se}{\PYZcb{}\PYZcb{}}\PY{l+s+s2}{\PYZca{}n(x\PYZus{}i\PYZhy{}}\PY{l+s+s2}{\PYZbs{}}\PY{l+s+s2}{bar}\PY{l+s+se}{\PYZob{}\PYZob{}}\PY{l+s+s2}{x}\PY{l+s+se}{\PYZcb{}\PYZcb{}}\PY{l+s+s2}{)(y\PYZus{}i\PYZhy{}}\PY{l+s+s2}{\PYZbs{}}\PY{l+s+s2}{bar}\PY{l+s+se}{\PYZob{}\PYZob{}}\PY{l+s+s2}{y}\PY{l+s+se}{\PYZcb{}\PYZcb{}}\PY{l+s+s2}{)}\PY{l+s+s2}{\PYZdq{}}
\PY{p}{)}\PY{p}{)}

\PY{n}{covariance\PYZus{}terms} \PY{o}{=} \PY{l+s+s2}{\PYZdq{}}\PY{l+s+s2}{ + }\PY{l+s+s2}{\PYZdq{}}\PY{o}{.}\PY{n}{join}\PY{p}{(}
    \PY{l+s+sa}{rf}\PY{l+s+s2}{\PYZdq{}}\PY{l+s+s2}{(}\PY{l+s+si}{\PYZob{}}\PY{n}{h}\PY{l+s+si}{:}\PY{l+s+s2}{g}\PY{l+s+si}{\PYZcb{}}\PY{l+s+s2}{\PYZhy{}}\PY{l+s+si}{\PYZob{}}\PY{n}{mean\PYZus{}practice\PYZus{}x}\PY{l+s+si}{:}\PY{l+s+s2}{g}\PY{l+s+si}{\PYZcb{}}\PY{l+s+s2}{)(}\PY{l+s+si}{\PYZob{}}\PY{n}{s}\PY{l+s+si}{:}\PY{l+s+s2}{g}\PY{l+s+si}{\PYZcb{}}\PY{l+s+s2}{\PYZhy{}}\PY{l+s+si}{\PYZob{}}\PY{n}{mean\PYZus{}practice\PYZus{}y}\PY{l+s+si}{:}\PY{l+s+s2}{g}\PY{l+s+si}{\PYZcb{}}\PY{l+s+s2}{)}\PY{l+s+s2}{\PYZdq{}}
    \PY{k}{for} \PY{n}{h}\PY{p}{,} \PY{n}{s} \PY{o+ow}{in} \PY{n+nb}{zip}\PY{p}{(}\PY{n}{practice\PYZus{}x}\PY{p}{,} \PY{n}{practice\PYZus{}y}\PY{p}{)}
\PY{p}{)}

\PY{n}{covariance\PYZus{}terms\PYZus{}after} \PY{o}{=} \PY{l+s+s2}{\PYZdq{}}\PY{l+s+s2}{ + }\PY{l+s+s2}{\PYZdq{}}\PY{o}{.}\PY{n}{join}\PY{p}{(}
    \PY{l+s+sa}{rf}\PY{l+s+s2}{\PYZdq{}}\PY{l+s+s2}{(}\PY{l+s+si}{\PYZob{}}\PY{p}{(}\PY{n}{h}\PY{o}{\PYZhy{}}\PY{n}{mean\PYZus{}practice\PYZus{}x}\PY{p}{)}\PY{l+s+si}{:}\PY{l+s+s2}{g}\PY{l+s+si}{\PYZcb{}}\PY{l+s+s2}{)(}\PY{l+s+si}{\PYZob{}}\PY{p}{(}\PY{n}{s}\PY{o}{\PYZhy{}}\PY{n}{mean\PYZus{}practice\PYZus{}y}\PY{p}{)}\PY{l+s+si}{:}\PY{l+s+s2}{g}\PY{l+s+si}{\PYZcb{}}\PY{l+s+s2}{)}\PY{l+s+s2}{\PYZdq{}}
    \PY{k}{for} \PY{n}{h}\PY{p}{,} \PY{n}{s} \PY{o+ow}{in} \PY{n+nb}{zip}\PY{p}{(}\PY{n}{practice\PYZus{}x}\PY{p}{,} \PY{n}{practice\PYZus{}y}\PY{p}{)}
\PY{p}{)}

\PY{n}{display}\PY{p}{(}\PY{n}{Math}\PY{p}{(}
    \PY{l+s+sa}{r}\PY{l+s+s2}{\PYZdq{}}\PY{l+s+s2}{\PYZbs{}}\PY{l+s+s2}{operatorname}\PY{l+s+si}{\PYZob{}Cov\PYZcb{}}\PY{l+s+s2}{(X,Y)}\PY{l+s+s2}{\PYZdq{}}
    \PY{l+s+sa}{r}\PY{l+s+s2}{\PYZdq{}}\PY{l+s+s2}{ = }\PY{l+s+s2}{\PYZbs{}}\PY{l+s+s2}{frac}\PY{l+s+s2}{\PYZob{}}\PY{l+s+s2}{\PYZbs{}}\PY{l+s+s2}{sum\PYZus{}}\PY{l+s+s2}{\PYZob{}}\PY{l+s+s2}{i=1\PYZcb{}\PYZca{}}\PY{l+s+si}{\PYZob{}n\PYZcb{}}\PY{l+s+s2}{(x\PYZus{}i\PYZhy{}}\PY{l+s+s2}{\PYZbs{}}\PY{l+s+s2}{bar}\PY{l+s+si}{\PYZob{}x\PYZcb{}}\PY{l+s+s2}{)(y\PYZus{}i\PYZhy{}}\PY{l+s+s2}{\PYZbs{}}\PY{l+s+s2}{bar}\PY{l+s+si}{\PYZob{}y\PYZcb{}}\PY{l+s+s2}{)\PYZcb{}}\PY{l+s+s2}{\PYZob{}}\PY{l+s+s2}{n\PYZhy{}1\PYZcb{}}\PY{l+s+s2}{\PYZdq{}}
\PY{p}{)}\PY{p}{)}

\PY{n}{products} \PY{o}{=} \PY{p}{(}\PY{n}{practice\PYZus{}x} \PY{o}{\PYZhy{}} \PY{n}{practice\PYZus{}x}\PY{o}{.}\PY{n}{mean}\PY{p}{(}\PY{p}{)}\PY{p}{)} \PY{o}{*} \PY{p}{(}\PY{n}{practice\PYZus{}y} \PY{o}{\PYZhy{}} \PY{n}{practice\PYZus{}x}\PY{o}{.}\PY{n}{mean}\PY{p}{(}\PY{p}{)}\PY{p}{)}
\PY{n}{manual\PYZus{}covariance} \PY{o}{=} \PY{n}{products}\PY{o}{.}\PY{n}{sum}\PY{p}{(}\PY{p}{)} \PY{o}{/} \PY{p}{(}\PY{n}{n\PYZus{}practice\PYZus{}x}\PY{o}{\PYZhy{}}\PY{l+m+mi}{1}\PY{p}{)}


\PY{n}{display}\PY{p}{(}\PY{n}{Math}\PY{p}{(}
    \PY{l+s+sa}{rf}\PY{l+s+s2}{\PYZdq{}}\PY{l+s+s2}{\PYZbs{}}\PY{l+s+s2}{begin}\PY{l+s+se}{\PYZob{}\PYZob{}}\PY{l+s+s2}{aligned}\PY{l+s+se}{\PYZcb{}\PYZcb{}}\PY{l+s+s2}{\PYZdq{}}
    \PY{l+s+sa}{rf}\PY{l+s+s2}{\PYZdq{}}\PY{l+s+s2}{\PYZbs{}}\PY{l+s+s2}{boxed}\PY{l+s+se}{\PYZob{}\PYZob{}}\PY{l+s+s2}{\PYZbs{}}\PY{l+s+s2}{operatorname}\PY{l+s+se}{\PYZob{}\PYZob{}}\PY{l+s+s2}{Cov}\PY{l+s+se}{\PYZcb{}\PYZcb{}}\PY{l+s+s2}{(}\PY{l+s+s2}{\PYZbs{}}\PY{l+s+s2}{mathrm}\PY{l+s+se}{\PYZob{}\PYZob{}}\PY{l+s+s2}{hours}\PY{l+s+se}{\PYZcb{}\PYZcb{}}\PY{l+s+s2}{,}\PY{l+s+s2}{\PYZbs{}}\PY{l+s+s2}{mathrm}\PY{l+s+se}{\PYZob{}\PYZob{}}\PY{l+s+s2}{scores}\PY{l+s+se}{\PYZcb{}\PYZcb{}}\PY{l+s+s2}{)}\PY{l+s+se}{\PYZcb{}\PYZcb{}}\PY{l+s+s2}{\PYZdq{}}
    \PY{l+s+sa}{rf}\PY{l+s+s2}{\PYZdq{}}\PY{l+s+s2}{ \PYZam{}= }\PY{l+s+s2}{\PYZbs{}}\PY{l+s+s2}{frac}\PY{l+s+se}{\PYZob{}\PYZob{}}\PY{l+s+si}{\PYZob{}}\PY{n}{covariance\PYZus{}terms}\PY{l+s+si}{\PYZcb{}}\PY{l+s+se}{\PYZcb{}\PYZcb{}}\PY{l+s+se}{\PYZob{}\PYZob{}}\PY{l+s+si}{\PYZob{}}\PY{n}{n\PYZus{}practice\PYZus{}x}\PY{l+s+si}{\PYZcb{}}\PY{l+s+s2}{\PYZhy{}1}\PY{l+s+se}{\PYZcb{}\PYZcb{}}\PY{l+s+s2}{ }\PY{l+s+se}{\PYZbs{}\PYZbs{}}\PY{l+s+s2}{[6pt]}\PY{l+s+s2}{\PYZdq{}}
    \PY{l+s+sa}{rf}\PY{l+s+s2}{\PYZdq{}}\PY{l+s+s2}{ \PYZam{}= }\PY{l+s+s2}{\PYZbs{}}\PY{l+s+s2}{frac}\PY{l+s+se}{\PYZob{}\PYZob{}}\PY{l+s+si}{\PYZob{}}\PY{n}{covariance\PYZus{}terms\PYZus{}after}\PY{l+s+si}{\PYZcb{}}\PY{l+s+se}{\PYZcb{}\PYZcb{}}\PY{l+s+se}{\PYZob{}\PYZob{}}\PY{l+s+si}{\PYZob{}}\PY{n}{n\PYZus{}practice\PYZus{}x}\PY{l+s+si}{\PYZcb{}}\PY{l+s+s2}{\PYZhy{}1}\PY{l+s+se}{\PYZcb{}\PYZcb{}}\PY{l+s+s2}{\PYZdq{}}
    \PY{l+s+sa}{rf}\PY{l+s+s2}{\PYZdq{}}\PY{l+s+s2}{ = }\PY{l+s+s2}{\PYZbs{}}\PY{l+s+s2}{frac}\PY{l+s+se}{\PYZob{}\PYZob{}}\PY{l+s+si}{\PYZob{}}\PY{n}{latex\PYZus{}sum}\PY{p}{(}\PY{n}{products}\PY{p}{)}\PY{l+s+si}{\PYZcb{}}\PY{l+s+se}{\PYZcb{}\PYZcb{}}\PY{l+s+se}{\PYZob{}\PYZob{}}\PY{l+s+si}{\PYZob{}}\PY{n}{n\PYZus{}practice\PYZus{}x}\PY{o}{\PYZhy{}}\PY{l+m+mi}{1}\PY{l+s+si}{\PYZcb{}}\PY{l+s+se}{\PYZcb{}\PYZcb{}}\PY{l+s+s2}{\PYZdq{}}
    \PY{l+s+sa}{rf}\PY{l+s+s2}{\PYZdq{}}\PY{l+s+s2}{ = }\PY{l+s+s2}{\PYZbs{}}\PY{l+s+s2}{frac}\PY{l+s+se}{\PYZob{}\PYZob{}}\PY{l+s+si}{\PYZob{}}\PY{n}{products}\PY{o}{.}\PY{n}{sum}\PY{p}{(}\PY{p}{)}\PY{l+s+si}{:}\PY{l+s+s2}{g}\PY{l+s+se}{\PYZcb{}\PYZcb{}}\PY{l+s+si}{\PYZcb{}}\PY{l+s+se}{\PYZob{}\PYZob{}}\PY{l+s+si}{\PYZob{}}\PY{n}{n\PYZus{}practice\PYZus{}x}\PY{o}{\PYZhy{}}\PY{l+m+mi}{1}\PY{l+s+si}{\PYZcb{}}\PY{l+s+se}{\PYZcb{}\PYZcb{}}\PY{l+s+s2}{\PYZdq{}}
    \PY{l+s+sa}{rf}\PY{l+s+s2}{\PYZdq{}}\PY{l+s+s2}{ = }\PY{l+s+s2}{\PYZbs{}}\PY{l+s+s2}{boxed}\PY{l+s+se}{\PYZob{}\PYZob{}}\PY{l+s+si}{\PYZob{}}\PY{n}{manual\PYZus{}covariance}\PY{l+s+si}{:}\PY{l+s+s2}{g}\PY{l+s+se}{\PYZcb{}\PYZcb{}}\PY{l+s+si}{\PYZcb{}}\PY{l+s+s2}{\PYZdq{}}
    \PY{l+s+sa}{rf}\PY{l+s+s2}{\PYZdq{}}\PY{l+s+s2}{\PYZbs{}}\PY{l+s+s2}{end}\PY{l+s+se}{\PYZob{}\PYZob{}}\PY{l+s+s2}{aligned}\PY{l+s+se}{\PYZcb{}\PYZcb{}}\PY{l+s+s2}{\PYZdq{}}
\PY{p}{)}\PY{p}{)}
\end{Verbatim}
\end{tcolorbox}

    \subsubsection{Exercice 2}\label{exercice-2}

    
    \begin{tabular}{lrrrrrr}
\toprule
 & X value & Deviation from mean X  & Squared deviation X & Y value & Deviation from mean Y & Squared deviation Y \\
\midrule
0 & 2 & -2 & 4 & 6 & 2 & 4 \\
1 & 4 & 0 & 0 & 4 & 0 & 0 \\
2 & 6 & 2 & 4 & 2 & -2 & 4 \\
\bottomrule
\end{tabular}


    
    $\displaystyle \bar{y} = \frac{6+4+2}{3} = 4.0$

    
    $\displaystyle s^2 = \frac{(6-4)^2+(4-4)^2+(2-4)^2}{3-1} = \frac{4+0+4}{2} = 4.0$

    
    $\displaystyle \operatorname{Var}_{\mathrm{NumPy}}(X) = 4.0$

    
    $\displaystyle s = \sqrt{s^2} = \sqrt{4.0} = 2.000$

    
    $\displaystyle \operatorname{Cov}(X,Y)=\frac{1}{n-1}\sum_{i=1}^n(x_i-\bar{x})(y_i-\bar{y})$

    
    $\displaystyle \operatorname{Cov}(X,Y) = \frac{\sum_{i=1}^{n}(x_i-\bar{x})(y_i-\bar{y})}{n-1}$

    
    $\displaystyle \begin{aligned}\boxed{\operatorname{Cov}(\mathrm{hours},\mathrm{scores})} &= \frac{(2-4)(6-4) + (4-4)(4-4) + (6-4)(2-4)}{3-1} \\[6pt] &= \frac{(-2)(2) + (0)(0) + (2)(-2)}{3-1} = \frac{-4 + 0 + -4}{2} = \frac{-8}{2} = \boxed{-4}\end{aligned}$

    
    \begin{tcolorbox}[breakable, size=fbox, boxrule=1pt, pad at break*=1mm,colback=cellbackground, colframe=cellborder]
\prompt{In}{incolor}{26}{\boxspacing}
\begin{Verbatim}[commandchars=\\\{\}]
\PY{n}{data} \PY{o}{=} \PY{n}{pd}\PY{o}{.}\PY{n}{DataFrame}\PY{p}{(}\PY{p}{\PYZob{}}\PY{l+s+s2}{\PYZdq{}}\PY{l+s+s2}{X}\PY{l+s+s2}{\PYZdq{}}\PY{p}{:} \PY{n}{practice\PYZus{}x}\PY{p}{,} \PY{l+s+s2}{\PYZdq{}}\PY{l+s+s2}{Y}\PY{l+s+s2}{\PYZdq{}}\PY{p}{:} \PY{n}{practice\PYZus{}y}\PY{p}{\PYZcb{}}\PY{p}{)}
\PY{n+nb}{print}\PY{p}{(}\PY{l+s+s2}{\PYZdq{}}\PY{l+s+s2}{Sample covariance matrix:}\PY{l+s+s2}{\PYZdq{}}\PY{p}{)}
\PY{n}{display}\PY{p}{(}\PY{n}{data}\PY{o}{.}\PY{n}{cov}\PY{p}{(}\PY{p}{)}\PY{p}{)}
\PY{n+nb}{print}\PY{p}{(}\PY{l+s+s2}{\PYZdq{}}\PY{l+s+s2}{Correlation matrix:}\PY{l+s+s2}{\PYZdq{}}\PY{p}{)}
\PY{n}{display}\PY{p}{(}\PY{n}{data}\PY{o}{.}\PY{n}{corr}\PY{p}{(}\PY{p}{)}\PY{p}{)}
\PY{n+nb}{print}\PY{p}{(}\PY{l+s+s2}{\PYZdq{}}\PY{l+s+s2}{np.cov agrees with pandas:}\PY{l+s+se}{\PYZbs{}n}\PY{l+s+s2}{\PYZdq{}}\PY{p}{,} \PY{n}{np}\PY{o}{.}\PY{n}{cov}\PY{p}{(}\PY{n}{practice\PYZus{}x}\PY{p}{,} \PY{n}{practice\PYZus{}y}\PY{p}{)}\PY{p}{)}
\PY{n+nb}{print}\PY{p}{(}\PY{l+s+s2}{\PYZdq{}}\PY{l+s+s2}{np.corrcoef agrees with pandas:}\PY{l+s+se}{\PYZbs{}n}\PY{l+s+s2}{\PYZdq{}}\PY{p}{,} \PY{n}{np}\PY{o}{.}\PY{n}{corrcoef}\PY{p}{(}\PY{n}{practice\PYZus{}x}\PY{p}{,} \PY{n}{practice\PYZus{}y}\PY{p}{)}\PY{p}{)}

\PY{n}{display}\PY{p}{(}\PY{n}{Markdown}\PY{p}{(}\PY{l+s+sa}{f}\PY{l+s+s2}{\PYZdq{}}\PY{l+s+s2}{**The covariance value calculated by numpy is }\PY{l+s+si}{\PYZob{}}\PY{n}{np}\PY{o}{.}\PY{n}{cov}\PY{p}{(}\PY{n}{practice\PYZus{}x}\PY{p}{,}\PY{+w}{ }\PY{n}{practice\PYZus{}y}\PY{p}{)}\PY{p}{[}\PY{l+m+mi}{0}\PY{p}{,}\PY{l+m+mi}{1}\PY{p}{]}\PY{l+s+si}{:}\PY{l+s+s2}{g}\PY{l+s+si}{\PYZcb{}}\PY{l+s+s2}{**}\PY{l+s+s2}{\PYZdq{}}\PY{p}{)}\PY{p}{)}
\end{Verbatim}
\end{tcolorbox}

    \begin{Verbatim}[commandchars=\\\{\}]
Sample covariance matrix:
    \end{Verbatim}

    
    \begin{Verbatim}[commandchars=\\\{\}]
     X    Y
X  4.0 -4.0
Y -4.0  4.0
    \end{Verbatim}

    
    \begin{Verbatim}[commandchars=\\\{\}]
Correlation matrix:
    \end{Verbatim}

    
    \begin{Verbatim}[commandchars=\\\{\}]
     X    Y
X  1.0 -1.0
Y -1.0  1.0
    \end{Verbatim}

    
    \begin{Verbatim}[commandchars=\\\{\}]
np.cov agrees with pandas:
 [[ 4. -4.]
 [-4.  4.]]
np.corrcoef agrees with pandas:
 [[ 1. -1.]
 [-1.  1.]]
    \end{Verbatim}

    \textbf{The covariance value calculated by numpy is -4}

    
    \begin{tcolorbox}[breakable, size=fbox, boxrule=1pt, pad at break*=1mm,colback=cellbackground, colframe=cellborder]
\prompt{In}{incolor}{27}{\boxspacing}
\begin{Verbatim}[commandchars=\\\{\}]
\PY{k+kn}{import}\PY{+w}{ }\PY{n+nn}{numpy}\PY{+w}{ }\PY{k}{as}\PY{+w}{ }\PY{n+nn}{np}
\PY{k+kn}{import}\PY{+w}{ }\PY{n+nn}{matplotlib}\PY{n+nn}{.}\PY{n+nn}{pyplot}\PY{+w}{ }\PY{k}{as}\PY{+w}{ }\PY{n+nn}{plt}

\PY{n}{practice\PYZus{}x} \PY{o}{=} \PY{n}{np}\PY{o}{.}\PY{n}{array}\PY{p}{(}\PY{p}{[}\PY{l+m+mf}{2.}\PY{p}{,} \PY{l+m+mf}{4.}\PY{p}{,} \PY{l+m+mf}{6.}\PY{p}{]}\PY{p}{)}
\PY{n}{practice\PYZus{}y} \PY{o}{=} \PY{n}{np}\PY{o}{.}\PY{n}{array}\PY{p}{(}\PY{p}{[}\PY{l+m+mf}{6.}\PY{p}{,} \PY{l+m+mf}{4.}\PY{p}{,} \PY{l+m+mf}{2.}\PY{p}{]}\PY{p}{)}

\PY{n}{fig}\PY{p}{,} \PY{n}{ax} \PY{o}{=} \PY{n}{plt}\PY{o}{.}\PY{n}{subplots}\PY{p}{(}\PY{n}{figsize}\PY{o}{=}\PY{p}{(}\PY{l+m+mi}{6}\PY{p}{,} \PY{l+m+mi}{5}\PY{p}{)}\PY{p}{)}

\PY{c+c1}{\PYZsh{} Plot observations}
\PY{n}{ax}\PY{o}{.}\PY{n}{scatter}\PY{p}{(}
    \PY{n}{practice\PYZus{}x}\PY{p}{,} \PY{n}{practice\PYZus{}y}\PY{p}{,}
    \PY{n}{color}\PY{o}{=}\PY{l+s+s2}{\PYZdq{}}\PY{l+s+s2}{royalblue}\PY{l+s+s2}{\PYZdq{}}\PY{p}{,}
    \PY{n}{s}\PY{o}{=}\PY{l+m+mi}{100}\PY{p}{,}
    \PY{n}{label}\PY{o}{=}\PY{l+s+s2}{\PYZdq{}}\PY{l+s+s2}{Observations}\PY{l+s+s2}{\PYZdq{}}\PY{p}{,}
    \PY{n}{zorder}\PY{o}{=}\PY{l+m+mi}{3}
\PY{p}{)}

\PY{c+c1}{\PYZsh{} Connect the observations}
\PY{n}{ax}\PY{o}{.}\PY{n}{plot}\PY{p}{(}
    \PY{n}{practice\PYZus{}x}\PY{p}{,} \PY{n}{practice\PYZus{}y}\PY{p}{,}
    \PY{n}{color}\PY{o}{=}\PY{l+s+s2}{\PYZdq{}}\PY{l+s+s2}{darkorange}\PY{l+s+s2}{\PYZdq{}}\PY{p}{,}
    \PY{n}{linestyle}\PY{o}{=}\PY{l+s+s2}{\PYZdq{}}\PY{l+s+s2}{\PYZhy{}\PYZhy{}}\PY{l+s+s2}{\PYZdq{}}\PY{p}{,}
    \PY{n}{label}\PY{o}{=}\PY{l+s+s2}{\PYZdq{}}\PY{l+s+s2}{Linear relationship}\PY{l+s+s2}{\PYZdq{}}
\PY{p}{)}

\PY{c+c1}{\PYZsh{} Label each observation}
\PY{k}{for} \PY{n}{x}\PY{p}{,} \PY{n}{y} \PY{o+ow}{in} \PY{n+nb}{zip}\PY{p}{(}\PY{n}{practice\PYZus{}x}\PY{p}{,} \PY{n}{practice\PYZus{}y}\PY{p}{)}\PY{p}{:}
    \PY{n}{ax}\PY{o}{.}\PY{n}{annotate}\PY{p}{(}
        \PY{l+s+sa}{f}\PY{l+s+s2}{\PYZdq{}}\PY{l+s+s2}{(}\PY{l+s+si}{\PYZob{}}\PY{n}{x}\PY{l+s+si}{:}\PY{l+s+s2}{g}\PY{l+s+si}{\PYZcb{}}\PY{l+s+s2}{, }\PY{l+s+si}{\PYZob{}}\PY{n}{y}\PY{l+s+si}{:}\PY{l+s+s2}{g}\PY{l+s+si}{\PYZcb{}}\PY{l+s+s2}{)}\PY{l+s+s2}{\PYZdq{}}\PY{p}{,}
        \PY{p}{(}\PY{n}{x}\PY{p}{,} \PY{n}{y}\PY{p}{)}\PY{p}{,}
        \PY{n}{xytext}\PY{o}{=}\PY{p}{(}\PY{l+m+mi}{8}\PY{p}{,} \PY{l+m+mi}{8}\PY{p}{)}\PY{p}{,}
        \PY{n}{textcoords}\PY{o}{=}\PY{l+s+s2}{\PYZdq{}}\PY{l+s+s2}{offset points}\PY{l+s+s2}{\PYZdq{}}
    \PY{p}{)}

\PY{n}{ax}\PY{o}{.}\PY{n}{set}\PY{p}{(}
    \PY{n}{xlabel}\PY{o}{=}\PY{l+s+s2}{\PYZdq{}}\PY{l+s+s2}{X}\PY{l+s+s2}{\PYZdq{}}\PY{p}{,}
    \PY{n}{ylabel}\PY{o}{=}\PY{l+s+s2}{\PYZdq{}}\PY{l+s+s2}{Y}\PY{l+s+s2}{\PYZdq{}}\PY{p}{,}
    \PY{n}{title}\PY{o}{=}\PY{l+s+s2}{\PYZdq{}}\PY{l+s+s2}{Negative Linear Correlation}\PY{l+s+s2}{\PYZdq{}}
\PY{p}{)}

\PY{n}{ax}\PY{o}{.}\PY{n}{set\PYZus{}xlim}\PY{p}{(}\PY{l+m+mi}{0}\PY{p}{,} \PY{l+m+mi}{8}\PY{p}{)}
\PY{n}{ax}\PY{o}{.}\PY{n}{set\PYZus{}ylim}\PY{p}{(}\PY{l+m+mi}{0}\PY{p}{,} \PY{l+m+mi}{8}\PY{p}{)}
\PY{n}{ax}\PY{o}{.}\PY{n}{set\PYZus{}aspect}\PY{p}{(}\PY{l+s+s2}{\PYZdq{}}\PY{l+s+s2}{equal}\PY{l+s+s2}{\PYZdq{}}\PY{p}{)}
\PY{n}{ax}\PY{o}{.}\PY{n}{grid}\PY{p}{(}\PY{n}{alpha}\PY{o}{=}\PY{l+m+mf}{0.2}\PY{p}{)}
\PY{n}{ax}\PY{o}{.}\PY{n}{legend}\PY{p}{(}\PY{p}{)}

\PY{n}{plt}\PY{o}{.}\PY{n}{show}\PY{p}{(}\PY{p}{)}
\end{Verbatim}
\end{tcolorbox}

    \begin{center}
    \adjustimage{max size={0.9\linewidth}{0.9\paperheight}}{week4_pdf_build_pca_files/week4_pdf_build_pca_44_0.png}
    \end{center}
    { \hspace*{\fill} \\}
    

    The covariance is \textbf{negative (-4)} because X and Y move in
opposite directions. As X increases, Y decreases, producing a
downward-sloping line in the scatter plot.

Mathematically, the deviations from their respective means have opposite
signs, resulting in negative products when calculating covariance.

Since the relationship is perfectly linear and decreasing, the Pearson
correlation coefficient is \textbf{-1}.



    \subsubsection{Exercice 3}\label{exercice-3}

\subsubsection{Calculating the eigenvalues of a
matrix}\label{calculating-the-eigenvalues-of-a-matrix}

Given the matrix:

\[
\mathbf{A} =
\begin{bmatrix}
4 & 0 \\
0 & 2
\end{bmatrix}
\]

We want to find the eigenvalues \(\lambda\) of \(\mathbf{A}\).

\paragraph{\texorpdfstring{Step 1: Calculate
\(\mathbf{A} - \lambda \mathbf{I}\)}{Step 1: Calculate \textbackslash mathbf\{A\} - \textbackslash lambda \textbackslash mathbf\{I\}}}\label{step-1-calculate-mathbfa---lambda-mathbfi}

The identity matrix is:

\[
\mathbf{I} =
\begin{bmatrix}
1 & 0 \\
0 & 1
\end{bmatrix}
\]

Multiplying the identity matrix by \(\lambda\):

\[
\lambda \mathbf{I} =
\begin{bmatrix}
\lambda & 0 \\
0 & \lambda
\end{bmatrix}
\]

Subtracting \(\lambda \mathbf{I}\) from \(\mathbf{A}\):

\[
\mathbf{A} - \lambda \mathbf{I}
=
\begin{bmatrix}
4 & 0 \\
0 & 2
\end{bmatrix}
-
\begin{bmatrix}
\lambda & 0 \\
0 & \lambda
\end{bmatrix}
\]

Therefore:

\[
\boxed{
\mathbf{A} - \lambda \mathbf{I}
=
\begin{bmatrix}
4-\lambda & 0 \\
0 & 2-\lambda
\end{bmatrix}
}
\]

\paragraph{Step 2: Compute the
determinant}\label{step-2-compute-the-determinant}

To find the eigenvalues, solve the \textbf{characteristic equation}:

\[
\det(\mathbf{A} - \lambda \mathbf{I}) = 0
\]

For a \(2 \times 2\) matrix:

\[
\det
\begin{bmatrix}
a & b \\
c & d
\end{bmatrix}
= ad - bc
\]

Applying this formula:

\[
\det
\begin{bmatrix}
4-\lambda & 0 \\
0 & 2-\lambda
\end{bmatrix}
= 0
\]

\[
(4-\lambda)(2-\lambda) - (0 \times 0) = 0
\]

Simplifying:

\[
(4-\lambda)(2-\lambda) = 0
\]

Expanding:

\[
8 - 4\lambda - 2\lambda + \lambda^2 = 0
\]

Rearranging into standard quadratic form:

\[
\boxed{\lambda^2 - 6\lambda + 8 = 0}
\]

\paragraph{Step 3: Solve the quadratic
equation}\label{step-3-solve-the-quadratic-equation}

We now solve:

\[
\lambda^2 - 6\lambda + 8 = 0
\]

The general quadratic equation is:

\[
ax^2 + bx + c = 0
\]

Its solutions are given by:

\[
x = \frac{-b \pm \sqrt{b^2 - 4ac}}{2a}
\]

In our case:

\begin{itemize}
\tightlist
\item
  \(a = 1\)
\item
  \(b = -6\)
\item
  \(c = 8\)
\end{itemize}

Substituting:

\[
\lambda =
\frac{-(-6) \pm \sqrt{(-6)^2 - 4(1)(8)}}{2(1)}
\]

Simplifying:

\[
\lambda =
\frac{6 \pm \sqrt{36-32}}{2}
\]

\[
\lambda =
\frac{6 \pm \sqrt{4}}{2}
\]

\[
\lambda =
\frac{6 \pm 2}{2}
\]

\textbf{First eigenvalue:}

\[
\lambda_1 = \frac{6+2}{2} = 4
\]

\textbf{Second eigenvalue:}

\[
\lambda_2 = \frac{6-2}{2} = 2
\]

Therefore, the eigenvalues are:

\[
\boxed{\lambda_1 = 4 \qquad \lambda_2 = 2}
\]

\begin{center}\rule{0.5\linewidth}{0.5pt}\end{center}

\paragraph{Shortcut: Eigenvalues of a diagonal
matrix}\label{shortcut-eigenvalues-of-a-diagonal-matrix}

Our original matrix is:

\[
\mathbf{A} =
\begin{bmatrix}
4 & 0 \\
0 & 2
\end{bmatrix}
\]

This is a \textbf{diagonal matrix}, meaning all elements outside the
main diagonal are zero.

For any diagonal matrix:

\[
\mathbf{D} =
\begin{bmatrix}
d_1 & 0 & 0 \\
0 & d_2 & 0 \\
0 & 0 & d_3
\end{bmatrix}
\]

The eigenvalues are simply its diagonal elements:

\[
\lambda_1 = d_1, \qquad
\lambda_2 = d_2, \qquad
\lambda_3 = d_3
\]

Therefore, for our matrix:

\[
\boxed{\lambda_1 = 4 \qquad \lambda_2 = 2}
\]

\textbf{Key takeaway:} For a diagonal matrix, you can read the
eigenvalues directly from the main diagonal without calculating the
determinant or solving a quadratic equation.

This shortcut applies to diagonal matrices of any size.

\subsubsection{Calculating the eigenvectors of a
matrix}\label{calculating-the-eigenvectors-of-a-matrix}

Given the matrix:

\[
\mathbf{A} =
\begin{bmatrix}
4 & 0 \\
0 & 2
\end{bmatrix}
\]

We have already calculated its eigenvalues:

\[
\lambda_1 = 4, \qquad \lambda_2 = 2
\]

Now we need to find the \textbf{eigenvector corresponding to each
eigenvalue}.

\paragraph{Step 1: The eigenvector
equation}\label{step-1-the-eigenvector-equation}

An eigenvector \(\mathbf{v}\) satisfies:

\[
\mathbf{A}\mathbf{v} = \lambda\mathbf{v}
\]

Rearranging:

\[
\mathbf{A}\mathbf{v} - \lambda\mathbf{v} = 0
\]

Factoring out \(\mathbf{v}\) using the identity matrix:

\[
\boxed{(\mathbf{A} - \lambda\mathbf{I})\mathbf{v} = \mathbf{0}}
\]

For a \(2 \times 2\) matrix, the eigenvector has two components:

\[
\mathbf{v} =
\begin{bmatrix}
x \\
y
\end{bmatrix}
\]

We solve this equation separately for each eigenvalue.

\begin{center}\rule{0.5\linewidth}{0.5pt}\end{center}

\paragraph{\texorpdfstring{Step 2: Find the eigenvector for
\(\lambda_1 = 4\)}{Step 2: Find the eigenvector for \textbackslash lambda\_1 = 4}}\label{step-2-find-the-eigenvector-for-lambda_1-4}

Substitute \(\lambda_1 = 4\):

\[
(\mathbf{A} - 4\mathbf{I})\mathbf{v}_1 = \mathbf{0}
\]

Expanding:

\[
\left(
\begin{bmatrix}
4 & 0 \\
0 & 2
\end{bmatrix}
-
\begin{bmatrix}
4 & 0 \\
0 & 4
\end{bmatrix}
\right)
\begin{bmatrix}
x \\
y
\end{bmatrix}
=
\begin{bmatrix}
0 \\
0
\end{bmatrix}
\]

Subtracting the matrices:

\[
\begin{bmatrix}
0 & 0 \\
0 & -2
\end{bmatrix}
\begin{bmatrix}
x \\
y
\end{bmatrix}
=
\begin{bmatrix}
0 \\
0
\end{bmatrix}
\]

Multiplying the matrix by the vector gives two equations:

\[
\begin{cases}
0x + 0y = 0 \\
0x - 2y = 0
\end{cases}
\]

From the second equation:

\[
-2y = 0 \implies y = 0
\]

The first equation is always satisfied, meaning \textbf{\(x\) can be any
real number except zero}.

We can choose \(x = 1\) for simplicity.

Therefore:

\[
\boxed{
\mathbf{v}_1 =
\begin{bmatrix}
1 \\
0
\end{bmatrix}
}
\]

Any nonzero scalar multiple of this vector is also an eigenvector:

\[
\mathbf{v}_1 =
\begin{bmatrix}
t \\
0
\end{bmatrix},
\qquad t \neq 0
\]

\paragraph{\texorpdfstring{Step 3: Find the eigenvector for
\(\lambda_2 = 2\)}{Step 3: Find the eigenvector for \textbackslash lambda\_2 = 2}}\label{step-3-find-the-eigenvector-for-lambda_2-2}

Substitute \(\lambda_2 = 2\):

\[
(\mathbf{A} - 2\mathbf{I})\mathbf{v}_2 = \mathbf{0}
\]

Expanding:

\[
\left(
\begin{bmatrix}
4 & 0 \\
0 & 2
\end{bmatrix}
-
\begin{bmatrix}
2 & 0 \\
0 & 2
\end{bmatrix}
\right)
\begin{bmatrix}
x \\
y
\end{bmatrix}
=
\begin{bmatrix}
0 \\
0
\end{bmatrix}
\]

Subtracting the matrices:

\[
\begin{bmatrix}
2 & 0 \\
0 & 0
\end{bmatrix}
\begin{bmatrix}
x \\
y
\end{bmatrix}
=
\begin{bmatrix}
0 \\
0
\end{bmatrix}
\]

Multiplying the matrix by the vector gives:

\[
\begin{cases}
2x + 0y = 0 \\
0x + 0y = 0
\end{cases}
\]

From the first equation:

\[
2x = 0 \implies x = 0
\]

The second equation is always satisfied, meaning \textbf{\(y\) can be
any real number except zero}.

We can choose \(y = 1\) for simplicity.

Therefore:

\[
\boxed{
\mathbf{v}_2 =
\begin{bmatrix}
0 \\
1
\end{bmatrix}
}
\]

Any nonzero scalar multiple is also an eigenvector:

\[
\mathbf{v}_2 =
\begin{bmatrix}
0 \\
t
\end{bmatrix},
\qquad t \neq 0
\]

\begin{center}\rule{0.5\linewidth}{0.5pt}\end{center}

\paragraph{Step 4: Verify the results}\label{step-4-verify-the-results}

We can verify each eigenvector using:

\[
\mathbf{A}\mathbf{v} = \lambda\mathbf{v}
\]

\textbf{For \(\lambda_1 = 4\):}

\[
\begin{bmatrix}
4 & 0 \\
0 & 2
\end{bmatrix}
\begin{bmatrix}
1 \\
0
\end{bmatrix}
=
\begin{bmatrix}
4 \\
0
\end{bmatrix}
=
4
\begin{bmatrix}
1 \\
0
\end{bmatrix}
\]

The equation is satisfied.

\textbf{For \(\lambda_2 = 2\):}

\[
\begin{bmatrix}
4 & 0 \\
0 & 2
\end{bmatrix}
\begin{bmatrix}
0 \\
1
\end{bmatrix}
=
\begin{bmatrix}
0 \\
2
\end{bmatrix}
=
2
\begin{bmatrix}
0 \\
1
\end{bmatrix}
\]

The equation is also satisfied.

\paragraph{Addendum: Why can we choose any nonzero
value?}\label{addendum-why-can-we-choose-any-nonzero-value}

An eigenvector satisfies:

\[
A\mathbf{v} = \lambda\mathbf{v}
\]

The eigenvalue is the \textbf{factor by which the vector is scaled}, not
a coordinate of the resulting vector.

For example, instead of choosing \(x = 1\), choose \(x = 3\):

\[
A =
\begin{bmatrix}
4 & 0 \\
0 & 2
\end{bmatrix},
\qquad
\mathbf{v} =
\begin{bmatrix}
3 \\
0
\end{bmatrix}
\]

Multiplying:

\[
A\mathbf{v}
=
\begin{bmatrix}
12 \\
0
\end{bmatrix}
=
4
\begin{bmatrix}
3 \\
0
\end{bmatrix}
=
4\mathbf{v}
\]

The eigenvalue is still \textbf{4}, not 12: the original coordinate was
3 and became 12, so the scaling factor is \(12/3 = 4\).

More generally, for any \(t \neq 0\):

\[
A
\begin{bmatrix}
t \\
0
\end{bmatrix}
=
\begin{bmatrix}
4t \\
0
\end{bmatrix}
=
4
\begin{bmatrix}
t \\
0
\end{bmatrix}
\]

Similarly, for the eigenvalue \(\lambda = 2\):

\[
A
\begin{bmatrix}
0 \\
t
\end{bmatrix}
=
\begin{bmatrix}
0 \\
2t
\end{bmatrix}
=
2
\begin{bmatrix}
0 \\
t
\end{bmatrix}
\]

Therefore, choosing a different nonzero value changes the eigenvector's
length, and possibly its orientation along the same line, but
\textbf{does not change its eigenvalue}.

We choose \(t = 1\) for convenience: it gives a vector of length 1. We
exclude \(t = 0\) because the zero vector is not an eigenvector.

\paragraph{Final results}\label{final-results}

The eigenvalues and their corresponding eigenvectors are:

\[
\boxed{
\lambda_1 = 4
\quad \longrightarrow \quad
\mathbf{v}_1 =
\begin{bmatrix}
1 \\
0
\end{bmatrix}
}
\]

\[
\boxed{
\lambda_2 = 2
\quad \longrightarrow \quad
\mathbf{v}_2 =
\begin{bmatrix}
0 \\
1
\end{bmatrix}
}
\]

\begin{center}\rule{0.5\linewidth}{0.5pt}\end{center}

\paragraph{Shortcut: Eigenvectors of a diagonal
matrix}\label{shortcut-eigenvectors-of-a-diagonal-matrix}

For a diagonal matrix with \textbf{distinct diagonal values}, each
eigenvalue corresponds to a standard basis vector.

For our matrix:

\[
\mathbf{A} =
\begin{bmatrix}
4 & 0 \\
0 & 2
\end{bmatrix}
\]

\begin{itemize}
\tightlist
\item
  The first diagonal element, \(4\), corresponds to the horizontal
  direction \(\begin{bmatrix}1\\0\end{bmatrix}\).
\item
  The second diagonal element, \(2\), corresponds to the vertical
  direction \(\begin{bmatrix}0\\1\end{bmatrix}\).
\end{itemize}

\textbf{Geometric interpretation:}

The matrix \(\mathbf{A}\) transforms vectors by:

\begin{itemize}
\tightlist
\item
  Stretching the horizontal (\(x\)) direction by a factor of \textbf{4}.
\item
  Stretching the vertical (\(y\)) direction by a factor of \textbf{2}.
\end{itemize}

The eigenvectors identify directions that remain unchanged under the
transformation, while the eigenvalues indicate how much vectors along
those directions are scaled.

\textbf{Important:} Eigenvectors must be nonzero vectors, but they do
not need to have length 1. For diagonal matrices with repeated
eigenvalues, the corresponding eigenvectors are not restricted to
individual standard basis vectors.


    \begin{tcolorbox}[breakable, size=fbox, boxrule=1pt, pad at break*=1mm,colback=cellbackground, colframe=cellborder]
\prompt{In}{incolor}{28}{\boxspacing}
\begin{Verbatim}[commandchars=\\\{\}]
\PY{n}{A} \PY{o}{=} \PY{n}{np}\PY{o}{.}\PY{n}{array}\PY{p}{(}\PY{p}{[}\PY{p}{[}\PY{l+m+mi}{4}\PY{p}{,} \PY{l+m+mi}{0}\PY{p}{]}\PY{p}{,} \PY{p}{[}\PY{l+m+mi}{0}\PY{p}{,} \PY{l+m+mi}{2}\PY{p}{]}\PY{p}{]}\PY{p}{)}

\PY{c+c1}{\PYZsh{} Calculate eigenvalues and eigenvectors directly}
\PY{n}{values\PYZus{}A}\PY{p}{,} \PY{n}{vectors\PYZus{}A} \PY{o}{=} \PY{n}{np}\PY{o}{.}\PY{n}{linalg}\PY{o}{.}\PY{n}{eigh}\PY{p}{(}\PY{n}{A}\PY{p}{)}

\PY{c+c1}{\PYZsh{} Sort eigenvalues in descending order}
\PY{n}{order\PYZus{}A} \PY{o}{=} \PY{n}{np}\PY{o}{.}\PY{n}{argsort}\PY{p}{(}\PY{n}{values\PYZus{}A}\PY{p}{)}\PY{p}{[}\PY{p}{:}\PY{p}{:}\PY{o}{\PYZhy{}}\PY{l+m+mi}{1}\PY{p}{]}
\PY{n}{values\PYZus{}A} \PY{o}{=} \PY{n}{values\PYZus{}A}\PY{p}{[}\PY{n}{order\PYZus{}A}\PY{p}{]}
\PY{n}{vectors\PYZus{}A} \PY{o}{=} \PY{n}{vectors\PYZus{}A}\PY{p}{[}\PY{p}{:}\PY{p}{,} \PY{n}{order\PYZus{}A}\PY{p}{]}

\PY{n+nb}{print}\PY{p}{(}\PY{l+s+s2}{\PYZdq{}}\PY{l+s+s2}{Eigenvalues (descending):}\PY{l+s+s2}{\PYZdq{}}\PY{p}{,} \PY{n}{values\PYZus{}A}\PY{p}{)}
\PY{n+nb}{print}\PY{p}{(}\PY{l+s+s2}{\PYZdq{}}\PY{l+s+s2}{Eigenvectors (columns):}\PY{l+s+se}{\PYZbs{}n}\PY{l+s+s2}{\PYZdq{}}\PY{p}{,} \PY{n}{vectors\PYZus{}A}\PY{p}{)}

\PY{c+c1}{\PYZsh{} Verify A @ v = lambda * v}
\PY{c+c1}{\PYZsh{} np.diag(values\PYZus{}A) converts the eigenvalues into a diagonal matrix}
\PY{k}{assert} \PY{n}{np}\PY{o}{.}\PY{n}{allclose}\PY{p}{(}\PY{n}{A} \PY{o}{@} \PY{n}{vectors\PYZus{}A}\PY{p}{,} \PY{n}{vectors\PYZus{}A} \PY{o}{@} \PY{n}{np}\PY{o}{.}\PY{n}{diag}\PY{p}{(}\PY{n}{values\PYZus{}A}\PY{p}{)}\PY{p}{)}

\PY{c+c1}{\PYZsh{} plot\PYZus{}pca\PYZus{}components(A, vectors\PYZus{}A, values\PYZus{}A, mode=\PYZdq{}vector\PYZdq{})}
\end{Verbatim}
\end{tcolorbox}

    \begin{Verbatim}[commandchars=\\\{\}]
Eigenvalues (descending): [4. 2.]
Eigenvectors (columns):
 [[1. 0.]
 [0. 1.]]
    \end{Verbatim}

    \begin{tcolorbox}[breakable, size=fbox, boxrule=1pt, pad at break*=1mm,colback=cellbackground, colframe=cellborder]
\prompt{In}{incolor}{29}{\boxspacing}
\begin{Verbatim}[commandchars=\\\{\}]
\PY{c+c1}{\PYZsh{} In [Section 6](\PYZsh{}section\PYZhy{}6), what fraction of the variance does PC1 explain? What happens if you reverse the signs of its eigenvector?}
\PY{n}{display}\PY{p}{(}\PY{n}{Markdown}\PY{p}{(}\PY{l+s+s2}{\PYZdq{}}\PY{l+s+s2}{\PYZsh{}\PYZsh{}\PYZsh{} Exercice 4}\PY{l+s+s2}{\PYZdq{}}\PY{p}{)}\PY{p}{)}

\PY{n}{display}\PY{p}{(}\PY{n}{Markdown}\PY{p}{(}\PY{l+s+sa}{f}\PY{l+s+s2}{\PYZdq{}}\PY{l+s+s2}{**Part 1: What fraction of the variance does PC1 explain?**}\PY{l+s+se}{\PYZbs{}n}\PY{l+s+se}{\PYZbs{}n}\PY{l+s+s2}{The first PCA component captures }\PY{l+s+si}{\PYZob{}}\PY{n}{explained\PYZus{}variance\PYZus{}ratio}\PY{p}{[}\PY{l+m+mi}{0}\PY{p}{]}\PY{l+s+si}{:}\PY{l+s+s2}{.1\PYZpc{}}\PY{l+s+si}{\PYZcb{}}\PY{l+s+s2}{ of the variance.}\PY{l+s+se}{\PYZbs{}n}\PY{l+s+se}{\PYZbs{}n}\PY{l+s+s2}{**Part 2: What happens if you reverse the signs of its eigenvector?**}\PY{l+s+se}{\PYZbs{}n}\PY{l+s+se}{\PYZbs{}n}\PY{l+s+s2}{If I reverse the value of the first eigenvector}\PY{l+s+s2}{\PYZdq{}}\PY{p}{)}\PY{p}{)}

\PY{c+c1}{\PYZsh{}first\PYZus{}eigen\PYZus{}vector = X\PYZus{}eigen\PYZus{}vectors[:, 0]}
\PY{c+c1}{\PYZsh{}first\PYZus{}eigen\PYZus{}vector\PYZus{}negative = \PYZhy{}first\PYZus{}eigen\PYZus{}vector}
\PY{c+c1}{\PYZsh{}new\PYZus{}eigen\PYZus{}matrix = np.column\PYZus{}stack((first\PYZus{}eigen\PYZus{}vector\PYZus{}negative, X\PYZus{}eigen\PYZus{}vectors[:, 1:]))}
\PY{n}{new\PYZus{}eigen\PYZus{}vectors} \PY{o}{=} \PY{n}{np}\PY{o}{.}\PY{n}{column\PYZus{}stack}\PY{p}{(}\PY{p}{(}
    \PY{o}{\PYZhy{}}\PY{n}{X\PYZus{}eigen\PYZus{}vectors}\PY{p}{[}\PY{p}{:}\PY{p}{,} \PY{l+m+mi}{0}\PY{p}{]}\PY{p}{,} \PY{c+c1}{\PYZsh{} The negative of the first eigenvector}
    \PY{n}{X\PYZus{}eigen\PYZus{}vectors}\PY{p}{[}\PY{p}{:}\PY{p}{,} \PY{l+m+mi}{1}\PY{p}{:}\PY{p}{]} \PY{c+c1}{\PYZsh{} The second eigenvector remains the same}
\PY{p}{)}\PY{p}{)}

\PY{n}{display}\PY{p}{(}\PY{n}{Markdown}\PY{p}{(}
    \PY{l+s+sa}{f}\PY{l+s+s2}{\PYZdq{}}\PY{l+s+s2}{The original eigenvectors are:}\PY{l+s+se}{\PYZbs{}n}\PY{l+s+se}{\PYZbs{}n}\PY{l+s+s2}{\PYZdq{}}
    \PY{l+s+sa}{f}\PY{l+s+s2}{\PYZdq{}}\PY{l+s+s2}{\PYZdl{}\PYZdl{}}\PY{l+s+si}{\PYZob{}}\PY{n}{latex}\PY{p}{(}\PY{n}{Matrix}\PY{p}{(}\PY{n}{X\PYZus{}eigen\PYZus{}vectors}\PY{p}{[}\PY{p}{:}\PY{p}{,}\PY{+w}{ }\PY{l+m+mi}{0}\PY{p}{]}\PY{p}{)}\PY{p}{)}\PY{l+s+si}{\PYZcb{}}\PY{l+s+s2}{\PYZdl{}\PYZdl{}}\PY{l+s+se}{\PYZbs{}n}\PY{l+s+se}{\PYZbs{}n}\PY{l+s+s2}{\PYZdq{}}
    \PY{l+s+sa}{f}\PY{l+s+s2}{\PYZdq{}}\PY{l+s+s2}{We extract the first eigenvector (the most relevant, as it has the largest eigenvalue).}\PY{l+s+se}{\PYZbs{}n}\PY{l+s+se}{\PYZbs{}n}\PY{l+s+s2}{\PYZdq{}}
    \PY{l+s+sa}{f}\PY{l+s+s2}{\PYZdq{}}\PY{l+s+s2}{The negative of the first eigenvector is:}\PY{l+s+se}{\PYZbs{}n}\PY{l+s+se}{\PYZbs{}n}\PY{l+s+s2}{\PYZdq{}}
    \PY{l+s+sa}{f}\PY{l+s+s2}{\PYZdq{}}\PY{l+s+s2}{\PYZdl{}\PYZdl{}}\PY{l+s+si}{\PYZob{}}\PY{n}{latex}\PY{p}{(}\PY{n}{Matrix}\PY{p}{(}\PY{o}{\PYZhy{}}\PY{n}{X\PYZus{}eigen\PYZus{}vectors}\PY{p}{[}\PY{p}{:}\PY{p}{,}\PY{+w}{ }\PY{l+m+mi}{0}\PY{p}{]}\PY{p}{)}\PY{p}{)}\PY{l+s+si}{\PYZcb{}}\PY{l+s+s2}{\PYZdl{}\PYZdl{}}\PY{l+s+se}{\PYZbs{}n}\PY{l+s+se}{\PYZbs{}n}\PY{l+s+s2}{\PYZdq{}}
    \PY{l+s+sa}{f}\PY{l+s+s2}{\PYZdq{}}\PY{l+s+s2}{The new eigenmatrix is \PYZdl{}\PYZdl{}}\PY{l+s+si}{\PYZob{}}\PY{n}{latex}\PY{p}{(}\PY{n}{Matrix}\PY{p}{(}\PY{n}{new\PYZus{}eigen\PYZus{}vectors}\PY{p}{)}\PY{p}{)}\PY{l+s+si}{\PYZcb{}}\PY{l+s+s2}{\PYZdl{}\PYZdl{}}\PY{l+s+s2}{\PYZdq{}}
\PY{p}{)}\PY{p}{)}

\PY{n}{new\PYZus{}smaller\PYZus{}size} \PY{o}{=} \PY{p}{(}\PY{l+m+mi}{5}\PY{p}{,} \PY{l+m+mi}{4}\PY{p}{)} \PY{c+c1}{\PYZsh{} To fit 2 in one page}
\PY{n}{display}\PY{p}{(}\PY{n}{Markdown}\PY{p}{(}\PY{l+s+s2}{\PYZdq{}}\PY{l+s+s2}{**Original eigenvectors**}\PY{l+s+s2}{\PYZdq{}}\PY{p}{)}\PY{p}{)}
\PY{n}{plot\PYZus{}pca\PYZus{}components}\PY{p}{(}\PY{n}{X}\PY{p}{,} \PY{n}{X\PYZus{}eigen\PYZus{}vectors}\PY{p}{,} \PY{n}{X\PYZus{}eigen\PYZus{}values}\PY{p}{,}
              \PY{n}{mode}\PY{o}{=}\PY{l+s+s2}{\PYZdq{}}\PY{l+s+s2}{vector}\PY{l+s+s2}{\PYZdq{}}\PY{p}{,} \PY{n}{components}\PY{o}{=}\PY{p}{[}\PY{l+m+mi}{0}\PY{p}{]}\PY{p}{,}
              \PY{n}{figsize}\PY{o}{=}\PY{n}{new\PYZus{}smaller\PYZus{}size}\PY{p}{,}
\PY{p}{)}

\PY{n}{display}\PY{p}{(}\PY{n}{Markdown}\PY{p}{(}\PY{l+s+s2}{\PYZdq{}}\PY{l+s+s2}{**Reversed eigenvectors**}\PY{l+s+s2}{\PYZdq{}}\PY{p}{)}\PY{p}{)}
\PY{n}{plot\PYZus{}pca\PYZus{}components}\PY{p}{(}\PY{n}{X}\PY{p}{,} \PY{n}{new\PYZus{}eigen\PYZus{}vectors}\PY{p}{,} \PY{n}{X\PYZus{}eigen\PYZus{}values}\PY{p}{,}
              \PY{n}{mode}\PY{o}{=}\PY{l+s+s2}{\PYZdq{}}\PY{l+s+s2}{vector}\PY{l+s+s2}{\PYZdq{}}\PY{p}{,} \PY{n}{components}\PY{o}{=}\PY{p}{[}\PY{l+m+mi}{0}\PY{p}{]}\PY{p}{,}
              \PY{n}{figsize}\PY{o}{=}\PY{n}{new\PYZus{}smaller\PYZus{}size}\PY{p}{,}
\PY{p}{)}

\PY{n}{display}\PY{p}{(}\PY{n}{Markdown}\PY{p}{(}\PY{l+s+sa}{f}\PY{l+s+s2}{\PYZdq{}}\PY{l+s+s2}{The first PCA component captures }\PY{l+s+si}{\PYZob{}}\PY{n}{explained\PYZus{}variance\PYZus{}ratio}\PY{p}{[}\PY{l+m+mi}{0}\PY{p}{]}\PY{l+s+si}{:}\PY{l+s+s2}{.1\PYZpc{}}\PY{l+s+si}{\PYZcb{}}\PY{l+s+s2}{ of the variance.}\PY{l+s+se}{\PYZbs{}n}\PY{l+s+s2}{If I reverse the value of the first eigenvector, the variance captured remains the same.}\PY{l+s+s2}{\PYZdq{}}\PY{p}{)}\PY{p}{)}

\PY{n}{original\PYZus{}scores} \PY{o}{=} \PY{n}{X\PYZus{}centered} \PY{o}{@} \PY{n}{X\PYZus{}eigen\PYZus{}vectors}\PY{p}{[}\PY{p}{:}\PY{p}{,} \PY{l+m+mi}{0}\PY{p}{]}
\PY{n}{reversed\PYZus{}scores} \PY{o}{=} \PY{n}{X\PYZus{}centered} \PY{o}{@} \PY{n}{new\PYZus{}eigen\PYZus{}vectors}\PY{p}{[}\PY{p}{:}\PY{p}{,} \PY{l+m+mi}{0}\PY{p}{]}

\PY{k}{assert} \PY{n}{np}\PY{o}{.}\PY{n}{allclose}\PY{p}{(}\PY{n}{reversed\PYZus{}scores}\PY{p}{,} \PY{o}{\PYZhy{}}\PY{n}{original\PYZus{}scores}\PY{p}{)}
\PY{n+nb}{print}\PY{p}{(}\PY{l+s+s2}{\PYZdq{}}\PY{l+s+s2}{PC1 scores reverse sign:}\PY{l+s+s2}{\PYZdq{}}\PY{p}{,} \PY{n}{np}\PY{o}{.}\PY{n}{allclose}\PY{p}{(}\PY{n}{reversed\PYZus{}scores}\PY{p}{,} \PY{o}{\PYZhy{}}\PY{n}{original\PYZus{}scores}\PY{p}{)}\PY{p}{)}
\end{Verbatim}
\end{tcolorbox}

    \subsubsection{Exercice 4}\label{exercice-4}

    
    \textbf{Part 1: What fraction of the variance does PC1 explain?}

The first PCA component captures 93.2\% of the variance.

\textbf{Part 2: What happens if you reverse the signs of its
eigenvector?}

If I reverse the value of the first eigenvector

    
    The original eigenvectors are:

\[\left[\begin{matrix}-0.805383657416694\\-0.592753881780718\end{matrix}\right]\]

We extract the first eigenvector (the most relevant, as it has the
largest eigenvalue).

The negative of the first eigenvector is:

\[\left[\begin{matrix}0.805383657416694\\0.592753881780718\end{matrix}\right]\]

The new eigenmatrix is
\[\left[\begin{matrix}0.805383657416694 & 0.592753881780718\\0.592753881780718 & -0.805383657416694\end{matrix}\right]\]

    
    \textbf{Original eigenvectors}

    
    \begin{center}
    \adjustimage{max size={0.9\linewidth}{0.9\paperheight}}{week4_pdf_build_pca_files/week4_pdf_build_pca_48_4.png}
    \end{center}
    { \hspace*{\fill} \\}
    
    \textbf{Reversed eigenvectors}

    
    \begin{center}
    \adjustimage{max size={0.9\linewidth}{0.9\paperheight}}{week4_pdf_build_pca_files/week4_pdf_build_pca_48_6.png}
    \end{center}
    { \hspace*{\fill} \\}
    
    The first PCA component captures 93.2\% of the variance. If I reverse
the value of the first eigenvector, the variance captured remains the
same.

    
    \begin{Verbatim}[commandchars=\\\{\}]
PC1 scores reverse sign: True
    \end{Verbatim}


    \subsubsection{Conclusions exercise 4}\label{conclusions-exercise-4}

\textbf{Part 1: What fraction of the variance does PC1 explain?}

The first principal component (PC1) explains the proportion of total
variance given by its eigenvalue divided by the sum of all eigenvalues:

\[
\text{Explained Variance Ratio}_{PC1}
=
\frac{\lambda_1}{\sum_{i=1}^{n}\lambda_i}
\]

\textbf{Part 2: What happens if you reverse the sign of its
eigenvector?}

Reversing the sign of the first eigenvector does not change the variance
explained by PC1 or its principal axis. It only reverses the orientation
of the eigenvector and the signs of the corresponding PCA scores.

The visual effect depends on the plotting mode:

\begin{itemize}
\tightlist
\item
  \textbf{Axis mode (\texttt{\textless{}-\textgreater{}}):} The plot
  remains identical because the principal component is represented by a
  double-headed arrow. Reversing the eigenvector simply swaps the
  endpoints.
\item
  \textbf{Vector mode (\texttt{-\textgreater{}}):} The arrow points in
  the opposite direction, making the sign reversal visible. However, its
  length and the variance explained remain unchanged.
\end{itemize}

\textbf{Conclusion:} Eigenvector signs are arbitrary in PCA. Both
\(\mathbf{v}\) and \(-\mathbf{v}\) represent the same principal
component axis and explain the same amount of variance.


    \begin{tcolorbox}[breakable, size=fbox, boxrule=1pt, pad at break*=1mm,colback=cellbackground, colframe=cellborder]
\prompt{In}{incolor}{30}{\boxspacing}
\begin{Verbatim}[commandchars=\\\{\}]
\PY{c+c1}{\PYZsh{} 5. Change the data covariance in [Section 6](\PYZsh{}section\PYZhy{}6) to `[[2, 0], [0, 2]]`. Why might the principal directions now be unstable across samples?}
\PY{n}{display}\PY{p}{(}\PY{n}{Markdown}\PY{p}{(}\PY{l+s+s2}{\PYZdq{}}\PY{l+s+s2}{\PYZsh{}\PYZsh{}\PYZsh{} Exercice 5}\PY{l+s+s2}{\PYZdq{}}\PY{p}{)}\PY{p}{)}
\PY{n}{S\PYZus{}modified} \PY{o}{=} \PY{n}{np}\PY{o}{.}\PY{n}{array}\PY{p}{(}\PY{p}{[}\PY{p}{[}\PY{l+m+mi}{2}\PY{p}{,} \PY{l+m+mi}{0}\PY{p}{]}\PY{p}{,} \PY{p}{[}\PY{l+m+mi}{0}\PY{p}{,} \PY{l+m+mi}{2}\PY{p}{]}\PY{p}{]}\PY{p}{)}

\PY{c+c1}{\PYZsh{} eigh returns eigenvalues in ascending order.}
\PY{n}{modified\PYZus{}eigenvalues}\PY{p}{,} \PY{n}{modified\PYZus{}eigenvectors} \PY{o}{=} \PY{n}{np}\PY{o}{.}\PY{n}{linalg}\PY{o}{.}\PY{n}{eigh}\PY{p}{(}\PY{n}{S\PYZus{}modified}\PY{p}{)}

\PY{n}{modified\PYZus{}order} \PY{o}{=} \PY{n}{np}\PY{o}{.}\PY{n}{argsort}\PY{p}{(}\PY{n}{modified\PYZus{}eigenvalues}\PY{p}{)}\PY{p}{[}\PY{p}{:}\PY{p}{:}\PY{o}{\PYZhy{}}\PY{l+m+mi}{1}\PY{p}{]}
\PY{n}{modified\PYZus{}eigenvalues} \PY{o}{=} \PY{n}{modified\PYZus{}eigenvalues}\PY{p}{[}\PY{n}{modified\PYZus{}order}\PY{p}{]}
\PY{n}{modified\PYZus{}eigenvectors} \PY{o}{=} \PY{n}{modified\PYZus{}eigenvectors}\PY{p}{[}\PY{p}{:}\PY{p}{,} \PY{n}{modified\PYZus{}order}\PY{p}{]}

\PY{n}{modified\PYZus{}explained\PYZus{}variance\PYZus{}ratio} \PY{o}{=} \PY{p}{(}
    \PY{n}{modified\PYZus{}eigenvalues} \PY{o}{/} \PY{n}{modified\PYZus{}eigenvalues}\PY{o}{.}\PY{n}{sum}\PY{p}{(}\PY{p}{)}
\PY{p}{)}

\PY{n+nb}{print}\PY{p}{(}\PY{l+s+s2}{\PYZdq{}}\PY{l+s+s2}{Means:}\PY{l+s+s2}{\PYZdq{}}\PY{p}{,} \PY{n}{X}\PY{o}{.}\PY{n}{mean}\PY{p}{(}\PY{n}{axis}\PY{o}{=}\PY{l+m+mi}{0}\PY{p}{)}\PY{p}{)} \PY{c+c1}{\PYZsh{} axis=0 (Column\PYZhy{}wise). Collapses the rows vertically, calculating the mean of each column.}
\PY{n+nb}{print}\PY{p}{(}\PY{l+s+s2}{\PYZdq{}}\PY{l+s+s2}{Covariance matrix:}\PY{l+s+se}{\PYZbs{}n}\PY{l+s+s2}{\PYZdq{}}\PY{p}{,} \PY{n}{S\PYZus{}modified}\PY{p}{)}
\PY{n+nb}{print}\PY{p}{(}\PY{l+s+s2}{\PYZdq{}}\PY{l+s+s2}{Eigenvalues (descending):}\PY{l+s+s2}{\PYZdq{}}\PY{p}{,} \PY{n}{modified\PYZus{}eigenvalues}\PY{p}{)}
\PY{n+nb}{print}\PY{p}{(}\PY{l+s+s2}{\PYZdq{}}\PY{l+s+s2}{Eigenvectors (columns):}\PY{l+s+se}{\PYZbs{}n}\PY{l+s+s2}{\PYZdq{}}\PY{p}{,} \PY{n}{modified\PYZus{}eigenvectors}\PY{p}{)}
\PY{n+nb}{print}\PY{p}{(}\PY{l+s+s2}{\PYZdq{}}\PY{l+s+s2}{Explained variance ratios:}\PY{l+s+s2}{\PYZdq{}}\PY{p}{,} \PY{n}{modified\PYZus{}explained\PYZus{}variance\PYZus{}ratio}\PY{p}{)}
\PY{n+nb}{print}\PY{p}{(}
    \PY{l+s+s2}{\PYZdq{}}\PY{l+s+s2}{Total variance = trace(S\PYZus{}modified) = sum(eigenvalues):}\PY{l+s+s2}{\PYZdq{}}\PY{p}{,}
    \PY{n}{np}\PY{o}{.}\PY{n}{trace}\PY{p}{(}\PY{n}{S\PYZus{}modified}\PY{p}{)}\PY{p}{,}
    \PY{n}{modified\PYZus{}eigenvalues}\PY{o}{.}\PY{n}{sum}\PY{p}{(}\PY{p}{)}\PY{p}{,}
\PY{p}{)}

\PY{k}{assert} \PY{n}{np}\PY{o}{.}\PY{n}{allclose}\PY{p}{(}
    \PY{n}{S\PYZus{}modified} \PY{o}{@} \PY{n}{modified\PYZus{}eigenvectors}\PY{p}{,}
    \PY{n}{modified\PYZus{}eigenvectors} \PY{o}{@} \PY{n}{np}\PY{o}{.}\PY{n}{diag}\PY{p}{(}\PY{n}{modified\PYZus{}eigenvalues}\PY{p}{)}\PY{p}{,}
\PY{p}{)}
\end{Verbatim}
\end{tcolorbox}

    \subsubsection{Exercice 5}\label{exercice-5}

    
    \begin{Verbatim}[commandchars=\\\{\}]
Means: [2.2273 3.1063]
Covariance matrix:
 [[2 0]
 [0 2]]
Eigenvalues (descending): [2. 2.]
Eigenvectors (columns):
 [[0. 1.]
 [1. 0.]]
Explained variance ratios: [0.5 0.5]
Total variance = trace(S\_modified) = sum(eigenvalues): 4 4.0
    \end{Verbatim}


    \subsubsection{Exercise 5 solution}\label{exercise-5-solution}

\textbf{What happens when the covariance matrix is changed to
\(\begin{bmatrix}2&0\\0&2\end{bmatrix}\)? Why might the principal
directions be unstable across samples?}

Both eigenvalues are equal:

\[
\lambda_1=\lambda_2=2
\]

Therefore, each principal component explains the same proportion of the
total variance:

\[
\frac{2}{2+2}=50\%
\]

Since the covariance matrix is a multiple of the identity matrix
(\(S=2I\)), every nonzero vector is an eigenvector. Consequently, there
is no uniquely preferred principal direction.

Note: for this exercise \(I = \begin{bmatrix}1&0\\0&1\end{bmatrix}\) and
\(S = \begin{bmatrix}2&0\\0&2\end{bmatrix}\). Therefore \(S=2I\)

Small differences between samples can produce slightly different
covariance matrices, causing the principal component directions to
change significantly.

\textbf{Conclusion:} PCA remains mathematically valid, but it cannot
identify a unique dominant direction. Reducing the dataset from two
dimensions to one would retain only 50\% of its variance, regardless of
the chosen direction.


    \begin{tcolorbox}[breakable, size=fbox, boxrule=1pt, pad at break*=1mm,colback=cellbackground, colframe=cellborder]
\prompt{In}{incolor}{31}{\boxspacing}
\begin{Verbatim}[commandchars=\\\{\}]
\PY{c+c1}{\PYZsh{} StandardScaler() uses the population standard deviation (ddof=0), while the covariance calculation in this notebook uses ddof=1. Therefore, the resulting sample variances will be \PYZbs{}(n/(n\PYZhy{}1)\PYZbs{}), not exactly 1.}
\PY{n}{use\PYZus{}standard\PYZus{}scaler\PYZus{}bool} \PY{o}{=} \PY{k+kc}{False}

\PY{n}{display}\PY{p}{(}\PY{n}{Markdown}\PY{p}{(}\PY{l+s+s2}{\PYZdq{}}\PY{l+s+s2}{\PYZsh{}\PYZsh{}\PYZsh{} Bonus exercise}\PY{l+s+s2}{\PYZdq{}}\PY{p}{)}\PY{p}{)}

\PY{c+c1}{\PYZsh{} Original dataset}
\PY{n}{S\PYZus{}original} \PY{o}{=} \PY{n}{np}\PY{o}{.}\PY{n}{cov}\PY{p}{(}\PY{n}{X}\PY{p}{,} \PY{n}{rowvar}\PY{o}{=}\PY{k+kc}{False}\PY{p}{,} \PY{n}{ddof}\PY{o}{=}\PY{l+m+mi}{1}\PY{p}{)}
\PY{n}{values\PYZus{}original}\PY{p}{,} \PY{n}{vectors\PYZus{}original} \PY{o}{=} \PY{n}{np}\PY{o}{.}\PY{n}{linalg}\PY{o}{.}\PY{n}{eigh}\PY{p}{(}\PY{n}{S\PYZus{}original}\PY{p}{)}

\PY{k}{if} \PY{n}{use\PYZus{}standard\PYZus{}scaler\PYZus{}bool}\PY{p}{:}
    \PY{k+kn}{from}\PY{+w}{ }\PY{n+nn}{sklearn}\PY{n+nn}{.}\PY{n+nn}{preprocessing}\PY{+w}{ }\PY{k+kn}{import} \PY{n}{StandardScaler}

    \PY{c+c1}{\PYZsh{} Standardise each column}
    \PY{n}{scaler} \PY{o}{=} \PY{n}{StandardScaler}\PY{p}{(}\PY{p}{)}
    \PY{n}{X\PYZus{}standardised} \PY{o}{=} \PY{n}{scaler}\PY{o}{.}\PY{n}{fit\PYZus{}transform}\PY{p}{(}\PY{n}{X}\PY{p}{)} \PY{c+c1}{\PYZsh{} Internally this uses ddof=0}
\PY{k}{else}\PY{p}{:}
    \PY{n}{X\PYZus{}standardised} \PY{o}{=} \PY{p}{(}\PY{n}{X} \PY{o}{\PYZhy{}} \PY{n}{X}\PY{o}{.}\PY{n}{mean}\PY{p}{(}\PY{n}{axis}\PY{o}{=}\PY{l+m+mi}{0}\PY{p}{)}\PY{p}{)} \PY{o}{/} \PY{n}{X}\PY{o}{.}\PY{n}{std}\PY{p}{(}\PY{n}{axis}\PY{o}{=}\PY{l+m+mi}{0}\PY{p}{,} \PY{n}{ddof}\PY{o}{=}\PY{l+m+mi}{1}\PY{p}{)}

\PY{c+c1}{\PYZsh{} Calculate the new covariance matrix}
\PY{n}{S\PYZus{}standardised} \PY{o}{=} \PY{n}{np}\PY{o}{.}\PY{n}{cov}\PY{p}{(}\PY{n}{X\PYZus{}standardised}\PY{p}{,} \PY{n}{rowvar}\PY{o}{=}\PY{k+kc}{False}\PY{p}{,} \PY{n}{ddof}\PY{o}{=}\PY{l+m+mi}{1}\PY{p}{)}
\PY{n}{values\PYZus{}standardised}\PY{p}{,} \PY{n}{vectors\PYZus{}standardised} \PY{o}{=} \PY{n}{np}\PY{o}{.}\PY{n}{linalg}\PY{o}{.}\PY{n}{eigh}\PY{p}{(}\PY{n}{S\PYZus{}standardised}\PY{p}{)}

\PY{c+c1}{\PYZsh{} Sort both sets of eigenvalues and eigenvectors}
\PY{n}{order\PYZus{}original} \PY{o}{=} \PY{n}{np}\PY{o}{.}\PY{n}{argsort}\PY{p}{(}\PY{n}{values\PYZus{}original}\PY{p}{)}\PY{p}{[}\PY{p}{:}\PY{p}{:}\PY{o}{\PYZhy{}}\PY{l+m+mi}{1}\PY{p}{]}
\PY{n}{values\PYZus{}original} \PY{o}{=} \PY{n}{values\PYZus{}original}\PY{p}{[}\PY{n}{order\PYZus{}original}\PY{p}{]}
\PY{n}{vectors\PYZus{}original} \PY{o}{=} \PY{n}{vectors\PYZus{}original}\PY{p}{[}\PY{p}{:}\PY{p}{,} \PY{n}{order\PYZus{}original}\PY{p}{]}

\PY{n}{order\PYZus{}standardised} \PY{o}{=} \PY{n}{np}\PY{o}{.}\PY{n}{argsort}\PY{p}{(}\PY{n}{values\PYZus{}standardised}\PY{p}{)}\PY{p}{[}\PY{p}{:}\PY{p}{:}\PY{o}{\PYZhy{}}\PY{l+m+mi}{1}\PY{p}{]}
\PY{n}{values\PYZus{}standardised} \PY{o}{=} \PY{n}{values\PYZus{}standardised}\PY{p}{[}\PY{n}{order\PYZus{}standardised}\PY{p}{]}
\PY{n}{vectors\PYZus{}standardised} \PY{o}{=} \PY{n}{vectors\PYZus{}standardised}\PY{p}{[}\PY{p}{:}\PY{p}{,} \PY{n}{order\PYZus{}standardised}\PY{p}{]}

\PY{c+c1}{\PYZsh{} Explained variance ratios}
\PY{n}{ratio\PYZus{}original} \PY{o}{=} \PY{n}{values\PYZus{}original} \PY{o}{/} \PY{n}{values\PYZus{}original}\PY{o}{.}\PY{n}{sum}\PY{p}{(}\PY{p}{)}
\PY{n}{ratio\PYZus{}standardised} \PY{o}{=} \PY{n}{values\PYZus{}standardised} \PY{o}{/} \PY{n}{values\PYZus{}standardised}\PY{o}{.}\PY{n}{sum}\PY{p}{(}\PY{p}{)}

\PY{n+nb}{print}\PY{p}{(}\PY{l+s+s2}{\PYZdq{}}\PY{l+s+s2}{Original covariance matrix:}\PY{l+s+se}{\PYZbs{}n}\PY{l+s+s2}{\PYZdq{}}\PY{p}{,} \PY{n}{S\PYZus{}original}\PY{p}{)}
\PY{n+nb}{print}\PY{p}{(}\PY{l+s+s2}{\PYZdq{}}\PY{l+s+s2}{Standardised covariance matrix:}\PY{l+s+se}{\PYZbs{}n}\PY{l+s+s2}{\PYZdq{}}\PY{p}{,} \PY{n}{S\PYZus{}standardised}\PY{p}{)}

\PY{n+nb}{print}\PY{p}{(}\PY{l+s+s2}{\PYZdq{}}\PY{l+s+se}{\PYZbs{}n}\PY{l+s+s2}{Original eigenvalues:}\PY{l+s+s2}{\PYZdq{}}\PY{p}{,} \PY{n}{values\PYZus{}original}\PY{p}{)}
\PY{n+nb}{print}\PY{p}{(}\PY{l+s+s2}{\PYZdq{}}\PY{l+s+s2}{Standardised eigenvalues:}\PY{l+s+s2}{\PYZdq{}}\PY{p}{,} \PY{n}{values\PYZus{}standardised}\PY{p}{)}

\PY{n+nb}{print}\PY{p}{(}\PY{l+s+s2}{\PYZdq{}}\PY{l+s+se}{\PYZbs{}n}\PY{l+s+s2}{Original eigenvectors:}\PY{l+s+se}{\PYZbs{}n}\PY{l+s+s2}{\PYZdq{}}\PY{p}{,} \PY{n}{vectors\PYZus{}original}\PY{p}{)}
\PY{n+nb}{print}\PY{p}{(}\PY{l+s+s2}{\PYZdq{}}\PY{l+s+s2}{Standardised eigenvectors:}\PY{l+s+se}{\PYZbs{}n}\PY{l+s+s2}{\PYZdq{}}\PY{p}{,} \PY{n}{vectors\PYZus{}standardised}\PY{p}{)}

\PY{n+nb}{print}\PY{p}{(}\PY{l+s+s2}{\PYZdq{}}\PY{l+s+se}{\PYZbs{}n}\PY{l+s+s2}{Original explained variance ratios:}\PY{l+s+s2}{\PYZdq{}}\PY{p}{,} \PY{n}{ratio\PYZus{}original}\PY{p}{)}
\PY{n+nb}{print}\PY{p}{(}\PY{l+s+s2}{\PYZdq{}}\PY{l+s+s2}{Standardised explained variance ratios:}\PY{l+s+s2}{\PYZdq{}}\PY{p}{,} \PY{n}{ratio\PYZus{}standardised}\PY{p}{)}
\end{Verbatim}
\end{tcolorbox}

    \subsubsection{Bonus exercise}\label{bonus-exercise}

    
    \begin{Verbatim}[commandchars=\\\{\}]
Original covariance matrix:
 [[2.6729 1.7538]
 [1.7538 1.5808]]
Standardised covariance matrix:
 [[1.     0.8532]
 [0.8532 1.    ]]

Original eigenvalues: [3.9636 0.29  ]
Standardised eigenvalues: [1.8532 0.1468]

Original eigenvectors:
 [[-0.8054  0.5928]
 [-0.5928 -0.8054]]
Standardised eigenvectors:
 [[ 0.7071 -0.7071]
 [ 0.7071  0.7071]]

Original explained variance ratios: [0.9318 0.0682]
Standardised explained variance ratios: [0.9266 0.0734]
    \end{Verbatim}

    \begin{tcolorbox}[breakable, size=fbox, boxrule=1pt, pad at break*=1mm,colback=cellbackground, colframe=cellborder]
\prompt{In}{incolor}{32}{\boxspacing}
\begin{Verbatim}[commandchars=\\\{\}]
\PY{n}{display}\PY{p}{(}\PY{n}{Markdown}\PY{p}{(}\PY{l+s+s2}{\PYZdq{}}\PY{l+s+s2}{**Original eigenvectors**}\PY{l+s+s2}{\PYZdq{}}\PY{p}{)}\PY{p}{)}
\PY{n}{centre} \PY{o}{=} \PY{n}{plot\PYZus{}pca\PYZus{}components}\PY{p}{(}
    \PY{n}{X}\PY{p}{,} \PY{n}{vectors\PYZus{}original}\PY{p}{,} \PY{n}{values\PYZus{}original}\PY{p}{,}
    \PY{n}{mode}\PY{o}{=}\PY{l+s+s2}{\PYZdq{}}\PY{l+s+s2}{axis}\PY{l+s+s2}{\PYZdq{}}\PY{p}{,}
    \PY{n}{debug\PYZus{}mode}\PY{o}{=}\PY{k+kc}{False}\PY{p}{,}
    \PY{n}{figsize}\PY{o}{=}\PY{n}{new\PYZus{}smaller\PYZus{}size}
\PY{p}{)}

\PY{n}{display}\PY{p}{(}\PY{n}{Markdown}\PY{p}{(}\PY{l+s+s2}{\PYZdq{}}\PY{l+s+s2}{**Standardised eigenvectors**}\PY{l+s+s2}{\PYZdq{}}\PY{p}{)}\PY{p}{)}
\PY{n}{centre} \PY{o}{=}  \PY{n}{plot\PYZus{}pca\PYZus{}components}\PY{p}{(}
    \PY{n}{X\PYZus{}standardised}\PY{p}{,} \PY{n}{vectors\PYZus{}standardised}\PY{p}{,} \PY{n}{values\PYZus{}standardised}\PY{p}{,}
    \PY{n}{mode}\PY{o}{=}\PY{l+s+s2}{\PYZdq{}}\PY{l+s+s2}{axis}\PY{l+s+s2}{\PYZdq{}}\PY{p}{,}
    \PY{n}{debug\PYZus{}mode}\PY{o}{=}\PY{k+kc}{False}\PY{p}{,}
    \PY{n}{figsize}\PY{o}{=}\PY{n}{new\PYZus{}smaller\PYZus{}size}
\PY{p}{)}
\end{Verbatim}
\end{tcolorbox}

    \textbf{Original eigenvectors}

    
    \begin{center}
    \adjustimage{max size={0.9\linewidth}{0.9\paperheight}}{week4_pdf_build_pca_files/week4_pdf_build_pca_53_1.png}
    \end{center}
    { \hspace*{\fill} \\}
    
    \textbf{Standardised eigenvectors}

    
    \begin{center}
    \adjustimage{max size={0.9\linewidth}{0.9\paperheight}}{week4_pdf_build_pca_files/week4_pdf_build_pca_53_3.png}
    \end{center}
    { \hspace*{\fill} \\}
    

    \subsubsection{Bonus: Effect of Standardisation on
PCA}\label{bonus-effect-of-standardisation-on-pca}

After standardising the columns of \texttt{X}, I observed the following
changes:

\begin{itemize}
\tightlist
\item
  \textbf{Covariance matrix:} The diagonal elements changed to 1, since
  both features now have unit variance.
\item
  \textbf{Eigenvalues:} They changed from approximately
  \([3.9636,\ 0.2900]\) to \([1.8532,\ 0.1468]\).
\item
  \textbf{Eigenvectors:} The principal directions rotated because both
  features now contribute equally in terms of their individual
  variances.
\item
  \textbf{Explained variance ratios:} PC1 decreased slightly from
  \textbf{93.18\% to 92.66\%}, a difference of only \textbf{0.52
  percentage points}.
\end{itemize}

The standardised covariance matrix also reveals a strong positive
correlation between the two features:

\[
\rho = 0.8532
\]

\textbf{Conclusion:} Standardisation produced only a modest rotation of
the principal component directions (approximately 8.7°) and a small
reduction in the variance explained by PC1 (0.52 percentage points).
This suggests that the original differences in feature scales had a
relatively limited effect on the PCA results. Both before and after
standardisation, PC1 captures approximately 93\% of the total variance.

This demonstrates that standardisation can change PCA results, although
the effect depends on the original variance and correlation structure of
the dataset.



    \section{Relationship Between Covariance, Eigenvectors and
PCA}\label{relationship-between-covariance-eigenvectors-and-pca}

There is an important relationship between \textbf{eigenvectors and
covariance}, particularly in Principal Component Analysis (PCA).

The key distinctions are:

\begin{itemize}
\tightlist
\item
  \textbf{Covariance} describes how variables vary together.
\item
  \textbf{Eigenvectors} of the covariance matrix identify the principal
  directions in which the data varies.
\item
  \textbf{Eigenvalues} quantify how much variance exists along each of
  those directions.
\end{itemize}

\subsection{1. What does the covariance matrix tell
us?}\label{what-does-the-covariance-matrix-tell-us}

For a dataset with two features, \(X\) and \(Y\), the covariance matrix
is:

\[
S =
\begin{bmatrix}
\operatorname{Var}(X) & \operatorname{Cov}(X,Y) \\
\operatorname{Cov}(X,Y) & \operatorname{Var}(Y)
\end{bmatrix}
\]

It contains two kinds of information:

\begin{itemize}
\tightlist
\item
  \textbf{Diagonal elements:} The variance of each feature.
\item
  \textbf{Off-diagonal elements:} How the features vary together.
\end{itemize}

For example:

\[
S =
\begin{bmatrix}
4 & 2 \\
2 & 4
\end{bmatrix}
\]

This tells us that:

\begin{itemize}
\tightlist
\item
  Both features have a variance of \(4\).
\item
  Their covariance is positive (\(2\)).
\item
  The features tend to increase and decrease together.
\end{itemize}

However, the covariance matrix does not directly identify the directions
of maximum and minimum variance. That is where \textbf{eigenvectors}
become useful.

\subsection{2. What do eigenvectors tell us about
covariance?}\label{what-do-eigenvectors-tell-us-about-covariance}

When we calculate the eigenvectors of a covariance matrix, we identify
new directions along which the data varies.

Consider:

\[
S =
\begin{bmatrix}
4 & 2 \\
2 & 4
\end{bmatrix}
\]

Its eigenvalues and corresponding unit eigenvectors are:

\textbf{First principal component (PC1):}

\[
\lambda_1 = 6,
\qquad
\mathbf{v}_1 =
\frac{1}{\sqrt{2}}
\begin{bmatrix}
1 \\
1
\end{bmatrix}
\]

\textbf{Second principal component (PC2):}

\[
\lambda_2 = 2,
\qquad
\mathbf{v}_2 =
\frac{1}{\sqrt{2}}
\begin{bmatrix}
1 \\
-1
\end{bmatrix}
\]

\subsubsection{Addendum: Calculating the eigenvalues step by
step}\label{addendum-calculating-the-eigenvalues-step-by-step}

Given the covariance matrix:

\[
S =
\begin{bmatrix}
4 & 2 \\
2 & 4
\end{bmatrix}
\]

The eigenvalues satisfy the \textbf{characteristic equation}:

\[
\det(S-\lambda I)=0
\]

\paragraph{\texorpdfstring{Step 1: Subtract
\(\lambda I\)}{Step 1: Subtract \textbackslash lambda I}}\label{step-1-subtract-lambda-i}

\[
S-\lambda I =
\begin{bmatrix}
4-\lambda & 2 \\
2 & 4-\lambda
\end{bmatrix}
\]

\paragraph{Step 2: Calculate the
determinant}\label{step-2-calculate-the-determinant}

For a \(2\times2\) matrix:

\[
\det
\begin{bmatrix}
a & b \\
c & d
\end{bmatrix}
=ad-bc
\]

Therefore:

\[
\det(S-\lambda I)
=
(4-\lambda)(4-\lambda)-(2)(2)
\]

Setting this equal to zero:

\[
(4-\lambda)^2-4=0
\]

\paragraph{\texorpdfstring{Step 3: Solve for
\(\lambda\)}{Step 3: Solve for \textbackslash lambda}}\label{step-3-solve-for-lambda}

\[
(4-\lambda)^2=4
\]

Taking square roots gives both possibilities:

\[
4-\lambda=\pm2
\]

Thus:

\begin{itemize}
\tightlist
\item
  If \(4-\lambda=2\), then \(\lambda=2\).
\item
  If \(4-\lambda=-2\), then \(\lambda=6\).
\end{itemize}

Ordering the eigenvalues from largest to smallest:

\[
\boxed{\lambda_1=6,\qquad\lambda_2=2}
\]

\paragraph{Alternative: Using the quadratic
formula}\label{alternative-using-the-quadratic-formula}

Expanding the characteristic equation:

\[
(4-\lambda)^2-4
=
16-8\lambda+\lambda^2-4
=
\lambda^2-8\lambda+12
\]

We therefore solve:

\[
\lambda^2-8\lambda+12=0
\]

Using \(a=1\), \(b=-8\), and \(c=12\):

\[
\lambda
=
\frac{-b\pm\sqrt{b^2-4ac}}{2a}
=
\frac{8\pm\sqrt{64-48}}{2}
=
\frac{8\pm\sqrt{16}}{2}
=
\frac{8\pm4}{2}
\]

Hence:

\[
\boxed{\lambda_1=6,\qquad\lambda_2=2}
\]

\textbf{Arithmetic check:} The discriminant is \(64-48=16\), and the
denominator divides both terms in the numerator:

\[
\frac{8\pm4}{2}
=
\frac{8}{2}\pm\frac{4}{2}
=
4\pm2
\]

\paragraph{Step 4: Verify using the
eigenvectors}\label{step-4-verify-using-the-eigenvectors}

An eigenvalue and its eigenvector must satisfy:

\[
S\mathbf{v}=\lambda\mathbf{v}
\]

For the first unit eigenvector:

\[
S\left(
\frac{1}{\sqrt{2}}
\begin{bmatrix}
1\\
1
\end{bmatrix}
\right)
=
\frac{1}{\sqrt{2}}
\begin{bmatrix}
6\\
6
\end{bmatrix}
=
6\left(
\frac{1}{\sqrt{2}}
\begin{bmatrix}
1\\
1
\end{bmatrix}
\right)
\]

For the second unit eigenvector:

\[
S\left(
\frac{1}{\sqrt{2}}
\begin{bmatrix}
1\\
-1
\end{bmatrix}
\right)
=
\frac{1}{\sqrt{2}}
\begin{bmatrix}
2\\
-2
\end{bmatrix}
=
2\left(
\frac{1}{\sqrt{2}}
\begin{bmatrix}
1\\
-1
\end{bmatrix}
\right)
\]

Both eigenvalues are confirmed.

Because \(S\) is a covariance matrix, these eigenvalues represent the
variances along the corresponding unit principal directions:

\begin{itemize}
\tightlist
\item
  \textbf{PC1:} variance of \(6\).
\item
  \textbf{PC2:} variance of \(2\).
\end{itemize}

\subsubsection{Addendum: Calculating the eigenvectors step by
step}\label{addendum-calculating-the-eigenvectors-step-by-step}

Given the covariance matrix:

\[
S =
\begin{bmatrix}
4 & 2 \\
2 & 4
\end{bmatrix}
\]

We have already calculated its eigenvalues:

\[
\lambda_1=6,\qquad\lambda_2=2
\]

An eigenvector is a \textbf{nonzero vector} satisfying:

\[
S\mathbf{v}=\lambda\mathbf{v}
\]

Equivalently:

\[
(S-\lambda I)\mathbf{v}=\mathbf{0}
\]

Write the unknown vector as:

\[
\mathbf{v}=
\begin{bmatrix}
x\\
y
\end{bmatrix}
\]

\paragraph{\texorpdfstring{Step 1: Find an eigenvector for
\(\lambda_1=6\)}{Step 1: Find an eigenvector for \textbackslash lambda\_1=6}}\label{step-1-find-an-eigenvector-for-lambda_16}

Substitute \(\lambda_1=6\):

\[
S-6I=
\begin{bmatrix}
-2 & 2\\
2 & -2
\end{bmatrix}
\]

Therefore:

\[
\begin{bmatrix}
-2 & 2\\
2 & -2
\end{bmatrix}
\begin{bmatrix}
x\\
y
\end{bmatrix}
=
\begin{bmatrix}
0\\
0
\end{bmatrix}
\]

This gives:

\[
\begin{cases}
-2x+2y=0\\
2x-2y=0
\end{cases}
\]

Both equations express the same condition:

\[
y=x
\]

Thus, every eigenvector associated with \(\lambda_1=6\) has the form:

\[
\mathbf{v}=
t\begin{bmatrix}
1\\
1
\end{bmatrix},
\qquad t\neq0
\]

Choosing \(t=1\) gives:

\[
\mathbf{u}_1=
\begin{bmatrix}
1\\
1
\end{bmatrix}
\]

\paragraph{Step 2: Normalize the first
eigenvector}\label{step-2-normalize-the-first-eigenvector}

For PCA, we use eigenvectors of length 1.

The length of \(\mathbf{u}_1\) is:

\[
\|\mathbf{u}_1\|
=
\sqrt{1^2+1^2}
=
\sqrt{2}
\]

Divide the vector by its length:

\[
\boxed{
\mathbf{v}_1
=
\frac{\mathbf{u}_1}{\|\mathbf{u}_1\|}
=
\frac{1}{\sqrt{2}}
\begin{bmatrix}
1\\
1
\end{bmatrix}
}
\]

This is the direction of the first principal component, PC1.

\paragraph{\texorpdfstring{Step 3: Find an eigenvector for
\(\lambda_2=2\)}{Step 3: Find an eigenvector for \textbackslash lambda\_2=2}}\label{step-3-find-an-eigenvector-for-lambda_22}

Substitute \(\lambda_2=2\):

\[
S-2I=
\begin{bmatrix}
2 & 2\\
2 & 2
\end{bmatrix}
\]

Therefore:

\[
\begin{bmatrix}
2 & 2\\
2 & 2
\end{bmatrix}
\begin{bmatrix}
x\\
y
\end{bmatrix}
=
\begin{bmatrix}
0\\
0
\end{bmatrix}
\]

This gives:

\[
\begin{cases}
2x+2y=0\\
2x+2y=0
\end{cases}
\]

Both equations imply:

\[
y=-x
\]

Thus, every eigenvector associated with \(\lambda_2=2\) has the form:

\[
\mathbf{v}=
t\begin{bmatrix}
1\\
-1
\end{bmatrix},
\qquad t\neq0
\]

Choosing \(t=1\) gives:

\[
\mathbf{u}_2=
\begin{bmatrix}
1\\
-1
\end{bmatrix}
\]

\paragraph{Step 4: Normalize the second
eigenvector}\label{step-4-normalize-the-second-eigenvector}

Its length is:

\[
\|\mathbf{u}_2\|
=
\sqrt{1^2+(-1)^2}
=
\sqrt{2}
\]

Therefore:

\[
\boxed{
\mathbf{v}_2
=
\frac{\mathbf{u}_2}{\|\mathbf{u}_2\|}
=
\frac{1}{\sqrt{2}}
\begin{bmatrix}
1\\
-1
\end{bmatrix}
}
\]

This is the direction of the second principal component, PC2.

\paragraph{Step 5: Check that the directions are
perpendicular}\label{step-5-check-that-the-directions-are-perpendicular}

The dot product of the unit eigenvectors is:

\[
\mathbf{v}_1^\mathsf{T}\mathbf{v}_2
=
\frac{1}{2}\left(1\cdot1+1\cdot(-1)\right)
=
0
\]

A zero dot product means the directions are perpendicular.

\paragraph{Final result}\label{final-result}

The eigenvalues and corresponding unit eigenvectors are:

\[
\boxed{
\lambda_1=6,
\qquad
\mathbf{v}_1=
\frac{1}{\sqrt{2}}
\begin{bmatrix}
1\\
1
\end{bmatrix}
}
\]

\[
\boxed{
\lambda_2=2,
\qquad
\mathbf{v}_2=
\frac{1}{\sqrt{2}}
\begin{bmatrix}
1\\
-1
\end{bmatrix}
}
\]

PC1 follows the line \(y=x\), along which the variance is 6. PC2 follows
the line \(y=-x\), along which the variance is 2.

\textbf{Sign convention:} Multiplying either eigenvector by \(-1\) gives
an equally valid unit eigenvector with the same eigenvalue. It reverses
the arrow along the same principal axis.

\subsubsection{Visualising the principal component
directions}\label{visualising-the-principal-component-directions}

\begin{figure}
\centering
\pandocbounded{\includesvg[keepaspectratio]{images/week_4/covariance_and_pca.svg}}
\caption{Covariance and PCA principal component directions}
\end{figure}

The diagram illustrates that:

\begin{itemize}
\tightlist
\item
  \textbf{PC1} follows the diagonal where both features increase
  together.
\item
  \textbf{PC2} follows the perpendicular diagonal where one feature
  increases as the other decreases.
\item
  \textbf{PC1 has the larger eigenvalue}, so more variance lies along
  its direction.
\end{itemize}

The eigenvectors effectively \textbf{rotate the coordinate system} to
align with the natural directions of variation in the dataset.

\subsubsection{Explained variance
ratios}\label{explained-variance-ratios}

The proportion of total variance explained by each principal component
is:

\[
\operatorname{EVR}_i =
\frac{\lambda_i}{\sum_{j=1}^{n}\lambda_j}
\]

For PC1:

\[
\operatorname{EVR}_1 =
\frac{6}{6+2}
= 0.75 = 75\%
\]

For PC2:

\[
\operatorname{EVR}_2 =
\frac{2}{6+2}
= 0.25 = 25\%
\]

Therefore, \textbf{PC1 captures 75\% of the total variance}, while PC2
captures the remaining 25\%.

\subsection{3. The mathematical
relationship}\label{the-mathematical-relationship}

The eigenvalue equation for a covariance matrix is:

\[
S\mathbf{v}_i = \lambda_i\mathbf{v}_i
\]

Where:

\begin{itemize}
\tightlist
\item
  \(S\) is the covariance matrix.
\item
  \(\mathbf{v}_i\) is an eigenvector.
\item
  \(\lambda_i\) is its corresponding eigenvalue.
\end{itemize}

For a \textbf{unit eigenvector} \(\mathbf{v}_i\), the variance of the
data projected onto that direction is:

\[
\boxed{
\operatorname{Var}(X_c\mathbf{v}_i)
=
\mathbf{v}_i^T S\mathbf{v}_i
=
\lambda_i
}
\]

Where \(X_c\) is the centred dataset.

This gives us the fundamental relationship:

\textbf{The eigenvalue of a covariance matrix equals the variance of the
data projected along its corresponding unit eigenvector.}

\subsubsection{Why does PCA remove covariance between
components?}\label{why-does-pca-remove-covariance-between-components}

Because covariance matrices are symmetric, their eigenvectors can be
chosen to be orthonormal.

If we collect these eigenvectors into a matrix \(V\), the covariance
matrix in the new coordinate system becomes:

\[
\boxed{
V^T S V = \Lambda
}
\]

Where \(\Lambda\) is a diagonal matrix containing the eigenvalues:

\[
\Lambda =
\begin{bmatrix}
\lambda_1 & 0 \\
0 & \lambda_2
\end{bmatrix}
\]

For our example:

\[
\Lambda =
\begin{bmatrix}
6 & 0 \\
0 & 2
\end{bmatrix}
\]

Notice that the \textbf{off-diagonal elements are zero}.

This means:

\begin{itemize}
\tightlist
\item
  The principal components are uncorrelated.
\item
  Each principal component captures variance along a different,
  perpendicular direction.
\item
  PCA transforms correlated features into uncorrelated principal
  components.
\end{itemize}

\subsection{4. Connecting this to Exercise
5}\label{connecting-this-to-exercise-5}

In Exercise 5, we modified the covariance matrix to:

\[
S =
\begin{bmatrix}
2 & 0 \\
0 & 2
\end{bmatrix}
\]

This matrix is a multiple of the identity matrix:

\[
S = 2I
\]

Both eigenvalues are equal:

\[
\lambda_1 = \lambda_2 = 2
\]

Therefore, the explained variance ratio for each component is:

\[
\operatorname{EVR}_1 =
\operatorname{EVR}_2 =
\frac{2}{2+2}
= 50\%
\]

\subsubsection{Why are the principal directions
unstable?}\label{why-are-the-principal-directions-unstable}

For any nonzero vector \(\mathbf{v}\):

\[
S\mathbf{v} = 2I\mathbf{v} = 2\mathbf{v}
\]

Therefore, \textbf{every nonzero vector is an eigenvector with
eigenvalue \(2\)}.

Consequently:

\begin{itemize}
\tightlist
\item
  There is no uniquely preferred principal direction.
\item
  Every unit direction has the same projected variance.
\item
  Small differences between samples can change which directions PCA
  identifies as PC1 and PC2.
\item
  Reducing from two dimensions to one retains only 50\% of the total
  variance, regardless of the selected direction.
\end{itemize}

This situation is called \textbf{eigenvalue degeneracy}, where repeated
eigenvalues prevent the individual eigenvector directions from being
uniquely determined.

\subsection{5. Key takeaways}\label{key-takeaways}

{\def\LTcaptype{none} % do not increment counter
\begin{longtable}[]{@{}
  >{\raggedright\arraybackslash}p{(\linewidth - 2\tabcolsep) * \real{0.5000}}
  >{\raggedright\arraybackslash}p{(\linewidth - 2\tabcolsep) * \real{0.5000}}@{}}
\toprule\noalign{}
\begin{minipage}[b]{\linewidth}\raggedright
Concept
\end{minipage} & \begin{minipage}[b]{\linewidth}\raggedright
Relationship to PCA
\end{minipage} \\
\midrule\noalign{}
\endhead
\bottomrule\noalign{}
\endlastfoot
Variance & Measures the spread of an individual feature \\
Covariance & Measures how two features vary together \\
Covariance matrix & Describes the variance and covariance structure of
the dataset \\
Eigenvectors & Identify orthogonal directions of variation \\
Eigenvalues & Quantify variance along the corresponding unit
eigenvectors \\
Principal components & New, uncorrelated coordinates obtained by
projecting onto the eigenvectors \\
Explained variance ratio & Measures the fraction of total variance
captured by each component \\
Equal eigenvalues & Indicate that the associated individual principal
directions are not uniquely determined \\
\end{longtable}
}

\subsubsection{Conclusion}\label{conclusion}

\textbf{Covariance describes the spread and relationships of the
original features. Eigenvectors of the covariance matrix identify the
principal directions of that spread, while their eigenvalues quantify
the variance along those directions.}

Eigenvectors are a general linear algebra concept. Their specific
relationship to data variance arises when we calculate the eigenvectors
of a \textbf{covariance matrix}, rather than an arbitrary matrix.

When eigenvalues are distinct, the corresponding principal directions
are uniquely determined up to sign. When eigenvalues are repeated, the
associated principal directions may not be uniquely determined.



    \subsection{10. Summary for your notes}\label{summary-for-your-notes}

{\def\LTcaptype{none} % do not increment counter
\begin{longtable}[]{@{}
  >{\raggedright\arraybackslash}p{(\linewidth - 4\tabcolsep) * \real{0.3333}}
  >{\raggedright\arraybackslash}p{(\linewidth - 4\tabcolsep) * \real{0.3333}}
  >{\raggedright\arraybackslash}p{(\linewidth - 4\tabcolsep) * \real{0.3333}}@{}}
\toprule\noalign{}
\begin{minipage}[b]{\linewidth}\raggedright
Concept
\end{minipage} & \begin{minipage}[b]{\linewidth}\raggedright
Key question
\end{minipage} & \begin{minipage}[b]{\linewidth}\raggedright
Key fact
\end{minipage} \\
\midrule\noalign{}
\endhead
\bottomrule\noalign{}
\endlastfoot
\textbf{Variance} & How much does one feature vary? & Non-negative;
measured in squared units \\
\textbf{Covariance} & How do two features vary together? & Sign shows
direction of linear co-movement \\
\textbf{Correlation} & How strongly do two variables move linearly
together? & Standardised covariance in \([-1,1]\) \\
\textbf{Eigenvector} & Which directions remain on the same line after a
matrix transform? & Satisfies \(A\mathbf v=\lambda\mathbf v\) \\
\textbf{Eigenvalue} & By what factor does a matrix scale an eigenvector?
& For covariance matrices, equals variance along a unit eigenvector \\
\textbf{PCA} & Which axes capture the most variance? & Eigenvectors of
covariance matrix, ordered by eigenvalues \\
\end{longtable}
}

\textbf{Practical note:} PCA is not a clustering method. It is an
unsupervised dimensionality-reduction method that may be used before
clustering algorithms such as DBSCAN or K-means.



    % Add a bibliography block to the postdoc
    
    
    
\end{document}
